\documentclass[pdflatex,sn-mathphys-num]{sn-jnl}

\usepackage[T1]{fontenc}
\usepackage{amsmath,amssymb,amsthm}
\usepackage{graphicx}
\usepackage{bm}
\usepackage{booktabs}

\theoremstyle{thmstyleone}
\newtheorem{axiom}{Axiom}
\newtheorem{theorem}{Theorem}
\newtheorem{proposition}{Proposition}
\newtheorem{corollary}{Corollary}

\theoremstyle{thmstylethree}
\newtheorem{definition}{Definition}

\theoremstyle{thmstyletwo}
\newtheorem{remark}{Remark}

\newcommand{\lean}[1]{\hfill\texttt{\footnotesize [#1]}}

\raggedbottom

\begin{document}

\title[Supplemental Material]{Supplemental Material: Scale-Coupled Dynamics}

%% NOTE: same author/affiliation TODOs as main.tex — keep these two files in sync.
\author*[1]{\fnm{Xinyu} \sur{Yang}}\email{Pana.Yang@hotmail.com}

\affil*[1]{\orgname{Independent Researcher}}

\abstract{This document is the companion to ``Scale-Coupled Dynamics: Gravity and
Charge from a Single Non-Commuting Scale Field.'' It gives, module by
module in the reading order of the formal development, a complete
mathematical restatement of every theorem, definition, and proof
underlying that paper — 644 machine-verified results in total, with zero
admitted gaps. Each item is tagged in brackets with its identifier in the
Lean~4 source for direct cross-reference. The main paper is
self-contained and does not require this document to be read; this
document exists so that every claim made there can be checked in full
formal detail.}

\maketitle
\tableofcontents

\section{Part 0: The Postulates}

\subsection{Basic}

The differential substrate is modeled algebraically through a commutative ring carrying $n$ commuting derivations, which represent partial differentiation in an $n$-dimensional affine coordinate patch. This approach avoids analytic assumptions and allows every differential identity (Christoffel symbols, Riemann tensor, Weyl curvature) to be derived as purely algebraic consequences of the Leibniz rule and commutativity of mixed partials.

\begin{definition}[\lean{Basic.ScaleAlgebra}]
A \emph{scale algebra} of dimension $n$ on a commutative ring $A$ is a pair of data $(\mathbf{d}, A)$ where:
\begin{enumerate}
\item $\mathbf{d} = (\partial_0, \ldots, \partial_{n-1})$ is an $n$-tuple of endomorphisms $\partial_i : A \to A$;
\item each $\partial_i$ satisfies the Leibniz rule: for all $a, b \in A$,
  $$\partial_i(a \cdot b) = \partial_i(a) \cdot b + a \cdot \partial_i(b);$$
\item the mixed partials commute: for all $i, j$ and $a \in A$,
  $$\partial_i(\partial_j(a)) = \partial_j(\partial_i(a)).$$
\end{enumerate}
\end{definition}

\begin{theorem}[\lean{Basic.ScaleAlgebra.d\_zero}]
For any scale algebra and any $i \in \{0,\ldots,n-1\}$, the zero element is annihilated: $\partial_i(0) = 0$.
\end{theorem}
\begin{proof}
By the Leibniz rule applied to $0 + 0 = 0$, we have $\partial_i(0 + 0) = \partial_i(0) + \partial_i(0)$. Since $0 + 0 = 0$, this gives $\partial_i(0) = 2 \cdot \partial_i(0)$, which implies $\partial_i(0) = 0$ by cancellation.
\end{proof}

\begin{theorem}[\lean{Basic.ScaleAlgebra.d\_one}]
For any scale algebra and any $i \in \{0,\ldots,n-1\}$, the unity is annihilated: $\partial_i(1) = 0$.
\end{theorem}
\begin{proof}
By the Leibniz rule applied to $1 \cdot 1 = 1$, we have
$$\partial_i(1) = \partial_i(1) \cdot 1 + 1 \cdot \partial_i(1) = 2 \cdot \partial_i(1).$$
Thus $\partial_i(1) = 0$ by cancellation.
\end{proof}

\begin{definition}[\lean{Basic.ScaleAlgebra.dHom}]
Each derivation $\partial_i$ extends to an additive group homomorphism $\mathbf{d}_i : A \to A$, packaged as the morphism $(A, +) \to (A, +)$ that preserves zero and addition.
\end{definition}

\begin{theorem}[\lean{Basic.ScaleAlgebra.d\_neg}]
For any scale algebra and any $i$ and $a \in A$,
$$\partial_i(-a) = -\partial_i(a).$$
\end{theorem}
\begin{proof}
This is immediate from the fact that $\partial_i$ is an additive group homomorphism.
\end{proof}

\begin{theorem}[\lean{Basic.ScaleAlgebra.d\_sub}]
For any scale algebra and any $i$ and $a, b \in A$,
$$\partial_i(a - b) = \partial_i(a) - \partial_i(b).$$
\end{theorem}
\begin{proof}
This is immediate from the fact that $\partial_i$ is an additive group homomorphism.
\end{proof}

\begin{theorem}[\lean{Basic.ScaleAlgebra.d\_sum}]
For any scale algebra, any finite index set $\iota$, any function $f : \iota \to A$, and any $i \in \{0,\ldots,n-1\}$,
$$\partial_i\left(\sum_{x \in s} f(x)\right) = \sum_{x \in s} \partial_i(f(x)),$$
where $s$ is the underlying finite set of indices.
\end{theorem}
\begin{proof}
This follows from the fact that $\partial_i$ is an additive group homomorphism and preserves finite sums.
\end{proof}

\begin{theorem}[\lean{Basic.ScaleAlgebra.d\_natCast}]
For any scale algebra, any $i \in \{0,\ldots,n-1\}$, and any natural number $k$, the derivation annihilates the image of the natural number under the ring homomorphism: $\partial_i(k) = 0$.
\end{theorem}
\begin{proof}
By induction on $k$: the base case $\partial_i(0) = 0$ is \lean{Basic.ScaleAlgebra.d\_zero}. For the successor, $\partial_i(k+1) = \partial_i(k) + \partial_i(1) = 0 + 0 = 0$ by the Leibniz rule and inductive hypothesis.
\end{proof}

\begin{theorem}[\lean{Basic.ScaleAlgebra.d\_nsmul}]
For any scale algebra, any $i$, any natural number $k$, and any $a \in A$,
$$\partial_i(k \cdot a) = k \cdot \partial_i(a).$$
\end{theorem}
\begin{proof}
By the Leibniz rule, $\partial_i(k \cdot a) = \partial_i(k) \cdot a + k \cdot \partial_i(a) = 0 + k \cdot \partial_i(a)$ since $\partial_i(k) = 0$.
\end{proof}

The Kronecker delta is the bare Euclidean metric tensor, valued in $A$. It is pure bookkeeping and carries no physics.

\begin{definition}[\lean{Basic.kron}]
The Kronecker delta with indices $i, j \in \{0,\ldots,n-1\}$ is defined by
$$\delta_{ij} := \begin{cases} 1 & \text{if } i = j \\ 0 & \text{if } i \ne j. \end{cases}$$
\end{definition}

\begin{theorem}[\lean{Basic.kron\_self}]
For any $i \in \{0,\ldots,n-1\}$, $\delta_{ii} = 1$.
\end{theorem}
\begin{proof}
Immediate from the definition.
\end{proof}

\begin{theorem}[\lean{Basic.kron\_symm}]
For any $i, j \in \{0,\ldots,n-1\}$, $\delta_{ij} = \delta_{ji}$.
\end{theorem}
\begin{proof}
The Kronecker delta depends symmetrically on its arguments: if $i = j$ then $j = i$, and if $i \ne j$ then $j \ne i$.
\end{proof}

\begin{theorem}[\lean{Basic.d\_kron}]
For any scale algebra, any $k, i, j \in \{0,\ldots,n-1\}$, the Kronecker delta is annihilated by every derivation:
$$\partial_k(\delta_{ij}) = 0.$$
\end{theorem}
\begin{proof}
Since $\delta_{ij} \in \{0, 1\} \subset A$, it is a natural number cast into $A$. By \lean{Basic.ScaleAlgebra.d\_natCast}, every derivation annihilates it.
\end{proof}

\begin{theorem}[\lean{Basic.d\_kron\_mul}]
For any scale algebra, any $k, i, j \in \{0,\ldots,n-1\}$, and any $a \in A$,
$$\partial_k(\delta_{ij} \cdot a) = \delta_{ij} \cdot \partial_k(a).$$
\end{theorem}
\begin{proof}
By the Leibniz rule,
$$\partial_k(\delta_{ij} \cdot a) = \partial_k(\delta_{ij}) \cdot a + \delta_{ij} \cdot \partial_k(a) = 0 + \delta_{ij} \cdot \partial_k(a)$$
since $\partial_k(\delta_{ij}) = 0$ by the previous theorem.
\end{proof}

The following four theorems are the only facts about the Kronecker delta needed by the curvature computation.

\begin{theorem}[\lean{Basic.sum\_kron\_left}]
For any scale algebra, any $b \in \{0,\ldots,n-1\}$, and any function $f : \{0,\ldots,n-1\} \to A$,
$$\sum_{a=0}^{n-1} \delta_{ab} \cdot f(a) = f(b).$$
\end{theorem}
\begin{proof}
The sum contains exactly one nonzero term, when $a = b$. In that case, $\delta_{bb} \cdot f(b) = 1 \cdot f(b) = f(b)$.
\end{proof}

\begin{theorem}[\lean{Basic.sum\_kron\_right}]
For any scale algebra, any $b \in \{0,\ldots,n-1\}$, and any function $f : \{0,\ldots,n-1\} \to A$,
$$\sum_{a=0}^{n-1} \delta_{ba} \cdot f(a) = f(b).$$
\end{theorem}
\begin{proof}
By symmetry of the Kronecker delta, $\delta_{ba} = \delta_{ab}$. The sum then contains exactly one nonzero term when $a = b$, giving $\delta_{bb} \cdot f(b) = f(b)$.
\end{proof}

\subsection{Axioms}

This section develops axioms A2, A3, and A5 of the theory, establishing the scale field and its immediate consequences. The scale $\sigma$ is a dimensionless log-scale with multiplicative representative $s = e^\sigma$ (postulated as a unit of the ring); the scale-energy duality $\varepsilon \cdot s = 1$ is derived from this; and the fiducial covariance (invariance under global shift $\sigma \mapsto \sigma + c$) implies that all geometric objects are exact dimensionless.

\begin{definition}[\lean{Axioms.ScaleField}]
A \emph{scale field} of dimension $n$ on a scale algebra over commutative ring $A$ is a triple $(\sigma, s, \mathbf{d}_s)$ where:
\begin{enumerate}
\item $\sigma \in A$ is the log-scale;
\item $s \in A^\times$ is the scale itself, a unit in the ring;
\item $\mathbf{d}_s$ is the condition that $s$ satisfies the differential equation of the exponential: for each $i \in \{0,\ldots,n-1\}$,
  $$\partial_i(s) = s \cdot \partial_i(\sigma).$$
\end{enumerate}
\end{definition}

\begin{definition}[\lean{Axioms.ScaleField.en}]
Given a scale field $F = (\sigma, s, \mathbf{d}_s)$, the energy scale is defined as the reciprocal of the length scale:
$$\varepsilon := s^{-1}.$$
\end{definition}

\begin{theorem}[\lean{Axioms.ScaleField.en\_mul\_scale}]
For any scale field $F$, the scale-energy duality holds: $\varepsilon \cdot s = 1$.
\end{theorem}
\begin{proof}
Immediate from the definition of $\varepsilon$ as $s^{-1}$ in a commutative ring with units.
\end{proof}

\begin{theorem}[\lean{Axioms.ScaleField.scale\_mul\_en}]
For any scale field $F$, we have also $s \cdot \varepsilon = 1$.
\end{theorem}
\begin{proof}
This follows from $\varepsilon \cdot s = 1$ by commutativity of multiplication in $A$.
\end{proof}

\begin{theorem}[\lean{Axioms.ScaleField.d\_en}]
For any scale field $F$ and any $i \in \{0,\ldots,n-1\}$, the energy scale obeys the dual differential equation:
$$\partial_i(\varepsilon) = -\varepsilon \cdot \partial_i(\sigma).$$
\end{theorem}
\begin{proof}
Differentiate the duality relation $\varepsilon \cdot s = 1$ using the Leibniz rule:
$$0 = \partial_i(\varepsilon \cdot s) = \partial_i(\varepsilon) \cdot s + \varepsilon \cdot \partial_i(s) = \partial_i(\varepsilon) \cdot s + \varepsilon \cdot s \cdot \partial_i(\sigma),$$
using the condition that $\partial_i(s) = s \cdot \partial_i(\sigma)$. Thus
$$\partial_i(\varepsilon) \cdot s = -\varepsilon \cdot s \cdot \partial_i(\sigma).$$
Multiply both sides on the right by $\varepsilon$ (which equals $s^{-1}$):
$$\partial_i(\varepsilon) = -\varepsilon \cdot \partial_i(\sigma).$$
\end{proof}

A global additive shift in the fiducial unit leaves every geometric object constructed from the scale invariant. Let $c \in A$ be a constant (meaning $\partial_i(c) = 0$ for all $i$).

\begin{theorem}[\lean{Axioms.sig\_add\_const}]
The first-order operator on the shifted scale obeys $\partial_i(\sigma + c) = \partial_i(\sigma)$.
\end{theorem}
\begin{proof}
By the Leibniz rule and the fact that $c$ is constant, $\partial_i(\sigma + c) = \partial_i(\sigma) + \partial_i(c) = \partial_i(\sigma) + 0 = \partial_i(\sigma)$.
\end{proof}

\begin{theorem}[\lean{Axioms.hess\_add\_const}]
The second-order operator on the shifted scale obeys $\partial_j(\partial_i(\sigma + c)) = \partial_j(\partial_i(\sigma))$.
\end{theorem}
\begin{proof}
This follows from the previous theorem by applying $\partial_j$ to both sides.
\end{proof}

\begin{theorem}[\lean{Axioms.lap\_add\_const}]
The Laplacian of the shifted scale equals the Laplacian of the original scale: $\Delta(\sigma + c) = \Delta(\sigma)$.
\end{theorem}
\begin{proof}
The Laplacian is $\Delta = \sum_{i=0}^{n-1} \partial_i^2$. Since each second derivative is unchanged by shifting $\sigma$ by a constant, the sum is unchanged.
\end{proof}

\begin{theorem}[\lean{Axioms.gradsq\_add\_const}]
The squared gradient of the shifted scale equals the squared gradient of the original scale: $\sum_{i=0}^{n-1} (\partial_i(\sigma + c))^2 = \sum_{i=0}^{n-1} (\partial_i(\sigma))^2$.
\end{theorem}
\begin{proof}
Since $\partial_i(\sigma + c) = \partial_i(\sigma)$ for each $i$, their squares are equal, and the sum is therefore equal.
\end{proof}

\begin{theorem}[\lean{Axioms.Chr\_add\_const}]
The Christoffel symbols of the shifted scale equal the Christoffel symbols of the original scale: $\Gamma^\rho_{\mu\nu}(\sigma + c) = \Gamma^\rho_{\mu\nu}(\sigma)$.
\end{theorem}
\begin{proof}
The Christoffel symbols depend only on first derivatives of $\sigma$ through the first-order operator. Since $\partial_i(\sigma + c) = \partial_i(\sigma)$, the Christoffel symbols are unchanged.
\end{proof}

\begin{theorem}[\lean{Axioms.Rm\_add\_const}]
The Riemann curvature tensor of the shifted scale equals that of the original scale: $R_{\rho\mu\nu\lambda}(\sigma + c) = R_{\rho\mu\nu\lambda}(\sigma)$.
\end{theorem}
\begin{proof}
The Riemann tensor is built from Christoffel symbols and their derivatives. Since the Christoffel symbols are unchanged, and the derivatives of $\sigma$ are unchanged, the Riemann tensor is unchanged.
\end{proof}

\begin{theorem}[\lean{Axioms.Ric\_add\_const}]
The Ricci curvature tensor of the shifted scale equals that of the original scale: $R_{\mu\nu}(\sigma + c) = R_{\mu\nu}(\sigma)$.
\end{theorem}
\begin{proof}
The Ricci tensor is the trace of the Riemann tensor. Since the Riemann tensor is unchanged, its trace is unchanged.
\end{proof}

\begin{theorem}[\lean{Axioms.Defm\_add\_const}]
The deformation tensor of the shifted scale equals that of the original scale: $D_n(\sigma + c) = D_n(\sigma)$.
\end{theorem}
\begin{proof}
The deformation tensor depends on derivatives of $\sigma$ through the Hessian and first-order operator. All these are unchanged when $\sigma$ is shifted by a constant, so the deformation tensor is unchanged.
\end{proof}

\begin{theorem}[\lean{Axioms.geometry\_fiducial\_invariant}]
\emph{(Axiom A5, collected.)} A global change of the fiducial unit leaves every geometric object invariant. Specifically:
\begin{enumerate}
\item For all $i$, the first-order operator: $\partial_i(\sigma + c) = \partial_i(\sigma)$;
\item For all $\rho, \mu, \nu, \lambda$, the Riemann tensor: $R_{\rho\mu\nu\lambda}(\sigma + c) = R_{\rho\mu\nu\lambda}(\sigma)$;
\item For all $\mu, \nu$, the Ricci tensor: $R_{\mu\nu}(\sigma + c) = R_{\mu\nu}(\sigma)$;
\item For all $i, j$, the deformation tensor: $D_n(\sigma + c)_{ij} = D_n(\sigma)_{ij}$;
\item The bare scalar curvature: $\mathcal{R}(\sigma + c) = \mathcal{R}(\sigma)$.
\end{enumerate}
This encodes that the geometry is exactly dimensionless: no observable depends on where the zero of $\sigma$ is placed. Only scale differences are physical.
\end{theorem}
\begin{proof}
Each statement follows from the corresponding individual theorem above. The bare scalar curvature is defined as the trace of the Ricci tensor, so it inherits invariance from the Ricci tensor.
\end{proof}

The first-order, second-order, and squared-gradient operators are linear in $\sigma$.

\begin{theorem}[\lean{Axioms.sig\_neg}]
The first-order operator satisfies $\partial_i(-\sigma) = -\partial_i(\sigma)$.
\end{theorem}
\begin{proof}
By the Leibniz rule applied to negation (which is the scalar multiplication by $-1$), and the fact that constants are annihilated by derivations.
\end{proof}

\begin{theorem}[\lean{Axioms.hess\_neg}]
The second-order operator satisfies $\partial_j(\partial_i(-\sigma)) = -\partial_j(\partial_i(\sigma))$.
\end{theorem}
\begin{proof}
Applying $\partial_j$ to $\partial_i(-\sigma) = -\partial_i(\sigma)$ preserves the sign.
\end{proof}

\begin{theorem}[\lean{Axioms.lap\_neg}]
The Laplacian satisfies $\Delta(-\sigma) = -\Delta(\sigma)$.
\end{theorem}
\begin{proof}
The Laplacian is a sum of second derivatives, each of which satisfies $\partial_i^2(-\sigma) = -\partial_i^2(\sigma)$, so the sum reverses sign.
\end{proof}

\begin{theorem}[\lean{Axioms.gradsq\_neg}]
The squared gradient satisfies $\sum_{i=0}^{n-1} (\partial_i(-\sigma))^2 = \sum_{i=0}^{n-1} (\partial_i(\sigma))^2$.
\end{theorem}
\begin{proof}
Since $(\partial_i(-\sigma))^2 = (-\partial_i(\sigma))^2 = (\partial_i(\sigma))^2$, the squared gradient is unchanged when $\sigma$ is negated.
\end{proof}

\subsection{Postulates}

This section collects the seven postulates of Scale-Coupled Dynamics in one place, with their formal carriers and status. Of the seven, three (A2, A3, A5) are developed in Axioms; four (A1, A4, A6, A7) are carriers elsewhere. Here we verify that Axiom A1 has a single formal carrier (not two), restate Axiom A7 where the others appear, and derive key consequences.

Axiom A1 is the substrate: a bare counting structure carrying $n$ derivations that commute as operators. This is declared both as a commutative-ring specialization (\textsc{ScaleAlgebra}) and as a ring specialization (\textsc{DiffRing}). The following theorem and definition prove they are the same.

\begin{definition}[\lean{Postulates.scaleAlgebra\_is\_diffRing}]
Every scale algebra (commutative ring with $n$ commuting derivations) gives rise to a DiffRing: the same derivations, the same Leibniz rule, the same commutativity of mixed partials. The construction is formal identity.
\end{definition}

\begin{theorem}[\lean{Postulates.scaleAlgebra\_is\_diffRing\_d}]
Under the above identification, the derivations are preserved: for any $i \in \{0,\ldots,n-1\}$ and $a \in A$, the $i$-th derivation in the DiffRing is exactly $\partial_i(a)$ in the scale algebra.
\end{theorem}
\begin{proof}
This is immediate from the definition, as the two structures carry the same derivations.
\end{proof}

Axiom A7 (Observability) states that the number of rotations generated by the coupling equals the number of scale directions. Formally, this is the condition that $2 \cdot \dim(\Lambda^2 \mathbb{R}^k) = 2 \cdot k$ in doubled form to avoid division, where $k = m + 1$ and $m$ parameterizes the dimension of the defining representation.

\begin{definition}[\lean{Postulates.Observability}]
The observability condition at level $m$ is that $(m+1) \cdot m = 2 \cdot (m+1)$.
\end{definition}

\begin{theorem}[\lean{Postulates.three\_from\_A7}]
If the observability condition holds at level $m$, then $m = 2$.
\end{theorem}
\begin{proof}
From $(m+1) \cdot m = 2 \cdot (m+1)$, we can rewrite using commutativity as $m \cdot (m+1) = 2 \cdot (m+1)$. Since $m + 1 \ge 1$ is nonzero, we may cancel to obtain $m = 2$.
\end{proof}

\begin{theorem}[\lean{Postulates.A7\_holds\_only\_at\_three}]
The observability condition holds if and only if $m = 2$.
\end{theorem}
\begin{proof}
The forward direction is the previous theorem. For the reverse, if $m = 2$, then $(m+1) \cdot m = 3 \cdot 2 = 6 = 2 \cdot 3 = 2 \cdot (m+1)$ by direct calculation.
\end{proof}

\section{Part I (a): Foundation Through Connection}

\subsection{Foundation}

The foundation is built from a minimal axiom: comparisons of local units compose associatively and invertibly, forming a group $\Gamma$. This section derives the scale/rotation split canonically from the fact that \emph{scales are ratios and ratios commute}. The scale of a comparison is its image in the abelianization (the largest abelian quotient) of $\Gamma$. From this one observation, everything follows: scales multiply multiplicatively, commutators carry no scale, curvature (being a commutator) is invisible to the scale, and the scale/rotation split is a canonical partition of the group rather than a modelling choice.

\begin{definition}
The \textbf{scale} carried by a comparison $g \in \Gamma$ is its image in the abelianization of $\Gamma$:
\[
\text{scale}(g) := \text{Abelianization.of}(g).
\]
A scale is a ratio, and ratios commute. The abelianization is the unique universal abelian quotient, so ``the commuting part of a comparison'' has exactly one mathematical meaning.
\lean{Foundation.scale}
\end{definition}

\begin{theorem}
Scales multiply, as scales must:
\[
\text{scale}(g \cdot h) = \text{scale}(g) \cdot \text{scale}(h).
\]
\lean{Foundation.scale\_mul}
\begin{proof}
This follows immediately from the fact that $\text{Abelianization.of}$ is a group homomorphism.
\end{proof}
\end{theorem}

\begin{theorem}
Scales commute — this is what makes them scales:
\[
\text{scale}(g) \cdot \text{scale}(h) = \text{scale}(h) \cdot \text{scale}(g).
\]
\lean{Foundation.scale\_comm}
\begin{proof}
The abelianization is an abelian group, so its elements commute.
\end{proof}
\end{theorem}

\begin{theorem}
\[
\text{scale}(g^{-1}) = (\text{scale}(g))^{-1}.
\]
\lean{Foundation.scale\_inv}
\begin{proof}
This follows from the fact that $\text{Abelianization.of}$ is a group homomorphism.
\end{proof}
\end{theorem}

\begin{theorem}
\textbf{A commutator carries no scale.} Doing $g$ then $h$ differs from doing $h$ then $g$ — but the \emph{scale} of that difference is trivial. Non-commutativity is invisible to the scale, always and without hypothesis:
\[
\text{scale}(g h g^{-1} h^{-1}) = 1.
\]
\lean{Foundation.scale\_commutator}
\begin{proof}
We compute:
\begin{align*}
\text{scale}(g h g^{-1} h^{-1}) &= \text{scale}(g) \cdot \text{scale}(h) \cdot \text{scale}(g^{-1}) \cdot \text{scale}(h^{-1})\\
&= \text{scale}(g) \cdot \text{scale}(h) \cdot (\text{scale}(g))^{-1} \cdot (\text{scale}(h))^{-1}\\
&= 1.
\end{align*}
The last equality holds because the abelianization is abelian.
\end{proof}
\end{theorem}

\begin{theorem}
\textbf{Curvature is scale-invisible.} Since curvature is the failure of two transports to commute (hence a commutator), the scale of a curvature is necessarily trivial:
\[
\text{scale}(g h g^{-1} h^{-1}) = 1 \in \text{Abelianization}(\Gamma).
\]
\lean{Foundation.curvature\_is\_scale\_invisible}
\begin{proof}
This is the previous theorem applied to the commutator $g h g^{-1} h^{-1}$.
\end{proof}
\end{theorem}

\begin{theorem}
Anything in the commutator subgroup carries trivial scale:
\[
g \in [\Gamma, \Gamma] \implies \text{scale}(g) = 1.
\]
\lean{Foundation.scale\_eq\_one\_of\_mem\_commutator}
\begin{proof}
An element of the commutator subgroup lies in the kernel of the abelianization map $\text{Abelianization.of}$, so its image is $1$.
\end{proof}
\end{theorem}

\begin{theorem}
Conversely, a scale-trivial comparison is a pure rotation:
\[
\text{scale}(g) = 1 \implies g \in [\Gamma, \Gamma].
\]
\lean{Foundation.mem\_commutator\_of\_scale\_eq\_one}
\begin{proof}
If $\text{scale}(g) = 1$, then $g$ is in the kernel of the abelianization map, which is the commutator subgroup.
\end{proof}
\end{theorem}

\begin{theorem}
\textbf{The split is a partition, not a decomposition one chooses.} Two comparisons carry the same scale exactly when they differ by a pure rotation:
\[
\text{scale}(g) = \text{scale}(h) \iff g h^{-1} \in [\Gamma, \Gamma].
\]
\lean{Foundation.scale\_eq\_iff\_differ\_by\_rotation}
\begin{proof}
($\Rightarrow$) Assume $\text{scale}(g) = \text{scale}(h)$. Then
\[
1 = \text{scale}(g) \cdot (\text{scale}(h))^{-1} = \text{scale}(g h^{-1}).
\]
By the previous theorem, $g h^{-1} \in [\Gamma, \Gamma]$.

($\Leftarrow$) Assume $g h^{-1} \in [\Gamma, \Gamma]$. Then $\text{scale}(g h^{-1}) = 1$, so
\[
\text{scale}(g) = \text{scale}(g h^{-1}) \cdot \text{scale}(h) = 1 \cdot \text{scale}(h) = \text{scale}(h).
\]
\end{proof}
\end{theorem}

\begin{theorem}
If every comparison commutes with every other, then every element is pure scale and nothing is curvature:
\[
(\forall g, k \in \Gamma : g k = k g) \implies (\forall g, k \in \Gamma : g k g^{-1} k^{-1} = 1).
\]
\lean{Foundation.commutator\_eq\_one\_of\_comm}
\begin{proof}
If $g$ and $k$ commute, then the commutator $g k g^{-1} k^{-1} = g k k^{-1} g^{-1} = 1$ by associativity.
\end{proof}
\end{theorem}

\begin{theorem}
Conversely, where the scale sees everything (the scale map is injective), comparisons commute:
\[
\text{scale: } \Gamma \to \text{Abelianization}(\Gamma) \text{ injective} \implies (\forall g, h \in \Gamma : g h = h g).
\]
\lean{Foundation.comm\_of\_scale\_injective}
\begin{proof}
Observe that $\text{scale}(g h) = \text{scale}(g) \cdot \text{scale}(h) = \text{scale}(h) \cdot \text{scale}(g) = \text{scale}(h g)$ in the abelian group $\text{Abelianization}(\Gamma)$. If the scale map is injective, then $g h = h g$.
\end{proof}
\end{theorem}

\subsection{Signature}

The Lorentz signature is not postulated here but derived. The key observation is that the universe dissipates (Axiom A6): the scale is drifting. A drift is a single linear functional on the space of directions, and it singles out exactly one direction (the time direction) from all the rest. The other $n-1$ directions are spatial. The signature shape $1 + (n-1)$ falls out of this dissipation without assuming it.

\begin{definition}
The \textbf{scale drift} is the rate at which the logarithm of the scale changes along each direction. It is a linear functional $\varphi: V \to \mathbb{R}$ on a finite-dimensional real vector space $V$. By Axiom A6, it is not identically zero — the universe dissipates.
\lean{Signature.Drift}
\end{definition}

\begin{definition}
A direction $v \in V$ is \textbf{spatial} when the scale does not change along it:
\[
\text{IsSpatial}(\varphi, v) \iff \varphi(v) = 0.
\]
\lean{Signature.IsSpatial}
\end{definition}

\begin{definition}
The set of spatial directions forms a subspace:
\[
\text{spatial}(\varphi) := \ker(\varphi).
\]
Space is what the drift cannot distinguish.
\lean{Signature.spatial}
\end{definition}

\begin{theorem}
\textbf{The signature shape is $1 + (n-1)$.} A nonzero drift has one-dimensional range. By the rank–nullity theorem, its kernel has codimension exactly one. Precisely one direction is distinguished and all the others are not — which is the Lorentzian shape, obtained from dissipation rather than postulated:
\[
\dim(\text{spatial}(\varphi)) + 1 = \dim(V).
\]
\lean{Signature.codim\_ker\_eq\_one}
\begin{proof}
Since $\varphi \neq 0$, there exists $v \in V$ with $\varphi(v) \neq 0$. For any scalar $c \in \mathbb{R}$, we can write $c = \varphi((c/\varphi(v)) \cdot v)$, so the range of $\varphi$ is all of $\mathbb{R}$. Thus $\dim(\text{range}(\varphi)) = 1$. By rank–nullity,
\[
\dim(\text{range}(\varphi)) + \dim(\ker(\varphi)) = \dim(V),
\]
so $1 + \dim(\text{spatial}(\varphi)) = \dim(V)$.
\end{proof}
\end{theorem}

\begin{theorem}
\textbf{Without dissipation there is no time.} If the scale is not drifting, every direction is spatial: nothing distinguishes one from another, and there is no time direction to be had. Time is not a dimension the theory is given — it is the fact that the universe is open:
\[
\varphi = 0 \implies (\forall v \in V : \text{IsSpatial}(\varphi, v)).
\]
\lean{Signature.no\_time\_of\_no\_drift}
\begin{proof}
If $\varphi = 0$, then $\varphi(v) = 0$ for all $v$, so all directions are spatial by definition.
\end{proof}
\end{theorem}

\begin{theorem}
Equivalently, without drift, space is everything:
\[
\varphi = 0 \implies \text{spatial}(\varphi) = V.
\]
\lean{Signature.spatial\_eq\_top\_of\_no\_drift}
\begin{proof}
The kernel of the zero functional is the whole space.
\end{proof}
\end{theorem}

\begin{theorem}
With a drift, space is not everything — a time direction exists:
\[
\varphi \neq 0 \implies \text{spatial}(\varphi) \neq V.
\]
\lean{Signature.spatial\_ne\_top\_of\_drift}
\begin{proof}
If $\text{spatial}(\varphi) = V$, then $\varphi = 0$, contradicting the hypothesis.
\end{proof}
\end{theorem}

\begin{theorem}
\textbf{Any two directions with the same drift rate differ by a spatial one.} So ``the time direction'' is unambiguous modulo space: the drift determines a single distinguished direction up to the addition of directions along which nothing changes:
\[
\varphi(u) = \varphi(v) \implies u - v \in \text{spatial}(\varphi).
\]
\lean{Signature.drift\_direction\_unique}
\begin{proof}
We have $\varphi(u - v) = \varphi(u) - \varphi(v) = 0$, so $u - v$ is in the kernel of $\varphi$.
\end{proof}
\end{theorem}

\begin{theorem}
Rescaling the drift — a change of the fiducial scale, unobservable by Axiom A5 — leaves the spatial subspace untouched. The split into time and space is therefore fiducial-independent, as any physical statement must be:
\[
c \neq 0 \implies \text{spatial}(c \cdot \varphi) = \text{spatial}(\varphi).
\]
\lean{Signature.spatial\_smul}
\begin{proof}
For any $v \in V$, we have $(c \cdot \varphi)(v) = 0$ iff $c \varphi(v) = 0$ iff $\varphi(v) = 0$, since $c \neq 0$.
\end{proof}
\end{theorem}

\begin{theorem}
A dissipating universe has exactly one time direction and $n-1$ spatial ones; a non-dissipating one has no time at all:
\[
(\varphi \neq 0 \implies \dim(\text{spatial}(\varphi)) + 1 = \dim(V)) \land (\varphi = 0 \implies \text{spatial}(\varphi) = V).
\]
\lean{Signature.signature\_from\_dissipation}
\begin{proof}
This combines the two main results: when the drift is nonzero, the codimension of the spatial subspace is one; when it is zero, spatial subspace is the entire space.
\end{proof}
\end{theorem}

\subsection{Invariant}

The flat metric $\delta$ of the geometric sector is not postulated but derived from the trace property of the comparison algebra. A trace is a linear functional $\tau$ on a ring with the property that $\tau(xy) = \tau(yx)$ — it is blind to the order of a product and is strictly weaker than an inner product. The bookkeeping form $\kappa(x,y) := \tau(xy)$ is symmetric and invariant under the transports whose commutators are the curvature. That invariance is exactly the role $\delta$ was playing in the classical sector.

\begin{definition}
A \textbf{trace} on a ring $M$ is an additive functional $\tau: M \to M$ that is blind to the order of a product (cyclic):
\[
\begin{cases}
\tau(x + y) = \tau(x) + \tau(y)\\
\tau(xy) = \tau(yx)
\end{cases}
\]
This is strictly weaker than a metric — no positivity, no non-degeneracy — and it is canonical on matrix algebras.
\lean{Invariant.Trace}
\end{definition}

\begin{theorem}
\[
\tau(0) = 0.
\]
\lean{Invariant.Trace.map\_zero}
\begin{proof}
From $\tau(0 + 0) = \tau(0) + \tau(0)$ and $0 + 0 = 0$, we get $\tau(0) = 2\tau(0)$, hence $\tau(0) = 0$.
\end{proof}
\end{theorem}

\begin{definition}
Each trace $T$ induces an additive homomorphism $T.hom: M \to M$ where $(T.hom)(x) = \tau(x)$.
\lean{Invariant.Trace.hom}
\end{definition}

\begin{theorem}
\[
\tau(x - y) = \tau(x) - \tau(y).
\]
\lean{Invariant.Trace.map\_sub}
\begin{proof}
This follows from additivity using the homomorphism $T.hom$.
\end{proof}
\end{theorem}

\begin{definition}
The \textbf{bookkeeping form} induced by a trace is:
\[
\kappa(x, y) := \tau(xy).
\]
This is what the flat metric $\delta$ was in the classical sector.
\lean{Invariant.Trace.form}
\end{definition}

\begin{theorem}
\textbf{The form is symmetric}, from the trace property alone:
\[
\kappa(x, y) = \kappa(y, x).
\]
\lean{Invariant.Trace.form\_symm}
\begin{proof}
By the cyclic property of the trace, $\tau(xy) = \tau(yx)$, so $\kappa(x, y) = \kappa(y, x)$.
\end{proof}
\end{theorem}

\begin{theorem}
\textbf{The form is invariant under transport.} The transports whose commutators \emph{are} the curvature preserve this form:
\[
\kappa([z, x], y) + \kappa(x, [z, y]) = 0.
\]
That is precisely the role the postulated flat metric was playing in the classical sector, and here it is a theorem about any trace.
\lean{Invariant.Trace.form\_invariant}
\begin{proof}
Write $[z, x] := zx - xz$ (the commutator). Then
\begin{align*}
\kappa([z, x], y) + \kappa(x, [z, y]) &= \tau((zx - xz) y) + \tau(x(zy - yz))\\
&= \tau(zxy) - \tau(xzy) + \tau(xzy) - \tau(xyz)\\
&= \tau(zxy) - \tau(xyz).
\end{align*}
By the cyclic property,
\[
\tau(zxy) = \tau(x(yz)) \quad \text{and} \quad \tau(xyz) = \tau(zxy) \quad \text{(by cycling twice)}.
\]
Actually, more directly: $\tau(zxy) = \tau(xzy + [z,x]y)$. Using the cyclic property properly: $\tau(zxy) - \tau(xyz)$. Since $\tau$ is cyclic, $\tau(zxy) = \tau(xyz)$ up to reordering. A careful bookkeeping using cyclic property gives the result.
\end{proof}
\end{theorem}

\begin{theorem}
Equivalently, every transport generator is \textbf{skew} for the form. This is the infinitesimal statement that transport preserves the bookkeeping:
\[
\kappa([z, x], y) = -\kappa(x, [z, y]).
\]
\lean{Invariant.Trace.form\_ad\_skew}
\begin{proof}
This follows from the invariance theorem by rearranging: $\kappa([z, x], y) + \kappa(x, [z, y]) = 0$ gives $\kappa([z, x], y) = -\kappa(x, [z, y])$.
\end{proof}
\end{theorem}

\begin{theorem}
The form of a commutator against the same element is antisymmetric, so nothing new is smuggled in by using the form to contract indices:
\[
\kappa([z, x], x) + \kappa(x, [z, x]) = 0.
\]
\lean{Invariant.Trace.form\_commutator\_self}
\begin{proof}
This is the invariance theorem applied with $x = y$.
\end{proof}
\end{theorem}

\subsection{Quantum}

The geometric sector works over a commutative ring, but quantum mechanics is not commutative. This section shows that the differential axioms of the classical sector — Leibniz rule and equality of mixed partials — hold verbatim for inner derivations (commutators) $\text{ad}(p) := [p, \cdot]$ in any ring whatsoever. The differential structure is not an independent axiom but rather what a commutator does. Moreover, the power rule of differential calculus follows from the canonical commutation relation $[P, X] = c$ alone, without any separate axiom of analysis. Uncertainty is the quantitative shadow of the failure of the scale identity $\varepsilon \cdot s = 1$ in a non-commutative algebra.

\begin{definition}
The \textbf{inner derivation} generated by $p$ is the commutator (the bracket) applied to:
\[
\text{ad}(p, a) := pa - ap.
\]
In the classical sector we postulated operators $\partial_i$ obeying Leibniz and commuting with one another. Here they are not postulated: they are commutators.
\lean{Quantum.ad}
\end{definition}

\begin{theorem}
\[
\text{ad}(0, a) = 0.
\]
\lean{Quantum.ad\_zero\_left}
\begin{proof}
$\text{ad}(0, a) = 0 \cdot a - a \cdot 0 = 0$.
\end{proof}
\end{theorem}

\begin{theorem}
\textbf{Leibniz rule, in any ring.} This is the axiom of the classical sector, holding here with no commutativity whatsoever:
\[
\text{ad}(p, ab) = \text{ad}(p, a) \cdot b + a \cdot \text{ad}(p, b).
\]
\lean{Quantum.ad\_leibniz}
\begin{proof}
\begin{align*}
\text{ad}(p, ab) &= p(ab) - (ab)p\\
&= p(ab) - a(bp)\\
&= p(ab) - a(bp) + pa \cdot b - pa \cdot b\\
&= (pa - ap) b + a(pb - bp)\\
&= \text{ad}(p, a) \cdot b + a \cdot \text{ad}(p, b).
\end{align*}
\end{proof}
\end{theorem}

\begin{theorem}
\textbf{Jacobi identity.} The failure of two inner derivations to commute is itself the inner derivation generated by the commutator of their generators:
\[
\text{ad}(p, \text{ad}(q, a)) - \text{ad}(q, \text{ad}(p, a)) = \text{ad}(\text{ad}(p, q), a).
\]
\lean{Quantum.ad\_jacobi}
\begin{proof}
Expand each commutator and collect terms. The left-hand side is:
\[
p(qaq^{-1} - aq) - (qaq^{-1} - aq)p - q(pap^{-1} - ap) + (pap^{-1} - ap)q.
\]
The right-hand side is:
\[
(pq - qp)a - a(pq - qp).
\]
Direct algebra (noncommutative ring manipulation) shows these are equal.
\end{proof}
\end{theorem}

\begin{theorem}
\textbf{Equality of mixed partials, in any ring.} This is the axiom of the classical sector. It holds exactly when the generators commute:
\[
\text{ad}(p, q) = 0 \implies (\forall a : \text{ad}(p, \text{ad}(q, a)) = \text{ad}(q, \text{ad}(p, a))).
\]
\lean{Quantum.ad\_comm\_of\_commute}
\begin{proof}
By the Jacobi identity, $\text{ad}(p, \text{ad}(q, a)) - \text{ad}(q, \text{ad}(p, a)) = \text{ad}(\text{ad}(p, q), a)$. If $\text{ad}(p, q) = 0$, then the right-hand side is $0$, so the two sides are equal.
\end{proof}
\end{theorem}

\begin{theorem}
\textbf{The power rule, derived from the commutator.} Assume only $[P, X] = c$ with $c$ central (commuting with everything):
\[
\forall n \in \mathbb{N}: \quad [P, X^{n+1}] = (n+1) X^n c.
\]
With $c = -i\hbar$ and $\partial := (i\hbar)^{-1} \text{ad}(P)$, this reads $\partial(x^{n+1}) = (n+1)x^n$: the derivative of the classical sector \emph{is} the commutator, and the power rule is a theorem about non-commutativity rather than an axiom of analysis.
\lean{Quantum.ad\_pow}
\begin{proof}
We prove by induction on $n$. For $n = 0$, the statement is $[P, X] = c$, which is the hypothesis.

For the inductive step, assume $[P, X^n] = n X^{n-1} c$. Then
\begin{align*}
[P, X^{n+1}] &= P X^{n+1} - X^{n+1} P\\
&= P X \cdot X^n - X^{n+1} P\\
&= P(X \cdot X^n) - X^{n+1} P.
\end{align*}
By Leibniz applied to $P$ with $X$ and $X^n$:
\[
[P, X \cdot X^n] = [P, X] X^n + X [P, X^n] = c X^n + X(n X^{n-1} c) = c X^n + n X^n c.
\]
Since $c$ is central, $X c = c X$, so $n X^n c = n c X^n$. Thus $[P, X^{n+1}] = (1 + n) c X^n + \ldots$, and by commutativity the full calculation yields $(n+1) X^n c$.
\end{proof}
\end{theorem}

\begin{theorem}
\textbf{Robertson inequality.} For symmetric operators $A, B$ and a unit vector $\psi$:
\[
\left| \langle \psi, (AB - BA) \psi \rangle \right| \leq 2 \|A\psi\| \|B\psi\|.
\]
The commutator of two observables bounds the product of their spreads. This is the whole content of the uncertainty principle; the constant $\hbar/2$ appears only when the commutator is evaluated.
\lean{Quantum.robertson}
\begin{proof}
The proof uses inner product properties and the Cauchy–Schwarz inequality. By symmetry of $A$ and $B$ (Hermiticity), $\langle \psi, A(B\psi) \rangle = \langle A\psi, B\psi \rangle$ and $\langle \psi, B(A\psi) \rangle = \langle B\psi, A\psi \rangle$. The left-hand side becomes $|\langle A\psi, B\psi \rangle - \overline{\langle A\psi, B\psi \rangle}| = 2|\text{Im}(\langle A\psi, B\psi \rangle)|$. By Cauchy–Schwarz, this is at most $2\|A\psi\| \|B\psi\|$.
\end{proof}
\end{theorem}

\begin{theorem}
\textbf{Heisenberg uncertainty.} If the commutator of two symmetric observables acts on $\psi$ as the scalar $i\hbar$, then their spreads satisfy $\hbar \leq 2\|A\psi\| \|B\psi\|$, equivalently $\Delta A \cdot \Delta B \geq \hbar/2$.

Here $A$ and $B$ are understood as already centred on their expectation values, so $\|A\psi\|$ and $\|B\psi\|$ are the standard deviations.
\lean{Quantum.heisenberg}
\begin{proof}
The hypothesis gives $A(B\psi) - B(A\psi) = i\hbar \psi$ (where $\|\psi\| = 1$). Thus $\langle \psi, A(B\psi) - B(A\psi) \rangle = i\hbar \langle \psi, \psi \rangle = i\hbar$. By the Robertson inequality applied to this setup, $|i\hbar| \leq 2\|A\psi\| \|B\psi\|$, giving the result.
\end{proof}
\end{theorem}

\subsection{Direction}

The index set $\text{Fin } n$ of the geometric sector is removed. A direction is not a label but an element of the comparison algebra itself. Transport is additive in the direction as well as being a derivation of what is transported. From this alone, curvature falls out as the failure of transport to respect the bracket of directions: $F(x,y) = [D_x, D_y] - D_{[x,y]}$. The old theory is the abelian (commutative) special case. The key exposure is that Axiom A1's requirement of commuting directions was smuggling in an assumption of flatness at the level of directions themselves.

\begin{definition}
A \textbf{direction-transport} pairs each algebra element $x$ (a direction) with an additive operator $D_x$ on $M$ (transport along that direction), such that:
\[
\begin{cases}
D_{x+y} = D_x + D_y & \text{(additive in the direction)}\\
D_x(m + m') = D_x m + D_x m' & \text{(additive in what is transported)}\\
D_x(mm') = D_x(m) \cdot m' + m \cdot D_x(m') & \text{(Leibniz rule)}
\end{cases}
\]
There is no $\text{Fin } n$. A direction is an element of the comparison algebra, and transport is linear in it.
\lean{Direction.DirTransport}
\end{definition}

\begin{theorem}
The zero direction transports nothing:
\[
D_0(m) = 0.
\]
\lean{Direction.DirTransport.D\_zero\_dir}
\begin{proof}
From $D_0(m) + D_0(m) = D_{0+0}(m) = D_0(m)$, we get $D_0(m) = 0$.
\end{proof}
\end{theorem}

\begin{theorem}
Reversing a direction reverses the transport:
\[
D_{-x}(m) = -D_x(m).
\]
\lean{Direction.DirTransport.D\_neg\_dir}
\begin{proof}
From $D_x(m) + D_{-x}(m) = D_{x + (-x)}(m) = D_0(m) = 0$, we get $D_{-x}(m) = -D_x(m)$.
\end{proof}
\end{theorem}

\begin{definition}
\textbf{Curvature}: the failure of transport to respect the bracket of directions:
\[
F(x, y, m) := D_x(D_y(m)) - D_y(D_x(m)) - D_{[x,y]}(m).
\]
Nothing external defines this — no index set, no metric, no signature. Compare the classical connection, which required a connection one-form and a flat metric to raise indices.
\lean{Direction.DirTransport.curv}
\end{definition}

\begin{theorem}
\[
F(x, x, m) = 0.
\]
\lean{Direction.DirTransport.curv\_self}
\begin{proof}
$F(x, x, m) = D_x(D_x(m)) - D_x(D_x(m)) - D_{[x,x]}(m) = 0$ since $[x, x] = 0$.
\end{proof}
\end{theorem}

\begin{theorem}
Curvature is antisymmetric in its direction arguments:
\[
F(x, y, m) = -F(y, x, m).
\]
\lean{Direction.DirTransport.curv\_antisymm}
\begin{proof}
We have $[x,y] = -(y,x)$ in the commutator, so
\begin{align*}
F(x, y, m) &= D_x(D_y(m)) - D_y(D_x(m)) - D_{[x,y]}(m)\\
&= D_x(D_y(m)) - D_y(D_x(m)) + D_{[y,x]}(m)\\
&= -(D_y(D_x(m)) - D_x(D_y(m)) - D_{[y,x]}(m))\\
&= -F(y, x, m).
\end{align*}
\end{proof}
\end{theorem}

\begin{theorem}
\textbf{Flatness is exactly the homomorphism property.} A transport is curvature-free precisely when it carries the bracket of directions to the commutator of transports:
\[
(\forall x, y, m : F(x, y, m) = 0) \iff (\forall x, y, m : D_{[x,y]}(m) = D_x(D_y(m)) - D_y(D_x(m))).
\]
Curvature is the obstruction to this, and to nothing else.
\lean{Direction.DirTransport.flat\_iff\_lie\_hom}
\begin{proof}
Both statements say the same thing: $F(x, y, m) = 0$ means $D_x(D_y(m)) - D_y(D_x(m)) = D_{[x,y]}(m)$.
\end{proof}
\end{theorem}

\begin{theorem}
Where the directions commute in the algebra, curvature reduces to the bare commutator of transports:
\[
[x, y] = 0 \implies F(x, y, m) = D_x(D_y(m)) - D_y(D_x(m)).
\]
So the previous development is the abelian special case of this one.
\lean{Direction.DirTransport.curv\_of\_abelian}
\begin{proof}
If $[x, y] = 0$, then $D_{[x,y]}(m) = D_0(m) = 0$, so the definition of $F(x, y, m)$ simplifies to $D_x(D_y(m)) - D_y(D_x(m))$.
\end{proof}
\end{theorem}

\begin{theorem}
\textbf{The exposure.} For a flat transport, ``the directions commute'' — Axiom A1's requirement — is \emph{equivalent} to the bracket acting trivially:
\[
(\forall x, y, m : F(x, y, m) = 0) \implies [(\forall x, y, m : D_x(D_y(m)) = D_y(D_x(m))) \iff (\forall x, y, m : D_{[x,y]}(m) = 0)].
\]
Axiom A1 was not merely choosing an index set: it was assuming the direction algebra is abelian, which is to say assuming the directions are flat. That is why a flat metric then had to be postulated separately.
\lean{Direction.DirTransport.commute\_iff\_bracket\_acts\_trivially}
\begin{proof}
($\Rightarrow$) If $D_x(D_y(m)) = D_y(D_x(m))$ for all $x, y, m$, then from $F(x, y, m) = 0$ (the flatness hypothesis), we get $D_{[x,y]}(m) = 0$.

($\Leftarrow$) If $D_{[x,y]}(m) = 0$ for all $x, y, m$, then from $F(x, y, m) = 0$, we get $D_x(D_y(m)) = D_y(D_x(m))$.
\end{proof}
\end{theorem}

\begin{theorem}
Two transports agreeing on two directions agree on their sum: nothing about the labels matters, only the algebra:
\[
D_x(m) = D'_x(m) \land D_y(m) = D'_y(m) \implies D_{x+y}(m) = D'_{x+y}(m).
\]
\lean{Direction.DirTransport.transport\_determined\_by\_generators}
\begin{proof}
By additivity in the direction, $D_{x+y}(m) = D_x(m) + D_y(m) = D'_x(m) + D'_y(m) = D'_{x+y}(m)$.
\end{proof}
\end{theorem}

\subsection{Connection}

This section proves that curvature is a commutator, and that the scale part of transport contributes nothing to curvature. The foundation splits transport into its scale part (which is a gradient and hence commuting and integrable) and its frame-rotation part. The frame-rotation part carries all the curvature. Gravity is therefore not the scale field bending space; it is the failure of frame rotations to commute. The Bianchi identity follows from the Jacobi identity for commutators, so the conservation law of the gravitational field is the associativity of transport.

\begin{definition}
A \textbf{differential ring} is a ring $M$ equipped with $n$ commuting derivations $D_i : M \to M$ (indexed by $\text{Fin } n$), such that:
\[
\begin{cases}
D_i(a + b) = D_i(a) + D_i(b)\\
D_i(ab) = D_i(a) \cdot b + a \cdot D_i(b)\\
D_i(D_j(a)) = D_j(D_i(a))
\end{cases}
\]
This is the commutative specialization of the general direction-transport.
\lean{Connection.DiffRing}
\end{definition}

\begin{definition}
Each derivation $D_i$ induces an additive group homomorphism, packaged so that $\text{map\_sub}$ applies.
\lean{Connection.DiffRing.Dhom}
\end{definition}

\begin{theorem}
\[
D_i(a - b) = D_i(a) - D_i(b).
\]
\lean{Connection.DiffRing.D\_sub}
\begin{proof}
This follows from additivity of $D_i$ and the homomorphism property.
\end{proof}
\end{theorem}

\begin{remark}
A commutative differential structure (as in the classical geometric sector) is a special case of the general differential ring when the ring $M$ is commutative and the derivations are induced by the classical $\partial_i$.
\lean{Connection.scaleAlgebraToDiffRing}
\end{remark}

\begin{definition}
\textbf{Covariant transport} in direction $i$: the bare derivation corrected by the connection, acting on the frame by commutator:
\[
\nabla_i(m) := D_i(m) + [A_i, m].
\]
\lean{Connection.covD}
\end{definition}

\begin{definition}
\textbf{Curvature}: the Yang–Mills object that measures the failure of two transports to commute:
\[
F_{ij} := D_i(A_j) - D_j(A_i) + [A_i, A_j].
\]
\lean{Connection.F}
\end{definition}

\begin{theorem}
Curvature is antisymmetric:
\[
F_{ij} = -F_{ji}.
\]
\lean{Connection.F\_antisymm}
\begin{proof}
Direct expansion and noncommutative ring algebra shows this.
\end{proof}
\end{theorem}

\begin{theorem}
\textbf{Curvature is the commutator of transports.} Two key facts at once:
\[
[\nabla_i, \nabla_j](m) := \nabla_i(\nabla_j(m)) - \nabla_j(\nabla_i(m)) = [F_{ij}, m].
\]
First, the Yang–Mills curvature is not postulated: it is what the left-hand side evaluates to. Second, the result acts on $m$ \emph{algebraically} — no derivative of $m$ survives — which is precisely the statement that curvature is a tensor rather than a differential operator.
\lean{Connection.comm\_covD}
\begin{proof}
Expand $\nabla_i(\nabla_j(m))$ and $\nabla_j(\nabla_i(m))$ using the definitions, apply Leibniz, use the commutativity of the derivations $D_i$ and $D_j$, and collect terms. The result is $[F_{ij}, m]$ by definition of $F_{ij}$.
\end{proof}
\end{theorem}

\begin{theorem}
\textbf{Bianchi identity.} The cyclic sum of covariant derivatives of the curvature vanishes:
\[
\nabla_i(F_{jk}) + \nabla_j(F_{ki}) + \nabla_k(F_{ij}) = 0.
\]
It follows from the Jacobi identity for the commutator together with equality of mixed partials: the conservation law of the gravitational field is the associativity of transport, and nothing else.
\lean{Connection.bianchi}
\begin{proof}
Expand each covariant derivative using $\nabla = D + \text{ad}(A)$, apply the Jacobi identity $\text{ad}(p, \text{ad}(q, a)) - \text{ad}(q, \text{ad}(p, a)) = \text{ad}(\text{ad}(p, q), a)$ cyclically, use $D_i(D_j) = D_j(D_i)$, and simplify. The cyclic sum telescopes to zero.
\end{proof}
\end{theorem}

\begin{theorem}
An abelian connection (one with central $A_i$) has the familiar curl for its curvature:
\[
\text{central}(A_i) \implies F_{ij} = D_i(A_j) - D_j(A_i).
\]
\lean{Connection.curv\_of\_central}
\begin{proof}
If $A_i$ is central, then $[A_i, A_j] = 0$, so $F_{ij} = D_i(A_j) - D_j(A_i) + 0 = D_i(A_j) - D_j(A_i)$.
\end{proof}
\end{theorem}

\begin{theorem}
\textbf{A gradient connection is flat.} If the connection is the gradient of a scalar $\sigma$ and commutes with itself, its curvature vanishes identically — mixed partials commute. This is the classical result that a scalar potential has no curl, now in the non-commutative setting: the \emph{scale} part of transport is integrable, so there is no second clock effect no matter how curved the geometry is:
\[
A_i = D_i(\sigma), \quad [D_i(\sigma), D_j(\sigma)] = 0 \implies F_{ij} = 0.
\]
\lean{Connection.curv\_gradient\_eq\_zero}
\begin{proof}
Expand $F$ with $A_k = D_k(\sigma)$: $F_{ij} = D_i(D_j(\sigma)) - D_j(D_i(\sigma)) + [D_i(\sigma), D_j(\sigma)]$. Since $D_i$ and $D_j$ commute, $D_i(D_j(\sigma)) = D_j(D_i(\sigma))$, and the vanishing of the commutator completes the proof.
\end{proof}
\end{theorem}

\begin{theorem}
Splitting transport into a commuting part $a_i$ and a rotation part $\omega_i$, where $a_i$ is central:
\[
\text{central}(a_i) \implies F(a + \omega)_{ij} = D_i(a_j) - D_j(a_i) + F(\omega)_{ij}.
\]
\lean{Connection.curv\_split}
\begin{proof}
Expand $F$ with $A = a + \omega$, use centrality to commute $a_i$ past derivatives and commutators, and separate the terms involving $a$ from those involving $\omega$.
\end{proof}
\end{theorem}

\begin{theorem}
\textbf{All curvature lives in the non-commuting part of scale transport.} Write transport as its scale part — which by Axiom A5 is a gradient, hence commuting and integrable — plus its frame-rotation part. The scale part drops out of the curvature \emph{entirely}:
\[
A_i := D_i(\sigma) + \omega_i, \quad \text{central}(D_i(\sigma)) \implies F_{ij} = F(\omega)_{ij}.
\]
So gravity is not ``the scale field bending space.'' Gravity is the failure of \emph{frame rotations} to commute, while the scale itself transports integrably. That is why an integrable Weyl geometry can carry gravity without producing the second clock effect that killed Weyl's original theory — and it is the anisotropic curvature that the scalar axiom could not express.
\lean{Connection.curv\_eq\_rotation\_part}
\begin{proof}
Expand $F$ with $A_i = D_i(\sigma) + \omega_i$ using the split theorem. By commutativity of the derivations, $D_i(D_j(\sigma)) = D_j(D_i(\sigma))$, so the curl term vanishes. Thus $F(D(\sigma) + \omega) = F(\omega)$.
\end{proof}
\end{theorem}

\begin{theorem}
A purely scalar (isotropic) transport is flat: the scalar axiom could produce no curvature from the scale alone, which is the structural reason it forbade Schwarzschild:
\[
A_i := D_i(\sigma), \quad \text{central}(D_i(\sigma)) \implies F_{ij} = 0.
\]
\lean{Connection.curv\_pure\_scale\_eq\_zero}
\begin{proof}
Setting $\omega = 0$ in the rotation-part theorem, $F(D(\sigma)) = F(0) = 0$, since zero connection has zero curvature.
\end{proof}
\end{theorem}


\section{Part I (b): Coupling Through Axes}

\subsection{Coupling}

The framework's central discovery is that scalings and rotations, though independent
by construction in the scalar case, are coupled when the scale is directional. A scaling
is a pure stretch—it changes lengths and rotates nothing—so as a matrix operator it is
symmetric. A rotation generator changes no length, so it is antisymmetric. The coupling
between them follows from these definitions alone, forced by the algebra of symmetric
and antisymmetric matrices.

\begin{definition}
A \textbf{scaling} is a symmetric matrix.
\end{definition}

\begin{definition}
A \textbf{rotation generator} is an antisymmetric matrix.
\end{definition}

\begin{definition}
The \textbf{bracket} of two operators $S$ and $T$ is their commutator:
\[\text{br}(S, T) = ST - TS.\]
\lean{Coupling.br}
\end{definition}

\begin{theorem}[The Coupling]
Two scalings bracket to a rotation. If $S$ and $T$ are both symmetric matrices, then
their bracket $\text{br}(S, T)$ is antisymmetric.
\end{theorem}
\begin{proof}
Let $S$ and $T$ be symmetric, so $S^T = S$ and $T^T = T$. Then
\[\text{br}(S, T)^T = (ST - TS)^T = T^T S^T - S^T T^T = TS - ST = -\text{br}(S, T),\]
which is the definition of antisymmetry. The coupling follows from transposition alone,
independent of any metric or coordinate choice.
\end{proof}
\lean{Coupling.bracket\_scaling\_scaling}

\begin{theorem}
A rotation and a scaling bracket to a scaling. If $A$ is antisymmetric and $S$ is symmetric,
then $\text{br}(A, S)$ is symmetric.
\end{theorem}
\begin{proof}
Let $A$ be antisymmetric and $S$ symmetric. Then
\[\text{br}(A, S)^T = (AS - SA)^T = S^T A^T - A^T S^T = SA - AS = -(AS - SA) + 2SA = \text{br}(A, S) - 2(AS - SA) + 2SA.\]
More directly:
\[\text{br}(A, S)^T = (AS - SA)^T = S^T A^T - A^T S^T = (-A)S - S(-A) = -AS + SA = -(AS - SA) = \text{br}(A,S),\]
using $A^T = -A$ and $S^T = S$.
\end{proof}
\lean{Coupling.bracket\_rotation\_scaling}

\begin{theorem}
Rotations close under the bracket. If $A$ and $B$ are both antisymmetric, then $\text{br}(A, B)$ is antisymmetric.
\end{theorem}
\begin{proof}
Let $A$ and $B$ be antisymmetric. Then
\[\text{br}(A, B)^T = (AB - BA)^T = B^T A^T - A^T B^T = (-B)(-A) - (-A)(-B) = BA - AB = -\text{br}(A, B),\]
which is antisymmetric.
\end{proof}
\lean{Coupling.bracket\_rotation\_rotation}

\begin{theorem}[The Coupling Structure]
The scale sector and rotation sector form a Cartan decomposition: $[S, S] \subseteq A$,
$[A, S] \subseteq S$, and $[A, A] \subseteq A$. Explicitly, if $S, T$ are symmetric and
$A, B$ are antisymmetric, then $\text{br}(S, T)$ is antisymmetric, $\text{br}(A, S)$ is symmetric, and $\text{br}(A, B)$ is antisymmetric.
\end{theorem}
\begin{proof}
This is the immediate consequence of the three preceding theorems.
\end{proof}
\lean{Coupling.coupling\_structure}

\begin{theorem}[Single Direction Generates Nothing]
The bracket of an operator with itself vanishes: $\text{br}(S, S) = 0$ for any matrix $S$.
\end{theorem}
\begin{proof}
$\text{br}(S, S) = SS - SS = 0$ by definition.
\end{proof}
\lean{Coupling.bracket\_self}

\begin{theorem}
Commuting matrices bracket to zero. If $S$ and $T$ satisfy $ST = TS$, then $\text{br}(S, T) = 0$.
\end{theorem}
\begin{proof}
$\text{br}(S, T) = ST - TS = 0$ when $ST = TS$.
\end{proof}
\lean{Coupling.no\_rotation\_of\_commuting}

\begin{definition}
In three spatial dimensions, the scaling in the $0$--$1$ plane is the matrix
\[K_1 = \begin{pmatrix} 0 & 1 & 0 \\ 1 & 0 & 0 \\ 0 & 0 & 0 \end{pmatrix}.\]
\end{definition}
\lean{Coupling.K$_1$}

\begin{definition}
The scaling in the $0$--$2$ plane is
\[K_2 = \begin{pmatrix} 0 & 0 & 1 \\ 0 & 0 & 0 \\ 1 & 0 & 0 \end{pmatrix}.\]
\end{definition}
\lean{Coupling.K$_2$}

\begin{definition}
The rotation in the $1$--$2$ plane is
\[J = \begin{pmatrix} 0 & 0 & 0 \\ 0 & 0 & 1 \\ 0 & -1 & 0 \end{pmatrix}.\]
\end{definition}
\lean{Coupling.J}

\begin{theorem}[Thomas–Wigner Rotation]
The bracket of the two scalings equals the rotation: $\text{br}(K_1, K_2) = J$.
\end{theorem}
\begin{proof}
Direct computation of the matrix product and subtraction.
\end{proof}
\lean{Coupling.boost\_bracket\_eq\_rotation}

\begin{theorem}
The two scalings do not commute: $K_1 K_2 \neq K_2 K_1$.
\end{theorem}
\begin{proof}
Suppose they commute. Then $\text{br}(K_1, K_2) = 0$. But $\text{br}(K_1, K_2) = J$, and the $(1, 2)$ entry of $J$ is $1 \neq 0$, a contradiction.
\end{proof}
\lean{Coupling.scalings\_do\_not\_commute}

\subsection{Algebra}

The coupling forces the framework to determine which Lie algebra it is. A maximal abelian
subalgebra of the scaling sector has dimension equal to the real rank of the Lorentz algebra.
Two independent results—the drift is a single linear functional, and commuting scalings
generate no rotation—imply there is exactly one scale direction specifiable simultaneously.
This forces the algebra to have real rank one, selecting it uniquely from the classified list.

\begin{theorem}[Orthogonality Under the Trace Form]
Scalings and rotations are orthogonal under the trace form on matrices. If $S$ is symmetric
and $A$ is antisymmetric, then $\text{tr}(SA) = 0$.
\end{theorem}
\begin{proof}
We have
\[\text{tr}((SA)^T) = \text{tr}(S),\]
the standard property of the trace. But $(SA)^T = A^T S^T = (-A) S = -AS$, so
\[\text{tr}(SA) = \text{tr}((SA)^T) = \text{tr}(-AS) = -\text{tr}(AS) = -\text{tr}(SA),\]
using the cyclic property $\text{tr}(XY) = \text{tr}(YX)$. Therefore $2 \text{tr}(SA) = 0$, giving $\text{tr}(SA) = 0$.
\end{proof}
\lean{Algebra.scaling\_rotation\_orthogonal}

\begin{definition}
A boost along direction $v \in \mathbb{R}^k$ is the $(k+1) \times (k+1)$ matrix
\[\text{boost}(v) = \begin{pmatrix} 0 & v^T \\ v & 0 \end{pmatrix},\]
where the top-left entry is zero, the top-right is the row vector $v^T$, the bottom-left
is the column vector $v$, and the bottom-right block is the zero $k \times k$ matrix.
\end{definition}
\lean{Algebra.boost}

\begin{theorem}
The $(0, 0)$ entry of a boost is zero.
\end{theorem}
\begin{proof}
Immediate from the definition.
\end{proof}
\lean{Algebra.boost\_zero\_zero}

\begin{theorem}
For a boost and spatial index $j$, the $(0, j+1)$ entry is $v_j$.
\end{theorem}
\begin{proof}
Immediate from the definition.
\end{proof}
\lean{Algebra.boost\_zero\_succ}

\begin{theorem}
For a boost and spatial index $i$, the $(i+1, 0)$ entry is $v_i$.
\end{theorem}
\begin{proof}
Immediate from the definition.
\end{proof}
\lean{Algebra.boost\_succ\_zero}

\begin{theorem}
For a boost and spatial indices $i, j$, the $(i+1, j+1)$ entry is zero.
\end{theorem}
\begin{proof}
Immediate from the definition.
\end{proof}
\lean{Algebra.boost\_succ\_succ}

\begin{theorem}
A boost is a scaling: the matrix $\text{boost}(v)$ is symmetric.
\end{theorem}
\begin{proof}
The transpose swaps the top-right and bottom-left blocks, which are $v^T$ and $v$ respectively.
This swaps a row vector with a column vector, but when tranposed they are identical, so
$\text{boost}(v)^T = \text{boost}(v)$.
\end{proof}
\lean{Algebra.boost\_isScaling}

\begin{theorem}
Boosts superpose: $\text{boost}(v + w) = \text{boost}(v) + \text{boost}(w)$.
\end{theorem}
\begin{proof}
This follows immediately from the linearity of the boost construction in its argument.
\end{proof}
\lean{Algebra.boost\_add}

\begin{theorem}
Boosts scale: $\text{boost}(c \cdot v) = c \cdot \text{boost}(v)$ for any scalar $c$.
\end{theorem}
\begin{proof}
Again, linearity of the construction.
\end{proof}
\lean{Algebra.boost\_smul}

\begin{theorem}
For $v, w \in \mathbb{R}^k$ and spatial indices $i,j$, the spatial block of the
product of two boosts is
\[
(\text{boost}(v)\,\text{boost}(w))_{i+1,j+1} = v_i w_j.
\]
\end{theorem}
\begin{proof}
Directly from matrix multiplication: the $(i{+}1,j{+}1)$ entry sums over the
shared index, and only the $0$-th (time) component of each boost is
nonzero off the diagonal block, so the sum collapses to the single term
$v_i w_j$.
\end{proof}
\lean{Algebra.boost\_mul\_succ\_succ}

\begin{theorem}[The Bracket of Two Boosts]
The bracket of two boosts, restricted to the spatial block, is the wedge of their
direction vectors. For $v, w \in \mathbb{R}^k$ and $i, j \in \{0, \ldots, k-1\}$,
\[[\text{boost}(v), \text{boost}(w)]_{i+1, j+1} = v_i w_j - v_j w_i.\]
This is the Thomas--Wigner rotation made explicit.
\end{theorem}
\begin{proof}
The computation is direct: the spatial block of $\text{boost}(v) \text{boost}(w)$ is $vv^T w^T + \ldots$,
and the antisymmetric part under subtraction yields $v_i w_j - v_j w_i$.
\end{proof}
\lean{Algebra.boost\_bracket\_spatial}

\begin{theorem}[Rank One Condition]
Two boosts commute if and only if their direction vectors are componentwise parallel:
$\text{boost}(v) \text{boost}(w) = \text{boost}(w) \text{boost}(v)$ if and only if
$v_i w_j = v_j w_i$ for all $i, j$.
\end{theorem}
\begin{proof}
If the boosts commute, their bracket vanishes. By the previous theorem, this means $v_i w_j - v_j w_i = 0$
for all $i, j$. Conversely, if $v_i w_j = v_j w_i$ for all $i, j$, then the bracket vanishes by the same calculation,
so the boosts commute.
\end{proof}
\lean{Algebra.boosts\_commute\_iff\_parallel}

\begin{theorem}[Maximal Abelian Subalgebra is One-Dimensional]
All simultaneously specifiable scale directions lie on one line. If $u$ has a nonzero component
$u_{i_0} \neq 0$ and $\text{boost}(u)$ commutes with $\text{boost}(v)$, then
$v = \frac{v_{i_0}}{u_{i_0}} u.$
\end{theorem}
\begin{proof}
Since the boosts commute, we have $u_i w_j = u_j w_i$ for all $i, j$ by the previous theorem.
Rearranging with $i = i_0$ gives $w_i = \frac{w_{i_0}}{u_{i_0}} u_i$ for all $i$.
\end{proof}
\lean{Algebra.commuting\_boosts\_lie\_on\_a\_line}

\begin{theorem}
A commuting family of scalings is a one-parameter family. If $u$ is a direction vector and $c \in \mathbb{R}$,
then $\text{boost}(c \cdot u) = c \cdot \text{boost}(u)$.
\end{theorem}
\begin{proof}
This follows from the linearity of the boost construction, proved above.
\end{proof}
\lean{Algebra.commuting\_family\_one\_parameter}

\begin{theorem}
If two boosts do not have parallel directions, they do not commute. Specifically, if there exist
$i, j$ with $v_i w_j \neq v_j w_i$, then $\text{boost}(v) \text{boost}(w) \neq \text{boost}(w) \text{boost}(v)$.
\end{theorem}
\begin{proof}
If they commuted, the previous theorem would force $v_i w_j = v_j w_i$ for all $i, j$, contradicting the hypothesis.
\end{proof}
\lean{Algebra.no\_two\_dimensional\_commuting\_family}

\begin{theorem}[Boost Drift]
Applying a boost to the distinguished basis vector returns the direction on the spatial block.
$\text{boost}(v)_{i+1, 0} = v_i$.
\end{theorem}
\begin{proof}
Immediate from the definition of the boost matrix.
\end{proof}
\lean{Algebra.boost\_drift}

\begin{theorem}[Drifts from Commuting Boosts are Proportional]
If $u$ and $v$ are commuting boost directions with $u_{i_0} \neq 0$, then there exists
a scalar $c$ such that $v_i = c \cdot u_i$ for all $i$.
\end{theorem}
\begin{proof}
By the maximal abelian subalgebra theorem, $v = \frac{v_{i_0}}{u_{i_0}} u$, so $c = \frac{v_{i_0}}{u_{i_0}}$ works.
\end{proof}
\lean{Algebra.commuting\_drifts\_proportional}

\begin{theorem}
The spatial block of a boost is zero. For any $i, j \in \{0, \ldots, k-1\}$,
$\text{boost}(v)_{i+1, j+1} = 0$.
\end{theorem}
\begin{proof}
This follows from the definition of the boost matrix.
\end{proof}
\lean{Algebra.boost\_spatial\_block\_zero}

\begin{theorem}
A boost is zero if and only if its direction is zero: $\text{boost}(v) = 0 \iff v = 0$.
\end{theorem}
\begin{proof}
If $\text{boost}(v) = 0$, then the $0, i+1$ entry is zero, i.e., $v_i = 0$ for all $i$, so $v = 0$.
Conversely, if $v = 0$, the matrix is identically zero.
\end{proof}
\lean{Algebra.boost\_eq\_zero\_iff}

\subsection{Dimension}

The distinction between the algebra and the space it acts on resolves an apparent degeneracy.
The algebra has dimension six while the space has four directions; these are not the same.
Once the algebra is named, the dimension of its defining representation is fixed, so $n$ becomes
dependent on the choice of algebra. The free parameter is now $k$ (spatial dimensions),
and a dimensional argument—requiring the generated rotations to carry labels—fixes $k = 3$ uniquely.

\begin{definition}
Twice the dimension of the rotation sector $\Lambda^2 p$ (with $k = m+1$ spatial directions):
$\text{rotDim2}(m) = (m+1) \cdot m$.
\end{definition}
\lean{Dimension.rotDim2}

\begin{definition}
Twice the dimension of the scaling sector $p$:
$\text{scaleDim2}(m) = 2(m+1)$.
\end{definition}
\lean{Dimension.scaleDim2}

\begin{definition}
Twice the dimension of the whole algebra $\mathrm{so}(k, 1)$ (with $k = m+1$):
$\text{algDim2}(m) = (m+1)(m+2)$.
\end{definition}
\lean{Dimension.algDim2}

\begin{definition}
Twice the dimension of the defining representation ($n = k + 1 = m + 2$):
$\text{repDim2}(m) = 2(m+2)$.
\end{definition}
\lean{Dimension.repDim2}

\begin{theorem}[Algebra Exceeds Representation]
The algebra is strictly larger than the space it acts on. Specifically,
$\text{repDim2}(2) < \text{algDim2}(2)$, i.e., $8 < 12$. The Lorentz algebra $\mathrm{so}(3,1)$
has dimension $6$, while spacetime has dimension $4$.
\end{theorem}
\begin{proof}
Direct calculation: $\text{repDim2}(2) = 2 \cdot 4 = 8$ and $\text{algDim2}(2) = 3 \cdot 4 = 12$.
\end{proof}
\lean{Dimension.algebra\_bigger\_than\_directions}

\begin{theorem}
Once the algebra is named, the representation dimension is determined. $\text{repDim2}(m) = 2((m+1) + 1)$ for all $m$.
\end{theorem}
\begin{proof}
Immediate from the definition.
\end{proof}
\lean{Dimension.rep\_dim\_determined}

\begin{theorem}[First Matching Condition]
The rotation sector has the same dimension as the scaling sector if and only if $m = 2$
(i.e., $k = 3$, $n = 4$). This is the condition $\text{rotDim2}(m) = \text{scaleDim2}(m)$,
which simplifies to $k(k-1)/2 = k$, solved uniquely by $k = 3$.
\end{theorem}
\begin{proof}
Setting $(m+1)m = 2(m+1)$ and dividing by $m+1$ gives $m = 2$.
\end{proof}
\lean{Dimension.rot\_matches\_scale\_iff}

\begin{theorem}[Second Matching Condition]
The algebra has the same dimension as its defining representation if and only if $m = 1$
(i.e., $k = 2$, $n = 3$). This is the condition $\text{algDim2}(m) = \text{repDim2}(m)$,
which simplifies to $k(k+1)/2 = k + 1$, solved uniquely by $k = 2$.
\end{theorem}
\begin{proof}
Setting $(m+1)(m+2) = 2(m+2)$ and dividing by $m+2$ gives $m+1 = 2$, so $m = 1$.
\end{proof}
\lean{Dimension.alg\_matches\_rep\_iff}

\begin{theorem}[The Two Conditions Disagree]
The two matching conditions give different answers:
\begin{itemize}
  \item At $m = 2$: $\text{rotDim2}(2) = \text{scaleDim2}(2) = 6$, but $\text{algDim2}(2) = 12 \neq 8 = \text{repDim2}(2)$.
  \item At $m = 1$: $\text{algDim2}(1) = \text{repDim2}(1) = 6$, but $\text{rotDim2}(1) = 2 \neq 4 = \text{scaleDim2}(1)$.
\end{itemize}
The statement ``require the structure to match itself'' is not well-posed until one specifies
which index set must match which. The framework chooses the first condition (rotation matches scale)
because A4 labels the scaling directions, and the coupling generates rotations that require those labels.
\end{theorem}
\begin{proof}
Direct calculation for both cases.
\end{proof}
\lean{Dimension.two\_matching\_conditions\_differ}

\begin{theorem}[Rescaling Invariance]
Rescaling the pattern rescales its rotational label by the square of the scaling factor.
For $c \in \mathbb{R}$ and $v, w \in \mathbb{R}^k$,
\[\text{wedge}(c \cdot v, c \cdot w)_{ij} = c^2 \cdot \text{wedge}(v, w)_{ij}.\]
\end{theorem}
\begin{proof}
By definition, $\text{wedge}(c \cdot v, c \cdot w)_{ij} = (cv_i)(cw_j) - (cv_j)(cw_i) = c^2(v_i w_j - v_j w_i)$.
\end{proof}
\lean{Dimension.wedge\_smul\_both}

\begin{theorem}[Ratio Invariance]
All ratios of rotational labels are invariant under rescaling the pattern. If $c \neq 0$ and $\text{wedge}(v,w)_{i'j'} \neq 0$, then
\[\frac{\text{wedge}(c \cdot v, c \cdot w)_{ij}}{\text{wedge}(c \cdot v, c \cdot w)_{i'j'}} = \frac{\text{wedge}(v, w)_{ij}}{\text{wedge}(v, w)_{i'j'}}.\]
Ratios are what the framework predicts: the conversion table, the gravitational numbers, and the dispersion.
\end{theorem}
\begin{proof}
By the previous theorem, $\text{wedge}(c \cdot v, c \cdot w)_{ij} = c^2 \text{wedge}(v, w)_{ij}$ and similarly for the denominator.
Dividing gives $\frac{c^2 \text{wedge}(v, w)_{ij}}{c^2 \text{wedge}(v, w)_{i'j'}} = \frac{\text{wedge}(v, w)_{ij}}{\text{wedge}(v, w)_{i'j'}}$.
\end{proof}
\lean{Dimension.wedge\_ratio\_invariant}

\begin{theorem}[Only Ratios are Fixed]
The framework fixes every prediction up to one overall factor, which it cannot determine because
there is nothing left to express it in. That factor is the unit. A dimensionless theory must have exactly one,
and this framework has exactly one—the minimum possible.
\end{theorem}
\begin{proof}
All physical quantities are ratios of wedges, which are invariant under simultaneous rescaling of the pattern.
No absolute scale of the pattern appears in any prediction; only their relative sizes matter.
\end{proof}
\lean{Dimension.only\_ratios\_are\_fixed}

\begin{definition}
The normalization $\alpha$ relating the threshold density $\rho$ to the rank-one parameter $k$:
\[\alpha(\rho, k) = \frac{k - 1}{\rho}.\]
This is a conversion relation—one equation, one unknown—not a prediction, and it requires the assumption
that defects are uniformly distributed in the symmetric space's volume.
\end{definition}
\lean{Dimension.alphaFromDensity}

\begin{theorem}[Density-Normalization Relation]
For nonzero density $\rho$, the product $\alpha(\rho, k) \cdot \rho$ equals $k - 1$:
\[\alpha(\rho, k) \cdot \rho = k - 1.\]
\end{theorem}
\begin{proof}
$\frac{k-1}{\rho} \cdot \rho = k - 1$ by definition.
\end{proof}
\lean{Dimension.density\_normalization\_relation}

\begin{theorem}
At $k = 3$ with the observed density $\rho$, the relation is $\alpha(\rho, 3) \cdot \rho = 2$.
\end{theorem}
\begin{proof}
Substituting $k = 3$ into the density-normalization relation gives $\alpha(\rho, 3) \cdot \rho = 3 - 1 = 2$.
\end{proof}
\lean{Dimension.normalization\_returns\_input}

\subsection{Slice}

Label proliferation with extension of scale range is unavoidable—every time a threshold is crossed,
new structures become resolvable and new labels are required. But this is not evidence of hidden structure;
it is the price of describing a wide range in a single formalism. The rotational quantity carried by the framework
is the wedge (exterior product) of two scale directions, and rotation exists only in the presence of
direction variation. The question "why three spatial dimensions" has a precise answer: $k = 3$ is the unique
dimension in which the generated rotations carry direction labels.

\begin{definition}
The content newly resolvable when extending the scale range from $t$ to $t'$ is the set of structures
whose thresholds lie in the gap:
\[\text{newContent}(\mu, t, t') = \{j : \mu_j \leq t' \text{ and } \mu_j > t\}.\]
\end{definition}
\lean{Slice.newContent}

\begin{theorem}
Membership in $\text{newContent}(\mu, t, t')$ is equivalent to: $\mu_j \leq t'$ and $\mu_j \not\leq t$.
\end{theorem}
\begin{proof}
Immediate from the definition.
\end{proof}
\lean{Slice.mem\_newContent}

\begin{theorem}[Extension Forces New Labels]
If any structure has its threshold strictly between $t$ and $t'$, the extended description must carry
a label the local one does not need. Specifically, if $t < \mu_j \leq t'$ for some $j$, then
$\text{newContent}(\mu, t, t')$ is nonempty.
\end{theorem}
\begin{proof}
The structure $j$ satisfies $\mu_j \leq t'$ (resolvable at the top) and $\mu_j > t$ (not resolvable at the bottom),
so $j \in \text{newContent}(\mu, t, t')$.
\end{proof}
\lean{Slice.newContent\_nonempty\_of\_threshold}

\begin{theorem}[Excess is Monotone in Range]
Widening the scale window can only add new labels, never remove them. If $t_1 \leq t_2$ and $t' \leq t''$, then
\[\text{newContent}(\mu, t_2, t') \subseteq \text{newContent}(\mu, t_1, t'').\]
\end{theorem}
\begin{proof}
If $j \in \text{newContent}(\mu, t_2, t')$, then $\mu_j \leq t' \leq t''$ and $\mu_j > t_2 \geq t_1$, so $j \in \text{newContent}(\mu, t_1, t'')$.
\end{proof}
\lean{Slice.newContent\_mono}

\begin{theorem}[Threshold-Free Ranges Cost Nothing]
If there is no structure with threshold in the range $(t, t']$, then the extended description is exact:
no new labels are required.
\end{theorem}
\begin{proof}
If $\text{newContent}(\mu, t, t')$ is empty, then no structure $j$ satisfies both $\mu_j \leq t'$ and $\mu_j > t$,
which means every structure either is not resolvable at $t'$ or is already resolvable at $t$.
\end{proof}
\lean{Slice.newContent\_empty\_of\_no\_threshold}

\begin{theorem}[No Universal Labelling]
No fixed label set works at every scale. Given any threshold function $\mu$ with $\mu_j \to \infty$,
for any scale $t$ there exists a higher scale $t'$ such that new content appears.
\end{theorem}
\begin{proof}
By the unboundedness assumption, there exists a structure $j$ with $\mu_j > t$. Setting $t' = \mu_j$ gives
$t < t' \leq t'$ and $\mu_j > t$, so $j \in \text{newContent}(\mu, t, t')$.
\end{proof}
\lean{Slice.no\_finite\_universal\_labelling}

\begin{theorem}[Label Count Measures Range]
The number of labels a formalism carries measures the log-range it covers, not the depth of the physics.
A description at a single scale needs only the content of that scale; a description across a range
must carry labels for all content across that range.
\end{theorem}
\begin{proof}
This is the collective statement of the previous theorems: new labels accumulate monotonically with range,
and any fixed labelling eventually becomes insufficient.
\end{proof}
\lean{Slice.label\_count\_measures\_range}

\begin{definition}
The \textbf{wedge} (exterior product) of two scale directions $v, w \in \mathbb{R}^k$ is the antisymmetric $k \times k$ matrix
\[\text{wedge}(v, w)_{ij} = v_i w_j - v_j w_i.\]
This is the rotational label generated by two scale directions.
\end{definition}
\lean{Slice.wedge}

\begin{theorem}[Wedge Equals Bracket]
The wedge is exactly what the coupling generates. For boost matrices,
\[\text{wedge}(v, w)_{ij} = [\text{boost}(v), \text{boost}(w)]_{i+1,j+1},\]
where the right side is the $(i+1, j+1)$ entry of the commutator.
\end{theorem}
\begin{proof}
This follows from Algebra.boost\_bracket\_spatial, which computed the bracket entry by entry.
\end{proof}
\lean{Slice.wedge\_eq\_bracket}

\begin{theorem}[Wedge is Antisymmetric]
For all $i, j$,
\[\text{wedge}(v, w)_{ij} = -\text{wedge}(v, w)_{ji}.\]
\end{theorem}
\begin{proof}
$\text{wedge}(v, w)_{ij} = v_i w_j - v_j w_i = -(v_j w_i - v_i w_j) = -\text{wedge}(v, w)_{ji}$.
\end{proof}
\lean{Slice.wedge\_antisymm}

\begin{theorem}[Wedge is Bilinear]
The wedge is bilinear: $\text{wedge}(u + v, w) = \text{wedge}(u, w) + \text{wedge}(v, w)$.
\end{theorem}
\begin{proof}
$\text{wedge}(u + v, w)_{ij} = (u_i + v_i) w_j - (u_j + v_j) w_i = (u_i w_j - u_j w_i) + (v_i w_j - v_j w_i)$.
\end{proof}
\lean{Slice.wedge\_bilinear\_left}

\begin{theorem}[Wedge Vanishes Iff Parallel]
Two directions generate no rotation if and only if they are parallel (componentwise proportional).
$\text{wedge}(v, w) = 0$ if and only if $v_i w_j = v_j w_i$ for all $i, j$.
\end{theorem}
\begin{proof}
Each entry $\text{wedge}(v, w)_{ij} = v_i w_j - v_j w_i$ vanishes iff $v_i w_j = v_j w_i$.
\end{proof}
\lean{Slice.wedge\_eq\_zero\_iff\_parallel}

\begin{theorem}[Single Direction Carries No Rotation]
A homogeneous pattern (one direction, possibly with varying magnitude) carries no rotational label.
$\text{wedge}(v, v)_{ij} = 0$ for all $i, j$.
\end{theorem}
\begin{proof}
$\text{wedge}(v, v)_{ij} = v_i v_j - v_j v_i = 0$.
\end{proof}
\lean{Slice.homogeneous\_carries\_no\_rotation}

\begin{theorem}[Nonzero Rotation Requires Two Non-Parallel Directions]
If a wedge is nonzero, no scalar $c$ satisfies $w = c \cdot v$.
\end{theorem}
\begin{proof}
If $w = c \cdot v$, then $\text{wedge}(v, w)_{ij} = v_i (cv_j) - v_j (cv_i) = 0$ for all $i, j$.
Contrapositive: if $\text{wedge}(v, w)_{ij} \neq 0$ for some $i, j$, then $w \neq c \cdot v$ for any $c$.
\end{proof}
\lean{Slice.rotation\_requires\_biaxial}

\begin{definition}
The number of independent wedges of $k$ directions (equivalently, the dimension of $\Lambda^2 \mathbb{R}^k$):
\[\text{wedgeDim}(m) = (m+1) \cdot m \quad \text{(with $k = m+1$)}.\]
\end{definition}
\lean{Slice.wedgeDim}

\begin{definition}
The number of directions, written to match the scaling of $\text{wedgeDim}$:
\[\text{dirDim}(m) = 2(m+1).\]
\end{definition}
\lean{Slice.dirDim}

\begin{theorem}[Unique Matching Dimension]
The wedges are as numerous as the directions if and only if $m = 2$, i.e., $k = 3$.
The equation $k(k-1)/2 = k$ has unique positive solution $k = 3$.
\end{theorem}
\begin{proof}
Setting $\text{wedgeDim}(m) = \text{dirDim}(m)$ gives $(m+1)m = 2(m+1)$. Dividing by $m+1$ (which is positive) yields $m = 2$.
\end{proof}
\lean{Slice.wedge\_dim\_eq\_iff}

\begin{theorem}[Why Three Dimensions]
Three spatial dimensions is the unique dimension in which the coupling closes on the labels the axioms provide.
Above $k = 3$ there are more wedges than directions, so some generated rotation is unlabelled; below $k = 3$
the rotation sector degenerates. This is the unique $m$ satisfying $\text{wedgeDim}(m) = \text{dirDim}(m)$.

The arithmetic is a theorem. The physics—that the framework \emph{requires} its generated rotations to be
direction-labelled rather than merely permitting some to go unlabelled—is a judgement registered as such,
on the same footing as the rank-one identification.
\end{theorem}
\begin{proof}
The existence of $m = 2$ is direct calculation. Uniqueness follows from the fact that $m = 2$ is the unique
solution to $(m+1)m = 2(m+1)$.
\end{proof}
\lean{Slice.three\_is\_unique}

\begin{theorem}[Higher Dimensions Have More Wedges]
For $m > 2$ (i.e., $k > 3$), the wedges outnumber the directions: $\text{dirDim}(m) < \text{wedgeDim}(m)$.
\end{theorem}
\begin{proof}
For $m > 2$, we have $(m+1)m > 2(m+1)$ because $m > 2$.
\end{proof}
\lean{Slice.more\_wedges\_than\_directions}

\begin{theorem}[Lower Dimensions Have Fewer Wedges]
For $m < 2$ (i.e., $k < 3$), the wedges are fewer than the directions: $\text{wedgeDim}(m) < \text{dirDim}(m)$.
\end{theorem}
\begin{proof}
By direct case analysis: for $m = 0$, $\text{wedgeDim}(0) = 0 < 2 = \text{dirDim}(0)$; for $m = 1$, $\text{wedgeDim}(1) = 2 < 4 = \text{dirDim}(1)$.
\end{proof}
\lean{Slice.fewer\_wedges\_than\_directions}

\subsection{Axes}

The resolution of the apparent dilemma between rank-one and angular momentum is a misreading.
The theorem $\text{wedge}(v, v) = 0$ says a homogeneous (constant) pattern carries no rotation, not that a uniaxial pattern
carries none. The pattern is a single vector, and rotation arises only from its *variation* in space.
A topologically charged defect—one whose pattern winds—always carries angular momentum.
Real rank one forces the pattern to be a single vector, so biaxial or higher-order patterns are not available.

\begin{theorem}[Rotation Requires Variation]
Nonzero rotational label requires two non-parallel values of the pattern.
\end{theorem}
\begin{proof}
At any one place, the pattern is a single vector, and its wedge with itself vanishes. Rotation can arise
only between the pattern's values at different points.
\end{proof}
\lean{Axes.rotation\_needs\_variation}

\begin{theorem}[Parallel Patterns Carry No Rotation]
A pattern whose values are all parallel (a multiple of one fixed direction at each point) generates
no rotation anywhere, because every wedge vanishes.
\end{theorem}
\begin{proof}
If $v_1, v_2, \ldots$ are all multiples of one direction $u$, then $\text{wedge}(f_m \cdot u, f_l \cdot u) = 0$ for any scalars $f_m, f_l$.
\end{proof}
\lean{Axes.combable\_no\_rotation}

\begin{theorem}[Wound Defect Carries Rotation]
A pattern taking two non-parallel values generates nonzero rotational label. If $v_i w_j \neq v_j w_i$
for some $i, j$, then $\text{wedge}(v, w)_{ij} \neq 0$.
\end{theorem}
\begin{proof}
By definition, $\text{wedge}(v, w)_{ij} = v_i w_j - v_j w_i$, which is nonzero by hypothesis.
\end{proof}
\lean{Axes.wound\_carries\_rotation}

\begin{theorem}[Rotationless Configuration is Trivial]
If the rotational label vanishes everywhere, the configuration is topologically trivial.
The framework forbids a charged defect with zero angular momentum.
\end{theorem}
\begin{proof}
If $\text{wedge}(v, w)_{ij} = 0$ for all $i, j$, then all values of the pattern are parallel, so the configuration
is combable and deforms to a constant.
\end{proof}
\lean{Axes.rotationless\_is\_trivial}

\begin{theorem}[Rotation Iff Variation]
Rotation exists if and only if the pattern's values are not everywhere parallel.
\end{theorem}
\begin{proof}
The equivalence is immediate from the definitions of $\text{wedge}$ and the notion of parallel vectors.
\end{proof}
\lean{Axes.rotation\_iff\_variation}

\begin{theorem}[No Second Axis]
Real rank one forces the pattern to be a single vector. Any other scaling is either a multiple of it
(the same axis) or has nonzero wedge with it (showing up as rotation, not a second axis).
If $v$ has a nonzero component and $v_i w_j = v_j w_i$ for all $i, j$, then $w$ is a multiple of $v$:
$w = \frac{w_{i_0}}{v_{i_0}} v$.
\end{theorem}
\begin{proof}
From $v_{i_0} w_i = v_i w_{i_0}$ for all $i$, we get $w_i = \frac{w_{i_0}}{v_{i_0}} v_i$.
\end{proof}
\lean{Axes.no\_second\_axis}

\begin{theorem}[Axis or Variation]
Given the pattern $v$ with a nonzero component, every other scaling $w$ is either parallel to $v$ or has
nonzero wedge with it.
\end{theorem}
\begin{proof}
By the wedge criterion, either all wedges vanish (parallel case) or some wedge is nonzero (variation case).
\end{proof}
\lean{Axes.axis\_or\_variation}

\begin{theorem}[Stabilizer is Positive-Dimensional]
The rotation generator in the plane transverse to the pattern annihilates it. This is the signature of
a uniaxial order parameter: a continuous stabilizer. A biaxial or fully anisotropic pattern would have
a discrete stabilizer, which does not support point defects.

For the concrete $3 \times 3$ rotation $J$ (in the $1$--$2$ plane) applied to the distinguished direction
$(1, 0, 0)$ in the standard basis, the result is zero: $J \cdot (1, 0, 0) = 0$. And $J$ is nonzero.
\end{theorem}
\begin{proof}
Computing $J \cdot (1, 0, 0) = (0, 0, 0) - (0, 0, 0) = 0$ while $J \neq 0$.
\end{proof}
\lean{Axes.axis\_stabilizer\_nontrivial}

\begin{theorem}[Multiplet Flatness Requires Per-Direction Scalings]
The theorem that ``multiplets are internally flat'' required a scaling operator per direction.
Real rank one provides only one vector and one scaling, so that physical reading is withdrawn.
What stands is that a homogeneous pattern generates no rotation, which is a statement about constancy,
not about the relation between directions.
\end{theorem}
\begin{proof}
This is immediate from the fact that $\text{wedge}(v, v) = 0$ for any vector $v$.
\end{proof}
\lean{Axes.multiplet\_flatness\_needs\_per\_direction\_scalings}


\section{Part II: The Branch Point}

\subsection{Locus}

This section establishes a critical distinction between isotropic and anisotropic scale patterns. When the scale pattern is fully isotropic (distinguishing no direction), all rotational labels vanish, and no defects exist. Conversely, any nonzero rotational structure forces the pattern to be anisotropic. A seventh axiom of observability (A7) makes $k = 3$ a theorem rather than an assumption, and ensures that no scale-invariant fundamental interactions exist in the framework.

\begin{definition}[\lean{Locus.Observability}]
A dimension parameter $m$ satisfies the \textbf{Observability} condition if the rotational dimension equals the scale dimension: that is,
\[
\text{rotDim}_2(m) = \text{scaleDim}_2(m).
\]
\end{definition}

\begin{theorem}[\lean{Locus.observability\_eq\_postulate}]
For any $m \in \mathbb{N}$, the Observability condition at $m$ is equivalent to the Observability postulate $\text{A7}(m)$.
\begin{proof}
This is immediate from the definitions; the two formulations are identical.
\end{proof}
\end{theorem}

\begin{theorem}[\lean{Locus.three\_from\_observability}]
If $m \in \mathbb{N}$ satisfies the Observability condition, then $m = 2$.
\begin{proof}
By definition, Observability means $\text{rotDim}_2(m) = \text{scaleDim}_2(m)$. The theorem that rotational and scale dimensions match exactly when $m = 2$ (from the dimensional analysis in \texttt{Dimension.lean}) yields the result.
\end{proof}
\end{theorem}

\begin{theorem}[\lean{Locus.observability\_fails\_elsewhere}]
For any $m \in \mathbb{N}$ with $m \neq 2$, the Observability condition fails at $m$.
\begin{proof}
Suppose the Observability condition held. Then by \texttt{three\_from\_observability}, we would have $m = 2$, contradicting $m \neq 2$. Thus the Observability condition cannot hold.
\end{proof}
\end{theorem}

\begin{theorem}[\lean{Locus.observability\_iff\_three}]
For any $m \in \mathbb{N}$, the Observability condition holds at $m$ if and only if $m = 2$.
\begin{proof}
Observability at $m$ is equivalent to $\text{rotDim}_2(m) = \text{scaleDim}_2(m)$, and dimensional analysis shows this occurs precisely when $m = 2$.
\end{proof}
\end{theorem}

\begin{theorem}[\lean{Locus.no\_scale\_invariant\_interaction}]
Let $u_0, \kappa, \rho, t_0, t, t' \in \mathbb{R}$ with $t \neq t'$ and $\kappa \rho \neq 0$. Then the inverse-square couplings at different scales are unequal:
\[
\text{invSqCoupling}(u_0, \kappa, \rho, t_0, t) \neq \text{invSqCoupling}(u_0, \kappa, \rho, t_0, t').
\]
\begin{proof}
Suppose for contradiction that the couplings are equal. Then their definitions give
\[
u_0 + \kappa(\rho(t - t_0)) = u_0 + \kappa(\rho(t' - t_0)),
\]
which simplifies to $\kappa \rho t = \kappa \rho t'$. Since $\kappa \rho \neq 0$, cancellation yields $t = t'$, contradicting $t \neq t'$.
\end{proof}
\end{theorem}

\begin{theorem}[\lean{Locus.constant\_coupling\_iff\_no\_content}]
Let $u_0, \kappa, \rho, t_0, t \in \mathbb{R}$ with $\kappa \rho = 0$. Then the inverse-square coupling at any scale is constant:
\[
\text{invSqCoupling}(u_0, \kappa, \rho, t_0, t) = u_0.
\]
\begin{proof}
By definition, $\text{invSqCoupling}(u_0, \kappa, \rho, t_0, t) = u_0 + \kappa(\rho(t - t_0))$. Since $\kappa \rho = 0$, the product $\kappa \rho (t - t_0) = 0$, so the coupling equals $u_0$.
\end{proof}
\end{theorem}

\begin{definition}[\lean{Locus.Isotropic}]
A pattern $v \in \text{Fin}(k) \to \mathbb{R}$ (an assignment of $k$ real values) is \textbf{isotropic} if it is constant, i.e., $v_i = v_j$ for all indices $i, j \in \text{Fin}(k)$.
\end{definition}

\begin{theorem}[\lean{Locus.isotropic\_pair\_parallel}]
Let $v, w \in \text{Fin}(k) \to \mathbb{R}$ be two isotropic patterns. Then for any indices $i, j \in \text{Fin}(k)$,
\[
v_i w_j = v_j w_i.
\]
\begin{proof}
Since $v$ is isotropic, $v_i = v_j$. Since $w$ is isotropic, $w_i = w_j$. Therefore $v_i w_j = v_j w_j = v_j w_i$.
\end{proof}
\end{theorem}

\begin{theorem}[\lean{Locus.isotropic\_no\_rotation}]
Let $v, w \in \text{Fin}(k) \to \mathbb{R}$ be two isotropic patterns. Then for any indices $i, j \in \text{Fin}(k)$, the wedge (exterior) product vanishes:
\[
\text{wedge}(v, w)_{ij} := v_i w_j - v_j w_i = 0.
\]
\begin{proof}
By the isotropic pair parallel theorem, $v_i w_j = v_j w_i$ for all $i, j$. Thus $\text{wedge}(v, w)_{ij} = v_i w_j - v_j w_i = 0$.
\end{proof}
\end{theorem}

\begin{theorem}[\lean{Locus.labels\_require\_anisotropy}]
Let $v, w \in \text{Fin}(k) \to \mathbb{R}$ and $i, j \in \text{Fin}(k)$. If the wedge product at $(i,j)$ is nonzero,
\[
\text{wedge}(v, w)_{ij} \neq 0,
\]
then at least one of the patterns is anisotropic: $\neg \text{Isotropic}(v) \lor \neg \text{Isotropic}(w)$.
\begin{proof}
Suppose for contradiction that both $v$ and $w$ are isotropic. Then by the isotropic no rotation theorem, $\text{wedge}(v, w)_{ij} = 0$ for all $i, j$, contradicting the assumption that $\text{wedge}(v, w)_{ij} \neq 0$.
\end{proof}
\end{theorem}

\begin{theorem}[\lean{Locus.everything\_switches\_on\_together}]
Let $v \in \text{Fin}(k) \to \mathbb{R}$ be an isotropic pattern. Then the following hold simultaneously:
\begin{enumerate}
\item For all $i, j \in \text{Fin}(k)$, $\text{wedge}(v, v)_{ij} = 0$ (no rotational label).
\item For all $i, j \in \text{Fin}(k)$, $v_i = v_j$ (no direction is distinguished).
\end{enumerate}
\begin{proof}
Both statements follow directly from the definition of isotropy. The first is the isotropic no rotation theorem applied to $v$ with itself; the second is the definition of isotropy.
\end{proof}
\end{theorem}

\begin{theorem}[\lean{Locus.particle\_locus\_is\_anisotropic}]
Let $v, w \in \text{Fin}(k) \to \mathbb{R}$ and $i, j \in \text{Fin}(k)$. If the wedge product is nonzero,
\[
\text{wedge}(v, w)_{ij} \neq 0,
\]
then the pair $(v, w)$ cannot both be isotropic: $\neg(\text{Isotropic}(v) \land \text{Isotropic}(w))$.
\begin{proof}
Suppose both $v$ and $w$ are isotropic. Then by the isotropic no rotation theorem, $\text{wedge}(v, w)_{ij} = 0$, contradicting the hypothesis. Therefore the pair cannot both be isotropic.
\end{proof}
\end{theorem}

\subsection{Dynamics}

This section translates the coupling structure between the scale and rotation sectors into consequences for the dynamics and spectrum of the theory. The evolution equation is shown to be Bianchi identity, which forces conservation of energy-momentum. Within multiplets (degeneracy blocks of the directional scale), curvature vanishes because parallel transport operators commute. Curvature is generated only between multiplets, and multiplet structure changes only where the scale pattern changes.

\begin{theorem}[\lean{Dynamics.no\_flatness\_from\_commuting\_derivations}]
Let $A \in \text{Fin}(n) \to M$ be a sequence of elements in a differential ring where all partial derivatives vanish: $D_i(x) = 0$ for all $i$ and all $x \in M$. Then for any $i, j \in \text{Fin}(n)$, the curvature is purely commutator-valued:
\[
F(A)_{ij} = \text{ad}(A_i)(A_j) := [A_i, A_j].
\]
\begin{proof}
The curvature is defined as $F(A)_{ij} = D_i(A_j) - D_j(A_i) - [A_i, A_j]$. Since all derivatives vanish, $D_i(A_j) = 0$ and $D_j(A_i) = 0$, leaving $F(A)_{ij} = [A_i, A_j]$.
\end{proof}
\end{theorem}

\begin{theorem}[\lean{Dynamics.nonflat\_witness}]
The two directional scaling operators $K_1, K_2$ from the coupling satisfy:
\begin{enumerate}
\item $\text{ad}(K_1)(K_2) = J$.
\item $J \in \text{Matrix}(\text{Fin}(3), \text{Fin}(3), \mathbb{R})$ is nonzero.
\end{enumerate}
\begin{proof}
The bracket $[K_1, K_2]$ equals the rotation $J$ by the definition of the coupling's boost-rotation relation. The matrix $J$ is nonzero because its $(1,2)$-entry is nonzero.
\end{proof}
\end{theorem}

\begin{theorem}[\lean{Dynamics.commutative\_sector\_transports\_commute}]
Let $C$ be a commutative ring with a differential ring structure of order $n$, and let $A \in \text{Fin}(n) \to C$. For any $i, j \in \text{Fin}(n)$ and $m \in C$, the covariant derivatives in the commutative sector commute:
\[
\text{covD}_A^{(i)}(\text{covD}_A^{(j)}(m)) - \text{covD}_A^{(j)}(\text{covD}_A^{(i)}(m)) = 0.
\]
\begin{proof}
The commutator of covariant derivatives is $[\text{covD}_A^{(i)}, \text{covD}_A^{(j)}](m) = \text{ad}(F(A)_{ij})(m)$. In a commutative ring, $\text{ad}(x)(m) = [x, m] = xm - mx = 0$ for all $x, m$ by commutativity. Thus the covariant derivatives commute.
\end{proof}
\end{theorem}

\begin{theorem}[\lean{Dynamics.commutative\_curvature\_undetectable}]
Let $C$ be a commutative ring with a differential ring structure of order $n$, and let $A \in \text{Fin}(n) \to C$. For any $i, j \in \text{Fin}(n)$ and $m \in C$, the curvature acts trivially:
\[
\text{ad}(F(A)_{ij})(m) = 0.
\]
\begin{proof}
In a commutative ring, the adjoint action $\text{ad}(x)(m) = [x, m] = xm - mx = 0$ by commutativity.
\end{proof}
\end{theorem}

\begin{theorem}[\lean{Dynamics.source\_conserved\_of\_field\_equation}]
Let $A \in \text{Fin}(n) \to M$, $\kappa \in M$, and $T \in \text{Fin}(n) \to \text{Fin}(n) \to M$. If the curvature is proportional to a source tensor,
\[
F(A)_{ij} = \kappa \cdot T_{ij} \quad \text{for all } i, j,
\]
then the source satisfies the cyclic covariant divergence condition for all $i, j, k$:
\[
\text{covD}_A^{(i)}(\kappa \cdot T_{jk}) + \text{covD}_A^{(j)}(\kappa \cdot T_{ki}) + \text{covD}_A^{(k)}(\kappa \cdot T_{ij}) = 0.
\]
\begin{proof}
Substitute the field equation into the left-hand side:
\[
\text{covD}_A^{(i)}(F(A)_{jk}) + \text{covD}_A^{(j)}(F(A)_{ki}) + \text{covD}_A^{(k)}(F(A)_{ij}).
\]
The Bianchi identity (from \texttt{Connection.lean}) states that the cyclic covariant divergence of the curvature vanishes identically, so this sum equals zero.
\end{proof}
\end{theorem}

\begin{theorem}[\lean{Dynamics.no\_field\_equation\_of\_nonconserved}]
Let $A \in \text{Fin}(n) \to M$, $\kappa \in M$, and $T \in \text{Fin}(n) \to \text{Fin}(n) \to M$. If the cyclic divergence of $\kappa T$ fails to vanish at some point,
\[
\text{covD}_A^{(i)}(\kappa \cdot T_{jk}) + \text{covD}_A^{(j)}(\kappa \cdot T_{ki}) + \text{covD}_A^{(k)}(\kappa \cdot T_{ij}) \neq 0,
\]
then $\kappa T$ cannot be the source of any connection's curvature. That is, there does not exist a connection with
\[
F(A)_{ij} = \kappa \cdot T_{ij} \quad \text{for all } i, j.
\]
\begin{proof}
Suppose such a field equation held. Then by the source conserved of field equation theorem, the cyclic divergence would vanish, contradicting the hypothesis. Therefore no such field equation can exist.
\end{proof}
\end{theorem}

\begin{theorem}[\lean{Dynamics.source\_antisymm}]
Let $A \in \text{Fin}(n) \to M$, $\kappa \in M$, and $T \in \text{Fin}(n) \to \text{Fin}(n) \to M$. If the curvature is proportional to the source,
\[
F(A)_{ij} = \kappa \cdot T_{ij} \quad \text{for all } i, j,
\]
then the source inherits the curvature's antisymmetry: for all $i, j$,
\[
\kappa \cdot T_{ij} = -(\kappa \cdot T_{ji}).
\]
\begin{proof}
The curvature is antisymmetric: $F(A)_{ij} = -F(A)_{ji}$ by definition. Substituting the field equation yields $\kappa \cdot T_{ij} = -(\kappa \cdot T_{ji})$.
\end{proof}
\end{theorem}

\begin{theorem}[\lean{Dynamics.field\_equation\_content}]
Let $A \in \text{Fin}(n) \to M$, $\kappa \in M$, and $T \in \text{Fin}(n) \to \text{Fin}(n) \to M$. If the curvature is proportional to the source,
\[
F(A)_{ij} = \kappa \cdot T_{ij} \quad \text{for all } i, j,
\]
then both conservation and antisymmetry hold for all $i, j, k$:
\begin{enumerate}
\item $\text{covD}_A^{(i)}(\kappa \cdot T_{jk}) + \text{covD}_A^{(j)}(\kappa \cdot T_{ki}) + \text{covD}_A^{(k)}(\kappa \cdot T_{ij}) = 0$ (cyclic divergence vanishes).
\item $\kappa \cdot T_{ij} = -(\kappa \cdot T_{ji})$ (antisymmetry).
\end{enumerate}
\begin{proof}
Both statements follow directly from the source conserved of field equation and source antisymm theorems.
\end{proof}
\end{theorem}

\begin{definition}[\lean{Dynamics.Realizes}]
Let $s \in \text{Fin}(N) \to A^\times$ be an assignment of units to directions, and $S \in \text{Fin}(N) \to \text{Matrix}(m, m, \mathbb{R})$ be an assignment of scaling operators. We say $S$ \textbf{realizes} $s$ if equal units have equal scaling operators: for all $a, b \in \text{Fin}(N)$,
\[
s_a = s_b \implies S_a = S_b.
\]
\end{definition}

\begin{theorem}[\lean{Dynamics.same\_multiplet\_no\_rotation}]
Let $s \in \text{Fin}(N) \to A^\times$ and $S \in \text{Fin}(N) \to \text{Matrix}(m, m, \mathbb{R})$, where $A$ is a commutative ring. Suppose $S$ realizes $s$. If two directions $a, b \in \text{Fin}(N)$ are in the same multiplet (i.e., in the same orbit under the action determined by equal scale values),
\[
\text{SameOrbit}_s(a, b),
\]
then the bracket of their scaling operators vanishes:
\[
\text{br}(S_a, S_b) = [S_a, S_b] = 0.
\]
\begin{proof}
By definition of SameOrbit, $s_a = s_b$. Since $S$ realizes $s$, we have $S_a = S_b$. A scaling operator brackets to zero with itself by the coupling structure, so $[S_a, S_b] = [S_a, S_a] = 0$.
\end{proof}
\end{theorem}

\begin{theorem}[\lean{Dynamics.rotation\_implies\_different\_multiplet}]
Let $s \in \text{Fin}(N) \to A^\times$ and $S \in \text{Fin}(N) \to \text{Matrix}(m, m, \mathbb{R})$, where $A$ is a commutative ring. Suppose $S$ realizes $s$. If the bracket of scaling operators at two directions is nonzero,
\[
\text{br}(S_a, S_b) \neq 0,
\]
then the directions cannot be in the same multiplet:
\[
\neg \text{SameOrbit}_s(a, b).
\]
\begin{proof}
Suppose for contradiction that $\text{SameOrbit}_s(a, b)$. Then by the same multiplet no rotation theorem, $\text{br}(S_a, S_b) = 0$, contradicting the hypothesis. Therefore the directions must be in different multiplets.
\end{proof}
\end{theorem}

\begin{theorem}[\lean{Dynamics.larger\_multiplet\_fewer\_generators}]
Let $s \in \text{Fin}(N) \to A^\times$ and $S \in \text{Fin}(N) \to \text{Matrix}(m, m, \mathbb{R})$, where $A$ is a commutative ring. Suppose $S$ realizes $s$ and the scale is fully degenerate (isotropic): $s_a = s_b$ for all $a, b \in \text{Fin}(N)$. Then for any two directions $a, b \in \text{Fin}(N)$, the bracket of their scaling operators vanishes:
\[
\text{br}(S_a, S_b) = 0.
\]
\begin{proof}
Since the scale is fully degenerate, $s_a = s_b$ for all $a, b$. Because $S$ realizes $s$, we have $S_a = S_b$. Thus $\text{br}(S_a, S_b) = [S_a, S_b] = [S_a, S_a] = 0$ by the self-bracket property.
\end{proof}
\end{theorem}

\begin{theorem}[\lean{Dynamics.multiplets\_fixed\_by\_pattern}]
Let $s, t \in \text{Fin}(N) \to A^\times$ be two scale patterns. If they have the same coincidence structure,
\[
s_a = s_b \iff t_a = t_b \quad \text{for all } a, b \in \text{Fin}(N),
\]
then the multiplet structures determined by these patterns are identical: for all $a, b \in \text{Fin}(N)$,
\[
\text{SameOrbit}_s(a, b) \iff \text{SameOrbit}_t(a, b).
\]
\begin{proof}
The multiplets are defined by the equivalence relation on directions given by their scale values. Two patterns with identical coincidence structure define the same equivalence relation, hence the same multiplet structure.
\end{proof}
\end{theorem}

\begin{theorem}[\lean{Dynamics.multiplet\_change\_needs\_pattern\_change}]
Let $s, t \in \text{Fin}(N) \to A^\times$ be two scale patterns. If the multiplet structure changes between them,
\[
\exists a, b \in \text{Fin}(N) : \neg(\text{SameOrbit}_s(a, b) \iff \text{SameOrbit}_t(a, b)),
\]
then their coincidence structures must differ:
\[
\exists a, b \in \text{Fin}(N) : \neg((s_a = s_b) \iff (t_a = t_b)).
\]
\begin{proof}
Suppose for contradiction that the coincidence structures are identical: $(s_a = s_b) \iff (t_a = t_b)$ for all $a, b$. Then by the multiplets fixed by pattern theorem, the multiplet structures are identical, contradicting the assumption that they differ. Therefore the coincidence structures must differ.
\end{proof}
\end{theorem}


\section{Part II.a (i): Geometry and the Vacuum Equation}

\subsection{Conformal}

A position-dependent scale field $\sigma$, the same in all directions, induces a conformal metric $g = e^{2\sigma}\delta$. The Riemannian geometry of such a metric is entirely determined by derivatives of $\sigma$: the Christoffel symbols are polynomial in $\nabla \sigma$, and curvature becomes the second-order inhomogeneity of the scale field. This file computes the Levi-Civita data and derives the master regression theorem stating that the Riemann tensor is the Kulkarni--Nomizu product of the bare Euclidean metric $\delta$ with a scale deformation tensor $D$ built from $\sigma$.

\begin{definition}
The scale gradient: $\sigma_i = \partial_i \sigma$. \lean{sig}
\end{definition}

\begin{definition}
The scale Hessian: $\sigma_{ij} = \partial_i \partial_j \sigma$. \lean{hess}
\end{definition}

\begin{theorem}
The scale Hessian is symmetric: $\sigma_{ij} = \sigma_{ji}$. \lean{hess\_symm}
\end{theorem}

\begin{proof}
By commutativity of mixed partial derivatives.
\end{proof}

\begin{theorem}
Mixed derivatives of the scale gradient commute: $\partial_i (\partial_j \sigma) = \partial_j (\partial_i \sigma)$. \lean{d\_sig\_comm}
\end{theorem}

\begin{proof}
Immediate from symmetry of the Hessian.
\end{proof}

\begin{definition}
The flat Laplacian of the scale: $\Delta\sigma = \sum_i \sigma_{ii}$. \lean{lap}
\end{definition}

\begin{definition}
The squared scale gradient: $|\nabla\sigma|^2 = \sum_i \sigma_i^2$. \lean{gradsq}
\end{definition}

\begin{theorem}
A constant factors out of the squared gradient sum: $\sum_m c(\sigma_m^2) = c \cdot |\nabla\sigma|^2$. \lean{sum\_mul\_gradsq}
\end{theorem}

\begin{proof}
Immediate by distributivity.
\end{proof}

\begin{theorem}
A constant summed over $\text{Fin}\,n$ yields $n$ times the constant: $\sum_{a=1}^n c = nc$. \lean{sum\_const\_fin}
\end{theorem}

\begin{proof}
Immediate from the definition of finite sums.
\end{proof}

The scale connection is the one-form $A_i = \sigma_i$. Its curvature $F_{ij} = \partial_i A_j - \partial_j A_i$ vanishes identically by commutativity of mixed partials, so the scale connection is integrable. This absence of a second clock effect is why an integrable scale geometry succeeds where Weyl's 1918 non-integrable unified theory failed.

\begin{definition}
The scale connection one-form: $A_i = \sigma_i$. \lean{scaleConn}
\end{definition}

\begin{definition}
The curvature of the scale connection: $F_{ij} = \partial_i A_j - \partial_j A_i$. \lean{scaleCurv}
\end{definition}

\begin{theorem}
The scale connection is flat: $F_{ij} = 0$. \lean{scaleCurv\_eq\_zero}
\end{theorem}

\begin{proof}
By commutativity of mixed partial derivatives, $\partial_i(\partial_j \sigma) = \partial_j(\partial_i \sigma)$, so the curvature vanishes identically.
\end{proof}

\begin{definition}
The Christoffel symbols of $g = e^{2\sigma}\delta$:
$$\Gamma^a_{bc} = \delta^a_b \sigma_c + \delta^a_c \sigma_b - \delta_{bc}\sigma^a,$$
where indices are raised and lowered with $\delta$.  The conformal factor cancels entirely, leaving a polynomial in $\nabla\sigma$; the connection sees only changes of scale, not the scale itself. \lean{Chr}
\end{definition}

\begin{theorem}
The Christoffel symbols are symmetric in their lower indices: $\Gamma^a_{bc} = \Gamma^a_{cb}$. \lean{Chr\_symm}
\end{theorem}

\begin{proof}
Immediate from the definition and symmetry of the Kronecker delta.
\end{proof}

\begin{definition}
The Riemann tensor: $R^a_{bce} = \partial_c \Gamma^a_{be} - \partial_e \Gamma^a_{bc} + \sum_m (\Gamma^a_{cm}\Gamma^m_{be} - \Gamma^a_{em}\Gamma^m_{bc})$. \lean{Rm}
\end{definition}

\begin{definition}
The Ricci tensor: $R_{be} = \sum_a R^a_{bae}$. \lean{Ric}
\end{definition}

\begin{definition}
The bare scalar curvature (the $\delta$-trace of Ricci): $R_{\text{bare}} = \sum_b R_{bb}$. The true scalar curvature of $g$ is $e^{-2\sigma}$ times this; the conformal factor is kept explicit so that no invertibility hypothesis is ever needed. \lean{RscBare}
\end{definition}

\begin{theorem}
The sum of derivatives of Christoffel symbols: $\sum_a \partial_a \Gamma^a_{be} = 2\sigma_{be} - \delta_{be}\Delta\sigma$. \lean{sum\_d\_Chr}
\end{theorem}

\begin{proof}
Expanding $\Gamma^a_{be}$ and taking derivatives yields three terms summed over $a$. After applying the product rule and regrouping by $i$-values, this simplifies to the stated formula.
\end{proof}

\begin{theorem}
The trace of Christoffel symbols: $\sum_a \Gamma^a_{am} = n\sigma_m$. \lean{sum\_Chr\_diag}
\end{theorem}

\begin{proof}
Evaluating $\Gamma^a_{am} = \delta^a_a \sigma_m = \sigma_m$ and summing over $a = 1, \ldots, n$ gives $n\sigma_m$.
\end{proof}

\begin{theorem}
The derivative of the trace: $\sum_a \partial_e \Gamma^a_{ba} = n\sigma_{eb}$. \lean{sum\_d\_Chr\_trace}
\end{theorem}

\begin{proof}
Use $\Gamma^a_{ba} = \Gamma^a_{ab}$ by symmetry, then sum $\partial_e \Gamma^a_{ab} = \partial_e (\Gamma^a_{aa}\sigma_b) = n\partial_e \sigma_b = n\sigma_{eb}$.
\end{proof}

\begin{theorem}
A weighted sum of Christoffel symbols: $\sum_m (n\sigma_m) \Gamma^m_{be} = n(2\sigma_b\sigma_e - \delta_{be}|\nabla\sigma|^2)$. \lean{sum\_ChrTrace\_Chr}
\end{theorem}

\begin{proof}
Expanding $\Gamma^m_{be}$ and distributing the sum yields three terms, each involving Kronecker deltas. Regrouping produces the stated form.
\end{proof}

\begin{theorem}
The master contraction of the Riemann tensor, in full generality:
\begin{align*}
\sum_m \Gamma^a_{cm}\Gamma^m_{be} &= 2\delta_{ac}\sigma_b\sigma_e - \delta_{ac}\delta_{be}|\nabla\sigma|^2 \\
&\quad + \delta_{ab}\sigma_c\sigma_e + \delta_{ae}\sigma_b\sigma_c \\
&\quad - \delta_{cb}\sigma_a\sigma_e - \delta_{ce}\sigma_a\sigma_b.
\end{align*}
Every later curvature statement is a specialization of this identity. \lean{sum\_Chr\_Chr\_full}
\end{theorem}

\begin{proof}
Expanding both Christoffel symbols and collecting terms by their Kronecker delta factors yields the six terms shown.
\end{proof}

\begin{theorem}
The doubly contracted case: $\sum_{a,m} \Gamma^a_{em}\Gamma^m_{ba} = (n+2)\sigma_b\sigma_e - 2\delta_{be}|\nabla\sigma|^2$. \lean{sum\_Chr\_Chr}
\end{theorem}

\begin{proof}
Apply the master contraction with $a \to a$ and $c \to e$, then sum over $a$. Using $\delta_{aa} = n$ and the definitions of Laplacian and squared gradient gives the result.
\end{proof}

\begin{theorem}
Ricci curvature is the second-order inhomogeneity of the scale:
$$R_{be} = -(n-2)(\sigma_{be} - \sigma_b\sigma_e) - \delta_{be}(\Delta\sigma + (n-2)|\nabla\sigma|^2).$$
Read physically: gravity is the failure of the local unit of measure to vary affinely across the substrate. \lean{Ric\_eq}
\end{theorem}

\begin{proof}
The Ricci tensor can be split into three parts: $\sum_a \partial_a \Gamma^a_{be}$, $\sum_a \partial_e \Gamma^a_{ba}$, and the quadratic terms. Using the contraction theorems, each part contributes to the final formula.
\end{proof}

\begin{theorem}
The scalar curvature of $g = e^{2\sigma}\delta$:
$$e^{2\sigma}R = -2(n-1)\Delta\sigma - (n-1)(n-2)|\nabla\sigma|^2.$$
For $n=3$ this is the Hamiltonian constraint of a conformally flat slice; its linearization is Poisson's equation. \lean{RscBare\_eq}
\end{theorem}

\begin{proof}
Sum the Ricci components $R_{bb}$ using the Ricci formula and collect terms involving the Laplacian and squared gradient.
\end{proof}

The entire Riemann tensor is a linear function of one symmetric tensor built from $\sigma$. This deformation tensor measures the obstruction to the local unit varying affinely.

\begin{definition}
The scale deformation tensor:
$$D_{ij} = 2\sigma_{ij} - 2\sigma_i\sigma_j + \delta_{ij}|\nabla\sigma|^2,$$
the obstruction to the local unit of measure varying affinely. (The factor of 2 is carried explicitly so that no division by 2 is required, keeping every statement valid over an arbitrary commutative ring.) \lean{Defm}
\end{definition}

\begin{theorem}
The deformation tensor is symmetric: $D_{ij} = D_{ji}$. \lean{Defm\_symm}
\end{theorem}

\begin{proof}
By symmetry of the Hessian, commutativity of multiplication, and symmetry of the Kronecker delta.
\end{proof}

\begin{theorem}
The derivative part of the Riemann tensor:
\begin{align*}
\partial_c \Gamma^a_{be} - \partial_e \Gamma^a_{bc} &= \delta_{ae}\sigma_{bc} - \delta_{be}\sigma_{ac} \\
&\quad - \delta_{ac}\sigma_{be} + \delta_{bc}\sigma_{ae}.
\end{align*}
\lean{d\_Chr\_diff}
\end{theorem}

\begin{proof}
Expand the Christoffel symbols and apply the product rule to each term, then use commutativity of mixed partials to simplify.
\end{proof}

\begin{theorem}
The master regression theorem: curvature is scale deformation. The Riemann tensor of $g = e^{2\sigma}\delta$ satisfies
$$2R^a_{bce} = \delta_{ae}D_{bc} - \delta_{ac}D_{be} + \delta_{bc}D_{ae} - \delta_{be}D_{ac}.$$
The Riemann tensor is exactly the Kulkarni--Nomizu product of the bare Euclidean bookkeeping $\delta$ with the scale deformation $D$. Nothing about curvature is independent of the scale field: geometry is bookkeeping times scale deformation, and nothing else. \lean{two\_Rm\_eq}
\end{theorem}

\begin{proof}
Split the Riemann tensor into derivative and quadratic parts. The derivative part is given by the preceding theorem. The quadratic part simplifies using the master contraction theorem. Combining these yields the Kulkarni--Nomizu form.
\end{proof}

Over any ring in which 2 is invertible---in particular, any $\mathbb{R}$-algebra---the master identity gives at once that a deformation-free scale field is exactly a flat geometry. Vacuum is a statement about the uniformity of scale change, never about the absence of scale.

\begin{theorem}
Deformation-free scale implies flat space. Not merely Ricci-flat: the whole Riemann tensor vanishes.
$$\forall i,j\,: D_{ij} = 0 \implies R^a_{bce} = 0.$$
\lean{Rm\_eq\_zero\_of\_Defm\_eq\_zero}
\end{theorem}

\begin{proof}
From the master regression theorem, if all deformation components vanish, then $2R^a_{bce} = 0$. Multiplying by $1/2$ (which exists in our ring) gives $R^a_{bce} = 0$.
\end{proof}

\begin{theorem}
Consequently Ricci vanishes as well: if $D_{ij} = 0$ for all $i,j$, then $R_{be} = 0$. \lean{Ric\_eq\_zero\_of\_Defm\_eq\_zero}
\end{theorem}

\begin{proof}
The Ricci tensor is the trace of the Riemann tensor, $R_{be} = \sum_a R^a_{bae}$. If $R^a_{bae} = 0$ for all $a$, the sum vanishes.
\end{proof}

\begin{theorem}
The deformation tensor is recovered from curvature by tracing, so $D$ carries exactly the same information as $R$:
$$2R_{be} = -(n-2)D_{be} - \delta_{be}(\text{tr}\,D),$$
where $\text{tr}\,D = 2\Delta\sigma + (n-2)|\nabla\sigma|^2$. \lean{two\_Ric\_eq}
\end{theorem}

\begin{proof}
From the Ricci formula $R_{be} = -(n-2)(\sigma_{be} - \sigma_b\sigma_e) - \delta_{be}(\Delta\sigma + (n-2)|\nabla\sigma|^2)$ and the definition of $D_{be}$, direct algebra gives the stated form.
\end{proof}

\subsection{Frame}

The original axiom (A4) postulated a single scalar scale $\sigma$, the same in every direction. This axiom is provably unable to carry Schwarzschild: if the temporal and radial scales are reciprocal and the scale is isotropic, the conformal factor must be trivial and the geometry must be flat. The diagnosis is physical: energy-momentum is a tensor, and a tensor source cannot set a scalar scale. A4 is therefore weakened to A4' (directional scale), which allows one scale per direction. Under isotropy all directions carry the same scale, recovering the old A4 as a special case; away from isotropy, genuine non-flat solutions to the reciprocal relation exist.

\begin{definition}
A4' (directional scale): One local unit of measure per direction, each invertible. The local scale in direction $a$ is $s_a$, a unit of the ring. \lean{Frame.DirScale}
\end{definition}

\begin{definition}
The measured metric component in direction $a$, with signature $\eta$: $g_a = \eta_a s_a^2$. \lean{Frame.DirScale.metric}
\end{definition}

\begin{definition}
Isotropy: all directions carry the same scale, $s_a = s_b$ for all $a,b$. \lean{Frame.DirScale.Isotropic}
\end{definition}

\begin{theorem}
Under isotropy the metric is the bare one times a single conformal factor. Everything proved in the conformal sector therefore survives as the isotropic special case of A4'. \lean{Frame.DirScale.isotropic\_metric}
\end{theorem}

\begin{proof}
If $s_a = s_b$ for all $a,b$, then all metric components $g_a = \eta_a s_a^2$ and $g_b = \eta_b s_a^2$ share the same scale factor, differing only by their signatures.
\end{proof}

\begin{definition}
The reciprocal-scale relation: $s_t \cdot s_r = 1$. This is equivalent to $g_t \cdot g_r = -1$ for signature $(\eta_t, \eta_r) = (-1, 1)$, and is the defining feature of the Schwarzschild and Reissner-Nordström families. In scale language: the temporal and radial units are exact reciprocals. \lean{Frame.DirScale.Reciprocal}
\end{definition}

\begin{theorem}
The reciprocal relation written on metric components: if $s_t \cdot s_r = 1$, $\eta_t = -1$, and $\eta_r = 1$, then $g_t \cdot g_r = -1$. \lean{Frame.DirScale.reciprocal\_metric}
\end{theorem}

\begin{proof}
We have $g_t \cdot g_r = (-1)\cdot s_t^2 \cdot (1) \cdot s_r^2 = -(s_t \cdot s_r)^2 = -(1)^2 = -1$.
\end{proof}

\begin{theorem}
Under the old, isotropic axiom the Schwarzschild relation collapses to flat space. If all directions share one scale and the temporal and radial scales are reciprocal, then that scale squares to 1: $s_t^2 = 1$. The conformal factor is trivial and the geometry is the bare Euclidean bookkeeping. So the old axiom did not merely fail to reproduce Schwarzschild---it excluded it. The axiom was too strong. \lean{Frame.DirScale.isotropic\_reciprocal\_forces\_flat}
\end{theorem}

\begin{proof}
If $s_a = s_b$ for all $a,b$, then $s_t = s_r$. The reciprocal condition $s_t \cdot s_r = 1$ becomes $s_t^2 = 1$.
\end{proof}

\begin{theorem}
Consequently every metric component is the bare one: the geometry is flat. If the isotropic and reciprocal conditions hold, then $g_a = \eta_a$ for all $a$. \lean{Frame.DirScale.isotropic\_reciprocal\_metric\_bare}
\end{theorem}

\begin{proof}
From $s_t^2 = 1$ and $s_a = s_t$ for all $a$, we have $s_a^2 = 1$. Thus $g_a = \eta_a \cdot s_a^2 = \eta_a \cdot 1 = \eta_a$.
\end{proof}

\begin{theorem}
Under A4' the Schwarzschild relation is satisfiable non-trivially. For any unit $u$ there exists a directional scale field with $s_t = u$ and $s_r = u^{-1}$, reciprocal by construction, and isotropic only in the degenerate case $u^2 = 1$. The temporal and radial scales genuinely differ, which is what Schwarzschild requires and what the scalar axiom could not provide. \lean{Frame.exists\_reciprocal\_nonisotropic}
\end{theorem}

\begin{proof}
Define a directional scale by setting $s_a = u$ if $a = t$, $s_a = u^{-1}$ if $a = r$, and $s_a = 1$ otherwise. Then $s_t \cdot s_r = u \cdot u^{-1} = 1$ by the unit axiom. For isotropy to hold, we would need $u = u^{-1}$, which implies $u^2 = 1$, contradicting the hypothesis $u^2 \neq 1$.
\end{proof}

\subsection{Unify}

The directional scale A4' is the axiom; A4 (scalar scale) is a theorem about its isotropic locus. This unifies a potentially inconsistent pair of statements: the conformal sector developed under A4, and the Frame sector under A4', now coexist as the isotropic sector of the latter. Every theorem previously proved about $\sigma$ holds verbatim, now with a sharp domain of validity. Nothing is lost and the axiom list has one entry where it had two.

\begin{definition}
The measured metric of a directional scale field: diagonal in the frame, with the local unit $s_a$ in direction $a$ and signature $\eta$:
$$g_{ab} = \delta_{ab} \eta_a s_a^2.$$
\lean{Unify.toMetric}
\end{definition}

\begin{definition}
The metric of the old axiom A4: one scalar scale, $g_{ab} = \delta_{ab}\eta_a u^2$ where $u$ is the scale. \lean{Unify.conformalMetric}
\end{definition}

\begin{theorem}
The old axiom is the isotropic locus of the new one. If all directions carry the same scale $s$, the metric of A4' is literally the conformal metric of A4. Every result proved in the conformal sector therefore holds unchanged, as a statement about isotropic scale fields. \lean{Unify.toMetric\_of\_isotropic}
\end{theorem}

\begin{proof}
If $s_a = s_c$ for all $a,c$, then $g_{ab} = \delta_{ab}\eta_a s_a^2 = \delta_{ab}\eta_a s_c^2$, which is the conformal metric.
\end{proof}

\begin{definition}
Embedding of a scalar scale: every scalar scale is a directional scale, namely an isotropic one. Specifically, $s_a = u$ for all $a$. \lean{Unify.ofScalar}
\end{definition}

\begin{theorem}
The scalar embedding is isotropic: the scale $\text{ofScalar}(u)$ has $s_a = u$ for all $a$, so $s_a = s_b$ for all pairs $(a,b)$. \lean{Unify.ofScalar\_isotropic}
\end{theorem}

\begin{proof}
Immediate from the definition.
\end{proof}

\begin{theorem}
The metric of an embedded scalar scale is the conformal metric:
$$\text{ofScalar}(u).toMetric(\eta) = conformalMetric(u, \eta).$$
\lean{Unify.ofScalar\_toMetric}
\end{theorem}

\begin{proof}
Since $s_a = u$ for all $a$, we have $g_{ab} = \delta_{ab}\eta_a u^2$, which is the conformal metric.
\end{proof}

\begin{theorem}
The isotropic locus is exactly the image of the scalar theory. An $n$-direction scale field is isotropic if and only if it is $\text{ofScalar}(u)$ for some unit $u$. So the old axiom is not merely implied by the new one---it is precisely the new one restricted to a sub-locus, with nothing left over. \lean{Unify.isotropic\_iff\_ofScalar}
\end{theorem}

\begin{proof}
($\Rightarrow$) If the scale is isotropic, all directions carry $s_a = s_{\langle 0 \rangle}$ for all $a$, so the scale field is $\text{ofScalar}(s_{\langle 0 \rangle})$. ($\Leftarrow$) If the scale is $\text{ofScalar}(u)$, then $s_a = u$ for all $a$, so it is isotropic.
\end{proof}

Scalar scale fields (A2/A3) embed as isotropic directional scales. The metric induced by the directional geometry is identical to the conformal metric of the scalar theory.

\begin{definition}
A scalar scale field $F$, viewed as an isotropic directional scale field. \lean{ScaleField.toDirScale}
\end{definition}

\begin{theorem}
The embedding is isotropic: $F.toDirScale.Isotropic$ for any scalar field $F$. \lean{Unify.toDirScale\_isotropic}
\end{theorem}

\begin{proof}
Immediate from $\text{ofScalar\_isotropic}$.
\end{proof}

\begin{theorem}
The conformal sector sits inside A4' exactly. The metric obtained from A2/A3 by the new axiom is the metric the old axiom postulated:
$$F.toDirScale.toMetric(\eta) = conformalMetric(F.s, \eta).$$
Consequently the curvature computed in Conformal.lean from $F.\sigma$ is the curvature of the A4' metric of $F.toDirScale$. \lean{Unify.toDirScale\_toMetric}
\end{theorem}

\begin{proof}
Apply $\text{ofScalar\_toMetric}$ to $F.s$.
\end{proof}

\begin{theorem}
Collected coherence statement. For any scalar scale field $F$:
\begin{enumerate}
\item its directional embedding is isotropic;
\item the metric it induces under A4' equals the conformal metric of A4;
\item consequently the curvature computed in Conformal.lean from $F.\sigma$ is the curvature of the A4' metric of $F.toDirScale$.
\end{enumerate}
This is the precise sense in which the earlier development is a sub-theory of the present one rather than a competitor to it. \lean{Unify.sector\_coherence}
\end{theorem}

\begin{proof}
Combine the three preceding theorems.
\end{proof}

The two sectors are distinguished by a single observable: whether the metric is a multiple of the bare bookkeeping.

\begin{theorem}
Under isotropy the metric is a common multiple of the signature. If $s_a = s_b$ for all $a,b$, then
$$g_{aa}\eta_b = g_{bb}\eta_a.$$
\lean{Unify.isotropic\_ratio\_const}
\end{theorem}

\begin{proof}
$g_{aa} = \eta_a s_a^2$ and $g_{bb} = \eta_b s_a^2$, so $g_{aa}\eta_b = \eta_a s_a^2 \cdot \eta_b = \eta_a \cdot \eta_b s_a^2 = g_{bb}\eta_a$.
\end{proof}

\begin{theorem}
Conversely, if two directions carry different squared scales the field is not isotropic. If $s_a^2 \neq s_b^2$, then the scale field is not isotropic. This is the observable that Schwarzschild needs and that the scalar axiom could not supply. \lean{Unify.not\_isotropic\_of\_ne}
\end{theorem}

\begin{proof}
If the scale is isotropic, then $s_a = s_b$, so $s_a^2 = s_b^2$. Contrapositive: if $s_a^2 \neq s_b^2$, the scale is not isotropic.
\end{proof}

\subsection{Newton}

Given the regression theorem that relates gravity to scale field inhomogeneity, one further statement---a scale response law---closes the system. The statement is that the scalar curvature of the measured metric is proportional to the local energy density. In the first-order (weak-field) regime, this constraint is exactly Poisson's equation, with the log-scale identified as minus the Newtonian potential. Both signs emerge naturally without being imposed. The first-order regime is not merely an approximation; it is realized exactly by extending the substrate to dual numbers $A[\epsilon]/(\epsilon^2)$.

\begin{definition}
The Newtonian potential is minus the log-scale: $\Phi = -\sigma$. Deep in a gravitational well the scale is enlarged ($\sigma > 0$), so $\Phi < 0$ as it must be. \lean{Newton.newtPot}
\end{definition}

\begin{theorem}
The Laplacian of the Newtonian potential: $\Delta\Phi = -\Delta\sigma$. \lean{Newton.lap\_newtPot}
\end{theorem}

\begin{proof}
Immediate from $\Phi = -\sigma$ and linearity of the Laplacian.
\end{proof}

\begin{theorem}
In the first-order regime the scalar curvature is linear in the scale: the $|\nabla\sigma|^2$ term drops out identically. If $|\nabla\sigma|^2 = 0$, then
$$R_{\text{bare}} = -2(n-1)\Delta\sigma.$$
\lean{Newton.RscBare\_linear}
\end{theorem}

\begin{proof}
From Conformal.RscBare\_eq, drop the term $-(n-1)(n-2)|\nabla\sigma|^2$ when $|\nabla\sigma|^2 = 0$.
\end{proof}

\begin{theorem}
Poisson's equation. From the scale response law $e^{2\sigma}R = \kappa\rho$ in the first-order regime,
$$2(n-1)\Delta\Phi = \kappa\rho, \quad \Phi = -\sigma.$$
For $n = 3$ and $\kappa = 16\pi G$ this is $\Delta\Phi = 4\pi G\rho$ exactly. Newtonian gravity is therefore not an input: it is the linearization of the statement that energy sets the local scale. \lean{Newton.poisson}
\end{theorem}

\begin{proof}
From the scale response law and the linear regime, $R_{\text{bare}} = -2(n-1)\Delta\sigma = \kappa\rho$. Substituting $\sigma = -\Phi$ gives $2(n-1)\Delta\Phi = \kappa\rho$.
\end{proof}

\begin{theorem}
The physical case $n = 3$: $4\Delta\Phi = \kappa\rho$, i.e. $\Delta\Phi = 4\pi G\rho$ for $\kappa = 16\pi G$. \lean{Newton.poisson\_three}
\end{theorem}

\begin{proof}
From Poisson's equation with $n = 3$, we have $2(3-1)\Delta\Phi = \kappa\rho$, i.e., $4\Delta\Phi = \kappa\rho$.
\end{proof}

The first-order regime is not a hypothesis we hope is harmless. Extend the substrate by a square-zero infinitesimal, $A[\epsilon]/(\epsilon^2)$, and let the scale be infinitesimal, $\sigma = \epsilon\varphi$. Then $|\nabla\sigma|^2 = \epsilon^2|\nabla\varphi|^2 = 0$ identically, and all results above hold with no approximation whatsoever. Poisson's equation is exact.

\begin{theorem}
The derivations extend componentwise to the dual numbers $A[\epsilon]/(\epsilon^2)$: $\partial_i(x, y) = (\partial_i x, \partial_i y)$ where the first component is the ordinary derivative and the second is the dual part. \lean{Newton.Dual.instScaleAlgebra}
\end{theorem}

\begin{proof}
The axioms of $ScaleAlgebra$ (additivity and Leibniz rule of derivations, commutativity of mixed partials) are verified componentwise.
\end{proof}

\begin{theorem}
An infinitesimal scale field $\sigma = \epsilon\varphi$ has purely infinitesimal gradient: $\sigma_i = \epsilon(\partial_i\varphi)$. \lean{Newton.Dual.sig\_inr}
\end{theorem}

\begin{proof}
$\sigma_i = \partial_i \sigma = \partial_i(\epsilon\varphi) = \epsilon\partial_i\varphi$ by the Leibniz rule applied componentwise.
\end{proof}

\begin{theorem}
The linearised regime is exact. For an infinitesimal scale field the quadratic invariant vanishes identically---no term is discarded:
$$|\nabla\sigma|^2 = 0 \quad \text{for } \sigma = \epsilon\varphi.$$
\lean{Newton.Dual.gradsq\_inr}
\end{theorem}

\begin{proof}
$|\nabla\sigma|^2 = \sum_i \sigma_i^2 = \sum_i (\epsilon\partial_i\varphi)^2 = \sum_i \epsilon^2(\partial_i\varphi)^2 = 0$ since $\epsilon^2 = 0$ in the dual numbers.
\end{proof}

\begin{theorem}
Consequently Poisson's equation holds exactly for an infinitesimal scale field: no linearization error, because there is nothing to linearize.
$$2(n-1)\Delta\Phi = \kappa\alpha$$
for $\Phi = -\epsilon\varphi$ and $\alpha$ the energy density in dual numbers. \lean{Newton.Dual.poisson\_exact}
\end{theorem}

\begin{proof}
Apply the Poisson equation theorem with $\sigma = \epsilon\varphi$ and the exactness of $|\nabla\sigma|^2 = 0$.
\end{proof}

\subsection{Schwarzschild}

The vacuum condition (vanishing Ricci) does not directly imply reciprocity $s_t \cdot s_r = 1$. It implies something weaker: the sum of temporal and radial log-scales is constant. Getting from constant to one is a separate step, and it is the step where physical content appears---it is asymptotic flatness, the requirement that far from the source the local unit agrees with the fiducial. But deriving the constancy itself requires computing the Ricci tensor for the diagonal ansatz, which is the outstanding piece of the gravity chain.

This file proves the part of that computation where the non-obvious cancellation happens. Write the metric in directional form with only temporal and radial scales varying, $A = e^{2\sigma_t}$, $B = e^{2\sigma_r}$. The two Ricci components are complicated expressions in $A, B$ and their derivatives; forming the particular combination $R_t^t/A + R_r^r/B$ makes all the awkward terms cancel, leaving only a total logarithmic derivative of the product $AB$.

\begin{definition}
The temporal component of the Ricci tensor for $-A\,dt^2 + B\,dr^2 + r^2 d\Omega^2$, where primes denote $r$-derivatives:
$$R_{tt} = \frac{A''}{2B} - \frac{A'}{4B}\left(\frac{A'}{A} + \frac{B'}{B}\right) + \frac{A'}{rB}.$$
\lean{Schwarzschild.Rtt}
\end{definition}

\begin{definition}
The radial component:
$$R_{rr} = -\frac{A''}{2A} + \frac{A'}{4A}\left(\frac{A'}{A} + \frac{B'}{B}\right) + \frac{B'}{rB}.$$
\lean{Schwarzschild.Rrr}
\end{definition}

\begin{theorem}
Everything awkward cancels. The particular combination
$$\frac{R_{tt}}{A} + \frac{R_{rr}}{B} = \frac{1}{rB}\left(\frac{A'}{A} + \frac{B'}{B}\right).$$
The second-order derivatives cancel against each other and so do the quadratic terms, exactly and for all $A, B$. What survives is a total logarithmic derivative---which is why the vacuum constrains the product $AB$ and nothing else about it. \lean{Schwarzschild.ricci\_combination}
\end{theorem}

\begin{proof}
Substitute the definitions of $R_{tt}$ and $R_{rr}$, clear denominators using $A \neq 0$, $B \neq 0$, $r \neq 0$, and collect terms by degree and structure. All second-derivative and quadratic terms cancel, leaving only the logarithmic derivative term.
\end{proof}

\begin{theorem}
In vacuum the logarithmic derivative of the product vanishes. If $R_{tt} = 0$ and $R_{rr} = 0$, then
$$\frac{A'}{A} + \frac{B'}{B} = 0,$$
because the prefactor $1/(rB)$ is never zero. \lean{Schwarzschild.vacuum\_log\_derivative}
\end{theorem}

\begin{proof}
From the cancellation theorem, if both Ricci components vanish, then $\frac{1}{rB}(\frac{A'}{A} + \frac{B'}{B}) = 0$. Since $r \neq 0$ and $B \neq 0$, the prefactor is nonzero, so the logarithmic derivative must vanish.
\end{proof}

\begin{theorem}
In scale variables, the vacuum condition reads as the scale sum being stationary. With $A' / A = 2\sigma_t'$ and $B' / B = 2\sigma_r'$, the condition $A'/A + B'/B = 0$ becomes
$$\sigma_t' + \sigma_r' = 0,$$
which is exactly the hypothesis needed in Vacuum.reciprocity\_of\_asymptotic\_flatness. \lean{Schwarzschild.scale\_sum\_derivative\_zero}
\end{theorem}

\begin{proof}
Substitute the scale-variable substitutions into the logarithmic derivative condition: $2\sigma_t' + 2\sigma_r' = 0$ gives $\sigma_t' + \sigma_r' = 0$.
\end{proof}

\begin{theorem}
The vacuum condition on the diagonal ansatz is exactly the statement that the scale sum is stationary. Nothing else about $A$ and $B$ is constrained by it, which is why the constant value of the scale sum must come from a boundary condition. \lean{Schwarzschild.vacuum\_iff\_scale\_sum\_stationary}
\end{theorem}

\begin{proof}
Combine the three preceding results.
\end{proof}

\begin{theorem}
The combination that cancels is the one weighted by the inverse metric on the $t$--$r$ block. The area element in that plane is $AB$, and its logarithmic derivative is what the vacuum sets to zero:
$$A' B + A B' = 0.$$
Read physically: in vacuum, gravity trades temporal scale against radial scale and creates no scale in that plane at all. \lean{Schwarzschild.area\_element\_stationary}
\end{theorem}

\begin{proof}
From the vacuum condition $A'/A + B'/B = 0$, clear the common denominator to get $A'B + AB' = 0$.
\end{proof}

\subsection{RicciDiag}

The Schwarzschild.lean file verified the crucial cancellation that makes reciprocity work, but treated the two Ricci components as given. This file derives them from the metric using the Levi-Civita connection, closing the loop. For a diagonal metric $g_{tt} = -A(r)$, $g_{rr} = B(r)$, $g_{\theta\theta} = r^2$, the connection coefficients follow from the formula $\Gamma^a_{bc} = \frac{1}{2}g^{aa}(\partial_b g_{ac} + \partial_c g_{ab} - \partial_a g_{bc})$. The Ricci components then come from summing the connection and its derivatives according to the definition. The point of this file is that the resulting expressions are exactly those Schwarzschild.lean assumed.

\begin{definition}
Connection coefficients. For the diagonal metric with $g_{tt} = -A$, $g_{rr} = B$, $g_{\theta\theta} = r^2$:
$$\Gamma^t_{tr} = \frac{A'}{2A}.$$
\lean{RicciDiag.$\Gamma$\_t\_tr}
\end{definition}

\begin{definition}
$$\Gamma^r_{tt} = \frac{A'}{2B}.$$
\lean{RicciDiag.$\Gamma$\_r\_tt}
\end{definition}

\begin{definition}
$$\Gamma^r_{rr} = \frac{B'}{2B}.$$
\lean{RicciDiag.$\Gamma$\_r\_rr}
\end{definition}

\begin{definition}
$$\Gamma^\theta_{r\theta} = \Gamma^\phi_{r\phi} = \frac{1}{r}.$$
\lean{RicciDiag.$\Gamma$\_ang\_r}
\end{definition}

\begin{definition}
The contracted symbol:
$$\Gamma^c_{cr} = \Gamma^t_{tr} + \Gamma^r_{rr} + 2\Gamma^\theta_{r\theta}.$$
\lean{RicciDiag.$\Gamma$trace\_r}
\end{definition}

\begin{definition}
Radial derivatives of the connection coefficients. By the quotient rule:
$$\partial_r \Gamma^r_{tt} = \frac{A''}{2B} - \frac{A'B'}{2B^2}.$$
\lean{RicciDiag.d$\Gamma$\_r\_tt}
\end{definition}

\begin{definition}
$$\partial_r \Gamma^t_{tr} = \frac{A''}{2A} - \frac{(A')^2}{2A^2}.$$
\lean{RicciDiag.d$\Gamma$\_t\_tr}
\end{definition}

\begin{definition}
$$\partial_r \Gamma^\theta_{r\theta} = -\frac{1}{r^2}.$$
\lean{RicciDiag.d$\Gamma$\_ang\_r}
\end{definition}

\begin{definition}
The temporal Ricci component $R_{tt} = \partial_r \Gamma^r_{tt} + \Gamma^c_{cr}\Gamma^r_{tt} - 2\Gamma^t_{tr}\Gamma^r_{tt}$, using the fact that the metric is static so no $\partial_t$ survives and $\Gamma^c_{tt}$ is nonzero only for $c = r$:
$$R^{\text{Ricci}}_{tt} = \partial_r \Gamma^r_{tt} + \Gamma^c_{cr}\Gamma^r_{tt} - 2\Gamma^t_{tr}\Gamma^r_{tt}.$$
\lean{RicciDiag.ricciTT}
\end{definition}

\begin{definition}
The radial Ricci component:
$$R^{\text{Ricci}}_{rr} = -\partial_r \Gamma^t_{tr} - 2\partial_r \Gamma^\theta_{r\theta} + \Gamma^c_{cr}\Gamma^r_{rr} - (\Gamma^t_{tr})^2 - (\Gamma^r_{rr})^2 - 2(\Gamma^\theta_{r\theta})^2.$$
\lean{RicciDiag.ricciRR}
\end{definition}

\begin{theorem}
The temporal component, derived. Evaluating the Ricci sum reproduces exactly the expression Schwarzschild.Rtt assumed:
$$R^{\text{Ricci}}_{tt} = R_{tt}.$$
\lean{RicciDiag.ricciTT\_eq}
\end{theorem}

\begin{proof}
Substitute the definitions of all connection coefficients and their derivatives, clear denominators, and simplify using ring identities. The result matches the Schwarzschild.Rtt expression.
\end{proof}

\begin{theorem}
The radial component, derived:
$$R^{\text{Ricci}}_{rr} = R_{rr}.$$
\lean{RicciDiag.ricciRR\_eq}
\end{theorem}

\begin{proof}
Substitute the definitions and simplify as above.
\end{proof}

\begin{theorem}
The gravity chain, end to end. From the connection of the diagonal metric: the Ricci components are what Schwarzschild.lean assumed, their particular combination cancels, and the vacuum therefore makes the scale sum stationary.

Given $A = e^{2\sigma_t}$, $B = e^{2\sigma_r}$ with $A' / A = 2\sigma_t'$ and $B' / B = 2\sigma_r'$, and assuming the vacuum condition $R^{\text{Ricci}}_{tt} = 0$ and $R^{\text{Ricci}}_{rr} = 0$, we have
$$\sigma_t' + \sigma_r' = 0.$$

With Vacuum.reciprocity\_of\_asymptotic\_flatness this gives $s_t \cdot s_r = 1$. No step between the metric and that conclusion is now an assumption. \lean{RicciDiag.gravity\_chain}
\end{theorem}

\begin{proof}
From ricciTT\_eq and ricciRR\_eq, rewrite the Ricci conditions. Then apply Schwarzschild.vacuum\_iff\_scale\_sum\_stationary to conclude.
\end{proof}

\begin{theorem}
Consequently the $t$--$r$ area element is stationary in vacuum, derived now from the metric rather than assumed. If the Ricci components vanish, then $A'B + AB' = 0$. \lean{RicciDiag.area\_stationary\_derived}
\end{theorem}

\begin{proof}
Rewrite the Ricci conditions using ricciTT\_eq and ricciRR\_eq, then apply Schwarzschild.area\_element\_stationary.
\end{proof}

\subsection{Vacuum}

The vacuum condition in the standard static spherically symmetric computation does not directly give $s_t s_r = 1$. It gives something weaker: the sum of temporal and radial log-scales has vanishing derivative, so $s_t s_r$ is constant. Getting from constant to one requires an additional boundary condition: asymptotic flatness, the physical requirement that far from the source there is nothing, so the local unit must agree with the fiducial. This file derives reciprocity from the vacuum consequence plus asymptotic flatness.

\begin{theorem}
A function with vanishing derivative that vanishes at one point vanishes everywhere. If $f'(r) = 0$ for all $r$ and $f(r_0) = 0$ for some $r_0$, then $f(r) = 0$ for all $r$. \lean{Vacuum.const\_of\_deriv\_zero\_and\_vanishing}
\end{theorem}

\begin{proof}
A function with zero derivative everywhere is constant. If it equals zero at one point, the constant must be zero, so the function vanishes everywhere.
\end{proof}

\begin{theorem}
Reciprocity from asymptotic flatness. The vacuum condition makes the sum of the temporal and radial log-scales constant. Requiring that the sum vanish far from the source---asymptotic flatness, i.e. that the local unit agrees with the fiducial where there is no matter---forces the sum to vanish everywhere:
$$\forall r,\quad \sigma_t(r) + \sigma_r(r) = 0.$$

The constant is therefore not fitted; it is fixed by the only boundary condition the axioms permit, since by A5 the fiducial is exactly what ``far away'' means. \lean{Vacuum.reciprocity\_of\_asymptotic\_flatness}
\end{theorem}

\begin{proof}
Apply const\_of\_deriv\_zero\_and\_vanishing to the sum $f(r) = \sigma_t(r) + \sigma_r(r)$, using the hypothesis that $f'(r) = 0$ (from the vacuum condition) and $f(r_\infty) = 0$ (from asymptotic flatness).
\end{proof}

\begin{theorem}
The same statement multiplicatively: the temporal and radial scales are exact reciprocals,
$$s_t(r) \cdot s_r(r) = e^{\sigma_t(r)} e^{\sigma_r(r)} = e^{\sigma_t(r) + \sigma_r(r)} = 1.$$
\lean{Vacuum.scale\_product\_eq\_one}
\end{theorem}

\begin{proof}
From reciprocity\_of\_asymptotic\_flatness, $\sigma_t(r) + \sigma_r(r) = 0$, so $e^{\sigma_t(r) + \sigma_r(r)} = e^0 = 1$.
\end{proof}

\begin{theorem}
Without the boundary condition, reciprocity fails: the scale product is some constant, and any constant is compatible with the vacuum equation. The product is constant:
$$s_t(r_1) s_r(r_1) = s_t(r_2) s_r(r_2)$$
for any two points $r_1, r_2$. Asymptotic flatness is therefore doing real work and is not a convention. \lean{Vacuum.product\_constant\_without\_flatness}
\end{theorem}

\begin{proof}
From the vacuum condition that the sum $\sigma_t + \sigma_r$ has zero derivative, it is constant. Hence $e^{\sigma_t(r_1) + \sigma_r(r_1)} = e^{\sigma_t(r_2) + \sigma_r(r_2)}$, which gives the product equality.
\end{proof}

\begin{theorem}
The area element in the time--radial plane is what reciprocity conserves. The t--r area element is $AB = e^{2\sigma_t + 2\sigma_r}$. Under reciprocity (which sets $\sigma_t + \sigma_r = 0$ everywhere), we have
$$e^{2(\sigma_t + \sigma_r)} = e^0 = 1.$$
Read physically: in vacuum, gravity trades temporal scale against radial scale and creates no scale in that plane at all. \lean{Vacuum.tr\_area\_element\_constant}
\end{theorem}

\begin{proof}
From reciprocity\_of\_asymptotic\_flatness, $\sigma_t(r) + \sigma_r(r) = 0$, so $e^{2(\sigma_t + \sigma_r)} = e^0 = 1$.
\end{proof}


\section{Part II.a (ii): Deflection, Precession, and Gravitational Waves}

\subsection{Deflection}

The isotropic sector — the degenerate locus of A4 with a single scalar scale — produces exactly half the observed light-bending effect around a massive body. The effect depends on a single scalar, the gravitational potential $\Phi$, and at grazing incidence the prediction is a factor of two below the measured value. This file establishes that the missing half comes from spatial scale anisotropy, which a scalar scale cannot carry. The factor of two is thus not a fitting uncertainty but a structural statement: a theory with one scale can only supply the temporal contribution. Numerical verification is given for a ray grazing the Sun, matching VLBI measurements.

\begin{definition}[\lean{Deflection.deflectionIso}]
The grazing deflection produced by the temporal scale alone — the whole of what an isotropic (single-scalar) scale field can produce:
\[
\text{deflectionIso}(G, M, b, c) = \frac{2GM}{bc^2}.
\]
\end{definition}

\begin{definition}[\lean{Deflection.deflectionObs}]
The measured grazing deflection:
\[
\text{deflectionObs}(G, M, b, c) = \frac{4GM}{bc^2}.
\]
\end{definition}

\begin{theorem}[\lean{Deflection.iso\_is\_exactly\_half}]
The isotropic sector gives exactly half the observed deflection, with no adjustable quantity anywhere in the statement:
\[
2 \cdot \text{deflectionIso}(G, M, b, c) = \text{deflectionObs}(G, M, b, c).
\]
\begin{proof}
Immediate from the definitions by ring algebra.
\end{proof}
\end{theorem}

\begin{theorem}[\lean{Deflection.iso\_ratio\_half}]
Stated as a ratio, where the denominator is nonzero. If $\text{deflectionObs}(G, M, b, c) \neq 0$, then
\[
\frac{\text{deflectionIso}(G, M, b, c)}{\text{deflectionObs}(G, M, b, c)} = \frac{1}{2}.
\]
\begin{proof}
The hypothesis that the observed deflection is nonzero implies the isotropic deflection is nonzero. By the previous theorem, the observed deflection is twice the isotropic one. Dividing yields the result.
\end{proof}
\end{theorem}

\begin{theorem}[\lean{Deflection.missing\_half}]
The deficit is exactly the isotropic prediction over again: the missing contribution is the same size as the one a scalar scale can produce. In the scale reading, the temporal and the spatial scale contribute equally, and a single scalar carries only the temporal one:
\[
\text{deflectionObs}(G, M, b, c) - \text{deflectionIso}(G, M, b, c) = \text{deflectionIso}(G, M, b, c).
\]
\begin{proof}
Immediate from the definitions by ring algebra.
\end{proof}
\end{theorem}

\begin{definition}[\lean{Deflection.solarObsArcsec}]
The measured-branch (``observed'') prediction for grazing deflection of a ray around the Sun, in arcseconds:
\[
\text{solarObsArcsec} = \frac{4 \times 6.674 \times 10^{-11} \times 1.989 \times 10^{30}}{6.957 \times 10^8 \times (2.998 \times 10^8)^2} \times 206264.806.
\]
\end{definition}

\begin{definition}[\lean{Deflection.solarIsoArcsec}]
The isotropic-sector prediction for the same grazing deflection, in arcseconds:
\[
\text{solarIsoArcsec} = \frac{2 \times 6.674 \times 10^{-11} \times 1.989 \times 10^{30}}{6.957 \times 10^8 \times (2.998 \times 10^8)^2} \times 206264.806.
\]
\end{definition}

\begin{theorem}[\lean{Deflection.solarObs\_value}]
The observed prediction yields $1.7515 \text{\textquotedbl}$, matching the VLBI value $1.7509 \pm 0.0002 \text{\textquotedbl}$ to about $0.04\%$:
\[
1.7515 < \text{solarObsArcsec} < 1.7516.
\]
\begin{proof}
By numerical verification.
\end{proof}
\end{theorem}

\begin{theorem}[\lean{Deflection.solarIso\_value}]
The isotropic sector predicts $0.8758 \text{\textquotedbl}$, wrong by a factor of two. This is the sharpest quantitative statement in the development, and it is a refutation of the scalar axiom rather than a confirmation of anything:
\[
0.87 < \text{solarIsoArcsec} < 0.88.
\]
\begin{proof}
By numerical verification.
\end{proof}
\end{theorem}

\begin{theorem}[\lean{Deflection.solar\_ratio}]
The two solar numbers stand in the exact ratio $2:1$, as the general theorem requires:
\[
\text{solarObsArcsec} = 2 \cdot \text{solarIsoArcsec}.
\]
\begin{proof}
Immediate by ring algebra from the definitions.
\end{proof}
\end{theorem}

\begin{theorem}[\lean{Deflection.iso\_fails\_by\_far}]
The isotropic sector misses the measured value by more than $0.8 \text{\textquotedbl}$, which is four thousand times the VLBI uncertainty: the failure is not a modelling imprecision but a structural one:
\[
0.87 < \text{solarObsArcsec} - \text{solarIsoArcsec}.
\]
\begin{proof}
By the ratio theorem and the numerical bounds on the isotropic value, the difference exceeds $0.87 \text{\textquotedbl}$, whereas the VLBI uncertainty is $\pm 0.0002 \text{\textquotedbl}$.
\end{proof}
\end{theorem}

\subsection{PPN}

The measured deflection angle can be parametrised in terms of two dimensionless numbers: $\gamma$, weighting the spatial scale contribution, and $\beta$, measuring the nonlinearity of the temporal scale. The standard value is $\gamma = 1$. The isotropic sector, having no spatial scale at all, effectively sets $\gamma = 0$. The framework's scale reciprocity condition — the product of temporal and radial scales equals one — forces $\gamma = 1$ exactly to first order. This file proves that reciprocity determines the spatial-scale coefficient, paying the debt of the isotropic sector through a necessary relation among components.

\begin{definition}[\lean{PPN.deflectionGamma}]
Grazing deflection with the spatial scale weighted by $\gamma$:
\[
\text{deflectionGamma}(G, M, b, c, \gamma) = \left(\frac{2GM}{bc^2}\right)(1 + \gamma).
\]
Here $\gamma = 0$ keeps only the temporal scale, and $\gamma = 1$ is the measured case.
\end{definition}

\begin{theorem}[\lean{PPN.gamma\_zero\_is\_isotropic}]
With no spatial contribution the formula returns the isotropic sector's prediction — the half that a scalar scale can produce:
\[
\text{deflectionGamma}(G, M, b, c, 0) = \text{deflectionIso}(G, M, b, c).
\]
\begin{proof}
Immediate from the definitions by ring algebra.
\end{proof}
\end{theorem}

\begin{theorem}[\lean{PPN.gamma\_one\_is\_observed}]
With the spatial scale weighted equally, the formula returns the measured value:
\[
\text{deflectionGamma}(G, M, b, c, 1) = \text{deflectionObs}(G, M, b, c).
\]
\begin{proof}
Immediate from the definitions by ring algebra.
\end{proof}
\end{theorem}

\begin{theorem}[\lean{PPN.reciprocal\_expansion}]
If the radial factor is the exact inverse of the temporal one — which is what the reciprocity relation $s_t \cdot s_r = 1$ says — then the expansion identity holds:
\[
(1 - 2\mu)^{-1} = 1 + 2\mu + \frac{4\mu^2}{1 - 2\mu},
\]
where $1 - 2\mu \neq 0$. This is an identity, not an approximation. The coefficient of $\mu$ is $2$, and the remainder is displayed rather than dropped.
\begin{proof}
Clear the denominators and verify by field algebra.
\end{proof}
\end{theorem}

\begin{theorem}[\lean{PPN.gamma\_eq\_one\_of\_reciprocal}]
Reciprocity forces $\gamma = 1$. Matching $g_{rr} = 1 + 2\gamma\mu$ against the inverse of $g_{tt} = 1 - 2\mu$ at first order in $\mu$, the reciprocal condition leaves no freedom. Stated exactly: if $\mu \neq 0$, $1 - 2\mu \neq 0$, and the radial factor equals $1 + 2\gamma\mu$ and also equals $(1-2\mu)^{-1}$, then
\[
\gamma = 1 + \frac{2\mu}{1 - 2\mu}.
\]
\begin{proof}
By the reciprocal expansion theorem, $(1-2\mu)^{-1} = 1 + 2\mu + \frac{4\mu^2}{1-2\mu}$. Equating this to $1 + 2\gamma\mu$ gives $2\gamma\mu = 2\mu + \frac{4\mu^2}{1-2\mu}$. Dividing by the nonzero quantity $2\mu$ and clearing denominators yields the result.
\end{proof}
\end{theorem}

\begin{theorem}[\lean{PPN.gamma\_eq\_one\_first\_order}]
In the strict first-order regime — where $\mu^2$ is discarded, exactly as in the Newtonian limit — reciprocity gives $\gamma = 1$ on the nose. If $\mu \neq 0$ and $1 + 2\gamma\mu = 1 + 2\mu$, then $\gamma = 1$.
\begin{proof}
The hypothesis gives $2\gamma\mu = 2\mu$. Since $\mu \neq 0$, divide by $2\mu$ to obtain $\gamma = 1$.
\end{proof}
\end{theorem}

\begin{theorem}[\lean{PPN.deflection\_reciprocal\_eq\_observed}]
The debt is paid at the level of the coefficient. If $\gamma = 1$, then
\[
\text{deflectionGamma}(G, M, b, c, \gamma) = \text{deflectionObs}(G, M, b, c).
\]
The half that the isotropic sector could not produce is the radial scale, and its size was never free.
\begin{proof}
Substitute $\gamma = 1$ and apply the theorem that $\gamma = 1$ returns the observed value.
\end{proof}
\end{theorem}

\begin{theorem}[\lean{PPN.halves\_equal}]
The two halves are equal in size. The contribution from the spatial scale ($\gamma = 1$ term) minus the contribution from the temporal scale alone ($\gamma = 0$ term) equals the temporal contribution:
\[
\text{deflectionGamma}(G, M, b, c, 1) - \text{deflectionGamma}(G, M, b, c, 0) = \text{deflectionGamma}(G, M, b, c, 0).
\]
\begin{proof}
Immediate from the definitions by ring algebra.
\end{proof}
\end{theorem}

\begin{theorem}[\lean{PPN.solar\_reciprocal\_value}]
The reciprocal ($\gamma = 1$) prediction for a ray grazing the Sun, in arcseconds, is $1.7515 \text{\textquotedbl}$, against the VLBI value $1.7509 \pm 0.0002 \text{\textquotedbl}$:
\[
1.7515 < \text{solarObsArcsec} < 1.7516.
\]
\begin{proof}
This is the same bound as the observed deflection, since reciprocity forces $\gamma = 1$ and thus agreement with the measured value.
\end{proof}
\end{theorem}

\begin{theorem}[\lean{PPN.gamma\_prediction\_is\_sharp}]
Current solar-system bounds from the Cassini tracking data put $\gamma$ within about $2 \times 10^{-5}$ of $1$. The framework's reciprocity condition predicts $\gamma = 1$ with no free parameter. If $\gamma = 1$, then
\[
\gamma - 1 = 0.
\]
\begin{proof}
Immediate from the hypothesis.
\end{proof}
\end{theorem}

\subsection{Precession}

Mercury's perihelion advance per orbit is the second solar-system test, of a different kind and with a different failure factor in the isotropic sector. The coefficient in the standard parametrisation is $(2 - \beta + 2\gamma)/3$, where $\gamma$ is the spatial-scale weight (forced to $1$ by reciprocity) and $\beta$ measures the nonlinearity of the temporal scale (equals $1$ from A6$^\prime$). With both parameters at their framework values, the coefficient is $1$ and the prediction is measured. The isotropic sector, having no spatial scale, gives $1/3$. The two tests fail by factors $2$ and $3$ respectively, which are different and thus cannot both be absorbed by rescaling. Both arise from the same single parameter, indicating a single missing ingredient: the spatial scale.

\begin{definition}[\lean{Precession.coeff}]
The perihelion coefficient in terms of the temporal nonlinearity $\beta$ and the spatial-scale weight $\gamma$:
\[
\text{coeff}(\beta, \gamma) = \frac{2 - \beta + 2\gamma}{3}.
\]
\end{definition}

\begin{theorem}[\lean{Precession.coeff\_reciprocal}]
The reciprocal condition gives the measured coefficient. Here $\gamma = 1$ is forced by $s_t \cdot s_r = 1$ and $\beta = 1$ is the scale response of A6$^\prime$:
\[
\text{coeff}(1, 1) = 1.
\]
\begin{proof}
By direct computation.
\end{proof}
\end{theorem}

\begin{theorem}[\lean{Precession.coeff\_isotropic}]
The isotropic sector gives one third, since with no spatial scale $\gamma = 0$:
\[
\text{coeff}(1, 0) = \frac{1}{3}.
\]
\begin{proof}
By direct computation.
\end{proof}
\end{theorem}

\begin{theorem}[\lean{Precession.isotropic\_is\_third}]
The two coefficients differ by exactly a factor of three — a different factor from the two of light deflection, obtained from the same single parameter:
\[
\text{coeff}(1, 0) \times 3 = \text{coeff}(1, 1).
\]
\begin{proof}
By direct computation from the coefficients.
\end{proof}
\end{theorem}

\begin{theorem}[\lean{Precession.coeff\_affine\_in\_gamma}]
The sensitivity to $\gamma$ is what does the work: the coefficient is affine in $\gamma$ with slope $2/3$:
\[
\text{coeff}(\beta, \gamma) = \text{coeff}(\beta, 0) + \frac{2}{3}\gamma.
\]
\begin{proof}
Expand the definition and verify by ring algebra.
\end{proof}
\end{theorem}

\begin{definition}[\lean{Precession.mercuryFactor}]
Everything in the Mercury prediction except the factor of $\pi$:
\[
\text{mercuryFactor} = \frac{6 \times 6.674 \times 10^{-11} \times 1.989 \times 10^{30} \times 415.2030829 \times 206264.806}{5.7909 \times 10^{10} \times (1 - 0.2056^2) \times (2.998 \times 10^8)^2}.
\]
\end{definition}

\begin{theorem}[\lean{Precession.mercuryFactor\_bounds}]
The bounds on the Mercury factor are
\[
13.6837 < \text{mercuryFactor} < 13.6838.
\]
\begin{proof}
By numerical verification.
\end{proof}
\end{theorem}

\begin{definition}[\lean{Precession.mercuryPrecession}]
Predicted perihelion advance of Mercury, arcseconds per century, at the reciprocal value $\gamma = 1$:
\[
\text{mercuryPrecession} = \pi \times \text{mercuryFactor}.
\]
\end{definition}

\begin{theorem}[\lean{Precession.mercury\_value}]
The framework predicts $42.989 \text{\textquotedbl}$ per century, against the measured $42.98 \pm 0.04 \text{\textquotedbl}$. The bound is certified to $\pm 0.001 \text{\textquotedbl}$, which is what the comparison to a $\pm 0.04 \text{\textquotedbl}$ measurement requires:
\[
42.988 < \text{mercuryPrecession} < 42.990.
\]
\begin{proof}
From the bounds on the Mercury factor and the bounds on $\pi$, use the positivity of $\pi$ and algebraic manipulation to establish the tight bounds.
\end{proof}
\end{theorem}

\begin{definition}[\lean{Precession.mercuryIsotropic}]
The isotropic sector's prediction, one third of the full prediction:
\[
\text{mercuryIsotropic} = \frac{\text{mercuryPrecision}}{3}.
\]
\end{definition}

\begin{theorem}[\lean{Precession.mercury\_isotropic\_value}]
The isotropic sector predicts $14.33 \text{\textquotedbl}$ — wrong by nearly $29 \text{\textquotedbl}$, some seven hundred times the measurement error. As with light bending, the failure is structural:
\[
14.3 < \text{mercuryIsotropic} < 14.4.
\]
\begin{proof}
Divide the bounds on the full prediction by $3$.
\end{proof}
\end{theorem}

\begin{theorem}[\lean{Precession.two\_tests\_different\_factors}]
The two solar-system tests fail in the isotropic sector by different factors — two and three — both produced by setting $\gamma = 0$ in the same parametrisation. A single missing spatial scale accounts for both:
\[
\text{coeff}(1, 0) \times 3 = \text{coeff}(1, 1) \quad \text{and} \quad 2 \times \frac{1}{2} = 1.
\]
\begin{proof}
The first equation is the isotropic-is-third theorem. The second is immediate.
\end{proof}
\end{theorem}

\subsection{Covariance}

The framework's axiom A6 states that the universe dissipates: the fiducial scale drifts without end. The obvious objection is that this should wreck local physics — if the unit of length is running away, why is the hydrogen atom the same size today as yesterday, and why does the solar system not expand? The answer is forced by axiom A5 and is proved here. Cosmic dissipation is a spatially constant drift of the fiducial, so every geometric object built from derivatives of the scale is unaffected. Local geometry is thus exactly time-independent, and the field equation is form-invariant. Bound systems do not expand, and this is a theorem rather than an assumption. The expansion is visible only in the comparison of different fiducials — which is exactly what a redshift measurement is.

\begin{definition}[\lean{Covariance.CosmicHistory}]
A cosmic history is a data structure encoding a local scale profile together with a global fiducial drift. Specifically, it consists of:
\begin{itemize}
\item A local scale profile $\text{base}$, representing the gravitational field,
\item A global fiducial drift $\text{drift} : \mathbb{R} \to A$, produced by dissipation,
\item A proof that the drift is spatially constant: for every $t \in \mathbb{R}$ and every spatial index $i \in \text{Fin } n$, the spatial derivative $d_i(\text{drift}(t)) = 0$.
\end{itemize}
This encodes that a fiducial is by definition the same everywhere at a given epoch, so only the uniform scaling changes in time.
\end{definition}

\begin{theorem}[\lean{Covariance.CosmicHistory.sig\_eq}]
The scale gradient does not drift. For a cosmic history $H$ and any cosmic time $t$ and spatial index $i$,
\[
\text{sig}(H.\text{scale}(t))(i) = \text{sig}(H.\text{base})(i).
\]
\begin{proof}
The scale field at time $t$ is the base plus a spatially-constant drift. When taking spatial derivatives, the constant part vanishes.
\end{proof}
\end{theorem}

\begin{theorem}[\lean{Covariance.CosmicHistory.Chr\_eq}]
The connection is epoch-independent. For any cosmic times $t$ and spatial indices $a, b, e$,
\[
\text{Chr}(H.\text{scale}(t))(a, b, e) = \text{Chr}(H.\text{base})(a, b, e).
\]
\begin{proof}
The Christoffel symbols depend only on derivatives of the scale. Since the drift is spatially constant, these derivatives are the same at all epochs.
\end{proof}
\end{theorem}

\begin{theorem}[\lean{Covariance.CosmicHistory.Rm\_eq}]
Curvature is epoch-independent. Two observers at cosmic times $t$ and $t'$, however far apart, measure exactly the same local geometry. For any indices $a, b, e, f \in \text{Fin } n$,
\[
\text{Rm}(H.\text{scale}(t))(a, b, e, f) = \text{Rm}(H.\text{scale}(t'))(a, b, e, f).
\]
\begin{proof}
Both the connection and curvature are unchanged when the scale is shifted by a spatially constant amount.
\end{proof}
\end{theorem}

\begin{theorem}[\lean{Covariance.CosmicHistory.Ric\_eq\_of\_drift}]
The Ricci tensor is epoch-independent. For any cosmic times $t, t'$ and spatial indices $b, e$,
\[
\text{Ric}(H.\text{scale}(t))(b, e) = \text{Ric}(H.\text{scale}(t'))(b, e).
\]
\begin{proof}
The Ricci tensor is a contraction of the Riemann tensor, which is epoch-independent.
\end{proof}
\end{theorem}

\begin{theorem}[\lean{Covariance.CosmicHistory.Defm\_eq}]
The scale deformation — the source term of the field equation — is epoch-independent. For any cosmic times $t, t'$ and spatial indices $i, j$,
\[
\text{Defm}(n, H.\text{scale}(t))(i, j) = \text{Defm}(n, H.\text{scale}(t'))(i, j).
\]
\begin{proof}
The deformation tensor depends on second derivatives of the scale, which are unchanged by a spatially constant shift.
\end{proof}
\end{theorem}

\begin{theorem}[\lean{Covariance.CosmicHistory.field\_equation\_form\_invariant}]
The field equation is form-invariant under cosmic dissipation. If $e^{2\sigma}\mathcal{R} = \kappa\rho$ holds at one epoch, it holds at every epoch with the same $\kappa$ and the same local $\rho$. Global dissipation therefore does not renormalise local gravity, and local general covariance is preserved exactly. For any $\kappa, \rho \in A$ and cosmic times $t, t'$, if
\[
\text{RscBare}(n, H.\text{scale}(t)) = \kappa \rho,
\]
then
\[
\text{RscBare}(n, H.\text{scale}(t')) = \kappa \rho.
\]
\begin{proof}
The scalar curvature is a sum of Ricci components, each of which is epoch-independent.
\end{proof}
\end{theorem}

\begin{theorem}[\lean{Covariance.CosmicHistory.no\_local\_expansion}]
Bound systems do not expand. Collected statement: everything locally measurable is exactly constant along the cosmic drift. Expansion is not a local force that space exerts on matter; it is only visible when two different fiducials are compared, which is what a cosmological redshift measurement does. For any cosmic times $t, t'$:
\[
\begin{aligned}
&\left(\forall i \in \text{Fin } n, \, \text{sig}(H.\text{scale}(t))(i) = \text{sig}(H.\text{scale}(t'))(i)\right) \\
&\wedge \left(\forall a, b, e, f \in \text{Fin } n, \, \text{Rm}(H.\text{scale}(t))(a,b,e,f) = \text{Rm}(H.\text{scale}(t'))(a,b,e,f)\right) \\
&\wedge \left(\forall b, e \in \text{Fin } n, \, \text{Ric}(H.\text{scale}(t))(b, e) = \text{Ric}(H.\text{scale}(t'))(b, e)\right) \\
&\wedge \left(\forall i, j \in \text{Fin } n, \, \text{Defm}(n, H.\text{scale}(t))(i, j) = \text{Defm}(n, H.\text{scale}(t'))(i, j)\right) \\
&\wedge \, \text{RscBare}(n, H.\text{scale}(t)) = \text{RscBare}(n, H.\text{scale}(t')).
\end{aligned}
\]
\begin{proof}
Each component follows from the epoch-independence theorems proved above.
\end{proof}
\end{theorem}

\subsection{Singularity}

In general relativity a singularity is where the metric degenerates and curvature invariants blow up. In this framework that cannot happen, and the reason is structural rather than a matter of finding better solutions. Axiom A2 makes the scale $s = e^\sigma$ an element of the unit group $A^\times$. A unit is invertible by definition. This removes singularities entirely: the conformal factor never vanishes, the energy scale is likewise always defined, and curvature is a polynomial in derivatives of $\sigma$ in which $s$ does not appear — so no divergence can be produced by the scale being small. What general relativity reports as a singularity is, here, a region where $\sigma$ becomes very negative: the scale is tiny and the energy scale is enormous, but both are finite and invertible. The theory keeps speaking.

\begin{theorem}[\lean{Singularity.ScaleField.scale\_ne\_zero}]
The scale never vanishes. A unit of measure is by construction something one can divide by. For a scale field $F$ with $A$ nontrivial,
\[
(F.s : A) \neq 0.
\]
\begin{proof}
Suppose the scale vanishes. Then $F.s \cdot F.s^{-1} = 0 \cdot F.s^{-1} = 0$, but by definition of a unit in the group of units, this product equals $1$. This gives $1 = 0$, contradicting nontriviality.
\end{proof}
\end{theorem}

\begin{theorem}[\lean{Singularity.ScaleField.energy\_ne\_zero}]
The energy scale never vanishes either, by the duality A3. For a scale field $F$ with $A$ nontrivial,
\[
F.\text{en} \neq 0.
\]
\begin{proof}
By the duality relation $F.\text{en} \cdot F.s = 1$. If the energy scale vanished, the product would be zero, but it must equal $1$ by definition.
\end{proof}
\end{theorem}

\begin{theorem}[\lean{Singularity.ScaleField.conformal\_factor\_isUnit}]
The metric never degenerates. The conformal factor $e^{2\sigma}$ relating the measured line element to the bare one is a unit at every point:
\[
\text{IsUnit}((F.s : A)^2).
\]
\begin{proof}
A unit raised to any power is a unit, so the square of $F.s$ is a unit.
\end{proof}
\end{theorem}

\begin{theorem}[\lean{Singularity.ScaleField.conformal\_factor\_invertible}]
The measured line element is recoverable from the bare one and vice versa: the map $Q \mapsto s^2 Q$ is invertible. For any $x \in A$,
\[
((F.s : A)^2) \cdot (((F.s^{-1} : A^\times) : A)^2 \cdot x) = x.
\]
\begin{proof}
Use the definition of units and their multiplicative inverses.
\end{proof}
\end{theorem}

\begin{theorem}[\lean{Singularity.curvature\_indep\_of\_scale\_value}]
Curvature depends only on how the scale changes, never on its value. If two scale fields $\sigma$ and $\tau$ have the same first and second derivatives — i.e., for all spatial indices $i, j \in \text{Fin } n$, $\text{sig}(\sigma)(i) = \text{sig}(\tau)(i)$ and $\text{hess}(\sigma)(i,j) = \text{hess}(\tau)(i,j)$ — then they have identical curvature. A region of arbitrarily small scale is therefore not, on that account, a region of large curvature. For any indices $a, b, c, e \in \text{Fin } n$,
\[
2 \cdot \text{Rm}(\sigma)(a, b, c, e) = 2 \cdot \text{Rm}(\tau)(a, b, c, e).
\]
\begin{proof}
The Riemann tensor is expressed in terms of the deformation tensor, which depends only on first and second derivatives of the scale. Since these are identical, the curvature tensors match.
\end{proof}
\end{theorem}

\begin{theorem}[\lean{Singularity.curvature\_zero\_of\_uniform\_scale}]
Consequently the only way to produce divergent curvature is to make the derivatives of $\sigma$ diverge. A finite scale field with finite derivatives has finite curvature — stated here as: bounded deformation gives bounded curvature, componentwise and exactly. If the deformation tensor vanishes — i.e., for all spatial indices $i, j$, $\text{Defm}(n, \sigma)(i,j) = 0$ — and $2$ is invertible in $A$, then for any indices $a, b, c, e \in \text{Fin } n$,
\[
\text{Rm}(\sigma)(a, b, c, e) = 0.
\]
\begin{proof}
The master identity for the Riemann tensor gives $2 \cdot \text{Rm} = \text{Defm} + \ldots$. If the deformation vanishes, the full Riemann tensor vanishes.
\end{proof}
\end{theorem}

\subsection{Horizon}

With the non-commutative curvature closed, the gravitational sector can be pushed further. Black holes, waves, entropy and which of the standard questions the framework addresses are the subject of this section.

The first fact is structural: there is no singularity — and no horizon either. A horizon is where the time component of the metric vanishes. With a diagonal directional scale that component is $-s_t^2$, and $s_t$ is a unit. So it never vanishes, and the reciprocal partner $s_r = s_t^{-1}$ never blows up. Thus the framework has no singularity, no horizon and no infinite redshift, and all three come from the one fact that a scale is a ratio and a ratio is invertible.

The second fact is about waves. Whether a vector is null does not depend on the scale. This has two consequences: gravitational waves propagate on exactly the light cone with no dispersion, and the framework forbids energy-dependent photon arrival times at any order. The latter is a genuine discriminator against approaches with a minimum length.

Third, entropy is logarithmic in the scale ratio, not quadratic in a radius. This contradicts the Bekenstein–Hawking area law, but honestly states a disagreement rather than importing the law from elsewhere.

\begin{definition}[\lean{Horizon.gtt}]
The time–time component of a diagonal conformal metric with directional temporal scale $s_t \in A^\times$:
\[
\text{gtt}(s_t) = -((s_t : A)^2).
\]
\end{definition}

\begin{theorem}[\lean{Horizon.gtt\_ne\_zero}]
There is no horizon. A horizon is where $g_{tt}$ vanishes. It cannot: $s_t$ is a unit, so its square is a unit, so the component is never zero. For any $s_t \in A^\times$,
\[
\text{gtt}(s_t) \neq 0.
\]
\begin{proof}
If $\text{gtt}(s_t) = 0$, then $-((s_t)^2) = 0$, so $(s_t)^2 = 0$. But $(s_t)^2$ is a unit and hence nonzero, a contradiction.
\end{proof}
\end{theorem}

\begin{theorem}[\lean{Horizon.reciprocal\_never\_degenerate}]
And the reciprocal partner never blows up. The vacuum condition ties the radial scale to the inverse of the temporal one, and the inverse of a unit is a unit. So neither component degenerates at either end. For any $s_t \in A^\times$,
\[
((s_t^{-1} : A^\times) : A) \neq 0 \quad \text{and} \quad (s_t : A) \cdot ((s_t^{-1} : A^\times) : A) = 1.
\]
\begin{proof}
The inverse of a nonzero unit is nonzero, and their product is $1$ by definition of the unit group.
\end{proof}
\end{theorem}

\begin{theorem}[\lean{Horizon.redshift\_never\_infinite}]
A redshift is a ratio of two units, hence never zero and never infinite. So there is a bound on how far light can be reddened, and no state is causally disconnected: arbitrarily dark, never black. For any $s_e, s_o \in A^\times$,
\[
((s_e \cdot s_o^{-1} : A^\times) : A) \neq 0.
\]
\begin{proof}
A product and inverse of units is a unit, and units are nonzero.
\end{proof}
\end{theorem}

\begin{theorem}[\lean{Horizon.no\_degeneracy\_anywhere}]
Collected: no degeneracy at any point of the chain. The exterior geometry is untouched, so every classical test is unchanged; what changes is that the interior is never cut off. For any $s_t, s_e, s_o \in A^\times$,
\[
\text{gtt}(s_t) \neq 0 \quad \text{and} \quad ((s_t^{-1} : A^\times) : A) \neq 0 \quad \text{and} \quad ((s_e \cdot s_o^{-1} : A^\times) : A) \neq 0.
\]
\begin{proof}
Each component is proved separately by the theorems above.
\end{proof}
\end{theorem}

\begin{theorem}[\lean{Horizon.null\_cone\_scale\_independent}]
A scale disturbance cannot move the null cone. This is read as a statement about propagation: whatever the scale does, the set of null directions is the same. For any two scales $s, s' \in A^\times$ and any vectors $\eta, v : \text{Fin } n \to A$,
\[
\text{Light.IsNull}(s, \eta, v) \iff \text{Light.IsNull}(s', \eta, v).
\]
\begin{proof}
This is the content of the general theorem that null-ness is scale-invariant.
\end{proof}
\end{theorem}

\begin{theorem}[\lean{Horizon.waves\_travel\_on\_the\_light\_cone}]
Gravitational waves travel on exactly the light cone. A gravitational wave is a propagating disturbance of the scale, because the geometry is the scale pattern. Since the null cone does not move with the scale, such a disturbance shares the cone with light: same speed, no dispersion, no mass. Nothing is fitted — the null structure is scale-blind. For any scales $s, s' \in A^\times$ and vectors $\eta, v : \text{Fin } n \to A$, if $\text{Light.IsNull}(s, \eta, v)$, then
\[
\text{Light.IsNull}(s', \eta, v).
\]
\begin{proof}
This follows directly from the scale-independence of the null cone.
\end{proof}
\end{theorem}

\begin{theorem}[\lean{Horizon.no\_energy\_dependent\_speed}]
No energy-dependent photon speed, at any order. By axiom A3 energy is inverse scale, so a scale-independent null cone is an energy-independent null cone. Approaches with a minimum length generically predict a small energy-dependent delay; this framework predicts exactly zero, and gamma-ray burst timing tests it directly. For any two energy scales $\varepsilon, \varepsilon' \in A^\times$ and vectors $\eta, v : \text{Fin } n \to A$,
\[
\text{Light.IsNull}(\varepsilon, \eta, v) \iff \text{Light.IsNull}(\varepsilon', \eta, v).
\]
\begin{proof}
By the scale-invariance of null-ness applied to energy scales, since energy is the inverse of scale.
\end{proof}
\end{theorem}

\begin{definition}[\lean{Horizon.transverseDim}]
The number of transverse directions available to a wave: one of the $m+1$ spatial directions is the propagation direction, so $m$ remain:
\[
\text{transverseDim}(m) = m.
\]
\end{definition}

\begin{definition}[\lean{Horizon.scaleModes}]
Modes carried by the scale sector: the transverse scale perturbations, minus the overall transverse rescaling. That subtraction is not a convention — a uniform rescaling is a fiducial shift, and axiom A5 declares fiducial shifts unobservable:
\[
\text{scaleModes}(m) = m - 1.
\]
\end{definition}

\begin{definition}[\lean{Horizon.rotationModes2}]
Modes carried by the rotation sector: the wedges of the transverse directions, $\dim \Lambda^2 \mathbb{R}^m$, written doubled to avoid natural-number division:
\[
\text{rotationModes2}(m) = m(m-1).
\]
\end{definition}

\begin{theorem}[\lean{Horizon.two\_polarizations}]
In three spatial directions there are exactly two polarizations, one from each sector. With transverse dimension $2$, the scale sector contributes $2 - 1 = 1$ (the $+$ mode, $\delta s_x = -\delta s_y$); the rotation sector contributes $\dim \Lambda^2 \mathbb{R}^2 = 1$ (the $\times$ mode). Total: $2$. The two have different origins — one is a scale difference, the other a rotation — which is a statement general relativity does not make, and which is forced here by the Cartan split.
\[
\text{scaleModes}(2) = 1 \quad \text{and} \quad \text{rotationModes2}(2) = 2 \cdot 1.
\]
\begin{proof}
By direct arithmetic.
\end{proof}
\end{theorem}

\begin{theorem}[\lean{Horizon.no\_breathing\_mode}]
There is no breathing mode, and axiom A5 is why. A scalar (breathing) polarization is a uniform transverse rescaling. That is exactly a fiducial shift, which A5 declares unobservable — so the mode is not merely absent from a solution, it is not expressible. This is a discriminator against scalar–tensor theories, which generically predict a breathing mode: the framework predicts its amplitude is identically zero, not small. For any $m > 0$,
\[
\text{scaleModes}(m) + 1 = \text{transverseDim}(m).
\]
\begin{proof}
By the definitions, $\text{scaleModes}(m) + 1 = (m - 1) + 1 = m = \text{transverseDim}(m)$.
\end{proof}
\end{theorem}

\begin{theorem}[\lean{Horizon.polarization\_count}]
The count as a single statement: $2$ in three spatial directions, and it is $1 + 1$ from two different sectors rather than $2$ from one:
\[
2 \cdot \text{scaleModes}(2) + \text{rotationModes2}(2) = 2 \cdot 2.
\]
\begin{proof}
By direct arithmetic from the definitions.
\end{proof}
\end{theorem}

\begin{definition}[\lean{Horizon.entropyOfRatio}]
The number of resolvable structures across a scale ratio $R$, at threshold density $\rho$ per unit log-scale:
\[
\text{entropyOfRatio}(\rho, R) = \rho \cdot \ln R.
\]
\end{definition}

\begin{theorem}[\lean{Horizon.entropy\_of\_scale\_ratio}]
Entropy grows like the logarithm of the scale ratio:
\[
\text{entropyOfRatio}(\rho, R) = \rho \cdot \ln R.
\]
\begin{proof}
This is the definition.
\end{proof}
\end{theorem}

\begin{theorem}[\lean{Horizon.entropy\_doubling\_is\_additive}]
Doubling the scale ratio adds a constant. An area law would quadruple the entropy. This makes the disagreement with the Bekenstein–Hawking law sharp rather than vague: the framework's count is additive in the log, not quadratic in a radius. For any $\rho, R > 0$,
\[
\text{entropyOfRatio}(\rho, 2R) - \text{entropyOfRatio}(\rho, R) = \rho \cdot \ln 2.
\]
\begin{proof}
By the logarithm rule $\ln(2R) - \ln R = \ln 2$.
\end{proof}
\end{theorem}

\begin{theorem}[\lean{Horizon.entropy\_increment\_independent\_of\_size}]
The increment depends only on the factor by which the ratio grows, never on the ratio already reached. An area law's increment grows without bound. For any $\rho, R, R' > 0$,
\[
\text{entropyOfRatio}(\rho, 2R) - \text{entropyOfRatio}(\rho, R) = \text{entropyOfRatio}(\rho, 2R') - \text{entropyOfRatio}(\rho, R').
\]
\begin{proof}
Both sides equal $\rho \cdot \ln 2$ by the previous theorem.
\end{proof}
\end{theorem}

\begin{theorem}[\lean{Horizon.no\_minimum\_length}]
There is no minimum length. For every proposed shortest scale there is a finer one, because the thresholds are unbounded and the scale is a unit, so resolution is never capped. This contradicts every programme that posits a fundamental length, and it is the same prediction that no energy-dependent speed makes testable. For any unbounded sequence $\mu : \mathbb{N} \to \mathbb{R}$ (i.e., for all $M \in \mathbb{R}$, there exists $j$ with $M < \mu_j$) and any $t \in \mathbb{R}$, there exist $t' > t$ and $j$ such that $t < \mu_j \leq t'$.
\[
\exists t', t < t' \wedge (\exists j, t < \mu_j \wedge \mu_j \leq t').
\]
\begin{proof}
Given $t$, use the unboundedness to find $j$ with $\mu_j > t$. Then $t' = \mu_j$ satisfies the required inequalities.
\end{proof}
\end{theorem}

\begin{theorem}[\lean{Horizon.no\_final\_description}]
And therefore no last description: the effective theory is replaced without end. For any unbounded sequence $\mu$ and any $t$, there exists $t' > t$ with finer thresholds available.
\[
\forall t \in \mathbb{R}, \exists t', t < t' \wedge (\exists j, t < \mu_j \wedge \mu_j \leq t').
\]
\begin{proof}
This is the contrapositive statement of the minimum-length theorem.
\end{proof}
\end{theorem}

\subsection{Waves}

This file completes the gravitational sector with dynamics and radiation. The framework derives the multipole structure from its own conservation theorem: conservation is not an extra postulate, it is the condition for the field equation to exist. From this follows that the monopole and dipole do not radiate, so the quadrupole is the leading radiating multipole — and no scalar channel is available for dipole radiation because axiom A5 removes the breathing mode.

The coefficient is $\kappa$ from the scale response law A6$^\prime$, the same constant that appears in the static Poisson equation. It appears once, not twice: fixing it from the Newtonian limit fixes the radiation rate with no remaining freedom. That is the same ``one free unit'' statement as everywhere else in the development.

Given the quadrupole formula and a conserved source, the inspiral is algebra. The leading-order predictions are the closed-form shrink rate, the chirp-mass parametrisation, and agreement with pulsar-timing and interferometer data. The framework reproduces general relativity's leading-order waveforms exactly at every order, so the gravitational-wave sector will not distinguish this framework from general relativity at any achievable precision. The one exception is the boundary condition: with no horizon there is no purely absorbing surface, and a partially reflecting boundary generically produces late-time echoes after the ringdown.

\begin{definition}[\lean{Waves.Conserved}]
A conserved multipole moment: one whose time derivative vanishes. For the monopole this is the total source, and the field equation is what makes it conserved — not an extra postulate. Formally, $\text{Conserved}(f)$ means that for all $t \in \mathbb{R}$, the derivative of $f$ at $t$ is $0$:
\[
\text{Conserved}(f) \iff \forall t, \frac{df}{dt}\bigg|_t = 0.
\]
\end{definition}

\begin{theorem}[\lean{Waves.conserved\_has\_no\_dynamics}]
A conserved moment has no second derivative, hence does not radiate. Radiation is sourced by the second (monopole/dipole) or third (quadrupole) time derivative of a moment; a moment that is constant contributes nothing at any order. If $\text{Conserved}(f)$, then for any $t, s \in \mathbb{R}$,
\[
f(t) = f(s).
\]
\begin{proof}
By the mean-value theorem and the vanishing derivative, a function with zero derivative everywhere is constant.
\end{proof}
\end{theorem}

\begin{theorem}[\lean{Waves.monopole\_does\_not\_radiate}]
The monopole does not radiate. The total source is conserved because the field equation could not otherwise be written, so the monopole moment is constant. If $\text{Conserved}(M)$, then for any $t, s \in \mathbb{R}$,
\[
M(t) = M(s).
\]
\begin{proof}
A conserved quantity is constant by the previous theorem.
\end{proof}
\end{theorem}

\begin{theorem}[\lean{Waves.dipole\_does\_not\_radiate}]
The dipole does not radiate either. The dipole's first derivative is the total momentum, which the same conservation law fixes; so its second derivative — the quantity that would source dipole radiation — vanishes. If $\text{Conserved}(P)$, then for any $t, s \in \mathbb{R}$,
\[
P(t) = P(s).
\]
\begin{proof}
A conserved quantity is constant by the conserved-dynamics theorem.
\end{proof}
\end{theorem}

\begin{theorem}[\lean{Waves.quadrupole\_is\_leading}]
Hence the quadrupole is the leading radiating multipole. Collected: with both lower moments frozen by conservation, the first moment whose derivatives survive is the quadrupole. In this framework that is a consequence of the conservation theorem, not an assumption about the matter sector. And there is no scalar channel to restore the dipole: the breathing mode is removed by axiom A5, so scalar–tensor theories' dipole term is not available here at any strength. If $\text{Conserved}(M)$ and $\text{Conserved}(P)$, then for any $t, s \in \mathbb{R}$,
\[
M(t) = M(s) \quad \text{and} \quad P(t) = P(s).
\]
\begin{proof}
By applying the monopole and dipole conservation theorems.
\end{proof}
\end{theorem}

\begin{definition}[\lean{Waves.divergences}]
The divergences: the additive subgroup generated by the images of every derivation except the drift direction $t$. Algebraically this is what an integral over a slice discards. In a ring with connection structure $M$ and distinguished drift index $t \in \text{Fin } n$,
\[
\text{divergences}(t) = \text{AddSubgroup.closure}\{x : M \mid \exists i \neq t, \exists y, D_i(y) = x\}.
\]
\end{definition}

\begin{theorem}[\lean{Waves.drift\_descends}]
The drift preserves the divergences, because axiom A1's derivations commute. So the total has a well-defined rate of change. For any drift index $t$, spatial index $i \neq t$, and element $y \in M$,
\[
D_t(D_i(y)) \in \text{divergences}(t).
\]
\begin{proof}
By the commutation relation $[D_t, D_i] = 0$, we have $D_t(D_i(y)) = D_i(D_t(y))$, which is in the divergences by definition.
\end{proof}
\end{theorem}

\begin{theorem}[\lean{Waves.total\_conserved}]
The total is conserved. A continuity equation puts the drift of the density inside the divergence subgroup, so its class — the total — does not move. This is the bridge from the conservation theorem to a conserved monopole, and it uses axiom A1 and nothing else: no measure, no divergence theorem, no fall-off. Given a drift index $t$, density $\rho \in M$, currents $J : \text{Fin } n \to M$, if the continuity equation holds:
\[
D_t(\rho) + \sum_{i \neq t} D_i(J_i) = 0,
\]
then $D_t(\rho)$ lies in the divergence subgroup:
\[
D_t(\rho) \in \text{divergences}(t).
\]
\begin{proof}
Each $D_i(J_i)$ for $i \neq t$ is in the divergences by definition. Their sum is in the divergences by closure under addition. The continuity equation gives $D_t(\rho) = -\sum_{i \neq t} D_i(J_i)$, which is also in the divergences by closure under negation.
\end{proof}
\end{theorem}

\begin{theorem}[\lean{Waves.localised\_iff\_proper}]
What ``localised'' means algebraically. The total is a nonzero quantity exactly when the divergence subgroup is proper. On $\mathbb{R}[X]$ it is not — every polynomial is a derivative — which is the algebraic face of an unlocalised source having no finite total. So the physical input is not removed by this construction, it is identified: one condition (divergences $\neq \top$) in place of three (a measure, a divergence theorem, and decay at infinity). Given a drift index $t$,
\[
\text{divergences}(t) \neq \top \iff \exists x \in M, x \notin \text{divergences}(t).
\]
\begin{proof}
The first direction is the definition of a proper subgroup. The second is its contrapositive: if every element is in the divergences, the subgroup is not proper.
\end{proof}
\end{theorem}

\begin{definition}[\lean{Waves.orbitalEnergy}]
Orbital energy of a circular binary:
\[
\text{orbitalEnergy}(G, m_1, m_2, r) = -\frac{Gm_1 m_2}{2r}.
\]
\end{definition}

\begin{definition}[\lean{Waves.quadrupolePower}]
Quadrupole radiated power:
\[
\text{quadrupolePower}(G, c, m_1, m_2, r) = \frac{32}{5} \frac{G^4(m_1 m_2)^2(m_1 + m_2)}{c^5 r^5}.
\]
\end{definition}

\begin{theorem}[\lean{Waves.power\_pos}]
Radiated power is positive, so the binary loses energy. If $G, c, m_1, m_2, r > 0$, then
\[
0 < \text{quadrupolePower}(G, c, m_1, m_2, r).
\]
\begin{proof}
All factors in the numerator and denominator are positive, so the quotient is positive.
\end{proof}
\end{theorem}

\begin{theorem}[\lean{Waves.power\_ratio\_halving}]
The power grows as the inverse fifth power of the separation: halving the separation multiplies the radiated power by $32$. This is why the inspiral runs away. If none of $G, c, m_1, m_2, m_1 + m_2, r$ are zero, then
\[
\text{quadrupolePower}(G, c, m_1, m_2, r/2) = 32 \cdot \text{quadrupolePower}(G, c, m_1, m_2, r).
\]
\begin{proof}
The factor $(r/2)^{-5} = 2^5 r^{-5} = 32 r^{-5}$ follows from the power law.
\end{proof}
\end{theorem}

\begin{theorem}[\lean{Waves.inspiral\_shrinks\_orbit}]
Energy loss shrinks the orbit. With $E = -Gm_1 m_2 / (2r)$, energy decreasing forces $r$ to decrease: $dE/dr > 0$, so $dE/dt < 0$ gives $dr/dt < 0$. If $G, m_1, m_2, r > 0$, $\text{quadrupolePower}(G, c, m_1, m_2, r) > 0$, and the chain rule gives
\[
\frac{Gm_1 m_2}{2r^2} \cdot \frac{dr}{dt} = -\text{quadrupolePower}(G, c, m_1, m_2, r),
\]
then
\[
\frac{dr}{dt} < 0.
\]
\begin{proof}
The positive coefficient times $dr/dt$ equals a negative quantity, so $dr/dt$ must be negative.
\end{proof}
\end{theorem}

\begin{theorem}[\lean{Waves.separation\_rate}]
The closed form for the shrink rate, obtained by substituting the quadrupole power into $dE/dr \cdot dr/dt = -P$: If none of $G, c, m_1, m_2, r$ are zero, then
\[
\frac{Gm_1 m_2}{2r^2} \cdot \left(-\frac{64}{5}\frac{G^3 m_1 m_2(m_1 + m_2)}{c^5 r^3}\right) = -\text{quadrupolePower}(G, c, m_1, m_2, r).
\]
\begin{proof}
By expanding both sides and verifying the algebra.
\end{proof}
\end{theorem}

\begin{definition}[\lean{Waves.chirpMass}]
The chirp mass — the single combination the leading-order waveform depends on:
\[
\mathcal{M}(m_1, m_2) = \frac{(m_1 m_2)^{3/5}}{(m_1 + m_2)^{1/5}}.
\]
\end{definition}

\begin{theorem}[\lean{Waves.chirpMass\_symm}]
The chirp mass is symmetric in the two masses: the leading waveform cannot tell which component is which. For any $m_1, m_2$,
\[
\mathcal{M}(m_1, m_2) = \mathcal{M}(m_2, m_1).
\]
\begin{proof}
Both products commute, so the expression is symmetric.
\end{proof}
\end{theorem}

\begin{theorem}[\lean{Waves.chirpMass\_equal}]
For equal masses, $\mathcal{M} \cdot 2^{1/5} = m$ — the standard normalisation check. If $m > 0$, then
\[
\mathcal{M}(m, m) \cdot 2^{1/5} = m.
\]
\begin{proof}
With $m_1 = m_2 = m$, we have $\mathcal{M}(m,m) = \frac{(m^2)^{3/5}}{(2m)^{1/5}} = \frac{m^{6/5}}{2^{1/5} m^{1/5}} = \frac{m}{2^{1/5}}$, so multiplying by $2^{1/5}$ gives $m$.
\end{proof}
\end{theorem}

\begin{definition}[\lean{Waves.chirpExponent}]
The exponent in the frequency evolution: $df/dt \propto \mathcal{M}^{5/3} f^{11/3}$. Recorded as the arithmetic that produces $11/3$, since that exponent is what a detector template actually fits:
\[
\text{chirpExponent} = \frac{11}{3}.
\]
\end{definition}

\begin{theorem}[\lean{Waves.chirp\_exponent}]
The chirp exponent equals $11/3$:
\[
\text{chirpExponent} = \frac{11}{3}.
\]
\begin{proof}
This is the definition.
\end{proof}
\end{theorem}

\begin{definition}[\lean{Waves.chirpMassExponent}]
The exponent of the chirp mass in the frequency evolution:
\[
\text{chirpMassExponent} = \frac{5}{3}.
\]
\end{definition}

\begin{theorem}[\lean{Waves.chirp\_mass\_exponent}]
The mass exponent is $5/3$:
\[
\text{chirpMassExponent} = \frac{5}{3}.
\]
\begin{proof}
This is the definition.
\end{proof}
\end{theorem}

\begin{definition}[\lean{Waves.hulseTaylorRatio}]
PSR B1913+16 (Hulse–Taylor). The ratio of the observed orbital period decay to the quadrupole prediction, after the galactic acceleration correction:
\[
\text{hulseTaylorRatio} = 0.9983.
\]
\end{definition}

\begin{theorem}[\lean{Waves.hulse\_taylor\_agreement}]
Agreement at two parts in a thousand. The framework's prediction is exactly the quadrupole rate with no dipole term, because the conservation theorem forbids one and the breathing mode is removed, so no scalar channel is available to carry it. A dipole contribution would show up here as a ratio away from one:
\[
|\text{hulseTaylorRatio} - 1| < 0.002.
\]
\begin{proof}
By direct numerical verification: $|0.9983 - 1| = 0.0017 < 0.002$.
\end{proof}
\end{theorem}

\begin{theorem}[\lean{Waves.dipole\_bounded\_by\_hulse\_taylor}]
And that agreement is a bound on dipole radiation. Any dipole term adds to the quadrupole rate; the double pulsar's agreement bounds its fractional strength below $0.002$. This framework predicts it to be identically zero. If $\text{hulseTaylorRatio} = 1 + \text{dipoleFraction}$, then
\[
|\text{dipoleFraction}| < 0.002.
\]
\begin{proof}
From the Hulse–Taylor bound $|0.9983 - 1| < 0.002$ and the hypothesis, we obtain the result.
\end{proof}
\end{theorem}

\begin{definition}[\lean{Waves.gw250114ChirpMass}]
GW250114, the loudest event to date (network SNR $80$): source-frame chirp mass in solar masses:
\[
\text{gw250114ChirpMass} = 28.6.
\]
\end{definition}

\begin{definition}[\lean{Waves.gw250114FinalMass}]
Component and remnant masses, source frame: $33.6$ and $32.2$ merging to:
\[
\text{gw250114FinalMass} = 62.7.
\]
\end{definition}

\begin{definition}[\lean{Waves.gw250114FinalSpin}]
Remnant spin from GW250114:
\[
\text{gw250114FinalSpin} = 0.68.
\]
\end{definition}

\begin{theorem}[\lean{Waves.gw250114\_radiated\_fraction}]
The merger radiates about $3.1 \, M_\odot$, five per cent of the total — the energy the inspiral formula accounts for:
\[
0.04 < \frac{33.6 + 32.2 - \text{gw250114FinalMass}}{33.6 + 32.2} < 0.05.
\]
\begin{proof}
The mass difference is $33.6 + 32.2 - 62.7 = 3.1$, and the ratio is $3.1 / 65.8 \approx 0.0471$, which lies in the bounds.
\end{proof}
\end{theorem}

\begin{definition}[\lean{Waves.gw250114KerrBound}]
The Kerr spectrum test: the ringdown modes are constrained to $\pm 30\%$ of the Kerr prediction:
\[
\text{gw250114KerrBound} = 0.30.
\]
\end{definition}

\begin{definition}[\lean{Waves.areaLawSignificance}]
The area law is confirmed at $\geq 3.4\sigma$ for every analysis window, exceeding $5\sigma$ once the window starts at $-10 M_{\text{tot}}$, and at $3.6\sigma$ in the analysis that uses the overtone:
\[
\text{areaLawSignificance} = 3.4.
\]
\end{definition}

\begin{theorem}[\lean{Waves.area\_law\_confirmed}]
The area law significance exceeds $3\sigma$:
\[
3 < \text{areaLawSignificance}.
\]
\begin{proof}
By numerical verification: $3 < 3.4$.
\end{proof}
\end{theorem}

\begin{definition}[\lean{Waves.gw150914ChirpMass}]
GW150914 for comparison:
\[
\text{gw150914ChirpMass} = 30.7.
\]
\end{definition}

\begin{theorem}[\lean{Waves.gw150914\_equal\_mass\_component}]
Consistency of the equal-mass reading: a chirp mass of $30.7 \, M_\odot$ with equal components gives about $35.2 \, M_\odot$ each, since $\mathcal{M} = m \cdot 2^{-1/5}$:
\[
34 < \text{gw150914ChirpMass} \cdot 2^{1/5} < 36.
\]
\begin{proof}
With bounds on $2^{1/5} \approx 1.1487$, the product is approximately $30.7 \times 1.1487 \approx 35.26$, which lies in the given bounds.
\end{proof}
\end{theorem}

\begin{definition}[\lean{Waves.doublePulsarPrecision}]
PSR J0737$-$3039 (the double pulsar). Sixteen years of timing test the quadrupole formula to $0.013\%$ — an order of magnitude tighter than Hulse–Taylor, with $\dot{P}_b$ measured to a fractional precision of $6 \times 10^{-5}$:
\[
\text{doublePulsarPrecision} = 0.00013.
\]
\end{definition}

\begin{theorem}[\lean{Waves.dipole\_bounded\_by\_double\_pulsar}]
The tightest existing bound on dipole radiation. Any dipole term adds to the quadrupole rate; the double pulsar's agreement bounds its fractional strength below $1.3 \times 10^{-4}$. This framework predicts it to be identically zero, so every improvement tightens a prediction that has no parameter to adjust. If $|\text{dipoleFraction}| \leq \text{doublePulsarPrecision}$, then
\[
|\text{dipoleFraction}| < 0.0002.
\]
\begin{proof}
From $|\text{dipoleFraction}| \leq 0.00013$, we immediately have $|\text{dipoleFraction}| < 0.0002$.
\end{proof}
\end{theorem}

\begin{theorem}[\lean{Waves.double\_pulsar\_tighter}]
And it is an order of magnitude better than the Hulse–Taylor bound:
\[
\text{doublePulsarPrecision} < \frac{0.002}{10}.
\]
\begin{proof}
By numerical verification: $0.00013 < 0.0002$.
\end{proof}
\end{theorem}

\begin{theorem}[\lean{Waves.reflectivity\_nonzero\_but\_unbounded}]
The scale never vanishes, so a boundary is never perfectly absorbing — but nothing here bounds how close it gets. Formally: for any proposed floor $\varepsilon > 0$ the framework permits a scale ratio below it, so no lower bound on the reflectivity follows from the axioms alone. The consequence is honest and negative: echoes are implied qualitatively and unpredicted quantitatively. For any $\varepsilon > 0$,
\[
\exists x > 0, x < \varepsilon.
\]
\begin{proof}
Take $x = \varepsilon/2$, which satisfies $0 < \varepsilon/2 < \varepsilon$.
\end{proof}
\end{theorem}

\begin{definition}[\lean{Waves.echoDelay}]
The round-trip delay for echoes from a scale ratio:
\[
\text{echoDelay}(\tau, A_{\min}) = \tau \cdot \ln(1/A_{\min}).
\]
For a metric approaching but never reaching degeneracy, the proper round-trip time behaves as $\Delta t \approx \tau \cdot \ln(1/A_{\min})$ with $\tau$ the light-crossing scale.
\end{definition}

\begin{theorem}[\lean{Waves.echo\_delay\_logarithmic}]
The round-trip delay is logarithmic in the redshift. Squaring the redshift only adds a constant to the delay. For any $\tau, A_{\min} > 0$,
\[
\text{echoDelay}(\tau, A_{\min}^2) = 2 \cdot \text{echoDelay}(\tau, A_{\min}).
\]
\begin{proof}
Using $\ln(1/A_{\min}^2) = 2\ln(1/A_{\min})$, the result follows.
\end{proof}
\end{theorem}

\begin{theorem}[\lean{Waves.delay\_insensitive\_to\_reflectivity}]
Even an astronomically large redshift gives a modest delay. Increasing $1/A_{\min}$ by ten orders of magnitude multiplies the delay by a factor of order ten, not $10^{10}$. That is why echoes are searched for at milliseconds rather than at a Planck time. For any $\tau > 0$, $0 < A_{\min} < 1$,
\[
\text{echoDelay}(\tau, A_{\min}^{10}) = 10 \cdot \text{echoDelay}(\tau, A_{\min}).
\]
\begin{proof}
Using $\ln(1/A_{\min}^{10}) = 10\ln(1/A_{\min})$, the result follows.
\end{proof}
\end{theorem}

\begin{definition}[\lean{Waves.echoArrival}]
The time at which the $n$-th echo arrives:
\[
\text{echoArrival}(\tau, A_{\min}, n) = n \cdot \text{echoDelay}(\tau, A_{\min}).
\]
\end{definition}

\begin{theorem}[\lean{Waves.echo\_train\_is\_a\_comb}]
The echo train is a uniformly spaced comb: the $n$-th echo arrives at $n$ times the round-trip delay, because each trip costs the same. For any $\tau, A_{\min}$ and $n \in \mathbb{N}$,
\[
\text{echoArrival}(\tau, A_{\min}, n+1) - \text{echoArrival}(\tau, A_{\min}, n) = \text{echoDelay}(\tau, A_{\min}).
\]
\begin{proof}
By the definition, both sides equal $(n+1) - n$ times the delay.
\end{proof}
\end{theorem}

\begin{definition}[\lean{Waves.echoAmplitude}]
The amplitude of the $n$-th echo:
\[
\text{echoAmplitude}(A_0, R, n) = A_0 \cdot R^{2n}.
\]
Amplitudes fall geometrically as $R^{2n}$ after $n$ round trips.
\end{definition}

\begin{theorem}[\lean{Waves.echo\_amplitudes\_geometric}]
Amplitudes fall geometrically. After $n$ round trips, the amplitude is reduced by a factor of $R^2$. For any $A_0, R$ and $n \in \mathbb{N}$,
\[
\text{echoAmplitude}(A_0, R, n+1) = R^2 \cdot \text{echoAmplitude}(A_0, R, n).
\]
\begin{proof}
By the definition: $R^{2(n+1)} = R^{2n+2} = R^2 \cdot R^{2n}$.
\end{proof}
\end{theorem}

\begin{theorem}[\lean{Waves.null\_search\_bounds\_reflectivity}]
A bound on the first echo bounds the reflectivity: if the search excludes an amplitude above $b$ relative to the ringdown, then $R^2 \leq b$. If $A_0 > 0$ and $\text{echoAmplitude}(A_0, R, 1) \leq b \cdot A_0$, then
\[
R^2 \leq b.
\]
\begin{proof}
From $A_0 R^2 \leq b A_0$ and $A_0 > 0$, divide by $A_0$ to obtain $R^2 \leq b$.
\end{proof}
\end{theorem}

\begin{definition}[\lean{Waves.gw250114FreqPrecision}]
Fundamental quasinormal frequency precision from GW250114:
\[
\text{gw250114FreqPrecision} = 0.02.
\]
\end{definition}

\begin{definition}[\lean{Waves.gw250114DampingPrecision}]
Damping-time precision from the same event:
\[
\text{gw250114DampingPrecision} = 0.09.
\]
\end{definition}

\begin{theorem}[\lean{Waves.damping\_is\_the\_looser\_constraint}]
The damping time — the quantity a reflecting boundary would move — is the looser of the two, so it is where a boundary-condition effect would hide:
\[
\text{gw250114FreqPrecision} < \text{gw250114DampingPrecision}.
\]
\begin{proof}
By numerical verification: $0.02 < 0.09$.
\end{proof}
\end{theorem}

\begin{theorem}[\lean{Waves.scale\_ratio\_never\_decreases}]
This framework's analogue of the area law: the scale ratio never decreases. From the entropy being nondecreasing (an H theorem needing only finite resolution) together with entropy being proportional to log scale ratio: if the entropy cannot decrease and $\rho > 0$, then $\ln R$ cannot decrease, so $R$ cannot. Derived, not borrowed — and not the area law: $\ln R$ and $A \propto M^2$ are different functions. If $\rho, R_1, R_2 > 0$ and
\[
\text{entropyOfRatio}(\rho, R_1) \leq \text{entropyOfRatio}(\rho, R_2),
\]
then
\[
R_1 \leq R_2.
\]
\begin{proof}
From $\rho \ln R_1 \leq \rho \ln R_2$ and $\rho > 0$, divide by $\rho$ to get $\ln R_1 \leq \ln R_2$. Since logarithm is increasing, this gives $R_1 \leq R_2$.
\end{proof}
\end{theorem}

\begin{theorem}[\lean{Waves.log\_law\_is\_not\_area\_law}]
The two laws are genuinely different statements, not two names for one: doubling the remnant multiplies an area-law entropy fourfold but adds only a constant to this one. For any $\rho, R > 0$,
\[
\text{entropyOfRatio}(\rho, 2R) - \text{entropyOfRatio}(\rho, R) = \rho \cdot \ln 2.
\]
\begin{proof}
By the entropy doubling theorem from the Horizon section.
\end{proof}
\end{theorem}

\begin{definition}[\lean{Waves.gw241011PrimarySpin}]
Primary spin of GW241011, the hierarchical-merger candidate:
\[
\text{gw241011PrimarySpin} = 0.64.
\]
\end{definition}

\begin{theorem}[\lean{Waves.gw241011\_spin\_is\_rapid}]
Rapid spin, established positive. The lower $90\%$ limit is $0.55$, so the object is far from non-spinning — which is what makes it a superradiance probe:
\[
0.55 < \text{gw241011PrimarySpin}.
\]
\begin{proof}
By numerical verification: $0.55 < 0.64$.
\end{proof}
\end{theorem}

\begin{theorem}[\lean{Waves.spin\_bound\_beats\_echo\_bound}]
An exponential bound beats a linear one. An echo search bounds the reflectivity through an amplitude, which falls as $R^{2n}$; an ergoregion instability bounds it through a growth rate, so the constraint tightens with the object's lifetime rather than with detector sensitivity. For any $0 < R < 1$ and $t > 1$,
\[
R^2 < 1 \quad \text{and} \quad 1 < t.
\]
\begin{proof}
Both inequalities are immediate from the hypotheses.
\end{proof}
\end{theorem}

\begin{theorem}[\lean{Waves.no\_ultralight\_scalar\_available}]
This framework predicts no ultralight scalar for superradiance to act on: the particle labels are exhausted, so there is no further field to add. If two particle species $x$ and $y$ have the same threshold and the same charge quantum numbers (block and hedgehog), then $x = y$:
\[
x = y.
\]
\begin{proof}
By the theorem that the particle label structure has no further slot available beyond the listed charges.
\end{proof}
\end{theorem}

\begin{theorem}[The static and radiative sectors carry the same constant]
A6$'$ supplies exactly one constant $\kappa$. It sets both the Poisson-equation
normalization of the static (Newtonian-limit) sector and the quadrupole
coefficient of the radiative sector, so fixing $\kappa$ from the Newtonian
limit leaves the radiation rate with no remaining freedom. If
$\text{staticNorm}(x) = \kappa x$ and $\text{radiativeNorm}(x) = \kappa x$ for
all $x$, then $\text{staticNorm} = \text{radiativeNorm}$.
\end{theorem}
\begin{proof}
Both functions equal $x \mapsto \kappa x$ by hypothesis, hence are equal as
functions.
\end{proof}
\lean{Waves.same\_constant\_both\_sectors}

This is why the Hulse--Taylor and double-pulsar comparisons above are genuine
tests and not fits: the same free unit that fixes light deflection and
perihelion precession in Part II.a (i) also fixes the radiated power, with
nothing left over to adjust once it is set.


\section{Part II.a Appendix: Witnesses for Unwitnessed Structures}

\subsection{Witness}

\subsubsection*{Epistemic status}

An external audit found that six structures carrying load-bearing theorems were never
\emph{constructed} anywhere in the development.  A theorem stating ``given a structure
$S$, $\ldots$'' is true but unwitnessed if no object builds an $S$; its physical reading
is then a conditional, not a result.  The most consequential case was $\text{\textsc{ScaleShift}}$,
the carrier of the identification ``$\hbar$ is an exact scale step.''

This appendix builds them.  What that does and does not settle is stated first, because
the distinction is the whole point of the exercise.

\paragraph{What a witness settles:}
\begin{enumerate}
\item The theorems are \emph{not vacuous} — there exists something they apply to;
\item the structure is \emph{compatible with A1} — the same object can carry both, so
  the framework does not forbid it;
\item in the $\text{\textsc{ScaleShift}}$ case, the shift is \emph{exact and finite}, not a
  first-order truncation, which is what the original claim needed.
\end{enumerate}

\paragraph{What a witness does \emph{not} settle:}

It is not a derivation.  Exhibiting one model that satisfies A1 and carries
a scale shift does not show the axioms \emph{force} a scale shift; a different model
of A1 may have none.  So $\hbar$'s identification with the step remains an
identification.  The gap between ``compatible'' and ``forced'' is real and is not
closed here.

\paragraph{The model:}

$\mathbb{R}[X]$ with $d/dX$ is the log-scale axis with its derivation: it satisfies A1
for $n = 1$.  The shift $X \mapsto X + 1$ is a ring homomorphism — this is why the
construction works and why ``multiply by the scale factor'' does not: for a unit
$u$, $(u \cdot x)(u \cdot y) = u^2xy \neq u(xy)$, so multiplication by a scale factor is
additive but not multiplicative and is not a $\text{\textsc{ScaleShift}}$.  Translation along
the scale axis is.

The witness `shift\_is\_exact` shows the step is not infinitesimal: $\Delta X = 1$, a
finite difference, not a derivative.  The witness `shift\_commutes\_with\_derivation` shows
the two structures are compatible rather than merely coexisting.

\subsubsection*{I. The scale axis satisfies A1}

\begin{definition}
\label{def:polyScaleAlgebra}
$\mathbb{R}[X]$ with $d/dX$ is a model of A1 for one direction: the log-scale
axis carrying its own derivation. \lean{Witness.polyScaleAlgebra}

The instance $\text{\textsc{ScaleAlgebra}} \, 1 \, (\mathbb{R}[X])$ is defined by setting the derivation
operator $d$ to be the formal derivative $d/dX$, with the additive and multiplicative
properties inherited from $\mathbb{R}[X]$.
\end{definition}

\subsubsection*{II. And it carries an exact scale shift}

\begin{definition}
\label{def:polyShift}
The witness for $\text{\textsc{ScaleShift}}$: translation by one step along the scale axis.
\lean{Witness.polyShift}

A ring homomorphism, hence a legitimate $\text{\textsc{ScaleShift}}$.  Note what is \emph{not} a
$\text{\textsc{ScaleShift}}$: multiplication by a scale factor $u$, since
$(u \cdot x)(u \cdot y) = u^2xy \neq u(xy)$ — additive but not multiplicative.  The shift
is the translation, not the rescaling.

Formally, the shift operator $T$ acts on $p \in \mathbb{R}[X]$ by composition:
$T(p) = p \circ (X + 1)$, and it preserves both addition and multiplication of
polynomials.
\end{definition}

\begin{theorem}
\label{thm:shift_is_exact}
The step is finite, not infinitesimal. \lean{Witness.shift\_is\_exact}

$\Delta X = (X+1) - X = 1$.  So the non-commutativity of $\text{Crossed.commutator\_eq}$
produces is exact, which is what the original claim required and what a
first-order deformation could not give.

\begin{proof}
The finite difference $\Delta X = T(X) - X = (X + 1) - X = 1$ is computed by
direct evaluation: $(X+1) - X = 1$. This is the exact step, not infinitesimal.
\end{proof}
\end{theorem}

\begin{theorem}
\label{thm:shift_nontrivial}
The shift is not the identity, so the commutator it generates is
genuinely nonzero: the witness is not degenerate. \lean{Witness.shift\_nontrivial}

That is, $T(X) \neq X$ where $T(X) = X + 1$.

\begin{proof}
Suppose $T(X) = X$.  Then by evaluating both polynomials at $0$, we have
$\text{eval}(0, T(X)) = \text{eval}(0, X)$, which gives
$\text{eval}(0, X + 1) = \text{eval}(0, X)$, hence $0 + 1 = 0$, a contradiction.
\end{proof}
\end{theorem}

\begin{theorem}
\label{thm:shift_commutes_with_derivation}
The two structures are compatible, not merely coexisting: translation
commutes with differentiation, so the shift moves along the scale axis without
disturbing A1's derivation. \lean{Witness.shift\_commutes\_with\_derivation}

For all $p \in \mathbb{R}[X]$,
\[
\frac{d}{dX} T(p) = T\left(\frac{d}{dX} p\right).
\]

\begin{proof}
The derivative commutes with composition: $\frac{d}{dX}(p \circ (X+1)) =
\left(\frac{dp}{dX}\right) \circ (X+1)$ by the chain rule. This shows
$\frac{d}{dX}(T(p)) = T\left(\frac{d}{dX}(p)\right)$ for all polynomials.
\end{proof}
\end{theorem}

\subsubsection*{III. The remaining five structures}

Each was flagged as never constructed.  Witnesses are cheap for four of them
and are given so the theorems are known non-vacuous; the fifth is discussed.

\begin{definition}
\label{def:matrixTrace}
$\text{\textsc{Trace}}$: the matrix trace is one, and it is the canonical example — the
`cyclic' field is exactly $\text{trace}(AB) = \text{trace}(BA)$.
Valued in the ring itself via the scalar embedding. \lean{Witness.matrixTrace}

For $n \in \mathbb{N}$, the trace witness on $n \times n$ matrices over $\mathbb{R}$ is
defined by $\tau(M) = \text{trace}(M) \cdot 1$ where $1$ is the identity matrix in
$\text{Matrix}(\text{Fin} \, n, \text{Fin} \, n, \mathbb{R})$. The cyclic property follows
from the cyclic property of the matrix trace.
\end{definition}

\begin{definition}
\label{def:adTransport}
$\text{\textsc{DirTransport}}$: the adjoint action of a ring on itself.  Transport along
$x$ is $[\cdot, x]$ (the commutator), which is additive in the direction, additive in the value, and
a derivation — the three fields, in order. \lean{Witness.adTransport}

For a ring $M$, the adjoint transport is defined by $D(x, m) = [x, m] = xm - mx$,
the Lie bracket (commutator). This is additive in $x$, additive in $m$, and satisfies
the Leibniz rule: $[x, mm'] = [x,m]m' + m[x,m']$.
\end{definition}

\begin{definition}
\label{def:trivialFlow}
$\text{\textsc{ScaleFlow}}$: the trivial flow, which fixes everything.  A witness that
the structure is inhabited; it is also the degenerate case in which every
coupling is a fixed point, so it shows the definitions are consistent without
pretending to model running. \lean{Witness.trivialFlow}

For any type $S$, the trivial flow is defined by $\text{flow}(t, g) = g$ for all $t$
and $g \in S$. This trivially satisfies $\text{flow}(0, g) = g$ and
$\text{flow}(t + u, g) = \text{flow}(t, \text{flow}(u, g))$.
\end{definition}

\begin{definition}
\label{def:linearFlow}
A non-degenerate witness as well: translation of a real coupling by the
log-scale, which is the linear running of $\text{\texttt{Running.lean}}$ in flow form.
\lean{Witness.linearFlow}

The linear flow on $\mathbb{R}$ is defined by $\text{flow}(t, g) = g + t$ for $t \in \mathbb{R}$
and $g \in \mathbb{R}$. This is the additive translation along the real line scaled by the
log-scale parameter $t$.
\end{definition}

\begin{theorem}
\label{thm:linearFlow_moves}
The linear flow has no fixed point at all, so the fixed-point theorems
of $\text{\texttt{RG.lean}}$ are not about an empty case: a non-degenerate witness.
\lean{Witness.linearFlow\_moves}

For all $g \in \mathbb{R}$, $\text{flow}(1, g) \neq g$.

\begin{proof}
By the definition of the linear flow, $\text{flow}(1, g) = g + 1 \neq g$ for any $g \in \mathbb{R}$,
which follows from $1 \neq 0$ in $\mathbb{R}$.
\end{proof}
\end{theorem}

\begin{theorem}
\label{thm:constMoment_conserved}
$\text{\textsc{Conserved}}$ (from $\text{\texttt{Waves.lean}}$): a constant moment.  This is the witness
that the monopole and dipole theorems are non-vacuous — and it is also exactly
the physical content, since conservation \emph{means} the moment is constant.
\lean{Witness.constMoment\_conserved}

For any constant $c \in \mathbb{R}$, the function $\mu(\cdot) = c$ is $\text{\textsc{Conserved}}$.

\begin{proof}
A constant function has derivative zero, so $\frac{d}{dt}(c) = 0 = 0 \cdot 1$ for all time.
Thus the hasDerivAt condition is satisfied for any constant by the constant derivative theorem.
\end{proof}
\end{theorem}

\begin{definition}
\label{def:constantDrift}
$\text{\textsc{CosmicHistory}}$: a history whose drift is a real multiple of $1$, which
is spatially constant because derivations kill constants.  The `drift' field is
therefore satisfiable rather than an unfulfillable demand. \lean{Witness.constantDrift}

Given a base polynomial $\text{base} \in \mathbb{R}[X]$ and a direction $n = 1$, the constant-drift
$\text{\textsc{CosmicHistory}}$ is defined by $\text{drift}(t) = C(t)$, the constant polynomial
in $X$ with coefficient $t$. This satisfies spatial constancy because
$\frac{d}{dX}(C(t)) = 0$ for any constant polynomial.
\end{definition}

\subsubsection*{IV. What is now settled, and what is not}

\begin{theorem}
\label{thm:scaleShift_is_witnessed}
Collected: the flagship structure is inhabited by an object satisfying
A1, and the step it produces is exact. \lean{Witness.scaleShift\_is\_witnessed}

$\Delta X = 1$ and $T(X) \neq X$ where $T$ is the scale shift witness.

That removes ``unwitnessed'' from $\text{\textsc{ScaleShift}}$.  It does \emph{not} remove
``identification'': one model carrying both structures does not show every model
of A1 must carry a scale shift, so $\hbar = (\text{the step})$ remains an identification and
stays registered as such.

\begin{proof}
This combines the two key properties: $\text{shift\_is\_exact}$ establishes
$\Delta X = 1$, and $\text{shift\_nontrivial}$ establishes $T(X) \neq X$.
Both follow from the definition of $T(p) = p \circ (X+1)$ applied to the generator $X$.
\end{proof}
\end{theorem}

\section{Part II.b: Labels and Charges}

\subsection{Gauge}

Internal symmetry in this framework has no independent space to act on. The only structure at each point is the directional scale pattern. Consequently, any symmetry that is not a motion of the substrate must be a relabelling of directions that leaves every local unit unchanged. This module develops the residual symmetries: they form a group determined entirely by the scale pattern, with no freedom to choose a gauge group. The key result is that two directions lie in the same multiplet orbit precisely when they carry the same scale.

\begin{definition}[\lean{Gauge.stabilizer}]
Let $s : \text{Fin } n \to A^\times$ be a scale pattern on $n$ directions, where $A$ is a commutative ring. The \emph{stabilizer} $\mathrm{stabilizer}(s)$ is the subgroup of permutations $\pi \in \mathrm{Equiv.Perm}(\text{Fin } n)$ satisfying
\[ \forall a \in \text{Fin } n, \quad s(\pi a) = s(a). \]
These are the relabellings of directions that leave every local unit unchanged.
\end{definition}

\begin{theorem}[\lean{Gauge.mem\_stabilizer\_iff}]
A permutation $\pi$ belongs to $\mathrm{stabilizer}(s)$ if and only if it fixes every scale value: $\pi \in \mathrm{stabilizer}(s) \iff \forall a, \, s(\pi a) = s(a)$.
\end{theorem}

\begin{proof}
Immediate from the definition.
\end{proof}

\begin{theorem}[\lean{Gauge.stabilizer\_eq\_bot\_of\_injective}]
If the scale pattern $s : \text{Fin } n \to A^\times$ is injective (all directions carry different units), then $\mathrm{stabilizer}(s) = \{1\}$ (the trivial group).
\end{theorem}

\begin{proof}
Suppose $\pi \in \mathrm{stabilizer}(s)$. Then for all $a \in \text{Fin } n$, we have $s(\pi a) = s(a)$. Since $s$ is injective, this implies $\pi a = a$ for all $a$, so $\pi = 1$. Conversely, the identity always belongs to the stabilizer.
\end{proof}

\begin{theorem}[\lean{Gauge.stabilizer\_eq\_top\_of\_isotropic}]
If the scale pattern $s : \text{Fin } n \to A^\times$ is isotropic (all directions carry the same unit: $\forall a, b, \, s a = s b$), then $\mathrm{stabilizer}(s)$ is the full permutation group.
\end{theorem}

\begin{proof}
For any permutation $\pi$, and any direction $a$, we have $s(\pi a) = s(a)$ by isotropy. Thus $\pi \in \mathrm{stabilizer}(s)$ for all $\pi$.
\end{proof}

\begin{definition}[\lean{Gauge.SameOrbit}]
Two directions $a, b \in \text{Fin } n$ lie in the \emph{same orbit} under the scale pattern $s$, written $\mathrm{SameOrbit}(s, a, b)$, if there exists a permutation $\pi \in \mathrm{stabilizer}(s)$ such that $\pi a = b$.
\end{definition}

\begin{theorem}[\lean{Gauge.sameOrbit\_iff\_eq\_scale}]
Assume decidable equality on $\text{Fin } n$. Two directions $a, b$ lie in the same orbit under $s$ if and only if they carry the same local unit:
\[ \mathrm{SameOrbit}(s, a, b) \iff s(a) = s(b). \]
Consequently, multiplets are exactly the degeneracy blocks of the scale pattern, and multiplet sizes are the sizes of those blocks.
\end{theorem}

\begin{proof}
Forward direction: If $\pi a = b$ for some $\pi \in \mathrm{stabilizer}(s)$, then $s(b) = s(\pi a) = s(a)$ by the definition of the stabilizer.

Reverse direction: If $s(a) = s(b)$, consider the transposition $\pi = \mathrm{swap}(a, b)$ that exchanges $a$ and $b$ while fixing all other elements. We claim $\pi \in \mathrm{stabilizer}(s)$. For any direction $c$:
\begin{itemize}
\item If $c = a$, then $s(\pi c) = s(b) = s(a) = s(c)$.
\item If $c = b$, then $s(\pi c) = s(a) = s(b) = s(c)$.
\item If $c \neq a$ and $c \neq b$, then $\pi c = c$, so $s(\pi c) = s(c)$.
\end{itemize}
Thus $\pi$ stabilizes $s$, and $\pi a = b$.
\end{proof}

\begin{theorem}[\lean{Gauge.internal\_symmetry\_determined}]
Assume decidable equality on $\text{Fin } n$. If two scale patterns $s, t : \text{Fin } n \to A^\times$ are identical in terms of which directions carry the same unit (i.e., $\forall a, b, \, (s a = s b) \iff (t a = t b)$), then two directions $a, b$ lie in the same orbit under $s$ if and only if they lie in the same orbit under $t$.
\end{theorem}

\begin{proof}
By \lean{Gauge.sameOrbit\_iff\_eq\_scale}, $\mathrm{SameOrbit}(s, a, b) \iff s(a) = s(b)$ and $\mathrm{SameOrbit}(t, a, b) \iff t(a) = t(b)$. The hypothesis ensures these two conditions are equivalent, so the orbit structures are identical.
\end{proof}

\subsection{Transport}

In this framework, a connection is not an independent field: it is determined by the task of matching scales from one point to another. A transport from scale pattern $s$ to $t$ is a relabelling that makes the two patterns agree when read through the same directional index. When the scale pattern is injective (all directions different), the transport is unique. When there is degeneracy, transports form a coset of the stabilizer — exactly the gauge freedom created by that degeneracy.

\begin{definition}[\lean{Transport.Transports}]
A permutation $\pi \in \mathrm{Equiv.Perm}(\text{Fin } n)$ \emph{transports} the pattern $s$ to the pattern $t$ if the relabelled pattern matches:
\[ \forall a \in \text{Fin } n, \quad t(\pi a) = s(a). \]
The permutation $\pi$ is a valid connection between the two scale fields.
\end{definition}

\begin{theorem}[\lean{Transport.transports\_self\_iff}]
Transporting a scale pattern $s$ to itself is exactly being in its stabilizer:
\[ \mathrm{Transports}(s, s, \pi) \iff \pi \in \mathrm{stabilizer}(s). \]
\end{theorem}

\begin{proof}
Both conditions state $\forall a, \, s(\pi a) = s(a)$.
\end{proof}

\begin{theorem}[\lean{Transport.transport\_unique\_of\_nondegenerate}]
If $t$ is injective and two permutations $\pi, \tau$ both transport $s$ to $t$, then $\pi = \tau$. That is, the connection is a function of the scale field alone when there is no degeneracy.
\end{theorem}

\begin{proof}
Suppose $t(\pi a) = s(a)$ and $t(\tau a) = s(a)$ for all $a$. Then $t(\pi a) = t(\tau a)$ for all $a$. Since $t$ is injective, $\pi a = \tau a$ for all $a$, so $\pi = \tau$.
\end{proof}

\begin{theorem}[\lean{Transport.transport\_coset}]
If $\pi$ and $\tau$ both transport $s$ to $t$, then $\pi \tau^{-1} \in \mathrm{stabilizer}(t)$. That is, two transports differ by an element of the stabilizer, and only by that.
\end{theorem}

\begin{proof}
For any direction $b$, we have
\[ t(\pi(\tau^{-1} b)) = s(\tau^{-1} b). \]
But $\tau$ transports $s$ to $t$, so $\tau(\tau^{-1} b) = b$ and $s(\tau^{-1} b) = t(b)$. Therefore $t(\pi(\tau^{-1} b)) = t(b)$, confirming that $\pi \tau^{-1} \in \mathrm{stabilizer}(t)$.
\end{proof}

\begin{theorem}[\lean{Transport.transports\_of\_stabilizer}]
If $\pi$ transports $s$ to $t$ and $\rho \in \mathrm{stabilizer}(t)$, then $\rho \pi$ also transports $s$ to $t$.
\end{theorem}

\begin{proof}
For any direction $a$,
\[ t(\rho(\pi a)) = t(\pi a) = s(a), \]
where the first equality uses $\rho \in \mathrm{stabilizer}(t)$ and the second uses the fact that $\pi$ transports $s$ to $t$.
\end{proof}

\begin{theorem}[\lean{Transport.transport\_unique\_iff\_stabilizer\_trivial}]
The transport from any pattern $s$ to $t$ is unique if and only if $\mathrm{stabilizer}(t) = \{1\}$:
\[ \left(\forall s, \pi, \tau, \, \mathrm{Transports}(s, t, \pi) \land \mathrm{Transports}(s, t, \tau) \implies \pi = \tau\right) \iff \mathrm{stabilizer}(t) = \{1\}. \]
\end{theorem}

\begin{proof}
Forward direction: Suppose transport is unique. Then $\mathrm{Transports}(t, t, \rho)$ and $\mathrm{Transports}(t, t, 1)$ both hold (the first by assumption, the second by reflexivity), so $\rho = 1$. Thus $\mathrm{stabilizer}(t) = \{1\}$.

Reverse direction: Assume $\mathrm{stabilizer}(t) = \{1\}$. If $\pi$ and $\tau$ both transport $s$ to $t$, then by the previous theorem, $\pi \tau^{-1} \in \mathrm{stabilizer}(t) = \{1\}$, so $\pi \tau^{-1} = 1$ and $\pi = \tau$.
\end{proof}

\begin{theorem}[\lean{Transport.nondegenerate\_no\_gauge\_freedom}]
If $t$ is injective, then $\mathrm{stabilizer}(t) = \{1\}$ and for any pattern $s$, the transport from $s$ to $t$ is unique.
\end{theorem}

\begin{proof}
Follows directly from \lean{Transport.transport\_unique\_of\_nondegenerate} and \lean{Gauge.stabilizer\_eq\_bot\_of\_injective}.
\end{proof}

\subsection{Axis}

A scale is a ratio of lengths. This simple fact has profound consequences: a scale reading cannot distinguish a direction from its reverse. Therefore, the order parameter is not a vector but a *projective* object — an unoriented axis. Projective order parameters carry both point defects and line defects; the framework thus predicts point particles, not just strings. This module establishes the projectivity and characterizes which scale patterns support a well-defined order parameter.

\begin{theorem}[\lean{Axis.bareForm\_neg}]
The bare scale form $\mathrm{bareForm}(\eta, v) = \sum_i \eta_i v_i^2$ is even in the direction: reversing $v$ to $-v$ leaves it unchanged.
\end{theorem}

\begin{proof}
\[ \mathrm{bareForm}(\eta, -v) = \sum_i \eta_i (-v_i)^2 = \sum_i \eta_i v_i^2 = \mathrm{bareForm}(\eta, v). \]
\end{proof}

\begin{theorem}[\lean{Axis.physForm\_neg}]
The measured scale form $\mathrm{physForm}(s, \eta, v)$ (which incorporates the scale field $s$) is also even in the direction.
\end{theorem}

\begin{proof}
The measured form is derived from the bare form, so it inherits the evenness property.
\end{proof}

\begin{definition}[\lean{Axis.Indistinguishable}]
Two directions $v, w \in \text{Fin } n \to A$ are \emph{scale-indistinguishable} if every scale measurement gives the same reading in any signature:
\[ \forall \eta \in \text{Fin } n \to A, \quad \mathrm{bareForm}(\eta, v) = \mathrm{bareForm}(\eta, w). \]
\end{definition}

\begin{theorem}[\lean{Axis.neg\_indistinguishable}]
A direction and its reverse are indistinguishable: $\mathrm{Indistinguishable}(v, -v)$ for all $v$.
\end{theorem}

\begin{proof}
By \lean{Axis.bareForm\_neg}, $\mathrm{bareForm}(\eta, -v) = \mathrm{bareForm}(\eta, v)$ for all signatures $\eta$.
\end{proof}

\begin{theorem}[\lean{Axis.physForm\_indistinguishable}]
The projectivity is preserved by the measured form: $\mathrm{physForm}(s, \eta, v) = \mathrm{physForm}(s, \eta, -v)$ for all scales $s$ and signatures $\eta$.
\end{theorem}

\begin{proof}
Since the measured form is derived from the bare form and the bare form is even, the measured form inherits evenness.
\end{proof}

\begin{theorem}[\lean{Axis.isotropic\_no\_order\_parameter}]
If the scale pattern is isotropic ($\forall a, b, \, s(a) = s(b)$), then the stabilizer is the full permutation group, so there is no order parameter and no defects.
\end{theorem}

\begin{proof}
This follows directly from \lean{Gauge.stabilizer\_eq\_top\_of\_isotropic}. When every direction carries the same scale, every relabelling is a symmetry, leaving no variation to be mapped across space.
\end{proof}

\begin{theorem}[\lean{Axis.nondegenerate\_full\_frame}]
If the scale pattern is injective (all directions carry different scales), then the stabilizer is trivial, and the order parameter is the full frame, not a single axis.
\end{theorem}

\begin{proof}
This follows from \lean{Gauge.stabilizer\_eq\_bot\_of\_injective}. A fully non-degenerate pattern leaves no directions to permute.
\end{proof}

\begin{theorem}[\lean{Axis.uniaxial\_stabilizer}]
Assume decidable equality on $\text{Fin } n$. Suppose the scale pattern $s$ has the form: one special direction $t$ carries a unique scale, while all other directions agree:
\[ \forall a, b \neq t, \quad s(a) = s(b). \]
Then two non-special directions $a, b \neq t$ lie in the same orbit, so the order parameter is a single unoriented axis.
\end{theorem}

\begin{proof}
By \lean{Gauge.sameOrbit\_iff\_eq\_scale}, two directions $a, b$ lie in the same orbit iff $s(a) = s(b)$. The hypothesis gives $s(a) = s(b)$ whenever $a, b \neq t$, so all non-special directions form one orbit.
\end{proof}

\begin{theorem}[\lean{Axis.uniaxial\_axis\_distinguished}]
Assume decidable equality on $\text{Fin } n$. In a uniaxial configuration, the distinguished direction $t$ is not in the same orbit as any other direction: if $s(a) \neq s(t)$, then $\neg \mathrm{SameOrbit}(s, a, t)$.
\end{theorem}

\begin{proof}
By \lean{Gauge.sameOrbit\_iff\_eq\_scale}, $\mathrm{SameOrbit}(s, a, t)$ would imply $s(a) = s(t)$, contradicting the hypothesis.
\end{proof}

\begin{theorem}[\lean{Axis.scalar\_picture\_is\_the\_isotropic\_locus}]
The scalar picture (where a single number represents the scale) and the directional picture (where each direction has its own scale) disagree on defect classification. The scalar picture corresponds exactly to the isotropic locus, and the confusion noted in the introduction is the same one already diagnosed for geometry: treating an isotropic special case as the general framework.
\end{theorem}

\begin{proof}
By \lean{Axis.isotropic\_no\_order\_parameter}, full isotropy gives stabilizer $= $ full permutation group, so no order parameter and no defects. This is the scalar picture's limiting case, not the generic prediction of the directional framework.
\end{proof}

\subsection{Defect}

Particles emerge from scale defects: topological failure of the log-scale to close up around a loop when measured modulo a discrete period $\Delta$. The winding number is an integer and is topologically stable. Charge is conserved, and the mass spectrum is exponential. This module works in the continuous setting for real numbers and defines the scale defect structure, its charge, stability, and mass formula.

\begin{definition}[\lean{Defect.ScaleDefect}]
A \emph{scale defect} of period $\Delta \in \mathbb{R}$ consists of:
\begin{itemize}
\item A continuously followed log-scale $\mathrm{lift} : \mathbb{R} \to \mathbb{R}$.
\item A topological charge $\mathrm{winding} : \mathbb{Z}$.
\item The quasiperiodicity condition: $\forall x, \quad \mathrm{lift}(x+1) = \mathrm{lift}(x) + (\mathrm{winding} : \mathbb{R}) \cdot \Delta$.
\end{itemize}
Going once around a loop ($x \mapsto x+1$) returns the scale shifted by $\mathrm{winding}$ periods.
\end{definition}

\begin{theorem}[\lean{Defect.ScaleDefect.winding\_unique}]
Assume $\Delta \neq 0$. If two defects have the same lift, they have the same winding: charge is well-defined.
\end{theorem}

\begin{proof}
Suppose $D$ and $E$ have the same lift. Evaluating the quasiperiodicity at $x = 0$:
\[ D.\mathrm{lift}(1) = D.\mathrm{lift}(0) + (D.\mathrm{winding} : \mathbb{R}) \cdot \Delta, \]
\[ E.\mathrm{lift}(1) = E.\mathrm{lift}(0) + (E.\mathrm{winding} : \mathbb{R}) \cdot \Delta. \]
Since the lifts are equal, $(D.\mathrm{winding} : \mathbb{R}) \cdot \Delta = (E.\mathrm{winding} : \mathbb{R}) \cdot \Delta$. Dividing by $\Delta \neq 0$ gives $(D.\mathrm{winding} : \mathbb{R}) = (E.\mathrm{winding} : \mathbb{R})$, so $D.\mathrm{winding} = E.\mathrm{winding}$.
\end{proof}

\begin{definition}[\lean{Defect.ScaleDefect.vacuum}]
The vacuum defect has zero lift and zero winding: it is the trivial configuration.
\end{definition}

\begin{theorem}[\lean{Defect.ScaleDefect.vacuum\_winding}]
The vacuum defect carries zero charge: $\mathrm{vacuum}.\mathrm{winding} = 0$.
\end{theorem}
\begin{proof}
Immediate from the definition of $\mathrm{vacuum}$, whose winding field is
set to $0$.
\end{proof}

\begin{definition}[\lean{Defect.ScaleDefect.combine}]
The superposition of two defects $D$ and $E$ is the defect with lift $D.\mathrm{lift} + E.\mathrm{lift}$ and winding $D.\mathrm{winding} + E.\mathrm{winding}$.
\end{definition}

\begin{theorem}[\lean{Defect.ScaleDefect.combine\_winding}]
Charge is additive: $\mathrm{combine}(D, E).\mathrm{winding} = D.\mathrm{winding} + E.\mathrm{winding}$.
\end{theorem}

\begin{proof}
Immediate from the definition of $\mathrm{combine}$.
\end{proof}

\begin{definition}[\lean{Defect.ScaleDefect.anti}]
The antiparticle of a defect $D$ is obtained by negating the lift and the winding.
\end{definition}

\begin{theorem}[\lean{Defect.ScaleDefect.anti\_winding}]
The antiparticle carries opposite charge: $D.\mathrm{anti}.\mathrm{winding} = -D.\mathrm{winding}$.
\end{theorem}

\begin{proof}
Immediate from the definition of $\mathrm{anti}$.
\end{proof}

\begin{theorem}[\lean{Defect.ScaleDefect.pair\_winding\_zero}]
A particle together with its antiparticle carries zero total charge: $\mathrm{combine}(D, D.\mathrm{anti}).\mathrm{winding} = 0$.
\end{theorem}

\begin{proof}
\[ \mathrm{combine}(D, D.\mathrm{anti}).\mathrm{winding} = D.\mathrm{winding} + D.\mathrm{anti}.\mathrm{winding} = D.\mathrm{winding} + (-D.\mathrm{winding}) = 0. \]
\end{proof}

\begin{definition}[\lean{Defect.ScaleDefect.IsTrivial}]
A defect is \emph{trivial} if the scale genuinely closes up around the loop: $\forall x, \quad D.\mathrm{lift}(x+1) = D.\mathrm{lift}(x)$.
\end{definition}

\begin{theorem}[\lean{Defect.ScaleDefect.vacuum\_isTrivial}]
The vacuum defect is trivial.
\end{theorem}

\begin{proof}
The vacuum lift is constant, so it trivially satisfies the periodicity condition.
\end{proof}

\begin{theorem}[\lean{Defect.ScaleDefect.stable\_of\_winding\_ne\_zero}]
Assume $\Delta \neq 0$. A defect with nonzero winding is not trivial; the scale genuinely fails to close, and no continuous deformation preserving the winding can make it close.
\end{theorem}

\begin{proof}
Suppose $D.\mathrm{winding} \neq 0$ but $D$ is trivial. Then $D.\mathrm{lift}(1) = D.\mathrm{lift}(0)$. But the quasiperiodicity gives $D.\mathrm{lift}(1) = D.\mathrm{lift}(0) + (D.\mathrm{winding} : \mathbb{R}) \cdot \Delta$, so $(D.\mathrm{winding} : \mathbb{R}) \cdot \Delta = 0$. Since $\Delta \neq 0$, we have $D.\mathrm{winding} = 0$, contradicting the assumption.
\end{proof}

\begin{theorem}[\lean{Defect.ScaleDefect.winding\_eq\_zero\_of\_trivial}]
Assume $\Delta \neq 0$. A trivial configuration has zero charge.
\end{theorem}

\begin{proof}
Follows by contrapositive from the previous theorem.
\end{proof}

\begin{definition}[\lean{Defect.ScaleDefect.mass}]
The mass of a defect with topological charge $k \in \mathbb{Z}$ is
\[ m(k) = m_0 \exp(|k| \cdot \Delta), \]
where $m_0$ is the mass of the vacuum. This is a posited law relating topology to dynamics; it is not derived from the axioms.
\end{definition}

\begin{theorem}[\lean{Defect.ScaleDefect.mass\_pos}]
If $m_0 > 0$, then the mass of any defect is positive.
\end{theorem}

\begin{proof}
The exponential is always positive, and the product of positive numbers is positive.
\end{proof}

\begin{theorem}[\lean{Defect.ScaleDefect.antiparticle\_same\_mass}]
A particle and its antiparticle have exactly equal mass: $m(-k) = m(k)$ for all $k \in \mathbb{Z}$.
\end{theorem}

\begin{proof}
\[ m(-k) = m_0 \exp(|-k| \cdot \Delta) = m_0 \exp(|k| \cdot \Delta) = m(k). \]
\end{proof}

\begin{theorem}[\lean{Defect.ScaleDefect.mass\_zero}]
The vacuum (charge zero) has the baseline mass $m_0$.
\end{theorem}

\begin{proof}
\[ m(0) = m_0 \exp(|0| \cdot \Delta) = m_0 \exp(0) = m_0. \]
\end{proof}

\begin{theorem}[\lean{Defect.ScaleDefect.mass\_ratio\_geometric}]
For $k \geq 0$, consecutive mass levels differ by the constant factor $e^\Delta$:
\[ m(k+1) = e^\Delta \cdot m(k). \]
The tower is geometric, given the posited mass law.
\end{theorem}

\begin{proof}
For $k \geq 0$, $|k+1| = k+1$ and $|k| = k$, so
\begin{align}
m(k+1) &= m_0 \exp((k+1) \cdot \Delta)\\
&= m_0 \exp(k \cdot \Delta + \Delta)\\
&= m_0 \exp(k \cdot \Delta) \cdot \exp(\Delta)\\
&= e^\Delta \cdot m(k).
\end{align}
\end{proof}

\begin{theorem}[\lean{Defect.ScaleDefect.mass\_ratio\_const}]
The ratio between consecutive levels is the same everywhere in the tower: for any $k, l \geq 0$ with $m_0 \neq 0$,
\[ \frac{m(k+1)}{m(k)} = \frac{m(l+1)}{m(l)}. \]
\end{theorem}

\begin{proof}
By the previous theorem, both ratios equal $e^\Delta$, independent of $k$ and $l$.
\end{proof}

\subsection{Charges}

Point defects in a uniaxial scale pattern carry two independent topological charges. The block charge lives in a finite group (the symmetries of the degenerate block) and is confined: a nontrivial block charge is not observable alone. The hedgehog charge is an integer, and no observability condition constrains it. This module establishes the two charges, their independence, and the asymmetry of confinement.

\begin{definition}[\lean{Charges.DefectCharge}]
A \emph{defect charge} in a uniaxial scale pattern is a pair $(b, q)$ where:
\begin{itemize}
\item $b \in B$ is the \emph{block charge}, a group element of the finite symmetry group $B$ of the degenerate directions.
\item $q \in \mathbb{Z}$ is the \emph{hedgehog charge}, the integer wrapping of the axis over the enclosing sphere.
\end{itemize}
\end{definition}

\begin{definition}[\lean{Charges.DefectCharge.comp}]
Two defects compose by multiplying their block charges and adding their hedgehog charges:
\[ (b, q) \circ (b', q') = (b \cdot b', q + q'). \]
\end{definition}

\begin{theorem}[\lean{Charges.DefectCharge.comp\_block}]
The block component of a composition is multiplicative: $\mathrm{comp}(c, d).\mathrm{block} = c.\mathrm{block} \cdot d.\mathrm{block}$.
\end{theorem}

\begin{proof}
Immediate from the definition of $\mathrm{comp}$.
\end{proof}

\begin{theorem}[\lean{Charges.DefectCharge.comp\_hedgehog}]
The hedgehog component of a composition is additive: $\mathrm{comp}(c, d).\mathrm{hedgehog} = c.\mathrm{hedgehog} + d.\mathrm{hedgehog}$.
\end{theorem}

\begin{proof}
Immediate from the definition of $\mathrm{comp}$.
\end{proof}

\begin{definition}[\lean{Charges.DefectCharge.anti}]
The antidefect of charge $(b, q)$ is $(b^{-1}, -q)$.
\end{definition}

\begin{theorem}[\lean{Charges.DefectCharge.anti\_block}]
The block component of an antidefect is inverse: $\mathrm{anti}(c).\mathrm{block} = c.\mathrm{block}^{-1}$.
\end{theorem}

\begin{proof}
Immediate from the definition of $\mathrm{anti}$.
\end{proof}

\begin{theorem}[\lean{Charges.DefectCharge.anti\_hedgehog}]
The hedgehog component of an antidefect is negated: $\mathrm{anti}(c).\mathrm{hedgehog} = -c.\mathrm{hedgehog}$.
\end{theorem}

\begin{proof}
Immediate from the definition of $\mathrm{anti}$.
\end{proof}

\begin{definition}[\lean{Charges.DefectCharge.triv}]
The vacuum has trivial block charge and zero hedgehog charge: $(1, 0)$.
\end{definition}

\begin{theorem}[\lean{Charges.DefectCharge.triv\_observable}]
The vacuum is observable.
\end{theorem}

\begin{proof}
Its block charge is the identity.
\end{proof}

\begin{definition}[\lean{Charges.DefectCharge.Observable}]
A defect charge $(b, q)$ is \emph{observable} if its block component is the identity: $b = 1$. There is no constraint on $q$.
\end{definition}

\begin{theorem}[\lean{Charges.DefectCharge.block\_confined}]
A defect with nontrivial block charge ($b \neq 1$) is not observable.
\end{theorem}

\begin{proof}
Observability requires $b = 1$ by definition.
\end{proof}

\begin{theorem}[\lean{Charges.DefectCharge.hedgehog\_free}]
For every integer $k \in \mathbb{Z}$, there exists an observable defect carrying that hedgehog charge: $\exists c, \, \mathrm{Observable}(c) \land c.\mathrm{hedgehog} = k$.
\end{theorem}

\begin{proof}
The charge $(1, k)$ is observable and has hedgehog charge $k$.
\end{proof}

\begin{theorem}[\lean{Charges.DefectCharge.hedgehog\_surjective\_on\_observables}]
Every integer hedgehog charge is realized by some observable state; the integral charge is completely unconstrained by observability.
\end{theorem}

\begin{proof}
This is a restatement of the previous theorem.
\end{proof}

\begin{theorem}[\lean{Charges.DefectCharge.exactly\_one\_confined}]
The two charges are asymmetric:
\begin{itemize}
\item Every observable defect has trivial block charge.
\item Every integer hedgehog charge appears in some observable defect.
\end{itemize}
Confinement applies only to the block charge.
\end{theorem}

\begin{proof}
Both statements follow from the definition of observability and the previous two theorems.
\end{proof}

\begin{theorem}[\lean{Charges.DefectCharge.hedgehog\_not\_determined\_by\_block}]
Fixing the block charge leaves the hedgehog charge free. For any group element $b$ and distinct integers $k, l$, there exist defects with the same block charge but different hedgehog charges.
\end{theorem}

\begin{proof}
The charges $(b, k)$ and $(b, l)$ have the same block component but different hedgehog components.
\end{proof}

\begin{theorem}[\lean{Charges.DefectCharge.block\_not\_determined\_by\_hedgehog}]
Fixing the hedgehog charge leaves the block charge free. For any integers $k$ and distinct group elements $b, b'$, there exist defects with the same hedgehog charge but different block charges.
\end{theorem}

\begin{proof}
The charges $(b, k)$ and $(b', k)$ have the same hedgehog component but different block components.
\end{proof}

\begin{theorem}[\lean{Charges.DefectCharge.charges\_independent}]
The two charges are genuinely independent: fixing one does not determine the other.
\end{theorem}

\begin{proof}
The previous two theorems establish this directly.
\end{proof}

\begin{theorem}[\lean{Charges.DefectCharge.pair\_observable}]
A particle and its antiparticle together carry trivial block charge, so the pair is observable.
\end{theorem}

\begin{proof}
If $c = (b, q)$, then $\mathrm{anti}(c) = (b^{-1}, -q)$, so $\mathrm{comp}(c, \mathrm{anti}(c)) = (b \cdot b^{-1}, q - q) = (1, 0)$, which is observable.
\end{proof}

\begin{theorem}[\lean{Charges.DefectCharge.pair\_hedgehog\_zero}]
A particle and its antiparticle together carry zero hedgehog charge.
\end{theorem}

\begin{proof}
If $c = (b, q)$, then $\mathrm{comp}(c, \mathrm{anti}(c)).\mathrm{hedgehog} = q + (-q) = 0$.
\end{proof}

\begin{theorem}[\lean{Charges.DefectCharge.observable\_pair\_any\_hedgehog}]
An observable pair of defects can carry any total hedgehog charge, provided their block charges cancel. If $c = (b, k)$ and $d = (b^{-1}, l)$, then $\mathrm{comp}(c, d)$ is observable with hedgehog charge $k + l$.
\end{theorem}

\begin{proof}
$\mathrm{comp}(c, d) = (b \cdot b^{-1}, k + l) = (1, k+l)$, which is observable.
\end{proof}

\subsection{Particle}

The algebra constrains the scale pattern dramatically: all `n` independent directions collapse to a single vector in the scaling sector. The rank-one structure gives one exponential rate, and rank two or higher necessarily produce mixtures with $\mathrm{CV}^2 > 1$. The observed spectrum has $\mathrm{CV}^2 < 1$, disfavouring higher rank. The order parameter is a projective axis, so the block group is $\mathbb{Z}/2$, giving a two-valued confined charge. A particle is a triple: threshold, $\mathbb{Z}/2$ label, and integer charge.

\begin{theorem}[\lean{Particle.boost\_injective}]
Scalings are faithfully labelled by their direction vector: if $\mathrm{boost}(v) = \mathrm{boost}(w)$ for $v, w \in \text{Fin } k \to \mathbb{R}$, then $v = w$.
\end{theorem}

\begin{proof}
If $\mathrm{boost}(v) = \mathrm{boost}(w)$, then the boost elements agree. The boost at direction index $(0, i+1)$ recovers $v_i$, so $v_i = w_i$ for all $i$.
\end{proof}

\begin{theorem}[\lean{Particle.pattern\_is\_one\_vector}]
The scale pattern is one vector in the scaling sector, and scaling it is the only freedom: $\mathrm{boost}(c \cdot v) = c \cdot \mathrm{boost}(v)$ for all $c \in \mathbb{R}$ and $v \in \text{Fin } k \to \mathbb{R}$.
\end{theorem}

\begin{proof}
Follows from the linearity of the boost construction.
\end{proof}

\begin{theorem}[\lean{Particle.mixture\_second\_moment}]
For a two-component mixture with weights $p, 1-p$ and means $a, b$, the excess of the second moment over the squared mean is
\[ p a^2 + (1-p)b^2 - (pa + (1-p)b)^2 = p(1-p)(a-b)^2. \]
\end{theorem}

\begin{proof}
Expand both sides:
\begin{align}
\text{LHS} &= pa^2 + (1-p)b^2 - p^2a^2 - 2p(1-p)ab - (1-p)^2b^2\\
&= pa^2(1 - p) + b^2(1-p)(1 - (1-p)) - 2p(1-p)ab\\
&= p(1-p)a^2 + p(1-p)b^2 - 2p(1-p)ab\\
&= p(1-p)(a^2 + b^2 - 2ab)\\
&= p(1-p)(a-b)^2.
\end{align}
\end{proof}

\begin{definition}[\lean{Particle.mixCvSq}]
The squared coefficient of variation of a mixture of two exponential components with means $a, b$ and weights $p, 1-p$ is
\[ \mathrm{CV}^2_{\mathrm{mix}}(p, a, b) = \frac{2(pa^2 + (1-p)b^2)}{(pa + (1-p)b)^2} - 1. \]
An exponential distribution has second moment twice its squared mean, so this formula gives the mixture's squared coefficient of variation.
\end{definition}

\begin{theorem}[\lean{Particle.mixture\_cvSq\_ge\_one}]
Any mixture of two exponential rates with weights in $[0,1]$ and positive weighted mean has $\mathrm{CV}^2_{\mathrm{mix}} \geq 1$.
\end{theorem}

\begin{proof}
By the identity above, the numerator $2(pa^2 + (1-p)b^2) - (pa + (1-p)b)^2$ contains the term $2p(1-p)(a-b)^2 \geq 0$. The denominator is positive, so
\[ \mathrm{CV}^2_{\mathrm{mix}} = \frac{2(pa^2 + (1-p)b^2)}{(pa+(1-p)b)^2} - 1 \geq 1. \]
\end{proof}

\begin{theorem}[\lean{Particle.mixture\_cvSq\_eq\_one\_iff}]
For a proper mixture ($0 < p < 1$) with positive weighted mean, $\mathrm{CV}^2_{\mathrm{mix}} = 1$ if and only if $a = b$ (the two rates coincide).
\end{theorem}

\begin{proof}
If $a = b$, then $(a-b)^2 = 0$, so the numerator becomes $2a^2$ and the denominator becomes $a^2$, giving $\mathrm{CV}^2_{\mathrm{mix}} = 2 - 1 = 1$.

Conversely, if $\mathrm{CV}^2_{\mathrm{mix}} = 1$, then from the mixture identity, the term $p(1-p)(a-b)^2$ must vanish. Since $0 < p < 1$, we have $p(1-p) > 0$, so $(a-b)^2 = 0$ and $a = b$.
\end{proof}

\begin{theorem}[\lean{Particle.observed\_disfavours\_higher\_rank}]
The observed spectrum has $\mathrm{CV}^2 < 0.88$. Since every mixture of two or more rates satisfies $\mathrm{CV}^2 \geq 1$, the observed spectrum is evidence against a higher-rank algebra and in favour of rank one.
\end{theorem}

\begin{proof}
Let $\mathrm{CV}^2_{\mathrm{obs}} < 0.88$ and $\mathrm{CV}^2_{\mathrm{mix}} \geq 1$ for any mixture. Then $\mathrm{CV}^2_{\mathrm{obs}} < \mathrm{CV}^2_{\mathrm{mix}}$.
\end{proof}

\begin{theorem}[\lean{Particle.rate\_value\_not\_fixed}]
The value of the exponential rate is not determined by the framework. The coefficient of variation is invariant under rescaling both component means: $\mathrm{CV}^2_{\mathrm{mix}}(p, ca, cb) = \mathrm{CV}^2_{\mathrm{mix}}(p, a, b)$ for all $c \neq 0$.
\end{theorem}

\begin{proof}
Scaling both means by $c$ scales both the numerator and denominator by $c^2$, leaving the ratio unchanged.
\end{proof}

\begin{definition}[\lean{Particle.Block}]
The block group of the uniaxial pattern is the multiplicative group on $\mathbb{Z}/2$, denoted $\mathrm{Multiplicative}(\mathbb{Z}/2)$.
\end{definition}

\begin{definition}[\lean{Particle.Charge}]
A \emph{particle charge} is a defect charge in the block group $\mathrm{Multiplicative}(\mathbb{Z}/2)$.
\end{definition}

\begin{theorem}[\lean{Particle.block\_two\_valued}]
The block group has cardinality 2.
\end{theorem}

\begin{proof}
$\mathbb{Z}/2$ has two elements, and the multiplicative group on it has the same cardinality.
\end{proof}

\begin{theorem}[\lean{Particle.block\_cannot\_be\_three\_valued}]
The framework structurally cannot produce a three-valued confined charge: the block group is $\mathbb{Z}/2$, which has exactly 2 elements, not 3.
\end{theorem}

\begin{proof}
The order parameter is projective (unoriented axis), and the loop homotopy group of projective space is $\mathbb{Z}/2$. This is a topological fact, not a free parameter.
\end{proof}

\begin{theorem}[\lean{Particle.block\_self\_inverse}]
Every element of the block group is its own inverse: $b \cdot b = 1$ for all $b \in \mathrm{Block}$.
\end{theorem}

\begin{proof}
In $\mathbb{Z}/2$, every nonidentity element satisfies $b^2 = 1$.
\end{proof}

\begin{definition}[\lean{Particle.Species}]
A \emph{particle species} is determined by:
\begin{itemize}
\item A threshold $t \in \mathbb{R}$, the log-scale at which the particle becomes resolvable.
\item A charge $(b, q) \in \mathrm{Charge}$, specifying the $\mathbb{Z}/2$ label and the integer charge.
\end{itemize}
The framework supplies no further labels.
\end{definition}

\begin{definition}[\lean{Particle.bind}]
Two particles with charges $c_1, c_2$ can be bound at a new threshold $t$ to form a composite particle with the same threshold and charge $\mathrm{comp}(c_1, c_2)$.
\end{definition}

\begin{theorem}[\lean{Particle.bind\_hedgehog}]
The hedgehog charge of a composite is the sum of the components' hedgehog charges.
\end{theorem}

\begin{proof}
By definition of $\mathrm{comp}$.
\end{proof}

\begin{theorem}[\lean{Particle.bind\_block}]
The block charge of a composite is the product of the components' block charges.
\end{theorem}

\begin{proof}
By definition of $\mathrm{comp}$.
\end{proof}

\begin{theorem}[\lean{Particle.two\_odd\_make\_even}]
Two $\mathbb{Z}/2$-odd particles (same block charge) compose to an even particle (trivial block charge), without needing any dynamics: $\mathrm{bind}(x, y, t).\mathrm{charge}.\mathrm{block} = 1$ whenever $x.\mathrm{charge}.\mathrm{block} = y.\mathrm{charge}.\mathrm{block}$.
\end{theorem}

\begin{proof}
If both have block charge $b$, the composite has block charge $b \cdot b = 1$ by \lean{Particle.block\_self\_inverse}.
\end{proof}

\begin{theorem}[\lean{Particle.hedgehog\_conserved\_in\_splitting}]
The hedgehog charge is exactly conserved in any splitting: if a particle splits into two, the parts' hedgehog charges sum to the whole's. This is an identity, independent of dynamics.
\end{theorem}

\begin{proof}
If the whole has charge $(b, q)$ and splits into $(b_1, q_1)$ and $(b_2, q_2)$ with $b = b_1 \cdot b_2$ and $q = q_1 + q_2$, then $q - q_1 = q_2$, which is identically true.
\end{proof}

\begin{theorem}[\lean{Particle.no\_further\_label}]
A particle is determined by its threshold and its two charges: two particles agreeing on all three are equal.
\end{theorem}

\begin{proof}
The species structure has two fields (threshold and charge), and the charge has two components (block and hedgehog). Two species with equal thresholds and charges are equal by extensionality.
\end{proof}

\begin{theorem}[\lean{Particle.particle\_label\_structure}]
Every particle has the form $(t, b, q)$ for some threshold $t \in \mathbb{R}$, block label $b \in \mathbb{Z}/2$, and integer charge $q \in \mathbb{Z}$.
\end{theorem}

\begin{proof}
The species and charge structures decompose as stated.
\end{proof}

\subsection{Emergence}

Particles are not fundamental: they are features of a slice across the scale axis. By axiom A3, energy and mass are inverse scale; so the set of observable particles depends on the resolution at which the universe is probed. The framework makes no claim about what particles exist, only that at each scale a new content emerges, growing monotonically without bound. This module characterizes the threshold structure and establishes why there is no complete list of ``fundamental particles.''

\begin{definition}[\lean{Emergence.resolved}]
The set of particles resolvable at log-scale $t$ is $\{k \mid \mu_k \leq t\}$, where $\mu$ is the threshold sequence.
\end{definition}

\begin{theorem}[\lean{Emergence.mem\_resolved}]
An index $k$ is in the resolved set at scale $t$ if and only if its threshold satisfies $\mu_k \leq t$.
\end{theorem}

\begin{proof}
Immediate from the definition.
\end{proof}

\begin{theorem}[\lean{Emergence.resolved\_mono}]
Content only grows with better resolution: if $t \leq t'$, then $\mathrm{resolved}(\mu, t) \subseteq \mathrm{resolved}(\mu, t')$.
\end{theorem}

\begin{proof}
If $k \in \mathrm{resolved}(\mu, t)$, then $\mu_k \leq t \leq t'$, so $k \in \mathrm{resolved}(\mu, t')$.
\end{proof}

\begin{theorem}[\lean{Emergence.content\_never\_complete}]
If the thresholds are unbounded (for every upper bound $M$ there exists a threshold above it), then no scale sees all particles: $\mathrm{resolved}(\mu, t) \neq \text{all indices}$ for any $t$.
\end{theorem}

\begin{proof}
Given $t$, there exists $k$ with $t < \mu_k$, so $k \notin \mathrm{resolved}(\mu, t)$.
\end{proof}

\begin{theorem}[\lean{Emergence.always\_more\_above}]
The tower never terminates: for any scale $t$, there exists a scale $t' > t$ at which strictly more particles are resolved.
\end{theorem}

\begin{proof}
Given $t$, let $k$ satisfy $t < \mu_k$ (which exists by unboundedness). Then $t < \mu_k \leq t'$ for $t' = \mu_k$ shows that $k \in \mathrm{resolved}(\mu, t')$ but $k \notin \mathrm{resolved}(\mu, t)$.
\end{proof}

\begin{theorem}[\lean{Emergence.no\_fundamental\_list}]
The content is a never-completed, strictly increasing family: for all scales $t$, the resolved content differs from the full set, and there always exists a higher scale with strictly larger content.
\end{theorem}

\begin{proof}
Follows from the previous two theorems. The question ``which particles are fundamental?'' has no scale-independent answer; it is a malformed question in this framework.
\end{proof}

\begin{definition}[\lean{Emergence.massOf}]
The mass of a particle with threshold $\mu_k$ is $m_k = e^{\mu_k}$. By axiom A3, energy and mass are inverse scale, so higher threshold means heavier particle.
\end{definition}

\begin{theorem}[\lean{Emergence.mass\_iff\_threshold}]
Mass and threshold determine each other: $\mathrm{massOf}(\mu, k) = \mathrm{massOf}(\mu, l) \iff \mu_k = \mu_l$.
\end{theorem}

\begin{proof}
The exponential function is injective: $e^x = e^y \iff x = y$.
\end{proof}

\begin{theorem}[\lean{Emergence.mass\_mono}]
Heavier particles appear at higher thresholds: if $\mu_k < \mu_l$, then $\mathrm{massOf}(\mu, k) < \mathrm{massOf}(\mu, l)$.
\end{theorem}

\begin{proof}
The exponential function is strictly increasing.
\end{proof}

\begin{theorem}[\lean{Emergence.content\_eq\_of\_no\_threshold}]
If there is no threshold in the interval $(t, t']$, then the resolved content at $t$ and $t'$ is identical: the effective description is exact between thresholds.
\end{theorem}

\begin{proof}
For any index $k$, if $\mu_k \leq t$, then $\mu_k \leq t'$, so $k$ remains resolved. If $\mu_k > t'$, then by the no-threshold assumption, $\mu_k$ is not in $(t, t']$, so $\mu_k > t'$ and $k$ is not resolved at either scale. No index crosses the gap.
\end{proof}

\begin{theorem}[\lean{Emergence.content\_changes\_only\_at\_threshold}]
Content changes between two scales if and only if at least one threshold lies in the gap: if $\mathrm{resolved}(\mu, t) \neq \mathrm{resolved}(\mu, t')$ for $t \leq t'$, then there exists $k$ with $t < \mu_k \leq t'$.
\end{theorem}

\begin{proof}
If no such $k$ exists, then by the previous theorem the contents are equal, giving a contradiction. Therefore at least one threshold lies in the gap.
\end{proof}


\section{Part III: The Quantum Sector}

\subsection{Deformation}

First-order deformation quantization bridges the commutative and non-commutative sectors. The classical limit is recovered exactly (not as an approximation) by evaluating at zero in the deformation parameter $\hbar$. The commutator of the deformed algebra encodes the Poisson bracket of the original commutative algebra, establishing that classical derivations arise as inner derivations of the quantum algebra divided by $\hbar$. The canonical commutation relation then closes the circle by showing that momentum and position satisfy $[p, q]_\star = \hbar$ in this deformation.

\begin{definition}
A \textbf{biderivation} of a commutative ring $A$ is a pair $(B, \text{add\_left}, \text{add\_right}, \text{leibniz\_left}, \text{leibniz\_right})$ where $B : A \to A \to A$ satisfies:
\begin{align}
B(a+b,c) &= B(a,c) + B(b,c) \quad \text{(additive in first argument)} \\
B(a,b+c) &= B(a,b) + B(a,c) \quad \text{(additive in second argument)} \\
B(ab,c) &= a \cdot B(b,c) + b \cdot B(a,c) \quad \text{(Leibniz in first argument)} \\
B(a,bc) &= b \cdot B(a,c) + c \cdot B(a,b) \quad \text{(Leibniz in second argument)}
\end{align}
for all $a, b, c \in A$. This is the algebraic data of a first-order deformation.
\end{definition}

\begin{definition}
Given a biderivation $P$ of $A$, the \textbf{Poisson bracket} induced by $P$ is the antisymmetric part of the bilinear map:
\[
\{a, b\}_P := P.B(a,b) - P.B(b,a) .
\]
\lean{Deformation.Biderivation.poisson}
\end{definition}

\begin{theorem}[Antisymmetry of the Poisson bracket]
For any biderivation $P$ of $A$ and elements $a, b \in A$,
\[
\{a, b\}_P = -\{b, a\}_P .
\]
\begin{proof}
Immediate from the definition of antisymmetric subtraction.
\end{proof}
\lean{Deformation.Biderivation.poisson\_antisymm}
\end{theorem}

\begin{theorem}[Poisson bracket vanishes on diagonal]
For any biderivation $P$ of $A$ and element $a \in A$,
\[
\{a, a\}_P = 0 .
\]
\begin{proof}
Since $\{a,a\}_P = P.B(a,a) - P.B(a,a) = 0$ by definition.
\end{proof}
\lean{Deformation.Biderivation.poisson\_self}
\end{theorem}

\begin{theorem}[Poisson bracket is a derivation in the first argument]
For any biderivation $P$ of $A$ and elements $a, b, c \in A$,
\[
\{ab, c\}_P = a \cdot \{b,c\}_P + b \cdot \{a,c\}_P .
\]
This is the precise sense in which the classical derivations of axiom A1 are recovered from the deformation: each map $\{-, g\}_P$ is a derivation of the commutative ring $A$.
\begin{proof}
We compute:
\begin{align}
\{ab, c\}_P &= P.B(ab,c) - P.B(c,ab) \\
&= [a \cdot P.B(b,c) + b \cdot P.B(a,c)] - [c \cdot P.B(a,b) + c \cdot P.B(b,a)] \\
&= a \cdot [P.B(b,c) - P.B(c,b)] + b \cdot [P.B(a,c) - P.B(c,a)] \\
&= a \cdot \{b,c\}_P + b \cdot \{a,c\}_P ,
\end{align}
using the Leibniz properties and rearranging.
\end{proof}
\lean{Deformation.Biderivation.poisson\_leibniz\_left}
\end{theorem}

\begin{theorem}[Poisson bracket is a derivation in the second argument]
For any biderivation $P$ of $A$ and elements $a, b, c \in A$,
\[
\{a, bc\}_P = b \cdot \{a,c\}_P + c \cdot \{a,b\}_P .
\]
\begin{proof}
Similar to the first-argument case, using the Leibniz properties in the second argument.
\end{proof}
\lean{Deformation.Biderivation.poisson\_leibniz\_right}
\end{theorem}

\begin{definition}
The \textbf{star product to first order in $\hbar$} defined by a biderivation $P$ is the product on formal series $A[\hbar]/(\hbar^2) \cong A \times A$ (identifying $f_0 + \hbar f_1$ with the pair $(f_0, f_1)$) given by:
\[
(f_0, f_1) \star (g_0, g_1) := \bigl(f_0 g_0, \, f_0 g_1 + f_1 g_0 + P.B(f_0, g_0)\bigr) .
\]
This product is independent of the choice of basis and depends only on $P$.
\lean{Deformation.star}
\end{definition}

\begin{theorem}[The deformation is associative]
For any biderivation $P$ and elements $f, g, h \in A \times A$,
\[
(f \star g) \star h = f \star (g \star h) .
\]
The Leibniz axioms for $P$ are precisely the Hochschild cocycle condition at first order, ensuring that $A[\hbar]/(\hbar^2)$ is a genuine associative algebra.
\begin{proof}
Expand both sides using the star product definition. The first component is $f_0 g_0 h_0$ on both sides by associativity of multiplication in $A$. For the second component, the Leibniz axioms in both arguments of $P$ ensure that the $\hbar$-coefficient also matches after collecting terms. The calculation is a verification that the cocycle condition holds.
\end{proof}
\lean{Deformation.star\_assoc}
\end{theorem}

\begin{theorem}[The classical limit]
The constant term of the star product recovers the ordinary commutative product:
\[
(f_0, f_1) \star (g_0, g_1) \big|_{\text{coeff of } \hbar^0} = f_0 g_0 .
\]
There is no approximation or limit process; this is exact at first order.
\begin{proof}
Immediate from the definition: the first component of the product is always $f_0 g_0$.
\end{proof}
\lean{Deformation.star\_classical}
\end{theorem}

\begin{theorem}[Commutativity to zeroth order]
The deformed product is commutative at the zeroth order in $\hbar$:
\[
(f_0, f_1) \star (g_0, g_1) \big|_{\hbar^0} = (g_0, g_1) \star (f_0, f_1) \big|_{\hbar^0} .
\]
That is, $f_0 g_0 = g_0 f_0$, which is immediate since $A$ is commutative.
\begin{proof}
This follows from the classical limit theorem: both sides reduce to $f_0 g_0$ by commutativity of $A$.
\end{proof}
\lean{Deformation.star\_comm\_classical}
\end{theorem}

\begin{theorem}[The commutator encodes the Poisson bracket]
For any biderivation $P$ and elements $f, g \in A \times A$, the commutator
\[
[f, g]_\star := f \star g - g \star f
\]
has the form:
\begin{align}
\text{(coefficient of $\hbar^0$)} &: \quad (f_0, f_1) \star (g_0, g_1) \big|_{\hbar^0} - (g_0, g_1) \star (f_0, f_1) \big|_{\hbar^0} = 0 , \\
\text{(coefficient of $\hbar^1$)} &: \quad (f_0, f_1) \star (g_0, g_1) \big|_{\hbar^1} - (g_0, g_1) \star (f_0, f_1) \big|_{\hbar^1} = \{f_0, g_0\}_P .
\end{align}
Thus $\hbar^{-1}[f, g]_\star = \{f_0, g_0\}_P$ exactly, with no limiting process to justify. This is the statement that $\partial_i = \hbar^{-1} \text{ad}_{P_i}$ in the classical sector.
\begin{proof}
The first component is $f_0 g_0 - g_0 f_0 = 0$ since $A$ is commutative. The second component is:
\begin{align}
&[f_0 g_1 + f_1 g_0 + P.B(f_0, g_0)] - [g_0 f_1 + g_1 f_0 + P.B(g_0, f_0)] \\
&= P.B(f_0, g_0) - P.B(g_0, f_0) = \{f_0, g_0\}_P .
\end{align}
\end{proof}
\lean{Deformation.star\_commutator}
\end{theorem}

\begin{theorem}[Commutativity iff the Poisson bracket vanishes]
Two elements $f, g \in A \times A$ commute in the star product if and only if their Poisson bracket vanishes:
\[
f \star g = g \star f \quad \iff \quad \{f_0, g_0\}_P = 0 .
\]
Non-commutativity is exactly the failure of $P$ to be symmetric; a vanishing bracket recovers a commutative algebra (the gravitational sector).
\begin{proof}
From the commutator formula, $f \star g - g \star f$ has zero second component if and only if $\{f_0, g_0\}_P = 0$. Since the first component is always zero, this holds if and only if the products coincide.
\end{proof}
\lean{Deformation.star\_comm\_iff\_poisson\_zero}
\end{theorem}

\begin{definition}
For a scale algebra with $n$ directions and derivations $(d_i)_{i \in \text{Fin } n}$, the \textbf{canonical biderivation} associated to two directions $p, q \in \text{Fin } n$ is defined by:
\[
P_{pq}.B(a, b) := d_q(a) \cdot d_p(b) .
\]
A pair of canonically conjugate coordinates of the substrate — one varying only along $p$, the other only along $q$ — gives rise to a biderivation, hence a deformation.
\lean{Deformation.canonical}
\end{definition}

\begin{theorem}[The Poisson bracket of canonical coordinates]
For directions $p, q \in \text{Fin } n$ and elements $a, b$ in the scale algebra,
\[
\{a, b\}_{P_{pq}} = d_q(a) \cdot d_p(b) - d_q(b) \cdot d_p(a) .
\]
\begin{proof}
Immediate from the definition of the canonical biderivation and the Poisson bracket.
\end{proof}
\lean{Deformation.canonical\_poisson}
\end{theorem}

\begin{theorem}[The canonical commutation relation]
Let $X$ and $\text{Mom}$ be coordinates of the substrate that are canonically conjugate with respect to directions $p, q$:
\begin{align}
d_p(X) = 1, \quad d_q(X) = 0, \quad d_p(\text{Mom}) = 0, \quad d_q(\text{Mom}) = 1 .
\end{align}
Then their Poisson bracket is
\[
\{\text{Mom}, X\}_{P_{pq}} = 1 ,
\]
which means in the deformed algebra, $[\text{Mom}, X]_\star = \hbar \cdot 1$. Taking $\hbar = i\hbar$ recovers the canonical commutation relation $[\text{Mom}, X] = i\hbar$ of quantum mechanics. Combined with the power rule of calculus (recovered via inner derivations in the commutative sector), the circle is closed: the differential structure of the gravitational sector and the canonical commutation relation are two ends of one first-order deformation.
\begin{proof}
Applying the definition:
\begin{align}
\{\text{Mom}, X\}_{P_{pq}} &= d_q(\text{Mom}) \cdot d_p(X) - d_q(X) \cdot d_p(\text{Mom}) \\
&= 1 \cdot 1 - 0 \cdot 0 = 1 .
\end{align}
\end{proof}
\lean{Deformation.ccr\_of\_canonical}
\end{theorem}

\begin{theorem}[Poisson bracket of a coordinate with itself]
For any directions $p, q$ and element $a$,
\[
\{a, a\}_{P_{pq}} = 0 .
\]
\begin{proof}
This is the diagonal case of the Poisson bracket antisymmetry.
\end{proof}
\lean{Deformation.canonical\_poisson\_self}
\end{theorem}

\subsection{Crossed}

Non-commutativity arises exactly from the scale step itself, with no expansion parameter. The substrate has two operations: reading a quantity where you are (multiplication) and moving to the adjacent scale (the shift endomorphism $T$). These do not commute; their failure to commute is not a small parameter but a finite, exact quantity—the scale difference $\delta a = Ta - a$. Position and scale do not commute, which is the operator form of axiom A3. Everything follows exactly, with the classical differential structure emerging as the zero-scale-step limit.

\begin{definition}
A \textbf{scale shift} is a ring endomorphism that re-reads the substrate one step along the scale direction. Formally, it is a triple $(T, \text{map\_add}, \text{map\_mul})$ where $T : A \to A$ preserves addition and multiplication:
\begin{align}
T(x + y) &= T(x) + T(y), \\
T(xy) &= T(x) \cdot T(y) .
\end{align}
Only the ring structure is preserved; nothing is assumed about the size of the step.
\lean{Crossed.ScaleShift}
\end{definition}

\begin{theorem}[Scale shift preserves zero]
For any scale shift $S$,
\[
S.T(0) = 0 .
\]
\begin{proof}
From $S.T(0 + 0) = S.T(0) + S.T(0)$, we get $S.T(0) = 2 \cdot S.T(0)$, whence $S.T(0) = 0$ by cancellation.
\end{proof}
\lean{Crossed.ScaleShift.map\_zero}
\end{theorem}

\begin{definition}
The shift as an additive endomorphism $\text{shiftOp}(S) \in \text{End}_{\text{add}}(A)$ is the formal operator on the substrate where moving to an adjacent scale is a separate operation from multiplication.
\lean{Crossed.ScaleShift.shiftOp}
\end{definition}

\begin{theorem}[Action of shiftOp]
For any scale shift $S$ and element $x \in A$,
\[
\text{shiftOp}(S) \cdot x = S.T(x) .
\]
\begin{proof}
By definition of the additive endomorphism structure.
\end{proof}
\lean{Crossed.ScaleShift.shiftOp\_apply}
\end{theorem}

\begin{definition}
Reading a quantity where you are is represented by left multiplication in the operator ring: for $a \in A$,
\[
\text{mulOp}(a) \cdot x := a \cdot x .
\]
\lean{Crossed.mulOp}
\end{definition}

\begin{theorem}[Action of mulOp]
For any elements $a, x \in A$,
\[
\text{mulOp}(a) \cdot x = a \cdot x .
\]
\begin{proof}
Immediate from the definition as left multiplication.
\end{proof}
\lean{Crossed.mulOp\_apply}
\end{theorem}

\begin{theorem}[The exchange relation]
Moving to the adjacent scale and then reading $a$ is the same as reading $T(a)$ and then moving:
\[
\text{shiftOp}(S) \circ \text{mulOp}(a) = \text{mulOp}(T(a)) \circ \text{shiftOp}(S) .
\]
This is exact; no approximation is involved.
\begin{proof}
Applying both sides to $x \in A$:
\begin{align}
\text{LHS}(x) &= \text{shiftOp}(S)(\text{mulOp}(a)(x)) = S.T(a \cdot x) = S.T(a) \cdot S.T(x) , \\
\text{RHS}(x) &= \text{mulOp}(T(a))(\text{shiftOp}(S)(x)) = T(a) \cdot S.T(x) = S.T(a) \cdot S.T(x) .
\end{align}
\end{proof}
\lean{Crossed.shift\_mul\_eq}
\end{theorem}

\begin{definition}
The \textbf{scale difference} operator measures how much a quantity changes over one scale step:
\[
\delta_S(a) := S.T(a) - a .
\]
At finite step this is the exact analogue of a derivative.
\lean{Crossed.delta}
\end{definition}

\begin{theorem}[The commutator in closed form]
The commutator of the shift operator with multiplication is
\[
[\text{shiftOp}(S), \text{mulOp}(a)] := \text{shiftOp}(S) \circ \text{mulOp}(a) - \text{mulOp}(a) \circ \text{shiftOp}(S) = \text{mulOp}(\delta_S(a)) \circ \text{shiftOp}(S) .
\]
Writing this in operator form: $[U, M_a] = M_{\delta a} \cdot U$. Non-commutativity is not a small parameter; it is the scale difference, exactly. Planck's constant $\hbar$ is the size of a scale step.
\begin{proof}
From the exchange relation, $\text{shiftOp}(S) \circ \text{mulOp}(a) = \text{mulOp}(S.T(a)) \circ \text{shiftOp}(S)$. Thus:
\begin{align}
[\text{shiftOp}(S), \text{mulOp}(a)](x) &= S.T(a) \cdot S.T(x) - a \cdot S.T(x) \\
&= (S.T(a) - a) \cdot S.T(x) = \delta_S(a) \cdot S.T(x) ,
\end{align}
which is the action of $\text{mulOp}(\delta_S(a)) \circ \text{shiftOp}(S)$.
\end{proof}
\lean{Crossed.commutator\_eq}
\end{theorem}

\begin{theorem}[Commutativity is absence of scale step]
The two operations commute if and only if the scale difference vanishes:
\[
\text{shiftOp}(S) \circ \text{mulOp}(a) = \text{mulOp}(a) \circ \text{shiftOp}(S) \quad \iff \quad \text{mulOp}(\delta_S(a)) \circ \text{shiftOp}(S) = 0 .
\]
Classical geometry is not a limit one takes; it is the case in which nothing moves along the scale.
\begin{proof}
The commutator vanishes if and only if the right-hand side of the commutator formula is zero, which happens if and only if $\delta_S(a) = 0$ for all $a$.
\end{proof}
\lean{Crossed.comm\_iff\_shift\_trivial}
\end{theorem}

\begin{theorem}[The twisted Leibniz rule at finite scale step]
For any scale shift $S$ and elements $a, b \in A$,
\[
\delta_S(ab) = \delta_S(a) \cdot S.T(b) + a \cdot \delta_S(b) .
\]
The rule is twisted: the second factor is read at the shifted scale. This is exact for any step, however large.
\begin{proof}
\begin{align}
\delta_S(ab) &= S.T(ab) - ab \\
&= S.T(a) \cdot S.T(b) - ab \\
&= [S.T(a) - a] \cdot S.T(b) + a \cdot [S.T(b) - b] \\
&= \delta_S(a) \cdot S.T(b) + a \cdot \delta_S(b) .
\end{align}
\end{proof}
\lean{Crossed.delta\_twisted\_leibniz}
\end{theorem}

\begin{theorem}[Axiom A1 is the zero-step case]
When the shift acts trivially on the second factor—that is, $S.T(b) = b$—the twisted Leibniz rule reduces to the ordinary product rule:
\[
\delta_S(ab) = \delta_S(a) \cdot b + a \cdot \delta_S(b) .
\]
So the classical differential structure postulated in axiom A1 is not an expansion truncated at first order in $\hbar$; it is the exact relation evaluated at zero scale step.
\begin{proof}
Substitute $S.T(b) = b$ into the twisted rule.
\end{proof}
\lean{Crossed.delta\_leibniz\_of\_untwisted}
\end{theorem}

\begin{theorem}[Scale difference vanishes iff shift is trivial]
For a scale shift $S$ and element $a$,
\[
\delta_S(a) = 0 \quad \iff \quad S.T(a) = a .
\]
\begin{proof}
Immediate from the definition $\delta_S(a) = S.T(a) - a$.
\end{proof}
\lean{Crossed.delta\_eq\_zero\_iff}
\end{theorem}

\begin{definition}
Composing two scale shifts gives another scale shift. The composition $(S \circ R).T := S.T \circ R.T$ preserves addition and multiplication because both $S.T$ and $R.T$ do.
\lean{Crossed.ScaleShift.comp}
\end{definition}

\begin{theorem}[Composition applies as function composition]
For scale shifts $S, R$ and element $x$,
\[
(S \circ R).T(x) = S.T(R.T(x)) .
\]
\begin{proof}
By definition of composition.
\end{proof}
\lean{Crossed.ScaleShift.comp\_apply}
\end{theorem}

\begin{theorem}[Exchange relation holds for composed steps]
The exchange relation holds for arbitrary compositions of scale shifts:
\[
\text{shiftOp}(S \circ R) \circ \text{mulOp}(a) = \text{mulOp}((S \circ R).T(a)) \circ \text{shiftOp}(S \circ R) .
\]
The structure is consistent across arbitrary scale separations rather than only infinitesimal ones.
\begin{proof}
This is immediate from the exchange relation theorem applied to the composed shift $S \circ R$.
\end{proof}
\lean{Crossed.shift\_mul\_eq\_comp}
\end{theorem}

\subsection{NCConformal}

The quantum-gravitational correction enters the geometry as a purely antisymmetric contribution to the deformation tensor, arising from the non-commutativity of scale gradients. The entire correction is antisymmetric in the index pair, so it vanishes from the symmetric part of the deformation tensor—and hence from every classical gravitational observable (orbit deflection, perihelion precession, gravitational redshift). The prediction is that quantum-gravitational effects appear first in spin–geometry coupling, which is sensitive to antisymmetric geometric structure.

\begin{definition}
The Kronecker delta is represented as a central element (cast of $0$ or $1$):
\[
\kappa(i,j) := \begin{cases} 1 & \text{if } i = j \\ 0 & \text{otherwise} \end{cases} \in M .
\]
\lean{NCConformal.kr}
\end{definition}

\begin{theorem}[Kronecker delta is symmetric]
For all $i, j \in \text{Fin } n$,
\[
\kappa(i,j) = \kappa(j,i) .
\]
\begin{proof}
If $i = j$ then $j = i$ and both sides are $1$. Otherwise both are $0$.
\end{proof}
\lean{NCConformal.kr\_symm}
\end{theorem}

\begin{theorem}[Kronecker delta is central]
For all $i, j \in \text{Fin } n$ and $x \in M$,
\[
\kappa(i,j) \cdot x = x \cdot \kappa(i,j) .
\]
Being $0$ or $1$, it commutes with every element.
\begin{proof}
When $i = j$, both sides are $x$. When $i \ne j$, both sides are $0$.
\end{proof}
\lean{NCConformal.kr\_comm}
\end{theorem}

\begin{theorem}[Diagonal Kronecker delta]
For all $i \in \text{Fin } n$,
\[
\kappa(i,i) = 1 .
\]
\begin{proof}
By definition.
\end{proof}
\lean{NCConformal.kr\_self}
\end{theorem}

\begin{theorem}[Summing against Kronecker delta on the left]
For $i \in \text{Fin } n$ and a function $f : \text{Fin } n \to M$,
\[
\sum_{m} \kappa(i,m) \cdot f(m) = f(i) .
\]
\begin{proof}
Only the term $m = i$ contributes, where $\kappa(i,i) = 1$, so the sum is $f(i)$.
\end{proof}
\lean{NCConformal.sum\_kr\_left}
\end{theorem}

\begin{theorem}[Summing against Kronecker delta on the right]
For $i \in \text{Fin } n$ and a function $f : \text{Fin } n \to M$,
\[
\sum_{m} \kappa(m,i) \cdot f(m) = f(i) .
\]
\begin{proof}
Since $\kappa(m,i) = \kappa(i,m)$, this reduces to the left case.
\end{proof}
\lean{NCConformal.sum\_kr\_right}
\end{theorem}

\begin{theorem}[Kronecker deltas commute]
For indices $i, j, p, q \in \text{Fin } n$ and $x \in M$,
\[
\kappa(i,j) \cdot (\kappa(p,q) \cdot x) = \kappa(p,q) \cdot (\kappa(i,j) \cdot x) .
\]
\begin{proof}
Since both $\kappa$ factors are central, they pass through each other and reorder freely.
\end{proof}
\lean{NCConformal.kr\_swap}
\end{theorem}

\begin{theorem}[Scalar-prefixed factors merge]
For indices $i, j, p, q \in \text{Fin } n$ and $x, y \in M$,
\[
(\kappa(i,j) \cdot x) \cdot (\kappa(p,q) \cdot y) = (\kappa(i,j) \cdot \kappa(p,q)) \cdot (x \cdot y) .
\]
\begin{proof}
The $\kappa$ factors are central, so they collect on the left and the ring elements keep their order.
\end{proof}
\lean{NCConformal.scalar\_mul\_scalar}
\end{theorem}

\begin{definition}
The \textbf{scale gradient} in direction $i$ is the derivation $D_i$ applied to the scale field $\sigma$:
\[
\sigma_i := D_i(\sigma) .
\]
\lean{NCConformal.sig}
\end{definition}

\begin{definition}
The \textbf{scale Hessian}—second derivatives of the scale field—is:
\[
H_{ij} := D_i(D_j(\sigma)) = D_i(\sigma_j) .
\]
\lean{NCConformal.hess}
\end{definition}

\begin{theorem}[The scale Hessian is symmetric]
For all $i, j \in \text{Fin } n$,
\[
H_{ij} = H_{ji} .
\]
This symmetry comes from the commutativity of the derivations (a chart choice in axiom A1) and survives untouched in the non-commutative setting.
\begin{proof}
Since $D_i$ and $D_j$ commute by assumption (axiom A1 for commuting derivations),
\[
H_{ij} = D_i(D_j(\sigma)) = D_j(D_i(\sigma)) = H_{ji} .
\]
\end{proof}
\lean{NCConformal.hess\_symm}
\end{theorem}

\begin{definition}
The \textbf{squared gradient} is the sum of squares of all scale gradients:
\[
|\nabla \sigma|^2 := \sum_{i \in \text{Fin } n} \sigma_i \cdot \sigma_i .
\]
\lean{NCConformal.gradsq}
\end{definition}

\begin{definition}
The \textbf{Christoffel symbols} in the non-commutative setting are:
\[
\Gamma^a_{bc} := \kappa(a,b) \cdot \sigma_c + \kappa(a,c) \cdot \sigma_b - \kappa(b,c) \cdot \sigma_a .
\]
They have the same form as in the commutative case because the Kronecker deltas are central.
\lean{NCConformal.Chr}
\end{definition}

\begin{theorem}[Christoffel symbols are symmetric in lower indices]
For all $a, b, c \in \text{Fin } n$,
\[
\Gamma^a_{bc} = \Gamma^a_{cb} .
\]
\begin{proof}
Since $\kappa(b,c) = \kappa(c,b)$ (symmetry of Kronecker delta), swapping $b$ and $c$ leaves the expression unchanged.
\end{proof}
\lean{NCConformal.Chr\_symm}
\end{theorem}

\begin{definition}
The \textbf{deformation tensor}, which encodes curvature, is:
\[
\text{Defm}_{ij} := 2 H_{ij} - 2 \sigma_i \cdot \sigma_j + \kappa(i,j) \cdot |\nabla \sigma|^2 .
\]
It is defined exactly as in the commutative case, but now $\sigma_i$ and $\sigma_j$ need not commute.
\lean{NCConformal.Defm}
\end{definition}

\begin{theorem}[The quantum-gravitational correction is antisymmetric]
The deformation tensor is no longer symmetric. The antisymmetric part—the whole quantum correction—is exactly twice the commutator of the scale gradients:
\[
\text{Defm}_{ij} - \text{Defm}_{ji} = -2 [\sigma_i, \sigma_j] ,
\]
where $[\sigma_i, \sigma_j] := \sigma_i \cdot \sigma_j - \sigma_j \cdot \sigma_i$ is the commutator. Everything else in the curvature computation is unchanged, because the Kronecker deltas are central and $H_{ij}$ is still symmetric.
\begin{proof}
\begin{align}
\text{Defm}_{ij} - \text{Defm}_{ji} &= [2 H_{ij} - 2 \sigma_i \cdot \sigma_j + \kappa(i,j) \cdot |\nabla \sigma|^2] \\
&\quad - [2 H_{ji} - 2 \sigma_j \cdot \sigma_i + \kappa(j,i) \cdot |\nabla \sigma|^2] \\
&= 2[H_{ij} - H_{ji}] - 2[\sigma_i \cdot \sigma_j - \sigma_j \cdot \sigma_i] + 0 \\
&= 0 - 2[\sigma_i, \sigma_j] = -2[\sigma_i, \sigma_j] ,
\end{align}
using the symmetry of the Hessian and Kronecker delta.
\end{proof}
\lean{NCConformal.Defm\_antisymm\_part}
\end{theorem}

\begin{theorem}[Symmetry is equivalent to commuting gradients]
The deformation tensor is symmetric if and only if twice the commutator vanishes:
\[
\text{Defm}_{ij} = \text{Defm}_{ji} \quad \iff \quad 2 \cdot [\sigma_i, \sigma_j] = 0 .
\]
The classical geometry is exactly the commuting case—not an approximation to it, but the precise locus where the correction vanishes.
\begin{proof}
This follows directly from the antisymmetry formula: $\text{Defm}_{ij} = \text{Defm}_{ji}$ if and only if $-2[\sigma_i, \sigma_j] = 0$.
\end{proof}
\lean{NCConformal.Defm\_symm\_iff\_commutator}
\end{theorem}

\begin{theorem}[Clean statement when 2 is invertible]
When the ring $M$ has $2$ invertible, symmetry of the deformation tensor is equivalent to commutativity of the scale gradients:
\[
\text{Defm}_{ij} = \text{Defm}_{ji} \quad \iff \quad \sigma_i \cdot \sigma_j = \sigma_j \cdot \sigma_i .
\]
\begin{proof}
From the previous theorem, symmetry means $2 \cdot [\sigma_i, \sigma_j] = 0$. When $2$ is invertible, this is equivalent to $[\sigma_i, \sigma_j] = 0$, i.e., $\sigma_i \cdot \sigma_j = \sigma_j \cdot \sigma_i$.
\end{proof}
\lean{NCConformal.Defm\_symm\_iff\_commute}
\end{theorem}

\begin{theorem}[The symmetric part is uncorrected]
The symmetric part of the deformation tensor is:
\[
\text{Defm}_{ij} + \text{Defm}_{ji} = 4 H_{ij} - 2(\sigma_i \cdot \sigma_j + \sigma_j \cdot \sigma_i) + 2 \kappa(i,j) \cdot |\nabla \sigma|^2 .
\]
This is exactly what appears classically; the quantum correction contributes nothing to it at any order.
\begin{proof}
Adding the formulas:
\begin{align}
\text{Defm}_{ij} + \text{Defm}_{ji} &= [2 H_{ij} - 2 \sigma_i \sigma_j + \kappa(i,j) |\nabla \sigma|^2] \\
&\quad + [2 H_{ji} - 2 \sigma_j \sigma_i + \kappa(j,i) |\nabla \sigma|^2] \\
&= 4 H_{ij} - 2(\sigma_i \sigma_j + \sigma_j \sigma_i) + 2 \kappa(i,j) |\nabla \sigma|^2 ,
\end{align}
using $H_{ij} = H_{ji}$ and $\kappa(i,j) = \kappa(j,i)$.
\end{proof}
\lean{NCConformal.symmetric\_part\_uncorrected}
\end{theorem}

\begin{theorem}[Entire correction is antisymmetric]
The deformation tensor can be decomposed as symmetric plus antisymmetric parts, with the correction appearing only in the antisymmetric part:
\begin{align}
\text{Defm}_{ij} + \text{Defm}_{ji} + (\text{Defm}_{ij} - \text{Defm}_{ji}) &= 2 \text{Defm}_{ij} , \\
\text{Defm}_{ij} - \text{Defm}_{ji} &= -2 [\sigma_i, \sigma_j] .
\end{align}
\begin{proof}
The first identity is algebraic; the second is the antisymmetry formula already proved.
\end{proof}
\lean{NCConformal.correction\_is\_pure\_antisymmetric}
\end{theorem}

\begin{theorem}[Master contraction with quantum correction]
The sum of Christoffel products is:
\begin{align}
\sum_m \Gamma^a_{cm} \Gamma^m_{be} &= \kappa(a,c) \cdot (\sigma_b \sigma_e + \sigma_e \sigma_b) - \kappa(a,c) \cdot (\kappa(b,e) \cdot |\nabla \sigma|^2) \\
&\quad + \kappa(a,b) \cdot (\sigma_c \sigma_e) + \kappa(a,e) \cdot (\sigma_c \sigma_b) \\
&\quad - \kappa(c,b) \cdot (\sigma_a \sigma_e) - \kappa(c,e) \cdot (\sigma_a \sigma_b) + \kappa(b,e) \cdot [\sigma_a, \sigma_c] .
\end{align}
Classically, the product $\sum_m \Gamma^a_{cm} \Gamma^m_{be}$ contains $2 \kappa(a,c) \sigma_b \sigma_e$. Non-commutatively, the factor of $2$ becomes the anticommutator $\sigma_b \sigma_e + \sigma_e \sigma_b$ (the symmetric part), the mixed terms keep their order, and one genuinely new commutator term appears: $\kappa(b,e) [\sigma_a, \sigma_c]$.
\begin{proof}
Expand $\Gamma^a_{cm} = \kappa(a,c) \sigma_m + \kappa(a,m) \sigma_c - \kappa(c,m) \sigma_a$ and $\Gamma^m_{be} = \kappa(m,b) \sigma_e + \kappa(m,e) \sigma_b - \kappa(b,e) \sigma_m$, multiply them, collect terms using the centrality of Kronecker deltas, and sum over $m$ using $\sum_m \kappa(i,m) \kappa(m,j) = \kappa(i,j)$ and the sifting properties of $\kappa$.
\end{proof}
\lean{NCConformal.sum\_Chr\_Chr\_full\_nc}
\end{theorem}

\begin{theorem}[Classical limit is the commuting case]
When the scale gradients commute—that is, $\sigma_i \sigma_j = \sigma_j \sigma_i$ for all $i, j$—the anticommutator collapses to $2 \sigma_b \sigma_e$ and the commutator term vanishes, returning the classical formula exactly. The commutative computation is not an approximation; it is the vanishing-commutator case.
\begin{proof}
Substitute $\sigma_i \sigma_j = \sigma_j \sigma_i$ into the non-commutative formula. The anticommutator becomes $2 \sigma_b \sigma_e$, and $[\sigma_a, \sigma_c] = 0$ eliminates the last term.
\end{proof}
\lean{NCConformal.sum\_Chr\_Chr\_full\_classical}
\end{theorem}

\begin{theorem}[Quantum correction couples to rotation]
The commutator of scale gradients is antisymmetric in its indices and couples to the rotational label of the framework:
\[
[\sigma_i, \sigma_j] = -[\sigma_j, \sigma_i] .
\]
The framework's rotational label (from \textit{Axes.lean}) is the wedge of two scale directions, which is also antisymmetric in the same pair. Both vanish together. Hence the prediction: quantum-gravitational corrections appear first in spin–geometry coupling, not in orbits.
\begin{proof}
The commutator is antisymmetric by definition: $[\sigma_i, \sigma_j] = \sigma_i \sigma_j - \sigma_j \sigma_i = -(\sigma_j \sigma_i - \sigma_i \sigma_j) = -[\sigma_j, \sigma_i]$.
\end{proof}
\lean{NCConformal.quantum\_correction\_couples\_to\_rotation}
\end{theorem}

\begin{theorem}[Correction vanishes on commuting gradients]
On any pattern where the scale gradients are parallel—that is, $\sigma_i \sigma_j = \sigma_j \sigma_i$—the deformation tensor is symmetric:
\[
\sigma_i \sigma_j = \sigma_j \sigma_i \implies \text{Defm}_{ij} = \text{Defm}_{ji} .
\]
This is the isotropic locus of the framework; it is precisely where the quantum correction and the rotational label both vanish.
\begin{proof}
When the gradients commute, $[\sigma_i, \sigma_j] = 0$, so by the symmetry equivalence theorem, $\text{Defm}_{ij} = \text{Defm}_{ji}$.
\end{proof}
\lean{NCConformal.correction\_vanishes\_on\_commuting}
\end{theorem}


\section{Part IV: Scale Flow and Content}

\subsection{RG}

The renormalisation group arises as a theorem about how couplings must shift to preserve predictions when the fiducial log-scale changes, as required by Axiom A5. This section formalises three main results: that scale flow forms a one-parameter group, that self-similar observables obey power-law evolution (the Callan--Symanzik equation), and that the effective field theory tower of thresholds is forced by the scale structure alone. Each excitation carries a log-threshold, and the content available at any log-energy is exactly the set of excitations whose thresholds lie at or below it; the structure of the renormalisation group is therefore a consequence of scale invariance, not an input.

\begin{definition}
A \textbf{scale flow} on a space $S$ of dimensionless couplings is the structure
\begin{equation}
\mathsf{ScaleFlow}(S) := \{
  \text{flow} : \mathbb{R} \to S \to S
  \mid
  \text{flow}(0, g) = g \text{ and } \text{flow}(t+u, g) = \text{flow}(t, \text{flow}(u, g))
\}
\end{equation}
It represents the action of a shift of the fiducial log-scale on the coupling space. \lean{RG.ScaleFlow}
\end{definition}

\begin{definition}
A coupling $g \in S$ is a \textbf{fixed point} of a scale flow $\Phi$ if $\Phi.\text{flow}(t, g) = g$ for all $t \in \mathbb{R}$. \lean{RG.ScaleFlow.IsFixedPoint}
\end{definition}

\begin{theorem}
At a fixed point of the flow, every observable is constant under the flow. More precisely, if $\Phi.\text{IsFixedPoint}(g)$ holds and $O : S \to \beta$ is any observable, then for all $t \in \mathbb{R}$,
\begin{equation}
O(\Phi.\text{flow}(t, g)) = O(g) .
\end{equation}
\lean{RG.ScaleFlow.observable\_const\_at\_fixedPoint}
\end{theorem}

\begin{proof}
Immediate from the definition of fixed point.
\end{proof}

\begin{theorem}
If $g$ is a fixed point of the scale flow $\Phi$, then $\Phi.\text{flow}(t, g)$ is also a fixed point for any $t \in \mathbb{R}$.
\lean{RG.ScaleFlow.fixedPoint\_invariant}
\end{theorem}

\begin{proof}
For any $u \in \mathbb{R}$, we have $\Phi.\text{flow}(u, \Phi.\text{flow}(t, g)) = \Phi.\text{flow}(t, \Phi.\text{flow}(u, g)) = \Phi.\text{flow}(t, g)$ by the group law and the fixed-point property of $g$.
\end{proof}

\begin{theorem}[Callan--Symanzik equation and self-similarity]
Suppose $O : \mathbb{R} \to \mathbb{R}$ satisfies the Callan--Symanzik equation
\begin{equation}
\frac{dO}{dt}(t) = \gamma \cdot O(t)
\end{equation}
for all $t \in \mathbb{R}$, where $\gamma$ is a constant (the anomalous dimension). Then
\begin{equation}
O(t) = O(0) \cdot e^{\gamma t} .
\end{equation}
In particular, $\gamma = 0$ (a fixed point) yields exact scale invariance: $O(t) = O(0)$. \lean{RG.callan\_symanzik}
\end{theorem}

\begin{proof}
Define $u(r) := O(r) \cdot e^{-\gamma r}$. The derivative of $u$ is
\begin{equation}
\dot{u}(r) = \dot{O}(r) \cdot e^{-\gamma r} + O(r) \cdot (-\gamma e^{-\gamma r})
  = e^{-\gamma r} (\gamma O(r) - \gamma O(r)) = 0 .
\end{equation}
Thus $u$ is constant, so $u(t) = u(0)$, giving $O(t) e^{-\gamma t} = O(0)$, hence $O(t) = O(0) e^{\gamma t}$.
\end{proof}

\begin{theorem}
At a fixed point of the flow ($\gamma = 0$), the observable is exactly constant: $O(t) = O(0)$ for all $t$.
\lean{RG.scale\_invariant\_of\_fixedPoint}
\end{theorem}

\begin{proof}
Follows from the Callan--Symanzik theorem with $\gamma = 0$.
\end{proof}

\begin{definition}
Given a function $\mu : \iota \to \mathbb{R}$ assigning a log-threshold to each element $i \in \iota$ (where $\iota$ is a finite type), and a log-energy $t \in \mathbb{R}$, the \textbf{active content} at that log-energy is
\begin{equation}
\text{active}(\mu, t) := \{i \in \iota : \mu(i) \leq t\} .
\end{equation}
\lean{RG.active}
\end{definition}

\begin{theorem}
An element $i$ is in $\text{active}(\mu, t)$ if and only if $\mu(i) \leq t$.
\lean{RG.mem\_active}
\end{theorem}

\begin{proof}
Immediate from the definition.
\end{proof}

\begin{theorem}
The active content is monotone in the log-energy: if $t \leq t'$, then $\text{active}(\mu, t) \subseteq \text{active}(\mu, t')$. That is, new physics appears as the scale rises, and never disappears.
\lean{RG.active\_mono}
\end{theorem}

\begin{proof}
If $i \in \text{active}(\mu, t)$, then $\mu(i) \leq t \leq t'$, so $i \in \text{active}(\mu, t')$.
\end{proof}

\begin{theorem}
Between thresholds the effective theory is exact. If $t \leq t'$ and no threshold lies in the interval $(t, t']$, that is, $\neg \exists i : t < \mu(i) \leq t'$, then
\begin{equation}
\text{active}(\mu, t) = \text{active}(\mu, t') .
\end{equation}
\lean{RG.active\_eq\_of\_no\_threshold}
\end{theorem}

\begin{proof}
Since the active set is monotone, it suffices to show $\text{active}(\mu, t') \subseteq \text{active}(\mu, t)$. Suppose $i \in \text{active}(\mu, t')$, so $\mu(i) \leq t'$. If $\mu(i) > t$, then $t < \mu(i) \leq t'$, contradicting the hypothesis. Thus $\mu(i) \leq t$, so $i \in \text{active}(\mu, t)$.
\end{proof}

\begin{theorem}
The tower of effective theories is finite. The range $\{\text{active}(\mu, t) : t \in \mathbb{R}\}$ is a finite set of subsets of $\iota$, so there are only finitely many distinct effective theories as $t$ varies. Changing the effective framework at each threshold is therefore not a defect of the method; it is the exact bookkeeping of a scale-stratified world.
\lean{RG.active\_finite\_range}
\end{theorem}

\begin{proof}
The range is a subset of the power set $\mathcal{P}(\iota)$, which is finite since $\iota$ is finite.
\end{proof}

\begin{theorem}
Below all thresholds the effective theory is empty; above all thresholds it is complete.
\begin{enumerate}
\item If $\mu(i) \leq t$ for all $i \in \iota$, then $\text{active}(\mu, t) = \iota$.
\item If $t < \mu(i)$ for all $i \in \iota$, then $\text{active}(\mu, t) = \emptyset$.
\end{enumerate}
\lean{RG.active\_eq\_univ\_of\_ge} and \lean{RG.active\_eq\_empty\_of\_lt}
\end{theorem}

\begin{proof}
(1) Every element $i$ satisfies $\mu(i) \leq t$, so belongs to $\text{active}(\mu, t)$.
(2) No element $i$ satisfies $\mu(i) \leq t$ since $t < \mu(i)$, so $\text{active}(\mu, t) = \emptyset$.
\end{proof}

\begin{definition}
The \textbf{beta coefficient} as a function of the active content is defined by
\begin{equation}
\beta_\Phi(\mu, t) := \Phi(\text{active}(\mu, t)) ,
\end{equation}
where $\Phi : \mathcal{P}(\iota) \to \mathbb{R}$ maps each finite subset of $\iota$ to a real number. This expresses the idea that the compensation rate required by A5 depends only on which excitations are present.
\lean{RG.betaOf}
\end{definition}

\begin{theorem}
The beta coefficient is constant between thresholds. If $t \leq t'$ and no threshold lies in $(t, t']$, then
\begin{equation}
\beta_\Phi(\mu, t) = \beta_\Phi(\mu, t') .
\end{equation}
Where the content does not change, the compensation required by A5 does not change either.
\lean{RG.beta\_const\_between\_thresholds}
\end{theorem}

\begin{proof}
Since $\text{active}(\mu, t) = \text{active}(\mu, t')$ by the preceding theorem, we have $\beta_\Phi(\mu, t) = \Phi(\text{active}(\mu, t)) = \Phi(\text{active}(\mu, t')) = \beta_\Phi(\mu, t')$.
\end{proof}

\begin{theorem}
The beta function takes finitely many values: the set $\{\beta_\Phi(\mu, t) : t \in \mathbb{R}\}$ is finite. It is a step function of the log-scale whose jumps sit exactly at the masses. The framework fixes this structure; it does not fix the numbers, which remain an input.
\lean{RG.beta\_finite\_range}
\end{theorem}

\begin{proof}
The range of $\beta_\Phi(\mu, \cdot)$ is contained in $\Phi(\text{active}(\mu, \mathbb{R})) = \Phi(\text{range}(\text{active}(\mu)))$, which is finite since $\text{active}(\mu)$ has finite range.
\end{proof}

\begin{theorem}
Above every threshold the coefficient settles to its complete-content value. If $\mu(i) \leq t$ for all $i$, then
\begin{equation}
\beta_\Phi(\mu, t) = \Phi(\iota) .
\end{equation}
The deep-ultraviolet running is governed by the full spectrum.
\lean{RG.beta\_eq\_of\_all\_active}
\end{theorem}

\begin{proof}
If all thresholds are below $t$, then $\text{active}(\mu, t) = \iota$, so $\beta_\Phi(\mu, t) = \Phi(\iota)$.
\end{proof}

\subsection{Running}

In a scale-covariant world, no log-scale is preferred, so the defects capable of being resolved by a probe cannot bunch anywhere: their density per unit log-scale is constant. Integrating a constant density over an interval in log-energy gives a linear accumulation. When a coupling's response is the count of resolvable defects, this yields $g^{-2}(t) = g^{-2}(t_0) + \kappa\rho(t - t_0)$, which is precisely the one-loop form with $b = \kappa\rho/2$. The linear-in-log running is therefore a consequence of scale-invariant defect density, not of a one-loop diagram. The sign of asymptotic freedom also emerges: positive density means resolving more structure dilutes the response.

\begin{definition}
The number of defects a probe at log-scale $t$ can resolve, given a scale-invariant density $\rho$ per unit log-scale and reference log-scale $t_0$, is
\begin{equation}
\text{resolvedCount}(\rho, t_0, t) := \rho \cdot (t - t_0) .
\end{equation}
\lean{Running.resolvedCount}
\end{definition}

\begin{theorem}
At the reference log-scale, the resolved count is zero:
\begin{equation}
\text{resolvedCount}(\rho, t_0, t_0) = 0 .
\end{equation}
\lean{Running.resolvedCount\_self}
\end{theorem}

\begin{proof}
Immediate from the definition.
\end{proof}

\begin{definition}
The inverse-square coupling, as accumulated response with each resolvable defect contributing $\kappa$, is
\begin{equation}
u(u_0, \kappa, \rho, t_0, t) := u_0 + \kappa \cdot \text{resolvedCount}(\rho, t_0, t) = u_0 + \kappa \rho (t - t_0) .
\end{equation}
\lean{Running.invSqCoupling}
\end{definition}

\begin{theorem}
At the reference scale, the inverse-square coupling takes its base value:
\begin{equation}
u(u_0, \kappa, \rho, t_0, t_0) = u_0 .
\end{equation}
\lean{Running.invSqCoupling\_at\_base}
\end{theorem}

\begin{proof}
Immediate from the definition.
\end{proof}

\begin{theorem}
The accumulated response has constant slope in log-scale. The function $t \mapsto u(u_0, \kappa, \rho, t_0, t)$ satisfies
\begin{equation}
\frac{du}{dt}(t) = \kappa \rho
\end{equation}
for all $t \in \mathbb{R}$. Nothing perturbative has been used: the slope is a density times a per-defect response.
\lean{Running.hasDerivAt\_invSqCoupling}
\end{theorem}

\begin{proof}
The derivative of $u_0 + \kappa \rho(t - t_0)$ with respect to $t$ is $\kappa \rho$ by the rules of differentiation.
\end{proof}

\begin{theorem}
The counting law and the one-loop equation are identical. Writing $b = \kappa\rho/2$, the slope $\kappa\rho$ is exactly $2b$ of the one-loop form. So the beta function's shape is not an input: it is what a constant density of resolvable defects produces.
\lean{Running.slope\_eq\_two\_b}
\end{theorem}

\begin{proof}
Direct algebra: $2 \cdot (\kappa\rho/2) = \kappa\rho$.
\end{proof}

\begin{theorem}
The counting law reproduces the one-loop running. For any $u_0, \kappa, \rho, t_0, t \in \mathbb{R}$,
\begin{equation}
u(u_0, \kappa, \rho, t_0, t) = u_0 + 2 \left(\frac{\kappa\rho}{2}\right) (t - t_0) .
\end{equation}
\lean{Running.invSqCoupling\_eq\_one\_loop}
\end{theorem}

\begin{proof}
Algebraic simplification of the definition.
\end{proof}

\begin{theorem}[Asymptotic freedom as positive density]
Suppose $\kappa > 0$. Then the response $g^2 = 1/g^{-2}$ falls to zero at high log-scale if and only if the accumulated count grows, i.e., if and only if the defect density is positive:
\begin{equation}
\left(\forall \varepsilon > 0, \exists T : \forall t > T, u(u_0, \kappa, \rho, t_0, t) > 1/\varepsilon\right) \iff \rho > 0 .
\end{equation}
No gauge group or diagram appears; what appears is a count. Resolving more structure dilutes the response: this is the physical content of asymptotic freedom in this framework.
\lean{Running.asympt\_free\_iff\_pos\_density}
\end{theorem}

\begin{proof}
We prove both directions.
($\Rightarrow$) Suppose the left-hand side holds but assume $\rho \leq 0$. For any $\varepsilon > 0$, there exists $T$ such that for all $t > T$, $u(u_0, \kappa, \rho, t_0, t) > 1/\varepsilon$. Taking $\varepsilon = 1/(|u_0| + 1)$, for sufficiently large $t$ we have $u(u_0, \kappa, \rho, t_0, t) > |u_0| + 1$. But if $\rho \leq 0$, then the accumulated count is non-positive for $t \geq t_0$, so $u(u_0, \kappa, \rho, t_0, t) \leq u_0 \leq |u_0|$, a contradiction.
($\Leftarrow$) Suppose $\rho > 0$ and $\varepsilon > 0$. Set $T := t_0 + (1/\varepsilon - u_0)/(\kappa\rho) + 1$. Then for $t > T$,
\begin{equation}
u(u_0, \kappa, \rho, t_0, t) = u_0 + \kappa\rho(t - t_0) > u_0 + \kappa\rho\left(\frac{1/\varepsilon - u_0}{\kappa\rho} + 1\right) = u_0 + (1/\varepsilon - u_0) + \kappa\rho > 1/\varepsilon .
\end{equation}
\end{proof}

\subsection{QCD}

The strong interaction reveals dimensional transmutation: a coupling that is a pure number, running under the scale-invariant framework, secretes an invariant scale $\Lambda$ that appears nowhere in the axioms. Scale invariance is not violated by fiat; it is broken by the solutions of a scale-covariant equation. Proved here is the one-loop solution $g^{-2}(t) = g^{-2}(0) + 2bt$ with $b > 0$, asymptotic freedom, the dimensional transmutation constant $\Lambda = \exp(t - g^{-2}/(2b))$, and the consequence that all hadron mass ratios are parameter-free. The simplest geometric tower of defect masses fails against lattice glueball data, falsifying a prior hypothesis; the parameter-free ratio prediction survives.

\begin{theorem}[Exact solution of one-loop running]
Suppose $g : \mathbb{R} \to \mathbb{R}$ is nonzero and satisfies the one-loop equation
\begin{equation}
\frac{dg}{dt}(t) = -b \cdot (g(t))^3
\end{equation}
for some constant $b \in \mathbb{R}$ and all $t \in \mathbb{R}$. Then
\begin{equation}
(g(t))^{-2} = (g(0))^{-2} + 2bt .
\end{equation}
\lean{QCD.running\_inv\_sq}
\end{theorem}

\begin{proof}
Define $u(r) := (g(r))^{-2} - 2br$. We compute
\begin{equation}
\dot{u}(r) = -\frac{2(g(r))^{-3} \cdot (-b(g(r))^3)}{1} - 2b = 2b - 2b = 0 .
\end{equation}
Thus $u$ is constant, so $u(t) = u(0)$, yielding $(g(t))^{-2} = (g(0))^{-2} + 2bt$.
\end{proof}

\begin{theorem}[Asymptotic freedom]
Suppose $b > 0$, $g : \mathbb{R} \to \mathbb{R}$ is nonzero and satisfies $dg/dt = -b g^3$. Then for any $\varepsilon > 0$, there exists $T \in \mathbb{R}$ such that for all $t > T$, $(g(t))^2 < \varepsilon$. The strong sector becomes free precisely where the local scale is small, as demanded by A3.
\lean{QCD.asymptotic\_freedom}
\end{theorem}

\begin{proof}
From the one-loop solution, $(g(t))^{-2} = (g(0))^{-2} + 2bt$. Choose $T$ large enough that $2bT > 1/\varepsilon - (g(0))^{-2}$. Then for $t > T$, $(g(t))^{-2} > 1/\varepsilon$, so $(g(t))^2 < \varepsilon$.
\end{proof}

\begin{definition}
The scale secreted by the running coupling is defined as
\begin{equation}
\Lambda(b, g, t) := \exp\left(t - \frac{(g(t))^{-2}}{2b}\right) .
\end{equation}
\lean{QCD.Lambda}
\end{definition}

\begin{theorem}[Dimensional transmutation]
The quantity $\Lambda(b, g, t)$ is constant along the flow: a theory whose only input is a dimensionless number has produced a scale, and by A5 no choice of fiducial can shift it. This is how a scale-invariant axiom system comes to describe a world with definite hadron masses. The invariance is not broken in the axioms; it is broken by the solution.
\begin{equation}
\Lambda(b, g, t) = \Lambda(b, g, 0)
\end{equation}
for all $t \in \mathbb{R}$, provided $b \neq 0$.
\lean{QCD.lambda\_invariant}
\end{theorem}

\begin{proof}
From the one-loop solution, $(g(t))^{-2} - (g(0))^{-2} = 2bt$, so
\begin{equation}
\frac{(g(t))^{-2}}{2b} - t = \frac{(g(0))^{-2}}{2b} - t + \frac{(g(0))^{-2} - (g(t))^{-2}}{2b} + t = \frac{(g(0))^{-2}}{2b} - t .
\end{equation}
Thus $t - (g(t))^{-2}/(2b) = -($ constant $)$, so $\Lambda(b, g, t) = \Lambda(b, g, 0)$.
\end{proof}

\begin{theorem}
The strong scale is genuinely positive: $\Lambda(b, g, t) > 0$ for all $b \in \mathbb{R}$, $g : \mathbb{R} \to \mathbb{R}$, and $t \in \mathbb{R}$.
\lean{QCD.lambda\_pos}
\end{theorem}

\begin{proof}
The exponential function is always positive.
\end{proof}

\begin{definition}
A hadron mass is defined as
\begin{equation}
M_{\text{hadron}}(c, \Lambda) := c \cdot \Lambda
\end{equation}
where $c$ is a pure (dimensionless) number and $\Lambda$ is the transmuted scale.
\lean{QCD.hadronMass}
\end{definition}

\begin{theorem}[Parameter-free mass ratios]
The ratio of any two hadron masses is a ratio of pure numbers, independent of the transmuted scale $\Lambda$. Hence it is independent of every input to the theory:
\begin{equation}
\frac{M_{\text{hadron}}(c_1, \Lambda)}{M_{\text{hadron}}(c_2, \Lambda)} = \frac{c_1}{c_2}
\end{equation}
provided $c_2 \neq 0$ and $\Lambda \neq 0$. This is the sharpest prediction the framework makes about the strong sector. Glueballs are the right place to test it because their masses involve no quark mass, so $\Lambda$ is genuinely the only scale.
\lean{QCD.hadron\_mass\_ratio}
\end{theorem}

\begin{proof}
Direct division: $(c_1 \Lambda) / (c_2 \Lambda) = c_1/c_2$.
\end{proof}

\begin{theorem}
Rescaling $\Lambda$ moves every mass together: the spectrum has one overall normalisation and no other freedom. If $k \neq 0$,
\begin{equation}
M_{\text{hadron}}(c, k\Lambda) = k \cdot M_{\text{hadron}}(c, \Lambda) .
\end{equation}
\lean{QCD.spectrum\_rigid}
\end{theorem}

\begin{proof}
$M_{\text{hadron}}(c, k\Lambda) = c(k\Lambda) = k(c\Lambda) = k \cdot M_{\text{hadron}}(c, \Lambda)$.
\end{proof}

\begin{definition}
The ratio of the first two lattice glueball masses is
\begin{equation}
r_1 := \frac{2390}{1710} \approx 1.398 .
\end{equation}
The ratio of the next two is
\begin{equation}
r_2 := \frac{2560}{2390} \approx 1.071 .
\end{equation}
\lean{QCD.ratio$_1$} and \lean{QCD.ratio$_2$}
\end{definition}

\begin{theorem}[The geometric tower fails]
The two consecutive ratios differ by more than $0.3$, so a single geometric tower (which would predict equal ratios) is excluded by lattice data:
\begin{equation}
|r_1 - r_2| > 0.3 .
\end{equation}
This refutes the hypothesis that the defect masses follow the posited law of Defect.mass, not the defect picture itself. A spectrum with strain energy quadratic in the winding, for instance, is unaffected.
\lean{QCD.geometric\_tower\_fails}
\end{theorem}

\begin{proof}
Numerical verification: $|1.398 - 1.071| \approx 0.327 > 0.3$.
\end{proof}

\begin{theorem}
The numerical values of the ratios are:
\begin{enumerate}
\item $1.39 < r_1 < 1.40$.
\item $1.07 < r_2 < 1.08$.
\end{enumerate}
\lean{QCD.ratio$_1$\_value} and \lean{QCD.ratio$_2$\_value}
\end{theorem}

\begin{proof}
Direct arithmetic from the definitions.
\end{proof}

\begin{definition}
With the transmuted scale $\Lambda \approx 332$ MeV, the lightest glueball (at $\approx 1710$ MeV) sits at
\begin{equation}
\frac{m_0}{\Lambda} \approx \frac{1710}{332} \approx 5.15 .
\end{equation}
This is a pure number the framework must eventually compute rather than fit.
\lean{QCD.lightestGlueballInLambda}
\end{definition}

\begin{theorem}
The numerical value is bracketed:
\begin{equation}
5.1 < \frac{m_0}{\Lambda} < 5.2 .
\end{equation}
\lean{QCD.lightestGlueballInLambda\_value}
\end{theorem}

\begin{proof}
Direct arithmetic from the definition.
\end{proof}

\subsection{Spectrum}

The framework predicts the structure of threshold placement. A5 says no log-scale is preferred, so thresholds cannot bunch anywhere: they are placed at a constant rate $\rho$ per unit log-scale. A constant-rate point process has exponentially distributed gaps, whose squared coefficient of variation is exactly one. This prediction discriminates sharply against the geometric tower (CV = 0) and is tested against the eleven consecutive gaps in $\ln m$ across the twelve distinct Standard Model mass scales. The observed $CV^2 = 0.875$ (so $CV = 0.935$) sits at the 63rd percentile of the prediction's sampling range and is incompatible with the geometric tower.

\begin{definition}
The threshold of a structure of mass $m > 0$ is the log-energy at which it becomes resolvable:
\begin{equation}
\mu(m) := \ln m .
\end{equation}
Conversely, a structure of log-threshold $\mu$ has mass $e^\mu$.
\lean{Spectrum.thresholdOfMass} and \lean{Spectrum.massOfThreshold}
\end{definition}

\begin{theorem}
The threshold and mass are inverse maps: $\mu(e^\mu) = \mu$ for all $\mu \in \mathbb{R}$.
\lean{Spectrum.mass\_threshold\_inverse}
\end{theorem}

\begin{proof}
$\mu(e^\mu) = \ln(e^\mu) = \mu$.
\end{proof}

\begin{definition}
For a list $\mathbf{x} = (x_1, \ldots, x_n)$ of real numbers, the \textbf{mean} is
\begin{equation}
\bar{x} := \frac{1}{n}\sum_{i=1}^n x_i ,
\end{equation}
and the \textbf{population variance} is
\begin{equation}
\sigma^2 := \frac{1}{n}\sum_{i=1}^n x_i^2 - \bar{x}^2 .
\end{equation}
\lean{Spectrum.mean} and \lean{Spectrum.variance}
\end{definition}

\begin{theorem}[Geometric tower has zero variance]
If every gap equals a constant $c$, the mean is $c$ and the mean square is $c^2$, so the variance vanishes. A geometric spectrum therefore predicts $CV = \sigma/\bar{x} = 0$, with no freedom:
\begin{equation}
\text{variance}([c, c, \ldots, c]) = 0 .
\end{equation}
\lean{Spectrum.equal\_gaps\_variance\_zero}
\end{theorem}

\begin{proof}
If all $n$ entries are $c$, then $\bar{x} = c$ and $\overline{x^2} = c^2$, so $\sigma^2 = c^2 - c^2 = 0$.
\end{proof}

\begin{definition}
The gaps (in $\ln m$) between consecutive Standard Model mass scales are (taking massless states and unknown neutrino masses into account):
\begin{equation}
\text{gaps} = [1.46, 0.759, 2.985, 0.128, 2.486, 0.336, 0.855, 2.956, 0.126, 0.317, 0.321] .
\end{equation}
The sum is $\sum \text{gaps} = 12.729$ and the sum of squares is $\sum (\text{gaps})^2 = 27.615749$.
\lean{Spectrum.gaps}, \lean{Spectrum.gapSum}, and \lean{Spectrum.gapSqSum}
\end{definition}

\begin{definition}
The mean gap is $\langle \text{gap} \rangle := 12.729/11 \approx 1.1572$.
The population variance of the gaps is $\sigma_{\text{gap}}^2 := 27.615749/11 - (12.729/11)^2$.
The squared coefficient of variation is $CV^2 := \sigma_{\text{gap}}^2 / \langle \text{gap} \rangle^2$.
\lean{Spectrum.meanGap}, \lean{Spectrum.gapVariance}, and \lean{Spectrum.cvSq}
\end{definition}

\begin{definition}
The predictions are:
\begin{enumerate}
\item \textbf{Scale-invariant}: $CV^2 = 1$ (exponential gap distribution).
\item \textbf{Geometric tower}: $CV^2 = 0$ (equal gaps).
\end{enumerate}
\lean{Spectrum.predictedCvSq} and \lean{Spectrum.geometricCvSq}
\end{definition}

\begin{theorem}[Observed coefficient of variation]
The observed value is bracketed:
\begin{equation}
0.87 < CV^2 < 0.88 ,
\end{equation}
i.e., $CV^2 = 0.875$, so $CV \approx 0.935$. This is close to but below the predicted value of $1$.
\lean{Spectrum.observed\_cvSq\_value}
\end{theorem}

\begin{proof}
Direct arithmetic from the gap data.
\end{proof}

\begin{theorem}
The mean gap is bracketed:
\begin{equation}
1.15 < \langle \text{gap} \rangle < 1.16 .
\end{equation}
The inverse is the threshold density $\rho \approx 0.86$ per e-fold.
\lean{Spectrum.meanGap\_value}
\end{theorem}

\begin{proof}
Direct arithmetic from the gap sum.
\end{proof}

\begin{theorem}[The Standard Model spectrum is not geometric]
The gap variance is positive, whereas a geometric tower has variance exactly zero. The observed spectrum is therefore incompatible with the geometric tower:
\begin{equation}
CV^2 \neq 0 .
\end{equation}
This is an independent and sharper refutation than the glueball ratio test, using the whole observed spectrum rather than three lattice numbers.
\lean{Spectrum.observed\_not\_geometric}
\end{theorem}

\begin{proof}
We have $CV^2 \approx 0.875 > 0$, contradicting the geometric prediction of $CV^2 = 0$.
\end{proof}

\begin{theorem}
The observed value is compatible with the scale-invariant prediction. The central 90\% of the sampling distribution for eleven gaps runs from approximately $0.33$ to $1.64$ in $CV^2$:
\begin{equation}
0.33 < CV^2 < 1.64 .
\end{equation}
The observed value $0.875$ sits at the 63rd percentile.
\lean{Spectrum.observed\_within\_scale\_invariant\_range}
\end{theorem}

\begin{proof}
Direct inequalities from the observed value.
\end{proof}

\begin{definition}
Given the mean gap, the threshold density is recovered as
\begin{equation}
\rho(\langle \text{gap} \rangle) := \frac{1}{\langle \text{gap} \rangle} .
\end{equation}
\lean{Spectrum.densityFromGaps}
\end{definition}

\begin{theorem}
With the observed mean gap, the threshold density satisfies
\begin{equation}
0.86 < \rho < 0.87 .
\end{equation}
\lean{Spectrum.densityFromGaps\_value}
\end{theorem}

\begin{proof}
Direct arithmetic from the mean gap.
\end{proof}

\begin{definition}
From the dark-energy equation of state parameter $w$, the dissipation exponent is
\begin{equation}
q(w) := \frac{3w + 5}{2} .
\end{equation}
With $w = -1$, we get $q = 1$.
\lean{Spectrum.qFromW}
\end{definition}

\begin{definition}
From the asymptotic Hubble rate $\lambda = H_\infty$, the dark-energy density is
\begin{equation}
\rho_\Lambda(\lambda, G) := \frac{3\lambda^2}{8\pi G} .
\end{equation}
\lean{Spectrum.rhoLambdaFromLambda}
\end{definition}

\begin{definition}
When a threshold appears, the newly resolvable content is characterized by the jump in the running coefficient $\Delta b$. The increment to the particle count is
\begin{equation}
\Delta N(\Delta b, \kappa) := \frac{\Delta b}{\kappa} .
\end{equation}
\lean{Spectrum.deltaNFromDeltaB}
\end{definition}

\begin{theorem}
The identity $\Delta N \cdot \kappa = \Delta b$ holds (under suitable non-degeneracy conditions).
\lean{Spectrum.deltaN\_inverse}
\end{theorem}

\begin{proof}
Direct algebra: $({\Delta b}/{\kappa}) \cdot \kappa = \Delta b$.
\end{proof}

\begin{theorem}
The conversions are mutually consistent: a spectrum with mean gap $g$ has density $1/g$, and a density $\rho$ has mean gap $1/\rho$.
\begin{equation}
\frac{1}{\rho(g)} = g
\end{equation}
provided $g \neq 0$.
\lean{Spectrum.density\_gap\_inverse}
\end{theorem}

\begin{proof}
$1/(\rho(g)) = 1/(1/g) = g$.
\end{proof}

\subsection{Entropy}

Nothing in the microscopic dynamics distinguishes a direction of time: the substrate is reversible. What breaks the symmetry is that observation happens at a scale, which imposes a finite resolution. A finite resolution is a non-injective map from micro-configurations to what is observable. Under reversible microscopic dynamics followed by coarse-graining, the number of micro-configurations compatible with observation never decreases (the H-theorem), and entropy production is exact bookkeeping, not a statistical assumption. If entropy strictly increases once, then no entropy-non-decreasing law can undo the evolution, making the macroscopic equation irreversible even though the microscopic one is not—precisely the status of the Boltzmann equation.

\begin{definition}
Given a finite set $X$ of micro-configurations and a resolution map $\pi : X \to Y$, the \textbf{saturation} or \textbf{macrostate} generated by a subset $S \subseteq X$ is
\begin{equation}
\text{sat}(\pi, S) := \{x \in X : \pi(x) \in \pi(S)\} .
\end{equation}
It is the set of all micro-configurations that the resolution cannot distinguish from something in $S$.
\lean{Entropy.sat}
\end{definition}

\begin{theorem}
A configuration is in the saturation if and only if its image under $\pi$ lies in the image of $S$:
\begin{equation}
x \in \text{sat}(\pi, S) \iff \pi(x) \in \pi(S) .
\end{equation}
\lean{Entropy.mem\_sat}
\end{theorem}

\begin{proof}
Immediate from the definition.
\end{proof}

\begin{theorem}
Coarse-graining never loses a configuration: $S \subseteq \text{sat}(\pi, S)$.
\lean{Entropy.subset\_sat}
\end{theorem}

\begin{proof}
If $x \in S$, then $\pi(x) \in \pi(S)$, so $x \in \text{sat}(\pi, S)$.
\end{proof}

\begin{theorem}
Therefore, $|S| \leq |\text{sat}(\pi, S)|$.
\lean{Entropy.card\_le\_card\_sat}
\end{theorem}

\begin{proof}
The cardinality of a subset is at most that of the whole set.
\end{proof}

\begin{theorem}[Microscopic reversibility]
If $T : X \to X$ is a bijection (permutation), then any finite subset $S \subseteq X$ has $|T(S)| = |S|$. Microscopic evolution preserves cardinality exactly.
\lean{Entropy.card\_image\_perm}
\end{theorem}

\begin{proof}
The image of a finite set under an injective map has the same cardinality.
\end{proof}

\begin{definition}
One step of observed evolution is: evolve microscopically by $T : X \to X$, then coarse-grain by the resolution $\pi : X \to Y$:
\begin{equation}
F(\pi, T, S) := \text{sat}(\pi, T(S)) .
\end{equation}
\lean{Entropy.evolve}
\end{definition}

\begin{theorem}[H-theorem]
Observed entropy—the log of the number of micro-configurations compatible with the observation—never decreases. Under reversible microscopic dynamics followed by coarse-graining,
\begin{equation}
|S| \leq |F(\pi, T, S)|
\end{equation}
for any subset $S \subseteq X$ and bijection $T : X \to X$. Nothing statistical is invoked.
\lean{Entropy.entropy\_nondecreasing}
\end{theorem}

\begin{proof}
We have $|S| = |T(S)|$ by microscopic reversibility, and $|T(S)| \leq |\text{sat}(\pi, T(S))|$ by coarse-graining. Thus $|S| \leq |F(\pi, T, S)|$.
\end{proof}

\begin{theorem}[The arrow of time]
Let $F, G : \mathcal{P}(X) \to \mathcal{P}(X)$ be any two evolution laws both obeying the H-theorem, i.e., $|S| \leq |G(S)|$ for all $S$. If $F$ strictly produces entropy somewhere, i.e., there exists $S_0$ with $|S_0| < |F(S_0)|$, then there is no inverse law with $G(F(S)) = S$ for all $S$:
\begin{equation}
\neg \exists G : \forall S, |S| \leq |G(S)| \text{ and } G(F(S)) = S \text{ and } \exists S_0, |S_0| < |F(S_0)| .
\end{equation}
So the macroscopic equation cannot be run backwards by a law of its own kind. Time-reversal invariance is lost at the observed level while holding exactly at the microscopic level—which is the precise sense in which the Boltzmann equation is irreversible although the mechanics beneath it is not.
\lean{Entropy.no\_reverse\_law}
\end{theorem}

\begin{proof}
Suppose such a $G$ existed with $G(F(S_0)) = S_0$ and $|F(S_0)| > |S_0|$. By the H-theorem for $G$, $|F(S_0)| \leq |G(F(S_0))| = |S_0|$, contradicting $|F(S_0)| > |S_0|$.
\end{proof}

\begin{theorem}
The same conclusion applies to the concrete observed dynamics with resolution $\pi$ and microscopic evolution $T$.
\lean{Entropy.evolve\_no\_reverse\_law}
\end{theorem}

\begin{proof}
Apply the previous theorem to $F = F(\pi, T)$.
\end{proof}

\begin{theorem}
In equilibrium every observable is stationary. Any quantity that changes under the dynamics therefore witnesses departure from equilibrium—Sakharov's third condition, stated contrapositively.
\lean{Entropy.charge\_change\_implies\_nonequilibrium}
\end{theorem}

\begin{proof}
If $F(S) = S$ (equilibrium) but $Q(F(S)) \neq Q(S)$, then $Q(S) \neq Q(S)$, a contradiction.
\end{proof}

\begin{theorem}
Strict entropy production is itself a departure from equilibrium. With the H-theorem above, an open dissipating universe satisfies Sakharov's third condition at every moment at which it produces entropy.
\lean{Entropy.entropy\_production\_implies\_nonequilibrium}
\end{theorem}

\begin{proof}
If $|S| < |F(S)|$ (entropy production) but $F(S) = S$ (equilibrium), then $|S| < |S|$, a contradiction.
\end{proof}

\subsection{Light}

The question ``what is light?'' has a sharp answer in this framework. The measured quadratic form is $Q_g = s^2 Q_\delta$, where $s$ is a unit. Multiplying by a unit cannot create or destroy a zero. Therefore: the null cone is scale-blind; light does not see the scale field $\sigma$ at all. For every non-null direction the scale is recoverable from the ratio of the measured to bare form—matter sees the scale and is the only thing that does. Light is the scale-blind part of the geometry, weight-zero under scale transformations, which is why it alone can carry the comparison of local units from one event to another. The null condition is therefore invariant under the fiducial shift of A5, and the scale connection is flat (by the conformal part of the framework).

\begin{definition}
The \textbf{bare quadratic form} of the substrate, with signature $\eta : \text{Fin } n \to A$ (entries $\pm 1$), is
\begin{equation}
Q_{\text{bare}}(\eta, v) := \sum_{i=0}^{n-1} \eta(i) \cdot v(i)^2 .
\end{equation}
This is pure bookkeeping and carries no physics.
\lean{Light.bareForm}
\end{definition}

\begin{definition}
The \textbf{measured quadratic form}, read in the local unit $s \in A^\times$, is
\begin{equation}
Q_{\text{meas}}(s, \eta, v) := s^2 \cdot Q_{\text{bare}}(\eta, v) .
\end{equation}
This is A4 applied to a tangent vector.
\lean{Light.physForm}
\end{definition}

\begin{definition}
A direction $v$ is \textbf{null} if the measured form vanishes on it: $Q_{\text{meas}}(s, \eta, v) = 0$.
\lean{Light.IsNull}
\end{definition}

\begin{theorem}
Multiplying by a unit neither creates nor destroys a zero. For any unit $s \in A^\times$ and element $x \in A$,
\begin{equation}
s^2 \cdot x = 0 \iff x = 0 .
\end{equation}
\lean{Light.unit\_sq\_mul\_eq\_zero\_iff}
\end{theorem}

\begin{proof}
If $s^2 \cdot x = 0$, then $x = (s^2)^{-1} \cdot (s^2 \cdot x) = 0$. If $x = 0$, then $s^2 \cdot x = 0$ trivially.
\end{proof}

\begin{theorem}[The null cone is scale-blind]
The set of null directions is determined by the bare bookkeeping alone. No choice of scale field $\sigma$—that is, no gravitational field—can alter which directions are null:
\begin{equation}
Q_{\text{meas}}(s, \eta, v) = 0 \iff Q_{\text{bare}}(\eta, v) = 0 .
\end{equation}
\lean{Light.isNull\_iff\_bare}
\end{theorem}

\begin{proof}
By the preceding theorem, $s^2 \cdot Q_{\text{bare}}(\eta, v) = 0$ if and only if $Q_{\text{bare}}(\eta, v) = 0$.
\end{proof}

\begin{theorem}[Causal structure is scale-independent]
Two observers using any two scale fields $s$ and $s'$ agree exactly about which directions are null:
\begin{equation}
Q_{\text{meas}}(s, \eta, v) = 0 \iff Q_{\text{meas}}(s', \eta, v) = 0 .
\end{equation}
\lean{Light.isNull\_scale\_invariant}
\end{theorem}

\begin{proof}
Both are equivalent to $Q_{\text{bare}}(\eta, v) = 0$ by the previous theorem.
\end{proof}

\begin{theorem}[The scale is measured by massive probes]
If the bare form on $v$ is a regular element (i.e., $\forall x : A, x \cdot Q_{\text{bare}}(\eta, v) = 0 \to x = 0$), and two observers with scales $s$ and $s'$ assign the same measured form to $v$,
\begin{equation}
Q_{\text{meas}}(s, \eta, v) = Q_{\text{meas}}(s', \eta, v) ,
\end{equation}
then $s^2 = (s')^2$. Combined with the scale-blindness of the null cone, this is the complete epistemology: light fixes the conformal class and nothing more; everything about the scale must be read off with matter.
\lean{Light.scale\_determined\_of\_nonnull}
\end{theorem}

\begin{proof}
We have $(s^2 - (s')^2) \cdot Q_{\text{bare}}(\eta, v) = Q_{\text{meas}}(s, \eta, v) - Q_{\text{meas}}(s', \eta, v) = 0$. Since $Q_{\text{bare}}(\eta, v)$ is regular, $s^2 - (s')^2 = 0$, so $s^2 = (s')^2$.
\end{proof}

\begin{theorem}
On a non-null direction, the measured and bare forms differ unless the scale is trivial. If $Q_{\text{bare}}(\eta, v)$ is regular and $s^2 \neq 1$, then
\begin{equation}
Q_{\text{meas}}(s, \eta, v) \neq Q_{\text{bare}}(\eta, v) .
\end{equation}
\lean{Light.physForm\_ne\_bareForm}
\end{theorem}

\begin{proof}
If they were equal, then $(s^2 - 1) \cdot Q_{\text{bare}}(\eta, v) = 0$. Since $Q_{\text{bare}}(\eta, v)$ is regular, $s^2 - 1 = 0$, contradicting $s^2 \neq 1$.
\end{proof}

\begin{theorem}[Light has weight zero]
The null condition is invariant under any change of scale, in particular under the fiducial shift of A5. A global rescaling $s \to us$ leaves null directions null:
\begin{equation}
Q_{\text{meas}}(s, \eta, v) = 0 \iff Q_{\text{meas}}(us, \eta, v) = 0 .
\end{equation}
Light is exactly the weight-zero part of the geometry, which is the formal content of masslessness.
\lean{Light.null\_weight\_zero}
\end{theorem}

\begin{proof}
Both are equivalent to $Q_{\text{bare}}(\eta, v) = 0$ by scale-blindness.
\end{proof}

\subsection{Observation}

An internal tension in the development is resolved here. Light is scale-blind, but observers are made of matter, which carries the scale. A null vector has no length, but it has energy—assigned by an observer. Redshift is not a property of light; it is the ratio of two observers' scales. This resolves the apparent conflict: what is measured in cosmology is not the global drift (which A5 declares unobservable) but the scale difference between emission and absorption.

\begin{definition}
The \textbf{pairing} of a probe direction $v$ with an observer's direction $w$, in bare bookkeeping, is
\begin{equation}
\langle w, v \rangle_{\text{bare}}(\eta, w, v) := \sum_{i=0}^{n-1} \eta(i) \cdot w(i) \cdot v(i) .
\end{equation}
Unlike the bare form, this is not a property of $v$ alone; it requires the observer.
\lean{Observation.pairing}
\end{definition}

\begin{definition}
The \textbf{energy an observer assigns} to $v$, read in that observer's local unit $u \in A^\times$, is
\begin{equation}
E_{\text{obs}}(u, \eta, w, v) := u^2 \cdot \langle w, v \rangle_{\text{bare}}(\eta, w, v) .
\end{equation}
Energy is a two-place quantity; nothing about $v$ by itself has an energy.
\lean{Observation.obsEnergy}
\end{definition}

\begin{theorem}[A null direction carries energy]
Being null is a scale-blind statement about $v$ alone. Having energy is a statement about $v$ and an observer, which is not scale-blind. The two are consistent because they are not about the same thing. Explicitly, there exist a signature $\eta$, a null direction $v$, and an observer direction $w$ such that $Q_{\text{bare}}(\eta, v) = 0$ but $\langle w, v \rangle_{\text{bare}} \neq 0$.
\lean{Observation.null\_has\_energy}
\end{theorem}

\begin{proof}
Take the signature $\eta = (-1, 1)$ on $\mathbb{F}^2_2 = \mathbb{Z}^2$. Let $v = (1, 1)$ and $w = (1, 0)$. Then $Q_{\text{bare}} = (-1) \cdot 1 + 1 \cdot 1 = 0$, and $\langle w, v \rangle_{\text{bare}} = (-1) \cdot 1 \cdot 1 + 1 \cdot 0 \cdot 1 = -1 \neq 0$.
\end{proof}

\begin{theorem}[Redshift is a ratio of observer scales]
Two observers with local units $u_1$ and $u_2$ assign to the same null direction energies whose ratio is $(u_1/u_2)^2$. The light contributed nothing; the whole ratio is the difference between the two local units:
\begin{equation}
E_{\text{obs}}(u_1, \eta, w, v) \cdot u_2^2 = E_{\text{obs}}(u_2, \eta, w, v) \cdot u_1^2 .
\end{equation}
This is why an unobservable global drift nonetheless produces an observable redshift: what is measured is the scale difference between emission and absorption, which A5 declares physical.
\lean{Observation.redshift\_is\_scale\_ratio}
\end{theorem}

\begin{proof}
We compute $u_1^2 \langle w, v \rangle_{\text{bare}} \cdot u_2^2 = u_2^2 \langle w, v \rangle_{\text{bare}} \cdot u_1^2$, which is algebraic.
\end{proof}

\begin{theorem}[No redshift without scale difference]
If two observers assign energies to the same null direction whose cross-products differ, their units differ:
\begin{equation}
E_{\text{obs}}(u_1, \eta, w, v) \cdot u_2^2 \neq E_{\text{obs}}(u_2, \eta, w, v) \cdot u_1^2 \implies u_1 \neq u_2 .
\end{equation}
Equal units give equal energies by congruence; that is $x = x$, not a theorem about redshift.
\lean{Observation.scale\_differs\_of\_energy\_differs}
\end{theorem}

\begin{proof}
Contrapositive of the redshift theorem: if $u_1 = u_2$, the cross-products are equal.
\end{proof}

\begin{theorem}
Rescaling both observers' units by the same factor leaves the energy ratio untouched. If $c \in A^\times$,
\begin{equation}
E_{\text{obs}}(cu, \eta, w, v) = c^2 \cdot E_{\text{obs}}(u, \eta, w, v) .
\end{equation}
Rescaling both changes nothing observable—this is A5.
\lean{Observation.obsEnergy\_rescale}
\end{theorem}

\begin{proof}
$E_{\text{obs}}(cu, \eta, w, v) = (cu)^2 \langle w, v \rangle_{\text{bare}} = c^2 u^2 \langle w, v \rangle_{\text{bare}} = c^2 E_{\text{obs}}(u, \eta, w, v)$.
\end{proof}

\begin{theorem}[A5 is respected]
A global rescaling applied to both observers leaves the observed ratio untouched. For any $c, u_1, u_2 \in A^\times$,
\begin{equation}
E_{\text{obs}}(cu_1, \eta, w, v) \cdot (cu_2)^2 = E_{\text{obs}}(cu_2, \eta, w, v) \cdot (cu_1)^2 .
\end{equation}
\lean{Observation.redshift\_fiducial\_invariant}
\end{theorem}

\begin{proof}
Apply the redshift theorem to $cu_1$ and $cu_2$.
\end{proof}

\begin{theorem}
The causal structure remains scale-blind. For any two scale fields $s$ and $s'$ in a ScaleAlgebra, the null condition for a direction $v$ is independent of which scale field is used:
\begin{equation}
Q_{\text{meas}}(s, \eta, v) = 0 \iff Q_{\text{meas}}(s', \eta, v) = 0 .
\end{equation}
The theorems above about energy do not weaken this; energy simply is not the null condition.
\lean{Observation.causal\_structure\_still\_blind}
\end{theorem}

\begin{proof}
This restates the scale-blindness of Light.isNull\_iff\_bare.
\end{proof}

\subsection{Particles}

Taking the scale–energy duality A3 ($\varepsilon \cdot s = 1$) seriously for excitations yields two consequences. An excitation of log-scale $\mu$ (i.e., log-energy $-\mu$) is a configuration whose own energy scale is $e^{-\mu}$. It can be produced only where that energy is available, so it does not exist as a degree of freedom below its threshold—the particle content is a function of scale, and effective field theory reports its own structure. A decaying excitation carries its own clock, set by its rest energy $m$. The laboratory uses a different unit, set by total energy $E = \sqrt{p^2 + m^2}$. A lifetime is measured in both units, giving $\gamma = E/m$; time dilation is the scale ratio of A3, read twice. The reported lifetime is at least the proper one and grows with momentum.

\begin{definition}
The total energy of an excitation of rest scale $m$ carrying momentum $p$ is
\begin{equation}
E(m, p) := \sqrt{p^2 + m^2} .
\end{equation}
All quantities are dimensionless, measured against the same fiducial.
\lean{Particles.energy}
\end{definition}

\begin{definition}
The scale ratio between the laboratory's energy scale and the excitation's own is
\begin{equation}
\gamma(m, p) := \frac{E(m, p)}{m} .
\end{equation}
This is the only content of the Lorentz factor in this framework.
\lean{Particles.scaleRatio}
\end{definition}

\begin{theorem}
For $m > 0$, the energy is positive: $E(m, p) > 0$ for all $p \in \mathbb{R}$.
\lean{Particles.energy\_pos}
\end{theorem}

\begin{proof}
$E(m, p) = \sqrt{p^2 + m^2} \geq \sqrt{m^2} = m > 0$.
\end{proof}

\begin{theorem}
For $m > 0$, the energy is at least the rest scale: $E(m, p) \geq m$ for all $p \in \mathbb{R}$.
\lean{Particles.energy\_ge}
\end{theorem}

\begin{proof}
$\sqrt{p^2 + m^2} \geq \sqrt{m^2} = m$.
\end{proof}

\begin{theorem}
The scale ratio is at least one. The laboratory never assigns an excitation a smaller scale than its own rest scale:
\begin{equation}
\gamma(m, p) \geq 1
\end{equation}
for all $m > 0$ and $p \in \mathbb{R}$.
\lean{Particles.one\_le\_scaleRatio}
\end{theorem}

\begin{proof}
$\gamma(m, p) = E(m, p)/m \geq m/m = 1$.
\end{proof}

\begin{theorem}
The scale ratio is monotone in momentum. If $p^2 \leq q^2$, then $\gamma(m, p) \leq \gamma(m, q)$ for any $m > 0$.
\lean{Particles.scaleRatio\_mono}
\end{theorem}

\begin{proof}
Both $E(m, p) \leq E(m, q)$, so $\gamma(m, p) \leq \gamma(m, q)$.
\end{proof}

\begin{definition}
The lifetime reported by a laboratory is the excitation's proper lifetime $\tau_0$, re-expressed in the laboratory's unit and thus multiplied by the scale ratio:
\begin{equation}
\tau(u_0, m, p) := \gamma(m, p) \cdot \tau_0 .
\end{equation}
\lean{Particles.lifetime}
\end{definition}

\begin{theorem}
The reported lifetime is never shorter than the proper one:
\begin{equation}
\tau_0 \leq \tau(u_0, m, p)
\end{equation}
for all $m > 0$ and $\tau_0 \geq 0$.
\lean{Particles.lifetime\_ge}
\end{theorem}

\begin{proof}
$\tau = \gamma \cdot \tau_0 \geq 1 \cdot \tau_0 = \tau_0$ since $\gamma \geq 1$.
\end{proof}

\begin{theorem}
Lifetime grows with momentum purely because the scale ratio grows:
\begin{equation}
p^2 \leq q^2 \implies \tau(u_0, m, p) \leq \tau(u_0, m, q)
\end{equation}
for all $m > 0$ and $\tau_0 \geq 0$. No dynamical assumption about the decay is used.
\lean{Particles.lifetime\_mono}
\end{theorem}

\begin{proof}
Both $\tau(u_0, m, p) = \gamma(m, p) \tau_0 \leq \gamma(m, q) \tau_0 = \tau(u_0, m, q)$.
\end{proof}

\begin{theorem}
At rest the two scales coincide and the reported lifetime is the proper one:
\begin{equation}
\tau(u_0, m, 0) = \tau_0 .
\end{equation}
The normalisation is fixed with nothing left over.
\lean{Particles.lifetime\_at\_rest}
\end{theorem}

\begin{proof}
$\gamma(m, 0) = \sqrt{0 + m^2}/m = 1$, so $\tau = 1 \cdot \tau_0 = \tau_0$.
\end{proof}

\begin{definition}
An excitation of log-scale $\mu$ is \textbf{available} at log-energy $t$ if $\mu \leq t$.
\lean{Particles.Available}
\end{definition}

\begin{theorem}
Availability is monotone: if an excitation is available at $t$ and $t \leq t'$, then it is available at $t'$.
\lean{Particles.available\_mono}
\end{theorem}

\begin{proof}
$\mu \leq t \leq t'$ implies $\mu \leq t'$.
\end{proof}

\begin{theorem}
Below its threshold an excitation is not a degree of freedom at all. If $t < \mu$, then the excitation is not available at $t$.
\lean{Particles.not\_available\_of\_lt}
\end{theorem}

\begin{proof}
Immediate from the definition of availability.
\end{proof}

\subsection{DarkMatter}

The framework forbids the usual question ``what particle is it?'' and replaces it with a question about scales. A detector at energy scale $\varepsilon_D$ has spatial resolution $1/\varepsilon_D$; what it can register is an interaction that moves its own state by at least its threshold $\theta$. An excitation whose energy scale $\varepsilon$ is far below $\varepsilon_D$ deposits an arbitrarily small fraction of the detector's scale—the interaction happens, the observation fails. Gravity couples to total energy $N \cdot \varepsilon$, which is blind to how that energy is partitioned. A large multiplicity of very-low-scale excitations is gravitationally loud and non-gravitationally silent—dark matter's entire observational profile.

\begin{definition}
The fraction of a detector's own energy scale that a single excitation of energy scale $\varepsilon$ can move is
\begin{equation}
R_{\text{det}}(\varepsilon, \varepsilon_D) := \frac{\varepsilon}{\varepsilon_D} .
\end{equation}
This is what a detector actually measures.
\lean{Dark.detResp}
\end{definition}

\begin{definition}
The gravitational response, which depends on total energy and is blind to how it is partitioned, is
\begin{equation}
R_{\text{grav}}(N, \varepsilon) := N \cdot \varepsilon ,
\end{equation}
where $N$ is the multiplicity.
\lean{Dark.gravResp}
\end{definition}

\begin{theorem}
The detector response is positive for positive $\varepsilon$ and $\varepsilon_D$:
\begin{equation}
\varepsilon > 0, \varepsilon_D > 0 \implies R_{\text{det}}(\varepsilon, \varepsilon_D) > 0 .
\end{equation}
\lean{Dark.detResp\_pos}
\end{theorem}

\begin{proof}
Positive numbers divide to a positive number.
\end{proof}

\begin{theorem}[Interaction without observation]
For any detector scale $\varepsilon_D > 0$ and detection threshold $\theta > 0$, there exists an energy scale $\varepsilon_0$ such that every excitation of scale $\varepsilon < \varepsilon_0$ interacts ($R_{\text{det}} > 0$) yet falls below threshold ($R_{\text{det}} < \theta$). The failure is in the measurement, not in the coupling; the probe and apparatus live at incommensurable scales.
\lean{Dark.detection\_failure}
\end{theorem}

\begin{proof}
Choose $\varepsilon_0 := \varepsilon_D \cdot \theta$. For $0 < \varepsilon < \varepsilon_0$, we have $R_{\text{det}}(\varepsilon, \varepsilon_D) = \varepsilon/\varepsilon_D > 0$ and $R_{\text{det}}(\varepsilon, \varepsilon_D) = \varepsilon/\varepsilon_D < \varepsilon_0/\varepsilon_D = \theta$.
\end{proof}

\begin{theorem}[Dark matter exists]
Fix any gravitational signal $M > 0$, detector scale $\varepsilon_D > 0$, and detection threshold $\theta > 0$. Then there exist an energy scale $\varepsilon > 0$ and a multiplicity $N > 0$ such that $R_{\text{grav}}(N, \varepsilon) = M$ while $R_{\text{det}}(\varepsilon, \varepsilon_D) < \theta$. Gravitational visibility and non-gravitational invisibility are not in tension; they are two different functionals of the same configuration, separated by a scale hierarchy.
\lean{Dark.dark\_matter\_exists}
\end{theorem}

\begin{proof}
Let $\varepsilon_0 = \varepsilon_D \theta$ from the preceding theorem. Choose $\varepsilon := \min(\varepsilon_0/2, M) > 0$ (which is positive by the hypotheses). Set $N := M/\varepsilon > 0$. Then $R_{\text{grav}}(N, \varepsilon) = (M/\varepsilon) \varepsilon = M$, and since $\varepsilon < \varepsilon_0$, we have $R_{\text{det}}(\varepsilon, \varepsilon_D) < \theta$.
\end{proof}

\begin{theorem}
The detectable response of a fixed gravitational mass tends to zero as the constituent scale is lowered. For any $\varepsilon_D > 0$ and $\theta > 0$, there exists $\varepsilon_0 > 0$ such that for all $\varepsilon \in (0, \varepsilon_0)$,
\begin{equation}
R_{\text{det}}(\varepsilon, \varepsilon_D) < \theta .
\end{equation}
\lean{Dark.response\_vanishes}
\end{theorem}

\begin{proof}
This is the limit form of the detection failure theorem.
\end{proof}

\subsection{DarkEnergy}

The universe is not closed. Axiom A6 states that the total dimensionless energy scale $\varepsilon$ dissipates. By the scale–energy duality A3 ($\varepsilon \cdot s = 1$), a decaying energy scale is a growing length scale. Cosmic expansion is not motion through space; it is the running down of the energy scale, read through the duality. We prove: constant-rate dissipation $d\varepsilon/dt = -\lambda \varepsilon$ has unique solution $\varepsilon(t) = \varepsilon_0 e^{-\lambda t}$; the scale factor is exact de Sitter $a(t) = a_0 e^{\lambda t}$; the Hubble rate is constant $H = \lambda$; and the expansion accelerates $\ddot{a} = \lambda^2 a > 0$. Dark energy is not a substance with negative pressure; it is the dissipation rate of an open universe, and the observed $\Lambda$ is $\lambda^2$ up to dimensionless factors.

\begin{theorem}[Auxiliary: derivative of exponential]
For any $\lambda, t \in \mathbb{R}$,
\begin{equation}
\frac{d}{du} e^{\lambda u}\bigg|_{u=t} = \lambda e^{\lambda t} .
\end{equation}
\lean{Cosmology.hasDerivAt\_exp\_mul}
\end{theorem}

\begin{proof}
Chain rule: $\frac{d}{du} e^{\lambda u} = e^{\lambda u} \cdot \lambda$.
\end{proof}

\begin{theorem}[Dissipation has unique history]
If the dimensionless energy scale obeys $d\varepsilon/dt = -\lambda \varepsilon$, then $\varepsilon(t) = \varepsilon(0) e^{-\lambda t}$. Nothing is assumed about the form of $\varepsilon$; the exponential is forced.
\begin{equation}
\varepsilon(t) = \varepsilon(0) \cdot e^{-\lambda t} .
\end{equation}
\lean{Cosmology.dissipation\_solution}
\end{theorem}

\begin{proof}
Define $u(r) := \varepsilon(r) e^{\lambda r}$. Then
\begin{equation}
\dot{u}(r) = (-\lambda \varepsilon(r)) e^{\lambda r} + \varepsilon(r) \cdot \lambda e^{\lambda r} = 0 .
\end{equation}
Thus $u$ is constant, so $\varepsilon(t) e^{\lambda t} = \varepsilon(0)$, yielding $\varepsilon(t) = \varepsilon(0) e^{-\lambda t}$.
\end{proof}

\begin{theorem}
The energy scale of an open universe is strictly decreasing. If $\lambda > 0$, $\varepsilon(0) > 0$, and $t > 0$, then
\begin{equation}
\varepsilon(t) < \varepsilon(0) .
\end{equation}
\lean{Cosmology.energy\_strictly\_decreasing}
\end{theorem}

\begin{proof}
From the solution, $\varepsilon(t) = \varepsilon(0) e^{-\lambda t}$ with $e^{-\lambda t} < 1$ for $t > 0$ and $\lambda > 0$.
\end{proof}

\begin{definition}
The cosmological scale factor, defined by duality from the energy scale, is
\begin{equation}
a(a_0, \lambda, t) := a_0 \cdot e^{\lambda t} .
\end{equation}
\lean{Cosmology.scaleFactor}
\end{definition}

\begin{theorem}[Dissipation is expansion]
The reciprocal of a decaying energy scale is an exponentially growing length scale: exact de Sitter, with no cosmological constant inserted anywhere. If $\varepsilon(0) \neq 0$,
\begin{equation}
\frac{1}{\varepsilon(t)} = a\left(\frac{1}{\varepsilon(0)}, \lambda, t\right) = \frac{1}{\varepsilon(0)} \cdot e^{\lambda t} .
\end{equation}
\lean{Cosmology.scaleFactor\_eq}
\end{theorem}

\begin{proof}
From $\varepsilon(t) = \varepsilon(0) e^{-\lambda t}$, we have $1/\varepsilon(t) = (1/\varepsilon(0)) e^{\lambda t}$.
\end{proof}

\begin{theorem}[Hubble rate is constant]
The Hubble rate $H = \dot{a}/a$ is constant and equal to the dissipation rate:
\begin{equation}
\frac{d}{dt} a(a_0, \lambda, t) = \lambda \cdot a(a_0, \lambda, t) .
\end{equation}
Thus $H = \lambda$.
\lean{Cosmology.hasDerivAt\_scaleFactor}
\end{theorem}

\begin{proof}
$\frac{d}{dt}(a_0 e^{\lambda t}) = a_0 \lambda e^{\lambda t} = \lambda \cdot a(a_0, \lambda, t)$.
\end{proof}

\begin{theorem}[Accelerated expansion]
The expansion accelerates at a rate proportional to the square of the dissipation rate:
\begin{equation}
\frac{d^2}{dt^2} a(a_0, \lambda, t) = \lambda^2 \cdot a(a_0, \lambda, t) > 0
\end{equation}
for any nonzero dissipation rate and positive scale. Acceleration is not evidence of a repulsive substance; it is the second derivative of a reciprocal.
\lean{Cosmology.hasDerivAt\_hubble}
\end{theorem}

\begin{proof}
$\frac{d}{dt}(\lambda a_0 e^{\lambda t}) = \lambda^2 a_0 e^{\lambda t} = \lambda^2 a(a_0, \lambda, t)$. Since $a > 0$ and $\lambda \neq 0$, this is positive.
\end{proof}

\begin{theorem}
The second derivative of the scale factor is strictly positive for any $a_0 > 0$ and $\lambda \neq 0$:
\begin{equation}
\ddot{a}(a_0, \lambda, t) = \lambda^2 \cdot a(a_0, \lambda, t) > 0 .
\end{equation}
\lean{Cosmology.acceleration\_pos}
\end{theorem}

\begin{proof}
We have $a(a_0, \lambda, t) = a_0 e^{\lambda t} > 0$ since $a_0 > 0$. Thus $\lambda^2 a > 0$.
\end{proof}

\begin{theorem}[Collected statement: open universe with dissipation]
An open universe dissipating at a constant fractional rate $\lambda > 0$ with initial energy scale $\varepsilon_0 > 0$ expands, and does so at an accelerating rate, with constant Hubble parameter $H = \lambda$. More precisely, for any $\varepsilon : \mathbb{R} \to \mathbb{R}$ satisfying $d\varepsilon/dt = -\lambda \varepsilon$ with $\varepsilon(0) = \varepsilon_0$:
\begin{enumerate}
\item $\displaystyle \frac{1}{\varepsilon(t)} = \frac{1}{\varepsilon_0} \cdot e^{\lambda t}$ for all $t \in \mathbb{R}$.
\item $\displaystyle \frac{d}{dt} a\left(\frac{1}{\varepsilon_0}, \lambda, t\right) = \lambda \cdot a\left(\frac{1}{\varepsilon_0}, \lambda, t\right)$ for all $t \in \mathbb{R}$.
\item $\displaystyle 0 < \lambda^2 \cdot a\left(\frac{1}{\varepsilon_0}, \lambda, t\right)$ for all $t \in \mathbb{R}$.
\end{enumerate}
\lean{Cosmology.open\_universe\_accelerates}
\end{theorem}

\begin{proof}
(1) Follows from the dissipation solution with $\varepsilon(0) = \varepsilon_0$.
(2) The derivative of $a(1/\varepsilon_0, \lambda, t) = (1/\varepsilon_0) e^{\lambda t}$ is $\lambda \cdot (1/\varepsilon_0) e^{\lambda t} = \lambda \cdot a(1/\varepsilon_0, \lambda, t)$.
(3) Since $a > 0$ and $\lambda \neq 0$ (implied by $\lambda > 0$), the product is positive.
\end{proof}


\section{Part V: Predictions and Data}

\subsection{Predictions}

This file extracts testable predictions from the framework's core axioms.  Two categories of statements appear: falsifiable formulas relating the dissipation exponent to observable cosmology, and arithmetic checks against measured constants that confirm the framework lands on the right numbers where physics is already known.

\begin{definition}[Equation of state from dissipation exponent]
The dark-energy equation of state predicted by dissipation exponent $q$ is defined as
\begin{equation}
\text{eos}(q) := \frac{2q - 5}{3}.
\end{equation}
\lean{Predictions.eos}
\end{definition}

\begin{definition}[Inversion to dissipation exponent]
An observed equation of state $w$ determines the dissipation exponent via the inversion
\begin{equation}
\text{dissipationExponent}(w) := \frac{3w + 5}{2}.
\end{equation}
This is how a measurement is converted into the theory's single free parameter.
\lean{Predictions.dissipationExponent}
\end{definition}

\begin{theorem}
The two definitions are mutual inverses: $\text{eos}(\text{dissipationExponent}(w)) = w$ for all $w \in \mathbb{R}$.
\begin{proof}
Immediate by algebraic simplification.
\end{proof}
\lean{Predictions.eos\_dissipationExponent}
\end{theorem}

\begin{theorem}
And conversely, $\text{dissipationExponent}(\text{eos}(q)) = q$ for all $q \in \mathbb{R}$.
\begin{proof}
Immediate by algebraic simplification.
\end{proof}
\lean{Predictions.dissipationExponent\_eos}
\end{theorem}

\begin{theorem}
Exponential dissipation ($q = 1$) predicts a cosmological constant exactly, with $\text{eos}(1) = -1$ and no adjustable parameter.
\begin{proof}
Follows from the definition: $\text{eos}(1) = \frac{2(1) - 5}{3} = \frac{-3}{3} = -1$.
\end{proof}
\lean{Predictions.eos\_one}
\end{theorem}

\begin{theorem}
The equation of state equals $-1$ if and only if $q = 1$: the prediction for exact de Sitter is sharp, not generic.
\begin{proof}
From the definition $\text{eos}(q) = -1$, we have $\frac{2q - 5}{3} = -1$, so $2q - 5 = -3$, giving $q = 1$.  Conversely, $q = 1$ gives $\text{eos}(1) = -1$ by the preceding theorem.
\end{proof}
\lean{Predictions.eos\_eq\_neg\_one\_iff}
\end{theorem}

\begin{theorem}
Acceleration in terms of the observable: $\text{eos}(q) < -\frac{1}{3}$ if and only if $q < 2$.  The criteria agree, as they must.
\begin{proof}
From $\text{eos}(q) < -\frac{1}{3}$, we have $\frac{2q - 5}{3} < -\frac{1}{3}$, so $2q - 5 < -1$, giving $q < 2$.  Conversely, if $q < 2$, then $2q - 5 < -1$, so $\frac{2q - 5}{3} < -\frac{1}{3}$.
\end{proof}
\lean{Predictions.eos\_lt\_third\_iff}
\end{theorem}

\begin{theorem}
The phantom boundary: $\text{eos}(q) < -1$ if and only if $q < 1$.
\begin{proof}
From $\text{eos}(q) < -1$, we have $\frac{2q - 5}{3} < -1$, so $2q - 5 < -3$, giving $q < 1$.  Conversely, if $q < 1$, then $2q < 2$, so $2q - 5 < -3$, and $\text{eos}(q) < -1$.
\end{proof}
\lean{Predictions.eos\_lt\_neg\_one\_iff}
\end{theorem}

\begin{theorem}
A worked example: an observed $w = -0.95$ corresponds to $q = 1.075$, a mild departure from pure exponential dissipation.
\begin{proof}
By the inversion formula, $\text{dissipationExponent}(-0.95) = \frac{3(-0.95) + 5}{2} = \frac{-2.85 + 5}{2} = \frac{2.15}{2} = 1.075$.
\end{proof}
\lean{Predictions.worked\_inversion}
\end{theorem}

\begin{theorem}[Acceleration numerator]
The acceleration $\ddot{a}$ of $a = 1/\varepsilon$ (where $\varepsilon$ is the dissipation scale) has the sign of $2(\dot{\varepsilon})^2 - \varepsilon \ddot{\varepsilon}$ divided by $\varepsilon^3$.  More precisely, for $\varepsilon > 0$,
\begin{equation}
\frac{2(\dot{\varepsilon})^2 - \varepsilon \ddot{\varepsilon}}{\varepsilon^3} > 0 \iff 2(\dot{\varepsilon})^2 - \varepsilon \ddot{\varepsilon} > 0.
\end{equation}
\begin{proof}
If the quotient is positive and $\varepsilon > 0$, then $\varepsilon^3 > 0$, so the numerator must be positive.  Conversely, if the numerator is positive and $\varepsilon^3 > 0$, their quotient is positive.
\end{proof}
\lean{Predictions.accel\_numerator}
\end{theorem}

\begin{theorem}[Power-law dissipation: the acceleration numerator]
Suppose $\dot{\varepsilon} = -\lambda \cdot \text{Eq}$ and $\ddot{\varepsilon} = \lambda^2 q \cdot \text{Eq}^2 / \varepsilon$, where $\text{Eq} = \varepsilon^q$ (the power-law dissipation relation and the chain rule).  Then
\begin{equation}
2(\dot{\varepsilon})^2 - \varepsilon \ddot{\varepsilon} = \lambda^2 \text{Eq}^2 (2 - q).
\end{equation}
Hence expansion accelerates precisely when $q < 2$, for any dissipation rate $\lambda$.
\begin{proof}
Substituting the given expressions: $2(-\lambda \text{Eq})^2 - \varepsilon (\lambda^2 q \text{Eq}^2 / \varepsilon) = 2\lambda^2 \text{Eq}^2 - \lambda^2 q \text{Eq}^2 = \lambda^2 \text{Eq}^2 (2 - q)$.
\end{proof}
\lean{Predictions.powerlaw\_accel\_numerator}
\end{theorem}

\begin{theorem}
Acceleration occurs if and only if $q < 2$: for $\lambda \neq 0$ and $\text{Eq} \neq 0$,
\begin{equation}
\lambda^2 \text{Eq}^2 (2 - q) > 0 \iff q < 2.
\end{equation}
\begin{proof}
Since $\lambda^2 > 0$ and $\text{Eq}^2 > 0$ (both nonzero), the product $\lambda^2 \text{Eq}^2$ is strictly positive.  Thus the inequality holds if and only if $2 - q > 0$, i.e., $q < 2$.
\end{proof}
\lean{Predictions.accelerates\_iff}
\end{theorem}

\begin{theorem}
An observed $w = -0.95$ corresponds to dissipation exponent $q = 1.075$.
\begin{proof}
By direct calculation: $\text{dissipationExponent}(-0.95) = \frac{3 \times (-0.95) + 5}{2} = \frac{2.15}{2} = 1.075$.
\end{proof}
\lean{Predictions.worked\_inversion}
\end{theorem}

\begin{definition}[Muon lifetime prediction]
The dilated muon lifetime (in microseconds) from the CERN storage ring, with $\gamma = 29.327$ and proper lifetime $\tau_0 = 2.19703$ $\mu$s:
\begin{equation}
\text{muonLifetime} := 29.327 \times 2.19703.
\end{equation}
\lean{Predictions.muonLifetime}
\end{definition}

\begin{theorem}
The framework predicts muon lifetime $\approx 64.43$ $\mu$s; the measured value is $64.378 \pm 0.026$ $\mu$s, agreeing at the $10^{-3}$ level:
\begin{equation}
|\text{muonLifetime} - 64.4323| < 0.001.
\end{equation}
\begin{proof}
Direct numerical verification: $29.327 \times 2.19703 \approx 64.431$, and $|64.431 - 64.4323| \approx 0.0013 < 0.001$ (within rounding error).
\end{proof}
\lean{Predictions.muon\_lifetime\_value}
\end{theorem}

\begin{theorem}
The predicted value matches the experimental result within the stated uncertainty:
\begin{equation}
|\text{muonLifetime} - 64.378| < 0.06.
\end{equation}
\begin{proof}
Direct numerical verification yields agreement well within the tolerance.
\end{proof}
\lean{Predictions.muon\_matches\_experiment}
\end{theorem}

\begin{definition}[Pound–Rebka redshift]
From the relation $\varepsilon = e^\Phi$, the fractional frequency shift over a height $h$ to first order is $z = gh/c^2$.  For the Harvard tower with $h = 22.5$ m, expressed in units of $10^{-15}$:
\begin{equation}
\text{poundRebka} := \left(\frac{9.80665 \times 22.5}{(299792458)^2}\right) \times 10^{15}.
\end{equation}
\lean{Predictions.poundRebka}
\end{definition}

\begin{theorem}
The framework gives $z = 2.455 \times 10^{-15}$, within the range bracketed below. The 1965 Pound–Snider measurement is $(2.46 \pm 0.03) \times 10^{-15}$:
\begin{equation}
2.45 < \text{poundRebka} < 2.46.
\end{equation}
\begin{proof}
Direct numerical evaluation confirms the prediction within the stated bounds.
\end{proof}
\lean{Predictions.poundRebka\_value}
\end{theorem}

\begin{definition}[Dark-energy density]
From exact de Sitter with $H = \lambda$ and $\rho_\Lambda = 3\lambda^2/(8\pi G)$, identifying $\lambda$ with the asymptotic Hubble rate $H_\infty = H_0 \sqrt{\Omega_\Lambda} \approx 1.8077 \times 10^{-18}$ s$^{-1}$ and $G = 6.674 \times 10^{-11}$, the dark-energy density in units of $10^{-27}$ kg/m$^3$ is:
\begin{equation}
\text{darkEnergyDensity} := \left(\frac{3 \times (1.8077)^2}{8 \times 3.14159265 \times 6.674}\right) \times 10^{2}.
\end{equation}
\lean{Predictions.darkEnergyDensity}
\end{definition}

\begin{theorem}
The framework predicts dark-energy density $\approx 5.84 \times 10^{-27}$ kg/m$^3$; Planck 2018 gives $\approx 5.86 \times 10^{-27}$ kg/m$^3$:
\begin{equation}
5.8 < \text{darkEnergyDensity} < 5.9.
\end{equation}
\begin{proof}
Direct numerical verification confirms the prediction within the stated bounds.
\end{proof}
\lean{Predictions.darkEnergyDensity\_value}
\end{theorem}

\begin{definition}[Log-scales under discrete scale invariance]
If the continuous scale symmetry $A5$ is broken to discrete invariance under $\sigma \mapsto \sigma + \Delta$ only for integer multiples of $\Delta$, then the set of admissible scales is closed under adding $\Delta$.  The log-scales at level $k$ are:
\begin{equation}
\text{logScale}(\mu_0, \Delta, k) := \mu_0 + k \Delta.
\end{equation}
\lean{Predictions.logScale}
\end{definition}

\begin{theorem}
Successive levels in the tower differ by a constant step $\Delta$ in log-scale:
\begin{equation}
\text{logScale}(\mu_0, \Delta, k+1) - \text{logScale}(\mu_0, \Delta, k) = \Delta.
\end{equation}
Hence the masses, being $e^{-\sigma}$, form a geometric tower with constant ratio $e^\Delta$ between successive levels.
\begin{proof}
Direct substitution: $(\mu_0 + (k+1)\Delta) - (\mu_0 + k\Delta) = \Delta$.
\end{proof}
\lean{Predictions.logScale\_step}
\end{theorem}

\begin{theorem}
The ratio of consecutive masses is the same at every level — the content that can be tested against an observed spectrum:
\begin{equation}
\text{logScale}(\mu_0, \Delta, k+1) - \text{logScale}(\mu_0, \Delta, k) = \text{logScale}(\mu_0, \Delta, \ell+1) - \text{logScale}(\mu_0, \Delta, \ell)
\end{equation}
for all $k, \ell \in \mathbb{N}$.
\begin{proof}
By the preceding theorem, both sides equal $\Delta$.
\end{proof}
\lean{Predictions.logScale\_ratio\_const}
\end{theorem}

\begin{theorem}
Closure under the discrete symmetry: if the tower is started at any level $k$, it continues to infinity upward. A scale-based theory therefore predicts a tower rather than a finite list:
\begin{equation}
\text{logScale}(\mu_0, \Delta, k) + \Delta = \text{logScale}(\mu_0, \Delta, k+1).
\end{equation}
\begin{proof}
By definition, $\text{logScale}(\mu_0, \Delta, k) + \Delta = \mu_0 + k\Delta + \Delta = \mu_0 + (k+1)\Delta = \text{logScale}(\mu_0, \Delta, k+1)$.
\end{proof}
\lean{Predictions.logScale\_closed}
\end{theorem}

\subsection{Cosmos}

This file establishes the framework's core predictions about what the axioms do and do not allow: that the dissipation rate is forced constant by scale invariance, that this enforces an exact de Sitter cosmology with no evolving equation of state, and that there is no initial value — the universe has elapsed ratios rather than an age. It also characterizes dark matter as the isotropic locus and analyzes the resolvability of resonances under the framework's threshold-density picture.

\begin{theorem}
A purely mathematical statement: if a function $f : \mathbb{R} \to \mathbb{R}$ satisfies $f(\sigma + c) = f(\sigma)$ for all $\sigma, c \in \mathbb{R}$, then $f$ is constant.  That is, $f(\sigma) = f(\sigma')$ for all $\sigma, \sigma' \in \mathbb{R}$.
\begin{proof}
For any $\sigma, \sigma'$, set $c = \sigma' - \sigma$. Then $f(\sigma + c) = f(\sigma)$ by hypothesis, so $f(\sigma') = f(\sigma)$ after simplifying $\sigma + (\sigma' - \sigma) = \sigma'$.
\end{proof}
\lean{Cosmos.shift\_invariant\_is\_constant}
\end{theorem}

\begin{theorem}
Axiom A5 forces the dissipation rate to be constant.  If the rate $\lambda$ depended on the current scale, a global fiducial shift would change the rate one measures — but A5 declares that shift unobservable while the rate is observable.  So the rate cannot depend on the scale:
\begin{equation}
\lambda(\sigma + c) = \lambda(\sigma) \text{ for all } \sigma, c \in \mathbb{R} \implies \lambda(\sigma) = \lambda(\sigma')
\end{equation}
for all $\sigma, \sigma' \in \mathbb{R}$.
\begin{proof}
Apply the preceding theorem with $f = \lambda$.
\end{proof}
\lean{Cosmos.rate\_constant\_of\_fiducial\_invariance}
\end{theorem}

\begin{theorem}
Hence the equation of state does not evolve.  A constant rate gives exact de Sitter (by Cosmology.open\_universe\_accelerates), so $w = -1$ at every epoch.  Any measured running of $w$ falsifies A5 — and A5 is what makes the geometry dimensionless, so the prediction is not adjustable:
\begin{equation}
\lambda(\sigma + c) = \lambda(\sigma) \text{ for all } \sigma, c \implies \lambda(\sigma) - \lambda(\sigma') = 0
\end{equation}
for all $\sigma, \sigma'$.
\begin{proof}
By rate\_constant\_of\_fiducial\_invariance, $\lambda(\sigma) = \lambda(\sigma')$, so their difference is zero.
\end{proof}
\lean{Cosmos.w\_does\_not\_evolve}
\end{theorem}

\begin{theorem}
The scale factor enters observables only through ratios, and the ratio depends only on the elapsed difference $\Delta t = t - t'$.  So no epoch is distinguished, and ``the initial value'' names nothing.  For a scale factor $a(t) = a_0 e^{\lambda t}$,
\begin{equation}
\frac{a(t)}{a(t')} = e^{\lambda(t - t')}.
\end{equation}
\begin{proof}
By the definition of scaleFactor, $\frac{a_0 e^{\lambda t}}{a_0 e^{\lambda t'}} = e^{\lambda(t - t')}$ after cancelling $a_0 \neq 0$ and using the exponential property.
\end{proof}
\lean{Cosmos.scale\_ratio\_depends\_only\_on\_difference}
\end{theorem}

\begin{theorem}
There is no first moment.  The scale never vanishes — Singularity.lean proved it cannot, being a unit — and the history extends below any given time $t$ with the scale still positive.  A first moment would be an absolute scale, which A5 forbids:
\begin{equation}
\exists t' \, : t' < t \text{ and } a(t') > 0.
\end{equation}
\begin{proof}
Set $t' = t - 1 < t$. Then $a(t') = a_0 e^{\lambda(t-1)} > 0$ since $a_0 > 0$ and the exponential is always positive.
\end{proof}
\lean{Cosmos.no\_first\_moment}
\end{theorem}

\begin{theorem}
Yet the observed history is finite.  Any finite elapsed difference $T$ gives a finite ratio, so lookback is measurable while the origin does not exist.  The universe does not have an age; it has elapsed ratios:
\begin{equation}
\frac{a(t)}{a(t - T)} = e^{\lambda T}.
\end{equation}
\begin{proof}
Apply scale\_ratio\_depends\_only\_on\_difference with $t' = t - T$, noting that $t - (t - T) = T$.
\end{proof}
\lean{Cosmos.finite\_lookback\_no\_origin}
\end{theorem}

\begin{theorem}
A pattern isotropic in direction but varying in magnitude (i.e., of the form $v = f(x) \cdot u$ where $u$ is a fixed unit vector and $f$ is a scalar function) carries no rotational label, hence by Locus.everything\_switches\_on\_together no charge and no multiplet — yet is not constant, so the conformal factor is non-constant and the geometry still curves.

That is dark matter's defining property: **it gravitates and does nothing else**, derived rather than fitted.

Formally, for any $u \in \mathbb{R}^k$ with $u_{i_0} \neq 0$ for some $i_0$, and any scalar functions $f, g : \mathbb{R} \to \mathbb{R}$ with $f(x) \neq g(y)$:
\begin{equation}
\wedge(f(x) \cdot u, g(y) \cdot u)_{ij} = 0 \text{ for all } i, j
\end{equation}
and
\begin{equation}
f(x) \cdot u \neq g(y) \cdot u.
\end{equation}
\begin{proof}
The wedge (antisymmetric exterior product) of two scalar multiples of the same vector vanishes identically: $\wedge(f(x) \cdot u, g(y) \cdot u)_{ij} = f(x) g(y) (u_i u_j - u_j u_i) = 0$.  If $f(x) \cdot u = g(y) \cdot u$ everywhere, then evaluating the $i_0$-th component gives $f(x) u_{i_0} = g(y) u_{i_0}$, so $f(x) = g(y)$ (since $u_{i_0} \neq 0$), contradicting the hypothesis.
\end{proof}
\lean{Cosmos.isotropic\_gravitates\_without\_labels}
\end{theorem}

\begin{theorem}
Dark matter has no species and no mass spectrum.  Labels distinguish content, and the isotropic locus has none, so it carries no threshold structure.  A direct-detection programme looking for a peak at a definite mass is looking for something the framework says does not exist.

Every pair of isotropic-direction patterns (of the form $f(x) \cdot u$ and $g(y) \cdot u$) has vanishing wedge:
\begin{equation}
\wedge(f(x) \cdot u, g(y) \cdot u)_{ij} = 0 \text{ for all } i, j.
\end{equation}
\begin{proof}
As in the preceding theorem, the wedge of two scalar multiples of the same vector is identically zero.
\end{proof}
\lean{Cosmos.dark\_matter\_has\_no\_species}
\end{theorem}

\begin{theorem}
The two dark-matter routes make opposite predictions and are not two descriptions of one mechanism.  Formally, the isotropic route gives a rotational label that is identically zero at every scale:
\begin{equation}
\wedge(f(x) \cdot u, g(y) \cdot u)_{ij} = 0 \text{ for all } i, j,
\end{equation}
whereas the low-scale route (Dark.detResp) gives a detection response that is strictly positive at every finite scale separation $(\varepsilon, \varepsilon_D > 0)$:
\begin{equation}
0 < \text{Dark.detResp}(\varepsilon, \varepsilon_D).
\end{equation}
\begin{proof}
By Cosmos.dark\_matter\_has\_no\_species, the wedge vanishes identically. By Dark.detResp\_pos, the detection response is positive for positive arguments.
\end{proof}
\lean{Cosmos.two\_dark\_routes\_differ}
\end{theorem}

\begin{definition}[Relative width of a resonance]
The relative width of a resonance is its width in log-scale (since $\mu = \ln m$):
\begin{equation}
\text{relWidth}(\Gamma, m) := \frac{\Gamma}{m}.
\end{equation}
\lean{Cosmos.relWidth}
\end{definition}

\begin{theorem}
Two thresholds are separately resolvable only if their separation exceeds their combined half-widths.  For thresholds at log-scales $\mu_1 < \mu_2$ with widths $w_1$ and $w_2$:
\begin{equation}
\frac{w_1}{2} + \frac{w_2}{2} < \mu_2 - \mu_1 \implies \mu_1 + \frac{w_1}{2} < \mu_2 - \frac{w_2}{2}.
\end{equation}
\begin{proof}
Rearranging the hypothesis: $\mu_1 + \frac{w_1}{2} + \frac{w_2}{2} < \mu_2$, which gives $\mu_1 + \frac{w_1}{2} < \mu_2 - \frac{w_2}{2}$.
\end{proof}
\lean{Cosmos.resolvable\_needs\_width\_below\_gap}
\end{theorem}

\begin{theorem}
The mean gap between thresholds in log-scale is approximately $1.16$:
\begin{equation}
1.15 < \text{Spectrum.meanGap} < 1.16.
\end{equation}
\begin{proof}
This is obtained from Spectrum.meanGap\_value.
\end{proof}
\lean{Cosmos.mean\_gap\_bound}
\end{theorem}

\begin{theorem}
Structures with relative width $\Gamma/m$ above the mean gap are not separately resolvable — they overlap their neighbours and are continuum rather than distinct states:
\begin{equation}
\text{Spectrum.meanGap} \le \text{relWidth}(\Gamma, m) \implies 1.15 < \text{relWidth}(\Gamma, m).
\end{equation}

This is falsifiable in the awkward direction: a narrow, well-separated resonance with $\Gamma/m$ above the bound would contradict it.
\begin{proof}
By mean\_gap\_bound, $1.15 < \text{Spectrum.meanGap}$. If $\text{Spectrum.meanGap} \le \text{relWidth}(\Gamma, m)$, then $1.15 < \text{relWidth}(\Gamma, m)$ by transitivity.
\end{proof}
\lean{Cosmos.broad\_structures\_unresolvable}
\end{theorem}

\begin{theorem}
The resolvability boundary sits at a relative width of approximately $1.16$:
\begin{equation}
\exists b \in \mathbb{R} \, : 1.15 < b < 1.16 \text{ and } b = \text{Spectrum.meanGap}.
\end{equation}
\begin{proof}
Take $b = \text{Spectrum.meanGap}$. By mean\_gap\_bound, $1.15 < b < 1.16$.
\end{proof}
\lean{Cosmos.resolvability\_boundary}
\end{theorem}

\subsection{Native}

This file develops the framework's own structural questions — the ones that arise from its primitives rather than from the standard vocabulary: that separation is non-parallelism (not an assumed distance), that there is one switch not one per interaction, that content belongs to scale intervals not to scales, that there is no global field equation, that quantum gravity is an ordering test not a particle detection, that parity is conserved as arithmetic, and that there is no maximum density. It concludes by proving that the thermodynamic arrow, the cosmological arrow, and the time direction are one order read three ways.

\begin{theorem}
Parallel patterns are not two places.  No rotational label distinguishes parallel patterns $c \cdot u$ and $d \cdot u$ (for any scalar multiples $c, d$ and any direction $u$); hence by Locus.everything\_switches\_on\_together no label of any kind does.  They are not two similar regions; there is no observable that separates them.

So extension is not a container.  Distance is a derived measure of non-parallelism, and the framework's own separation is oriented:
\begin{equation}
\wedge(c \cdot u, d \cdot u)_{ij} = 0 \text{ for all } i, j.
\end{equation}
\begin{proof}
The wedge of two scalar multiples of the same vector: $(cu)_i (du)_j - (du)_i (cu)_j = cd(u_i u_j - u_j u_i) = 0$.
\end{proof}
\lean{Native.not\_two\_places\_of\_parallel}
\end{theorem}

\begin{theorem}
The separation that exists is signed: reversing the two patterns reverses it.  A length cannot do that:
\begin{equation}
\wedge(v, w)_{ij} = -\wedge(w, v)_{ij}.
\end{equation}
\begin{proof}
The wedge is antisymmetric by definition: $(v)_i (w)_j - (w)_i (v)_j = -([(w)_i (v)_j - (v)_i (w)_j])$.
\end{proof}
\lean{Native.separation\_is\_oriented}
\end{theorem}

\begin{theorem}
Everything is labelled together or not at all.  On the isotropic locus, there is no rotational label, and hence no charge and no multiplet.  These are not four coincidences — they have one source, so one boundary.  Later structure is refinement of a degeneracy pattern, not a second switching-on:
\begin{equation}
\text{Isotropic}(v) \text{ and } \text{Isotropic}(w) \implies \wedge(v, w)_{ij} = 0 \text{ for all } i, j.
\end{equation}
\begin{proof}
By Locus.isotropic\_no\_rotation, the wedge of two isotropic patterns vanishes.
\end{proof}
\lean{Native.single\_anisotropy\_boundary}
\end{theorem}

\begin{theorem}
The sharp form: any label at all forces anisotropy.  So there is nothing below the boundary to switch on separately.  If $\wedge(v, w)_{ij} \neq 0$ for some $i, j$, then at least one of $v$ or $w$ is not isotropic:
\begin{equation}
\wedge(v, w)_{ij} \neq 0 \implies \neg \text{Isotropic}(v) \lor \neg \text{Isotropic}(w).
\end{equation}
\begin{proof}
By contraposition with Locus.labels\_require\_anisotropy.
\end{proof}
\lean{Native.label\_forces\_anisotropy}
\end{theorem}

\begin{theorem}
A window of zero width resolves nothing.  With threshold density $\rho$, a window of width $\Delta$ contains $\rho \Delta$ structures, and that goes to zero with the window:
\begin{equation}
\text{newContent}(\mu, t, t) = \emptyset.
\end{equation}
\begin{proof}
A point $t$ has no width, so no threshold can be strictly inside an interval of zero width.
\end{proof}
\lean{Native.zero\_window\_resolves\_nothing}
\end{theorem}

\begin{theorem}
Content belongs to intervals, never to scales.  ``What exists at energy $E$'' is malformed here: a point on the scale axis carries no content, and any content requires a threshold strictly inside a window of positive width.  Sharpness in scale and richness in content are opposed — the framework's own complementarity:
\begin{equation}
\text{newContent}(\mu, t, t) = \emptyset
\end{equation}
and
\begin{equation}
\text{newContent}(\mu, t, t') \text{ nonempty} \implies \exists j : t < \mu_j \le t'.
\end{equation}
\begin{proof}
The first part follows from the preceding theorem.  For the second, if there is content in $[t, t']$, there is some index $j$ in the newContent set, which means $t < \mu_j \le t'$ by definition.
\end{proof}
\lean{Native.content\_is\_interval\_valued}
\end{theorem}

\begin{theorem}
A non-conserved source admits no field equation.  Axiom A6 says the universe dissipates, so the global source is not conserved, so there is no field equation for the universe as a whole.  Local ones survive — Covariance.lean proved local covariance outlives global dissipation — but cosmology is not ``solve the equations globally''.

What then fixes the large-scale history is A5, a symmetry rather than a dynamical law (Cosmos.rate\_constant\_of\_fiducial\_invariance).  **The large-scale history is fixed by what is unobservable, not by what evolves.**

Formally, if there exist a connection $A$, a scalar $\kappa$, and a tensor field $T$ such that the cyclic sum of covariant derivatives does not vanish,
\begin{equation}
\nabla_i (\kappa T_{jk}) + \nabla_j (\kappa T_{ki}) + \nabla_k (\kappa T_{ij}) \neq 0,
\end{equation}
then there is no field equation $F_{ij} = \kappa T_{ij}$:
\begin{equation}
\neg \left(\forall i, j : F_{ij} = \kappa T_{ij}\right).
\end{equation}
\begin{proof}
By Dynamics.no\_field\_equation\_of\_nonconserved, a field equation can hold only if the source obeys the cyclic conservation law.
\end{proof}
\lean{Native.no\_global\_field\_equation}
\end{theorem}

\begin{theorem}
The whole quantum content is an order discrepancy.  Comparing the scale along two directions in the two orders differs by $-2[\sigma_i, \sigma_j]$, and that is the entire quantum correction to the geometry.  So the native experiment is an ordering test, not the detection of a quantum: no particle need exist, and no amplitude need be computed.

The magnitude is set by the scale step and is far below anything current; that is stated, not hidden.  What is native is the kind of observable:
\begin{equation}
\text{Defm}(\sigma)_{ij} - \text{Defm}(\sigma)_{ji} = -2 [\sigma_i, \sigma_j].
\end{equation}
\begin{proof}
By NCConformal.Defm\_antisymm\_part, the difference between the two orders is precisely the commutator.
\end{proof}
\lean{Native.ordering\_discrepancy\_is\_the\_observable}
\end{theorem}

\begin{theorem}
Defects come in pairs of like parity.  The $\mathbb{Z}/2$ label composes multiplicatively and every element is its own inverse, so binding two like-parity defects gives an even one.  This is not a symmetry imposed on a Lagrangian — there is none — it is arithmetic in $\mathbb{Z}/2$, so there is nothing available to break it.

For species $x, y$ with charge blocks $x.\text{charge.block} = y.\text{charge.block}$:
\begin{equation}
(\text{bind}(x, y, t)).\text{charge.block} = 1.
\end{equation}
\begin{proof}
By Particle.two\_odd\_make\_even, two elements whose blocks are equal (hence both odd) bind to an even element.
\end{proof}
\lean{Native.parity\_conserved\_under\_binding}
\end{theorem}

\begin{theorem}
The integer charge is likewise conserved by identity, so no process can violate either charge:
\begin{equation}
(\text{bind}(x, y, t)).\text{charge.hedgehog} = x.\text{charge.hedgehog} + y.\text{charge.hedgehog}.
\end{equation}
\begin{proof}
Immediate by the definition of the bind operation.
\end{proof}
\lean{Native.both\_labels\_conserved}
\end{theorem}

\begin{theorem}
Ordering by entropy and ordering by scale ratio are the same order.  This is not ``entropy increase implies expansion'' but an equivalence — the two arrows are the same relation rather than two that agree.

For threshold density $\rho > 0$ and scale ratios $R_1, R_2 > 0$, the entropy of a region spanning a scale ratio $R$ is $H(R) = \rho \ln R$, so:
\begin{equation}
H(R_1) \le H(R_2) \iff R_1 \le R_2.
\end{equation}
\begin{proof}
Dividing $H(R_1) = \rho \ln R_1 \le \rho \ln R_2 = H(R_2)$ by $\rho > 0$, we get $\ln R_1 \le \ln R_2$. Since $\ln$ is strictly monotone on positive reals, this is equivalent to $R_1 \le R_2$.
\end{proof}
\lean{Native.entropy\_order\_is\_scale\_order}
\end{theorem}

\begin{theorem}
Strictly: entropy strictly increases exactly when the scale ratio does.  So neither arrow can move while the other stands still:
\begin{equation}
H(R_1) < H(R_2) \iff R_1 < R_2.
\end{equation}
\begin{proof}
By strict monotonicity of $\ln$ on positive reals: $\rho \ln R_1 < \rho \ln R_2 \iff \ln R_1 < \ln R_2 \iff R_1 < R_2$.
\end{proof}
\lean{Native.entropy\_strict\_iff\_scale\_strict}
\end{theorem}

\begin{theorem}
The three-way identification, collected.  Entropy order, scale-ratio order, and — since Signature.lean makes time the drift direction and A6 makes the drift increase the scale — the time order.  One relation, three names:
\begin{equation}
(H(R_1) \le H(R_2) \iff R_1 \le R_2) \text{ and } (H(R_1) < H(R_2) \iff R_1 < R_2).
\end{equation}

The framework's own question this answers: not ``why are the arrows aligned'' but ``why did anyone think there were two''.
\begin{proof}
By the two preceding theorems.
\end{proof}
\lean{Native.arrows\_are\_one\_order}
\end{theorem}

\begin{theorem}
A ratio of scales is a unit: never zero, never infinite, unbounded in magnitude.  Matter here is anisotropy of the scale, so there is no maximum matter density, no Planck density, and no scale at which the description must be replaced by a different kind of thing.  This is Horizon.no\_minimum\_length seen from the matter side.

For any two units (invertible elements) $s_a, s_b$ in a commutative ring $A$:
\begin{equation}
\left(\frac{s_a}{s_b}\right) \neq 0.
\end{equation}
\begin{proof}
A unit in a ring is never zero, and the quotient of two nonzero elements is nonzero.
\end{proof}
\lean{Native.no\_maximum\_anisotropy}
\end{theorem}

\begin{theorem}
The two statements are the same fact: no cutoff at either end. For units $s_a, s_b, s_t$,
\begin{equation}
\left(\frac{s_a}{s_b}\right) \neq 0 \text{ and } g_{tt}(s_t) \neq 0.
\end{equation}
\begin{proof}
By the preceding theorem and Horizon.gtt\_ne\_zero, both statements hold.
\end{proof}
\lean{Native.no\_cutoff\_either\_end}
\end{theorem}

\subsection{Data}

This file confronts the predictions with observations.  One prediction is in tension with DESI: the framework's requirement that $w = -1$ exactly, enforced by A5.  The resonance width bound sits between two flagship determinations and makes a live falsifiable claim about which is correct.  Photon dispersion is tightly bounded and consistent with zero.  The charge count is undecided pending neutrinoless double-beta decay.  Parity is conserved but carries no discriminating power.  The particle spectrum shows no second switch in label types across fourteen orders of magnitude, consistent with the framework's structural requirement.  A scorecard collects the status.

\begin{theorem}
The DESI DR2 preference for an evolving equation of state, combining BAO with CMB and supernovae, exceeds three sigma.  On the strongest combination it reaches $4.2\sigma$, where the framework predicts $w = -1$ exactly and immovable (since A5 is load-bearing and not adjustable):
\begin{equation}
4.2 > 3.
\end{equation}
\begin{proof}
Direct numerical comparison.
\end{proof}
\lean{Data.desi\_tension\_exceeds\_three\_sigma}
\end{theorem}

\begin{theorem}
But the significance depends on which supernova compilation is used, by exactly a factor of $1.5$.  That is the reason the preference is not yet a falsification: the signal lives in the combination, and independent analyses argue the datasets are mutually inconsistent:
\begin{equation}
\frac{4.2}{2.8} = 1.5.
\end{equation}
\begin{proof}
Direct division.
\end{proof}
\lean{Data.desi\_significance\_dataset\_dependent}
\end{theorem}

\begin{theorem}
The two combinations disagree with each other about $w_0$.  BAO+CMB gives $w_0 \approx -0.42 \pm 0.21$; BAO+CMB+SN gives $w_0 \approx -0.838$.  That is two standard deviations apart on the same parameter, and DESI BAO alone is fully consistent with $\Lambda$CDM — so the signal lives entirely in the combination, which is the step independent analyses question:
\begin{equation}
\frac{0.838 - 0.42}{0.21} > 1.9.
\end{equation}
\begin{proof}
Direct numerical evaluation: $(0.838 - 0.42) / 0.21 \approx 1.99 > 1.9$.
\end{proof}
\lean{Data.desi\_combinations\_disagree\_internally}
\end{theorem}

\begin{theorem}
The threshold at which the framework would be dead.  Stated so that it cannot be quietly moved later.  If the significance of evolving $w$ reaches five sigma with consistent datasets, then A5 is falsified:
\begin{equation}
5 > 4.2.
\end{equation}
\begin{proof}
Direct numerical comparison.
\end{proof}
\lean{Data.falsification\_threshold}
\end{theorem}

\begin{definition}[Relative width abbreviation]
Abbreviation for the relative width to keep numerical statements concise:
\begin{equation}
\text{relW}(\Gamma, m) := \frac{\Gamma}{m}.
\end{equation}
\lean{Data.relW}
\end{definition}

\begin{definition}[Resolvability bound]
The bound: the mean threshold spacing in log-scale, $1/\rho$:
\begin{equation}
\text{widthBound} := 1.157.
\end{equation}
\lean{Data.widthBound}
\end{definition}

\begin{theorem}
CCL (Bern, using Roy equations and chiral-perturbation-theory input) places the $f_0(500)$ pole at $441 - i \cdot 272$, giving relative width $\Gamma/m = 544/441 = 1.234$ — **above the bound**:
\begin{equation}
\text{widthBound} < \text{relW}(544, 441).
\end{equation}
\begin{proof}
Numerical calculation: $544/441 \approx 1.234 > 1.157$.
\end{proof}
\lean{Data.ccl\_above\_bound}
\end{theorem}

\begin{theorem}
GKPY (Madrid, using once-subtracted dispersion relations and $\pi\pi$ data alone, without chiral perturbation theory) places the pole at $457 - i \cdot 249$, giving relative width $\Gamma/m = 498/457 = 1.090$ — **below the bound**:
\begin{equation}
\text{relW}(498, 457) < \text{widthBound}.
\end{equation}
\begin{proof}
Numerical calculation: $498/457 \approx 1.090 < 1.157$.
\end{proof}
\lean{Data.gkpy\_below\_bound}
\end{theorem}

\begin{theorem}
The bound falls between the two flagship determinations.  So it is not violated; it sits at the resolution limit of the best measurement of the most contested resonance, and the framework makes a live falsifiable claim about which side the true pole is on:
\begin{equation}
\text{relW}(498, 457) < \text{widthBound} < \text{relW}(544, 441).
\end{equation}
\begin{proof}
By the two preceding theorems.
\end{proof}
\lean{Data.determinations\_straddle\_bound}
\end{theorem}

\begin{theorem}
The updated Madrid values, $(445-450) - i(220-245)$, are entirely inside: the widest corner is $\Gamma/m = 490/445 = 1.101$:
\begin{equation}
\text{relW}(490, 445) < \text{widthBound}.
\end{equation}
\begin{proof}
Numerical calculation: $490/445 \approx 1.101 < 1.157$.
\end{proof}
\lean{Data.gkpy\_updated\_inside}
\end{theorem}

\begin{theorem}
$K_0^*(700)$, the second-broadest established state, has pole $648 - i \cdot 280$, giving relative width $\Gamma/m = 560/648 \approx 0.863$ — comfortably inside:
\begin{equation}
\text{relW}(560, 648) < \text{widthBound}.
\end{equation}
\begin{proof}
Numerical calculation confirms this is well below the bound.
\end{proof}
\lean{Data.kappa\_below\_bound}
\end{theorem}

\begin{theorem}
$\rho(770)$ is comfortably inside:
\begin{equation}
\text{relW}(149, 775) < \text{widthBound}.
\end{equation}
\begin{proof}
Numerical calculation: $149/775 \approx 0.192 < 1.157$.
\end{proof}
\lean{Data.rho\_below\_bound}
\end{theorem}

\begin{theorem}
$\Delta(1232)$ is comfortably inside:
\begin{equation}
\text{relW}(117, 1232) < \text{widthBound}.
\end{equation}
\begin{proof}
Numerical calculation: $117/1232 \approx 0.095 < 1.157$.
\end{proof}
\lean{Data.delta\_below\_bound}
\end{theorem}

\begin{theorem}
The electroweak states (W, Z, top, Higgs) are orders of magnitude inside:
\begin{equation}
\text{relW}(2.50, 91.19) < 0.03.
\end{equation}
\begin{proof}
Numerical calculation: $2.50/91.19 \approx 0.0274 < 0.03$.
\end{proof}
\lean{Data.z\_boson\_far\_inside}
\end{theorem}

\begin{theorem}
The bound still lands **between** the two broadest states, which is what makes it a test rather than a description:
\begin{equation}
\text{relW}(560, 648) < \text{widthBound} < \text{relW}(544, 441).
\end{equation}
\begin{proof}
By the theorems on $K_0^*$ and CCL above.
\end{proof}
\lean{Data.bound\_separates\_kappa\_from\_sigma}
\end{theorem}

\begin{theorem}
The Breit–Wigner ranges remain useless as a test: the PDG quotes $m \in [400, 800]$ and $\Gamma \in [100, 800]$, spanning $\Gamma/m$ from $0.125$ to $2$ — directly bracketing the bound itself.  So the prediction must be read on the pole:
\begin{equation}
\text{relW}(100, 800) < \text{widthBound} < \text{relW}(800, 400).
\end{equation}
\begin{proof}
Numerical calculation: $100/800 = 0.125 < 1.157 < 2 = 800/400$.
\end{proof}
\lean{Data.breit\_wigner\_spans\_the\_bound}
\end{theorem}

\begin{definition}[Fermi-LAT photon dispersion limit]
The Fermi-LAT limit on linear dispersion from GRB 090510, in units of the Planck energy:
\begin{equation}
E_{QG,1} \gtrsim 7.6 \, E_{\text{Pl}}.
\end{equation}

Expressed as a number:
\begin{equation}
\text{fermiLimitInPlanck} := 7.6.
\end{equation}
\lean{Data.fermiLimitInPlanck}
\end{definition}

\begin{theorem}
The limit is already above the Planck energy, where minimum-length programmes generically place the effect.  The framework predicts exactly zero dispersion, so every tightening moves toward it:
\begin{equation}
\text{fermiLimitInPlanck} > 1.
\end{equation}
\begin{proof}
Direct numerical comparison: $7.6 > 1$.
\end{proof}
\lean{Data.dispersion\_limit\_above\_planck}
\end{theorem}

\begin{theorem}
Stated as the divergence it is: a finite $E_{QG}$ is a prediction of the rival programmes and is excluded further with each measurement, while this framework predicts the limit can be pushed without end.  For any energy $E$, there exists a higher energy $E'$:
\begin{equation}
\exists E' : E < E'.
\end{equation}
\begin{proof}
Set $E' = E + 1$, which is strictly greater than $E$.
\end{proof}
\lean{Data.framework\_predicts\_no\_finite\_scale}
\end{theorem}

\begin{definition}[Framework's count of unconfined exactly conserved charges]
The framework predicts exactly one unconfined, exactly conserved integer charge:
\begin{equation}
\text{predictedExactCharges} := 1.
\end{equation}
\lean{Data.predictedExactCharges}
\end{definition}

\begin{definition}[Nature's candidate charges]
In the Standard Model, the candidates for an unconfined, exactly conserved integer charge are electric charge and $B - L$:
\begin{equation}
\text{candidateExactCharges} := 2.
\end{equation}
\lean{Data.candidateExactCharges}
\end{definition}

\begin{theorem}
The counts disagree unless $B - L$ is violated.  Neutrinoless double-beta decay violates $B - L$ by two units, so its observation removes the second candidate and confirms the count; its continued absence, with $B - L$ looking exact, refutes it:
\begin{equation}
\text{predictedExactCharges} \neq \text{candidateExactCharges}.
\end{equation}
\begin{proof}
Direct comparison: $1 \neq 2$.
\end{proof}
\lean{Data.one\_charge\_prediction}
\end{theorem}

\begin{theorem}
What $0\nu\beta\beta$ would settle, stated arithmetically so the logic is explicit.  Observation of neutrinoless double-beta decay would reduce the candidate count by one:
\begin{equation}
\text{candidateExactCharges} - 1 = \text{predictedExactCharges}.
\end{equation}
\begin{proof}
Direct calculation: $2 - 1 = 1$.
\end{proof}
\lean{Data.zero\_nu\_beta\_beta\_decides}
\end{theorem}

\begin{definition}[Label-type count]
Exactly-conserved label types observed at any probed scale: electric charge, colour, particle–antiparticle doubling, and $B - L$ if exact:
\begin{equation}
\text{labelTypeCount} := 4.
\end{equation}
\lean{Data.labelTypeCount}
\end{definition}

\begin{definition}[Bottom of probed range]
The neutrino mass scale, the lowest energy probed:
\begin{equation}
\text{probedLow} := 5 \times 10^{-11} \text{ GeV}.
\end{equation}
\lean{Data.probedLow}
\end{definition}

\begin{definition}[Top of probed range]
The LHC's collision energy:
\begin{equation}
\text{probedHigh} := 1.3 \times 10^{4} \text{ GeV}.
\end{equation}
\lean{Data.probedHigh}
\end{definition}

\begin{theorem}
**Fourteen orders of magnitude**, over which the label-type count does not change:
\begin{equation}
\frac{\text{probedHigh}}{\text{probedLow}} > 10^{14}.
\end{equation}
\begin{proof}
Direct calculation: $\frac{1.3 \times 10^4}{5 \times 10^{-11}} = \frac{1.3}{5} \times 10^{15} = 0.26 \times 10^{15} = 2.6 \times 10^{14} > 10^{14}$.
\end{proof}
\lean{Data.probed\_range\_orders}
\end{theorem}

\begin{theorem}
The count is the same at both ends of the probed range.  Stated as the equality it is: the inventory at the neutrino scale and at LHC energies is the same set of label types.  A second switch would break it:
\begin{equation}
\text{labelTypeCount} = \text{labelTypeCount}.
\end{equation}
\begin{proof}
Immediate by reflexivity.
\end{proof}
\lean{Data.no\_second\_switch\_observed}
\end{theorem}

\begin{definition}[Electroweak multiplet size]
Electroweak breaking takes multiplets of size 2 (doublets) to singlets of size 1:
\begin{equation}
\text{ewMultipletBefore} := 2, \quad \text{ewMultipletAfter} := 1.
\end{equation}
\lean{Data.ewMultipletBefore}
\lean{Data.ewMultipletAfter}
\end{definition}

\begin{definition}[Colour multiplet size]
Confinement leaves colour triplets as triplets:
\begin{equation}
\text{colourMultipletBefore} := 3, \quad \text{colourMultipletAfter} := 3.
\end{equation}
\lean{Data.colourMultipletBefore}
\lean{Data.colourMultipletAfter}
\end{definition}

\begin{theorem}
**Both breakings are refinements: sizes never increase.**  This is the content of the test.  A size going $1 \to 2$ — an exact degeneracy appearing where there was none — would be a second switch and would falsify Locus.everything\_switches\_on\_together.  Neither breaking does that:
\begin{equation}
\text{ewMultipletAfter} \le \text{ewMultipletBefore} \text{ and } \text{colourMultipletAfter} \le \text{colourMultipletBefore}.
\end{equation}
\begin{proof}
Direct comparison: $1 \le 2$ and $3 \le 3$.
\end{proof}
\lean{Data.breakings\_are\_refinements}
\end{theorem}

\begin{theorem}
The electroweak case is a **strict** refinement, so the test is not passed by both breakings being trivial:
\begin{equation}
\text{ewMultipletAfter} < \text{ewMultipletBefore}.
\end{equation}
\begin{proof}
Direct comparison: $1 < 2$.
\end{proof}
\lean{Data.electroweak\_strictly\_refines}
\end{theorem}

\begin{theorem}
**Scorecard.** One live tension and one live open question, neither fatal.  The DESI tension reaches $4.2\sigma$ but is contested and the signal depends on dataset combination.  The resonance width bound is open, with the two methods straddling it:
\begin{equation}
4.2 < 5 \text{ and } \text{relW}(498, 457) < \text{widthBound}.
\end{equation}
\begin{proof}
By desi\_tension\_exceeds\_three\_sigma (showing the tension exists) and gkpy\_below\_bound (placing GKPY inside the bound).
\end{proof}
\lean{Data.scorecard}
\end{theorem}


\section{Part VI: The Register — Retracted, Corrected, and Audited}

\subsection{Color}

A degenerate block in a scale pattern — a set of directions carrying the same local unit — has no preferred labelling within it.  When a defect sits in such a block and a circuit is made around it, the scale winds and the block labels may be permuted.  Asking when such a configuration is globally consistent (when the labelling closes up after the loop) yields three striking results: defects with genuine permutations cannot exist alone (confinement), bound states must contain exactly as many constituents as the block size, and charge comes in units of $1/n$ for a block of size $n$.  These theorems are mathematically correct but the identification with colour has been withdrawn — the geometry of the scale gives a block of size two, not three, as registered in \texttt{ColorAudit.lean}.

\begin{definition}
The \textbf{monodromy} of a defect in a degenerate block is a pair: a permutation of the block labels and an integer winding per direction.  Going once around a defect, the scale winds (forced by discrete scale invariance) and the block labels may be permuted (forced by the absence of preferred labelling).
\begin{align*}
\text{BlockMonodromy}(n) = \{\text{perm} : \text{Equiv.Perm}(\text{Fin } n), \text{ wind} : \text{Fin } n \to \mathbb{Z}\}
\end{align*}
\lean{Color.BlockMonodromy}
\end{definition}

\begin{definition}
The composition of two monodromies combines their permutations multiplicatively and their windings additively:
\begin{align*}
M \circ N &: \text{perm} = M.\text{perm} \cdot N.\text{perm}, \quad \text{wind} = \lambda a. M.\text{wind}(a) + N.\text{wind}(a)
\end{align*}
\lean{Color.BlockMonodromy.comp}
\end{definition}

\begin{definition}
The opposite defect, obtained by reversing the circuit direction, has inverted permutation and negated winding:
\begin{align*}
M^{-1} &: \text{perm} = M.\text{perm}^{-1}, \quad \text{wind} = -M.\text{wind}
\end{align*}
\lean{Color.BlockMonodromy.anti}
\end{definition}

\begin{definition}
The vacuum state is the trivial monodromy — identity permutation and zero winding everywhere.
\begin{align*}
\text{triv}(n) &: \text{perm} = 1, \quad \text{wind}(a) = 0
\end{align*}
\lean{Color.BlockMonodromy.triv}
\end{definition}

\begin{definition}
The total scale winding — the charge — is the sum of windings over all directions:
\begin{align*}
\text{charge}(M) = \sum_{a=0}^{n-1} M.\text{wind}(a)
\end{align*}
\lean{Color.BlockMonodromy.charge}
\end{definition}

\begin{proposition}
Composition acts on the permutation part by multiplication:
\begin{align*}
(M \circ N).\text{perm} = M.\text{perm} \cdot N.\text{perm}
\end{align*}
\end{proposition}
\begin{proof}
Immediate from the definition of $\circ$.
\end{proof}
\lean{Color.BlockMonodromy.comp\_perm}

\begin{proposition}
The opposite monodromy inverts the permutation part:
\begin{align*}
M^{-1}.\text{perm} = M.\text{perm}^{-1}
\end{align*}
\end{proposition}
\begin{proof}
Immediate from the definition of the opposite monodromy.
\end{proof}
\lean{Color.BlockMonodromy.anti\_perm}

\begin{proposition}
The vacuum monodromy has trivial permutation: $\text{triv}(n).\text{perm} = 1$.
\end{proposition}
\begin{proof}
Immediate from the definition of $\text{triv}$.
\end{proof}
\lean{Color.BlockMonodromy.triv\_perm}

\begin{proposition}
Composing two monodromies gives the sum of charges.
\begin{align*}
\text{charge}(M \circ N) = \text{charge}(M) + \text{charge}(N)
\end{align*}
\begin{proof}
Both sides are sums over the finite index set of directions: $\text{charge}(M\circ N)=\sum_a (M\circ N).\text{wind}(a) = \sum_a (M.\text{wind}(a)+N.\text{wind}(a))$ by definition of composition, and this splits by additivity of the finite sum into $\sum_a M.\text{wind}(a) + \sum_a N.\text{wind}(a) = \text{charge}(M)+\text{charge}(N)$.
\end{proof}
\lean{Color.BlockMonodromy.charge\_comp}
\end{proposition}

\begin{proposition}
The opposite defect has opposite charge.
\begin{align*}
\text{charge}(M^{-1}) = -\text{charge}(M)
\end{align*}
\begin{proof}
$\text{charge}(M^{-1}) = \sum_a M^{-1}.\text{wind}(a) = \sum_a (-M.\text{wind}(a))$ by definition of the opposite monodromy, and a finite sum of negatives is the negative of the sum, giving $-\sum_a M.\text{wind}(a) = -\text{charge}(M)$.
\end{proof}
\lean{Color.BlockMonodromy.charge\_anti}
\end{proposition}

\begin{definition}
A configuration is \textbf{observable} when its block labelling closes up after the circuit — the permutation is the identity.
\begin{align*}
\text{Observable}(M) :\Leftrightarrow M.\text{perm} = 1
\end{align*}
\lean{Color.BlockMonodromy.Observable}
\end{definition}

\begin{proposition}
The vacuum monodromy is observable: $\text{Observable}(\text{triv}(n))$.
\end{proposition}
\begin{proof}
The vacuum has trivial permutation, which is the identity by definition of
observability.
\end{proof}
\lean{Color.BlockMonodromy.triv\_observable}

\begin{theorem}[Confinement]
A defect whose monodromy permutes the block genuinely cannot exist on its own.  This is not a force growing with separation; the isolated configuration is simply not single-valued.
\begin{align*}
\text{If } M.\text{perm} \neq 1, \text{ then } \neg \text{Observable}(M)
\end{align*}
\begin{proof}
Immediate from the definitions: $\text{Observable}(M)$ unfolds to $M.\text{perm}=1$, so the hypothesis $M.\text{perm}\neq1$ is literally the negation of the conclusion.
\end{proof}
\lean{Color.BlockMonodromy.confinement}
\end{theorem}

\begin{theorem}
A defect paired with its opposite always closes up observably, whatever the block permutation.  No condition on $n$, no fine-tuning.
\begin{align*}
\text{Observable}(M \circ M^{-1})
\end{align*}
\begin{proof}
Unfolding $\text{Observable}$ and composition, the goal is $(M\circ M^{-1}).\text{perm}=1$, i.e.\ $M.\text{perm}\cdot M.\text{perm}^{-1}=1$, which holds for any group element by the cancellation law $g\cdot g^{-1}=1$.
\end{proof}
\lean{Color.BlockMonodromy.pair\_observable}
\end{theorem}

\begin{theorem}
The pair carries zero charge.
\begin{align*}
\text{charge}(M \circ M^{-1}) = 0
\end{align*}
\begin{proof}
By the charge-of-composition law, $\text{charge}(M\circ M^{-1}) = \text{charge}(M) + \text{charge}(M^{-1})$, and by the charge-of-opposite law $\text{charge}(M^{-1})=-\text{charge}(M)$, so the sum is $\text{charge}(M) + (-\text{charge}(M)) = 0$.
\end{proof}
\lean{Color.BlockMonodromy.pair\_charge}
\end{theorem}

\begin{definition}
Bringing together $k$ copies of the same defect is defined recursively:
\begin{align*}
\text{rep}(M, 0) &= \text{triv}(n) \\
\text{rep}(M, k+1) &= M \circ \text{rep}(M, k)
\end{align*}
\lean{Color.BlockMonodromy.rep}
\end{definition}

\begin{proposition}
The permutation of $k$ copies is the $k$-fold power of the single permutation.
\begin{align*}
\text{rep}(M, k).\text{perm} = M.\text{perm}^k
\end{align*}
\begin{proof}
By induction on $k$. At $k=0$, $\text{rep}(M,0)=\text{triv}(n)$ has trivial permutation, and $M.\text{perm}^0=1$, so both sides agree. For the step, assume $\text{rep}(M,m).\text{perm}=M.\text{perm}^m$; then $\text{rep}(M,m+1).\text{perm} = (M\circ\text{rep}(M,m)).\text{perm} = M.\text{perm}\cdot\text{rep}(M,m).\text{perm} = M.\text{perm}\cdot M.\text{perm}^m = M.\text{perm}^{m+1}$, using the composition-permutation law and the inductive hypothesis.
\end{proof}
\lean{Color.BlockMonodromy.rep\_perm}
\end{proposition}

\begin{proposition}
The charge of $k$ copies is $k$ times the single charge.
\begin{align*}
\text{charge}(\text{rep}(M, k)) = (k : \mathbb{Z}) \cdot \text{charge}(M)
\end{align*}
\begin{proof}
By induction on $k$. At $k=0$, $\text{rep}(M,0)=\text{triv}(n)$ has charge $0$ (a sum of zero windings), matching $(0:\mathbb{Z})\cdot\text{charge}(M)=0$. For the step, assume $\text{charge}(\text{rep}(M,m)) = m\cdot\text{charge}(M)$; then by the charge-of-composition law, $\text{charge}(\text{rep}(M,m+1)) = \text{charge}(M) + \text{charge}(\text{rep}(M,m)) = \text{charge}(M) + m\cdot\text{charge}(M) = (m+1)\cdot\text{charge}(M)$, using elementary ring arithmetic on the cast integers.
\end{proof}
\lean{Color.BlockMonodromy.rep\_charge}
\end{proposition}

\begin{theorem}
$k$ copies close up observably exactly when the permutation raised to the $k$-th power equals the identity.
\begin{align*}
\text{If } M.\text{perm}^k = 1, \text{ then } \text{Observable}(\text{rep}(M, k))
\end{align*}
\begin{proof}
By definition $\text{Observable}(\text{rep}(M,k))$ unfolds to
$\text{rep}(M,k).\text{perm}=1$, and $\text{rep}(M,k).\text{perm}=M.\text{perm}^k$
by the representation-power law above, so the goal reduces exactly to the
hypothesis $M.\text{perm}^k=1$.
\end{proof}
\lean{Color.BlockMonodromy.bound\_state\_observable}
\end{theorem}

\begin{theorem}[Bound state size]
The size of a bound state is the multiplicative order of its monodromy permutation.  $k$ copies close up exactly when $\pi^k = 1$.  For a defect that cycles a block of size $n$ once, the smallest consistent number of constituents is $n$: \textbf{baryon number is the block size}, not an extra quantum number.
\begin{align*}
\text{If } \text{ord}(M.\text{perm}) = n, \text{ then } \text{Observable}(\text{rep}(M, n))
\end{align*}
\begin{proof}
Any group element raised to its own order is the identity:
$M.\text{perm}^{\text{ord}(M.\text{perm})}=1$. Substituting $\text{ord}(M.\text{perm})=n$
gives $M.\text{perm}^n=1$, and the previous theorem (bound state observable)
applied with $k=n$ gives $\text{Observable}(\text{rep}(M,n))$.
\end{proof}
\lean{Color.BlockMonodromy.bound\_state\_size}
\end{theorem}

\begin{theorem}[No partial hadrons]
Fewer constituents than the block size never closes up observably.  With $n = 3$: a single constituent is unobservable, two are unobservable, three are observable.  The absence of two-quark hadrons is not an extra rule but a mathematical consequence.
\begin{align*}
\text{If } \text{ord}(M.\text{perm}) = n, \, 0 < k < n, \text{ then } \neg \text{Observable}(\text{rep}(M, k))
\end{align*}
\begin{proof}
Suppose for contradiction that $\text{Observable}(\text{rep}(M,k))$ holds, i.e.
$M.\text{perm}^k=1$ (using $\text{rep}(M,k).\text{perm}=M.\text{perm}^k$ and
$0<k$). Then the order of $M.\text{perm}$ must divide, and in particular be
$\leq k$: $\text{ord}(M.\text{perm}) \leq k$. Substituting the hypothesis
$\text{ord}(M.\text{perm})=n$ gives $n \leq k$, contradicting $k<n$.
\end{proof}
\lean{Color.BlockMonodromy.no\_partial\_state}
\end{theorem}

\begin{theorem}[Charge in $n$ths]
The bound state's charge is $n$ times a constituent's.  Therefore if observable charge is normalised to integers, constituent charges lie in $(1/n)\mathbb{Z}$.
\begin{align*}
\text{charge}(\text{rep}(M, n)) = (n : \mathbb{Z}) \cdot \text{charge}(M)
\end{align*}
\begin{proof}
This is exactly the representation-charge law ($\text{rep}(M,k).\text{charge} = k\cdot M.\text{charge}$) instantiated at $k=n$.
\end{proof}
\lean{Color.BlockMonodromy.constituent\_charge\_fraction}
\end{theorem}

\begin{theorem}
Stated as a fraction: the constituent carries $1/n$ of the observable charge.
\begin{align*}
(M.\text{charge} : \mathbb{Q}) = \frac{1}{n} \cdot (\text{rep}(M, n).\text{charge} : \mathbb{Q})
\end{align*}
when $0 < n$.
\begin{proof}
By the previous theorem, $\text{rep}(M,n).\text{charge} = n\cdot M.\text{charge}$
(cast into $\mathbb{Q}$). Since $n>0$, $(n:\mathbb{Q})\neq0$, so dividing both
sides by $n$ (equivalently, cancelling $n^{-1}\cdot n=1$) gives
$M.\text{charge} = \tfrac1n\cdot\text{rep}(M,n).\text{charge}$.
\end{proof}
\lean{Color.BlockMonodromy.constituent\_charge\_rat}
\end{theorem}

\begin{theorem}
An observable bound state of $n$ identical constituents has charge divisible by $n$.
\begin{align*}
n \mid \text{charge}(\text{rep}(M, n))
\end{align*}
\begin{proof}
Immediate from the charge-in-$n$ths theorem: $\text{rep}(M,n).\text{charge} = n\cdot M.\text{charge}$ exhibits $M.\text{charge}$ itself as the witnessing quotient.
\end{proof}
\lean{Color.BlockMonodromy.bound\_charge\_divisible}
\end{theorem}

\subsection{ColorAudit}

The identification of the degenerate block from the scale pattern with the colour block fails by a factor of $3/2$ in baryon number and is excluded by the charge quantum — halves are not thirds.  This file records the refutation.  The mathematics of \texttt{Color.lean} is untouched; it always said $n$ was an input.  What is excluded is the route that would derive $n$ from the spatial scale pattern.  That route gives $n = 2$ and is refuted.

\begin{definition}
The size of the degenerate block in a uniaxial pattern in three spatial dimensions is two (one distinguished direction, the other two are not).
\begin{align*}
\text{uniaxialBlockSize} = 2
\end{align*}
\lean{ColorAudit.uniaxialBlockSize}
\end{definition}

\begin{definition}
The number of constituents \texttt{Color.lean} predicts for a degenerate block of size $n$.
\begin{align*}
\text{boundStateSize}(n) = n
\end{align*}
\lean{ColorAudit.boundStateSize}
\end{definition}

\begin{definition}
The charge quantum \texttt{Color.lean} predicts for a block of size $n$.
\begin{align*}
\text{chargeQuantum}(n) = \frac{1}{n}
\end{align*}
\lean{ColorAudit.chargeQuantum}
\end{definition}

\begin{definition}
The observed baryon number in nature.
\begin{align*}
\text{observedBaryonNumber} = 3
\end{align*}
\lean{ColorAudit.observedBaryonNumber}
\end{definition}

\begin{definition}
The observed quark charge quantum.
\begin{align*}
\text{observedChargeQuantum} = \frac{1}{3}
\end{align*}
\lean{ColorAudit.observedChargeQuantum}
\end{definition}

\begin{theorem}
Under the identification, bound states would have two constituents.
\begin{align*}
\text{boundStateSize}(\text{uniaxialBlockSize}) = 2
\end{align*}
\begin{proof}
Immediate from the definitions: $\text{boundStateSize}(n)=n$ and $\text{uniaxialBlockSize}=2$, so $\text{boundStateSize}(\text{uniaxialBlockSize})=2$ by unfolding both.
\end{proof}
\lean{ColorAudit.identification\_predicts\_two}
\end{theorem}

\begin{theorem}
And charges would come in halves.
\begin{align*}
\text{chargeQuantum}(\text{uniaxialBlockSize}) = \frac{1}{2}
\end{align*}
\begin{proof}
Unfold $\text{chargeQuantum}(n)=1/n$ and $\text{uniaxialBlockSize}=2$ to get $1/2=1/2$, a numerical identity.
\end{proof}
\lean{ColorAudit.identification\_predicts\_halves}
\end{theorem}

\begin{theorem}
The identification is excluded by baryon number.
\begin{align*}
\text{boundStateSize}(\text{uniaxialBlockSize}) \neq \text{observedBaryonNumber}
\end{align*}
\begin{proof}
Unfolding the definitions reduces the claim to $2 \neq 3$, which holds by arithmetic on natural numbers.
\end{proof}
\lean{ColorAudit.identification\_fails\_baryon}
\end{theorem}

\begin{theorem}
And by the charge quantum.  Halves are not thirds, and not nearly.
\begin{align*}
\text{chargeQuantum}(\text{uniaxialBlockSize}) \neq \text{observedChargeQuantum}
\end{align*}
\begin{proof}
Unfolding the definitions reduces the claim to $1/2 \neq 1/3$, a numerical fact about rationals.
\end{proof}
\lean{ColorAudit.identification\_fails\_charge}
\end{theorem}

\begin{theorem}
The spatial degeneracy block cannot be the colour block.
\begin{align*}
\text{boundStateSize}(\text{uniaxialBlockSize}) \neq \text{observedBaryonNumber} \,\wedge\, \text{chargeQuantum}(\text{uniaxialBlockSize}) \neq \text{observedChargeQuantum}
\end{align*}
\begin{proof}
The conjunction of the two preceding refutations: pair the baryon-number failure with the charge-quantum failure.
\end{proof}
\lean{ColorAudit.spatial\_block\_is\_not\_colour\_block}
\end{theorem}

\begin{theorem}
Only $n = 3$ matches the observed baryon number.
\begin{align*}
\text{If } \text{boundStateSize}(n) = \text{observedBaryonNumber}, \text{ then } n = 3
\end{align*}
\begin{proof}
Since $\text{boundStateSize}(n)=n$ and $\text{observedBaryonNumber}=3$ by definition, the hypothesis $\text{boundStateSize}(n)=\text{observedBaryonNumber}$ literally is the equation $n=3$.
\end{proof}
\lean{ColorAudit.three\_is\_forced\_by\_data}
\end{theorem}

\begin{theorem}
And $n = 3$ does reproduce the observed charge quantum, so the two observations are consistent with a single block size — just not the spatial one.
\begin{align*}
\text{chargeQuantum}(3) = \text{observedChargeQuantum}
\end{align*}
\begin{proof}
Unfold both sides: $\text{chargeQuantum}(3)=1/3$ and $\text{observedChargeQuantum}=1/3$ by definition, so they are equal.
\end{proof}
\lean{ColorAudit.three\_matches\_charge}
\end{theorem}

\begin{theorem}
The gap, stated as a number: the geometry offers two where the data requires three.
\begin{align*}
\text{uniaxialBlockSize} = 2 \,\wedge\, \text{observedBaryonNumber} = 3 \,\wedge\, \text{uniaxialBlockSize} \neq \text{observedBaryonNumber}
\end{align*}
\begin{proof}
The first two conjuncts unfold directly from the definitions $\text{uniaxialBlockSize}=2$ and $\text{observedBaryonNumber}=3$; the third is the numerical fact $2\neq3$.
\end{proof}
\lean{ColorAudit.geometry\_offers\_two\_data\_needs\_three}
\end{theorem}

\subsection{Codimension}

A loop in three-dimensional space does not enclose a point; it encloses a line.  What a loop detects through the scale (which is circle-valued in this file) is a locus of codimension two.  So a circle-valued scale in three dimensions produces strings, not point particles.  The scale is one number so no second scale offers a way to detect anything else.  Every theorem below is true of the circle-valued case and stands as a record of the error corrected in \texttt{Axis.lean}, where the directional scale of axiom A4 is shown to restore point defects.

\begin{definition}
The dimension of the locus a loop can wind around in $d$ spatial dimensions.  A loop detects codimension two; nothing about a circle-valued scale offers a way to detect anything else.
\begin{align*}
\text{defectDim}(d) = d - 2
\end{align*}
\lean{Codimension.defectDim}
\end{definition}

\begin{proposition}
In two spatial dimensions, winding defects are zero-dimensional (pointlike).
\begin{align*}
\text{defectDim}(2) = 0
\end{align*}
\begin{proof}
Immediate from the definition $\text{defectDim}(d)=d-2$: substituting $d=2$ gives $2-2=0$.
\end{proof}
\lean{Codimension.defectDim\_two}
\end{proposition}

\begin{proposition}
In three spatial dimensions, winding defects are one-dimensional (strings).
\begin{align*}
\text{defectDim}(3) = 1
\end{align*}
\begin{proof}
Immediate from the definition: $\text{defectDim}(3) = 3-2 = 1$.
\end{proof}
\lean{Codimension.defectDim\_three}
\end{proposition}

\begin{proposition}
In four spatial dimensions, defects are two-dimensional.
\begin{align*}
\text{defectDim}(4) = 2
\end{align*}
\begin{proof}
Immediate from the definition: $\text{defectDim}(4) = 4-2 = 2$.
\end{proof}
\lean{Codimension.defectDim\_four}
\end{proposition}

\begin{theorem}
Point-like winding defects require exactly two spatial dimensions.  The picture of a point with a loop around it is available only in $d = 2$.
\begin{align*}
\text{If } 2 \le d, \text{ then } (\text{defectDim}(d) = 0 \Leftrightarrow d = 2)
\end{align*}
\begin{proof}
Both directions use $d\geq2$. ($\Rightarrow$) Unfolding $\text{defectDim}(d)=d-2$ (truncated natural subtraction), $d-2=0$ together with $2\leq d$ forces $d=2$ by elementary arithmetic on $\mathbb{N}$. ($\Leftarrow$) Substituting $d=2$ gives $\text{defectDim}(2)=2-2=0$ by the base computation above.
\end{proof}
\lean{Codimension.pointlike\_iff\_two\_dim}
\end{theorem}

\begin{theorem}
In three spatial dimensions a winding defect is a string — a genuinely one-dimensional object, not a point with internal structure.
\begin{align*}
\text{defectDim}(3) = 1
\end{align*}
\begin{proof}
Immediate from the definition, as above: $\text{defectDim}(3)=3-2=1$.
\end{proof}
\lean{Codimension.three\_dim\_gives\_strings}
\end{theorem}

\begin{definition}
Spacetime dimension: spatial dimension plus the single drift direction that the scale supplies.
\begin{align*}
\text{spacetimeDim}(d) = d + 1
\end{align*}
\lean{Codimension.spacetimeDim}
\end{definition}

\begin{theorem}
If the fundamental excitations are point-like, the world is $2+1$.
\begin{align*}
\text{If } 2 \le d \text{ and } \text{defectDim}(d) = 0, \text{ then } \text{spacetimeDim}(d) = 3
\end{align*}
\begin{proof}
By the point-like-iff-two-dimensional theorem, $\text{defectDim}(d)=0$ together with $2\leq d$ gives $d=2$. Substituting into the definition $\text{spacetimeDim}(d)=d+1$ gives $\text{spacetimeDim}(2)=3$.
\end{proof}
\lean{Codimension.pointlike\_forces\_three}
\end{theorem}

\begin{theorem}
In a $3+1$ world the excitations are one-dimensional.  The framework does not get to have both point-like defects and our spacetime.
\begin{align*}
\text{spacetimeDim}(3) = 4 \,\wedge\, \text{defectDim}(3) = 1
\end{align*}
\begin{proof}
Both conjuncts are direct computations from the definitions: $\text{spacetimeDim}(3)=3+1=4$ and $\text{defectDim}(3)=3-2=1$.
\end{proof}
\lean{Codimension.our\_world\_gives\_strings}
\end{theorem}

\begin{theorem}
The two possibilities are mutually exclusive: no spatial dimension gives both point-like defects and a $3+1$ spacetime.
\begin{align*}
\text{If } 2 \le d \text{ and } \text{spacetimeDim}(d) = 4, \text{ then } \text{defectDim}(d) \neq 0
\end{align*}
\begin{proof}
Suppose for contradiction that $\text{defectDim}(d)=0$. By the point-like-iff-two-dimensional theorem (using $2\leq d$) this gives $d=2$. Substituting into $\text{spacetimeDim}(d)=d+1$ gives $\text{spacetimeDim}(d)=3$, contradicting the hypothesis $\text{spacetimeDim}(d)=4$ (since $3\neq4$).
\end{proof}
\lean{Codimension.no\_pointlike\_in\_four}
\end{theorem}

\begin{definition}
The mass of a one-dimensional defect is tension times length.
\begin{align*}
\text{stringMass}(\text{tension}, \text{length}) = \text{tension} \times \text{length}
\end{align*}
\lean{Codimension.stringMass}
\end{definition}

\begin{theorem}
Winding does not determine mass.  Two defects of the same winding but different length have different mass.  So no function of the winding alone can be the mass law — independently confirming \texttt{MassAudit.mass\_law\_underdetermined} and explaining its cause: the defect is extended and topology never measured the extension.
\begin{align*}
\text{If } 0 < \text{tension} \text{ and } L_1 \neq L_2, \text{ then } \text{stringMass}(\text{tension}, L_1) \neq \text{stringMass}(\text{tension}, L_2)
\end{align*}
\begin{proof}
Suppose for contradiction $\text{stringMass}(\text{tension},L_1)=\text{stringMass}(\text{tension},L_2)$, i.e.\ $\text{tension}\cdot L_1 = \text{tension}\cdot L_2$ by definition of $\text{stringMass}$. Since $\text{tension}>0$, it is nonzero, so left-cancellation gives $L_1=L_2$, contradicting the hypothesis $L_1\neq L_2$.
\end{proof}
\lean{Codimension.string\_mass\_not\_topological}
\end{theorem}

\begin{theorem}
Mass grows strictly with length at fixed tension, so the spectrum of a string defect is continuous in length — not a tower indexed by an integer.
\begin{align*}
\text{If } 0 < \text{tension} \text{ and } L_1 < L_2, \text{ then } \text{stringMass}(\text{tension}, L_1) < \text{stringMass}(\text{tension}, L_2)
\end{align*}
\begin{proof}
Unfolding $\text{stringMass}(\text{tension},L)=\text{tension}\cdot L$, the claim is $\text{tension}\cdot L_1 < \text{tension}\cdot L_2$, which follows from $L_1<L_2$ by multiplying a strict inequality on the left by the positive constant $\text{tension}$.
\end{proof}
\lean{Codimension.stringMass\_strictMono}
\end{theorem}

\begin{proposition}
A closed string of zero length is massless.
\begin{align*}
\text{stringMass}(\text{tension}, 0) = 0
\end{align*}
\begin{proof}
Immediate from the definition: $\text{stringMass}(\text{tension},0) = \text{tension}\times0 = 0$.
\end{proof}
\lean{Codimension.stringMass\_zero}
\end{proposition}

\begin{theorem}
For any spatial dimension at least two, either the defects are extended or the world is $2+1$.
\begin{align*}
\text{If } 2 \le d, \text{ then } (\text{defectDim}(d) \neq 0 \,\vee\, \text{spacetimeDim}(d) = 3)
\end{align*}
\begin{proof}
Case split on whether $d=2$ or $2<d$ (the trichotomy available from $2\leq d$). If $d=2$, take the right disjunct: $\text{spacetimeDim}(2)=2+1=3$. If $2<d$, take the left disjunct: unfolding $\text{defectDim}(d)=d-2$ (truncated subtraction), $2<d$ gives $d-2\neq0$ by elementary arithmetic on $\mathbb{N}$.
\end{proof}
\lean{Codimension.extended\_or\_low\_dimensional}
\end{theorem}

\subsection{Expressive}

The framework's vocabulary is three operations: comparing scales, failing to commute, and counting.  This file proves these really are the whole vocabulary and that they are mutually independent.  The scale channel expresses differences and nothing else, rotation says nothing about scale and scale says nothing about rotation, so the two channels are disjoint.  A cross-sector test would need them to constrain each other, and the axioms guarantee they do not.  These theorems hold of the scalar scale sector; the directional scale of A4 couples the channels, as shown in \texttt{Coupling.lean}.

\begin{definition}
A quantity is \textbf{observable} when a global change of the fiducial (an additive shift of the scale field everywhere) does not move it.  This is axiom A5 stated for an arbitrary function of the scale field.
\begin{align*}
\text{FiducialInvariant}(f) :\Leftrightarrow \forall (\sigma : X \to \mathbb{R})(c : \mathbb{R}), \, f(\lambda x. \sigma x + c) = f(\sigma)
\end{align*}
\lean{Expressive.FiducialInvariant}
\end{definition}

\begin{definition}
The difference pattern of a scale field relative to a base point $x_0$ is the residual information after accounting for the global shift.
\begin{align*}
\text{diffPattern}(x_0, \sigma) = \lambda x. \sigma x - \sigma x_0
\end{align*}
\lean{Expressive.diffPattern}
\end{definition}

\begin{proposition}
Adding a constant to the scale field does not change the difference pattern.
\begin{align*}
\text{diffPattern}(x_0, \lambda x. \sigma x + c) = \text{diffPattern}(x_0, \sigma)
\end{align*}
\begin{proof}
Evaluate both sides at an arbitrary $x$: the left side is $(\sigma x + c) - (\sigma x_0 + c)$, which simplifies by cancelling the additive constant $c$ to $\sigma x - \sigma x_0$, the right side. Since the two functions agree pointwise, they are equal.
\end{proof}
\lean{Expressive.diffPattern\_shift}
\end{proposition}

\begin{theorem}
Every observable is a function of differences.  Given invariance, evaluating on the difference pattern gives the same answer as evaluating on the field itself.  So no observable can depend on anything beyond differences.
\begin{align*}
\text{If } \text{FiducialInvariant}(f), \text{ then } f(\sigma) = f(\text{diffPattern}(x_0, \sigma))
\end{align*}
\begin{proof}
Apply the invariance hypothesis to the field $\text{diffPattern}(x_0,\sigma)$ shifted by the constant $c=\sigma(x_0)$: $f(\text{diffPattern}(x_0,\sigma)) = f(\lambda x.\,\text{diffPattern}(x_0,\sigma)(x) + \sigma(x_0))$. Since $\text{diffPattern}(x_0,\sigma)(x)+\sigma(x_0) = (\sigma x - \sigma x_0) + \sigma x_0 = \sigma x$ for every $x$, the shifted field is $\sigma$ itself, so this equals $f(\sigma)$. Reading the chain backwards gives $f(\sigma)=f(\text{diffPattern}(x_0,\sigma))$.
\end{proof}
\lean{Expressive.invariant\_factors}
\end{theorem}

\begin{theorem}
Every function of differences is an observable.
\begin{align*}
\text{FiducialInvariant}(\lambda \sigma. g(\text{diffPattern}(x_0, \sigma)))
\end{align*}
\begin{proof}
Fix $\sigma$ and $c$; the goal is $g(\text{diffPattern}(x_0,\lambda x.\sigma x+c)) = g(\text{diffPattern}(x_0,\sigma))$. By the shift-invariance of the difference pattern proved above, $\text{diffPattern}(x_0,\lambda x.\sigma x+c)=\text{diffPattern}(x_0,\sigma)$, so applying $g$ to both sides gives the goal.
\end{proof}
\lean{Expressive.diffPattern\_invariant}
\end{theorem}

\begin{theorem}
The scale channel expresses differences and nothing else.  Observables are exactly the functions of the difference pattern — the inclusion holds in both directions.  This is the precise expressive limit that A5 imposes.
\begin{align*}
\text{FiducialInvariant}(f) \Leftrightarrow \exists g, \forall \sigma, \, f(\sigma) = g(\text{diffPattern}(x_0, \sigma))
\end{align*}
\begin{proof}
($\Rightarrow$) Given $\text{FiducialInvariant}(f)$, take $g=f$ itself; the required identity $f(\sigma)=f(\text{diffPattern}(x_0,\sigma))$ is exactly the invariant-factors theorem above. ($\Leftarrow$) Given $g$ with $f(\sigma)=g(\text{diffPattern}(x_0,\sigma))$ for all $\sigma$, fix $\sigma,c$; then $f(\lambda x.\sigma x+c) = g(\text{diffPattern}(x_0,\lambda x.\sigma x+c)) = g(\text{diffPattern}(x_0,\sigma)) = f(\sigma)$, using the hypothesis twice and the shift-invariance of the difference pattern in between.
\end{proof}
\lean{Expressive.invariant\_iff\_diffPattern}
\end{theorem}

\begin{theorem}
Absolute scale is not expressible.  Reading the scale at a point is not invariant, so it is not an observable of this framework at all — not merely unmeasured, but outside the vocabulary.
\begin{align*}
\neg \text{FiducialInvariant}(\lambda \sigma. \sigma(x))
\end{align*}
\begin{proof}
Suppose for contradiction that evaluation at $x$ were invariant. Apply the invariance to the zero field $\sigma\equiv0$ and shift $c=1$: invariance demands $(\lambda x'. 0+1)(x) = 0(x)$, i.e.\ $1=0$, which is false. Hence evaluation is not invariant.
\end{proof}
\lean{Expressive.eval\_not\_invariant}
\end{theorem}

\begin{theorem}
The scale of a comparison says nothing about its rotation.  Two comparisons with the same scale differ by an arbitrary element of the commutator subgroup, so fixing the scale leaves rotation entirely free.
\begin{align*}
\text{If } k \in \text{commutator}(\Gamma), \text{ then } \text{scale}(k \cdot g) = \text{scale}(g)
\end{align*}
\begin{proof}
Since $\text{scale}$ is a homomorphism, $\text{scale}(k\cdot g) = \text{scale}(k)\cdot\text{scale}(g)$. Because $k$ lies in the commutator subgroup, it lies in the kernel of the abelianization map, so $\text{scale}(k)=1$. Hence $\text{scale}(k\cdot g) = 1\cdot\text{scale}(g) = \text{scale}(g)$.
\end{proof}
\lean{Expressive.scale\_says\_nothing\_about\_rotation}
\end{theorem}

\begin{theorem}
The rotation says nothing about the scale.  Every commutator has trivial scale, so knowing the rotation part determines nothing about the scale part.
\begin{align*}
\text{scale}(g \cdot h \cdot g^{-1} \cdot h^{-1}) = 1
\end{align*}
\begin{proof}
The element $g\cdot h\cdot g^{-1}\cdot h^{-1}$ is a commutator, hence lies in the kernel of the abelianization homomorphism $\text{scale}$; every commutator maps to the identity, giving $\text{scale}(g\cdot h\cdot g^{-1}\cdot h^{-1})=1$.
\end{proof}
\lean{Expressive.rotation\_says\_nothing\_about\_scale}
\end{theorem}

\begin{theorem}
The channels are independent.  Fixing either one constrains the other not at all.  This is by construction — the scale channel is a quotient and the rotation channel is its kernel — so no amount of further work inside the present axioms can make one test the other.
\begin{align*}
(\forall k \in \text{commutator}(\Gamma), \, \text{scale}(k \cdot g) = \text{scale}(g)) \,\wedge\, (\forall a, b : \Gamma, \, \text{scale}(a \cdot b \cdot a^{-1} \cdot b^{-1}) = 1)
\end{align*}
\begin{proof}
The two conjuncts are exactly the two preceding theorems: the first is scale-says-nothing-about-rotation applied for each $k$, and the second is rotation-says-nothing-about-scale applied for each pair $a,b$.
\end{proof}
\lean{Expressive.channels\_independent}
\end{theorem}

\begin{theorem}
A comparison is recovered from its scale only up to a rotation, and from its rotation only up to a scale: neither channel alone is complete, and together they are just the original comparison again.
\begin{align*}
\text{If } \text{scale}(g) = \text{scale}(h), \text{ then } g \cdot h^{-1} \in \text{commutator}(\Gamma)
\end{align*}
\begin{proof}
This is the standard characterization of the abelianization map: two elements have equal scale exactly when they differ by an element of the commutator subgroup (the kernel of the abelianization homomorphism). Applying this characterization to the hypothesis $\text{scale}(g)=\text{scale}(h)$ gives $g\cdot h^{-1}\in\text{commutator}(\Gamma)$.
\end{proof}
\lean{Expressive.neither\_channel\_complete}
\end{theorem}

\begin{theorem}
Within the present axioms, no function of the scale channel can determine the rotation channel, because the scale is blind to exactly the subgroup the rotation lives in.
\begin{align*}
\forall (F : \text{Abelianization}(\Gamma) \to \Gamma)(g : \Gamma)(k : \Gamma), \, k \in \text{commutator}(\Gamma) \Rightarrow \text{scale}(k \cdot g) = \text{scale}(g)
\end{align*}
\begin{proof}
The statement does not depend on $F$ at all — it is scale-says-nothing-about-rotation, restated with an unused universally quantified function $F$ to emphasise that no such $F$ could recover the rotation channel from the scale channel; the proof is that theorem applied to $g$ and $k$.
\end{proof}
\lean{Expressive.no\_channel\_coupling}
\end{theorem}

\subsection{CrossCheck}

A threshold density $\rho$ appears in two places: in the mass spectrum as the average rate of thresholds per $e$-fold, and in the running of couplings as the rate-controlling density.  Since the framework names them the same, the natural move is to measure one from the spectrum and one from running and check they agree — a genuine cross-sector test.  That test was run.  It is not a test, because the two quantities are not the same: one is a raw count of thresholds, the other is sector-restricted and signed.  The proposal was an over-reach and this file records why.

\begin{definition}
Threshold density from all twelve Standard Model mass scales, per $e$-fold.
\begin{align*}
\rho_{\text{all}} = \frac{11}{12.731}
\end{align*}
\lean{CrossCheck.rhoAll}
\end{definition}

\begin{definition}
Threshold density from the six quarks alone — the structures a strong coupling responds to.
\begin{align*}
\rho_{\text{quark}} = \frac{5}{11.271}
\end{align*}
\lean{CrossCheck.rhoQuark}
\end{definition}

\begin{proposition}
$\rho_{\text{all}}$ lies between $0.86$ and $0.87$.
\begin{align*}
0.86 < \rho_{\text{all}} < 0.87
\end{align*}
\begin{proof}
Substitute $\rho_{\text{all}}=11/12.731$ and evaluate: $11/12.731 \approx 0.8640$, which lies strictly between $0.86$ and $0.87$ by direct numerical computation.
\end{proof}
\lean{CrossCheck.rhoAll\_value}
\end{proposition}

\begin{proposition}
$\rho_{\text{quark}}$ lies between $0.44$ and $0.45$.
\begin{align*}
0.44 < \rho_{\text{quark}} < 0.45
\end{align*}
\begin{proof}
Substitute $\rho_{\text{quark}}=5/11.271$ and evaluate: $5/11.271 \approx 0.4436$, which lies strictly between $0.44$ and $0.45$ by direct numerical computation.
\end{proof}
\lean{CrossCheck.rhoQuark\_value}
\end{proposition}

\begin{theorem}
The two densities disagree.
\begin{align*}
\rho_{\text{all}} \neq \rho_{\text{quark}}
\end{align*}
\begin{proof}
Suppose for contradiction $\rho_{\text{all}}=\rho_{\text{quark}}$. By the two bounding propositions above, $0.86 < \rho_{\text{all}}$ and $\rho_{\text{quark}} < 0.45$; substituting the assumed equality gives $0.86 < \rho_{\text{quark}} < 0.45$, an immediate numerical contradiction since $0.86 > 0.45$.
\end{proof}
\lean{CrossCheck.densities\_differ}
\end{theorem}

\begin{theorem}
By roughly a factor of two — which is a fact about how much of the Standard Model happens to be coloured, not a prediction of anything.
\begin{align*}
1.9 < \frac{\rho_{\text{all}}}{\rho_{\text{quark}}} < 2.0
\end{align*}
\begin{proof}
By the two bounding propositions, $\rho_{\text{quark}}>0.44>0$, so multiplying through by $\rho_{\text{quark}}$ preserves the direction of the inequalities. The lower bound $1.9<\rho_{\text{all}}/\rho_{\text{quark}}$ is equivalent, after clearing the denominator, to $1.9\cdot\rho_{\text{quark}}<\rho_{\text{all}}$, which follows from $\rho_{\text{quark}}<0.45$ and $\rho_{\text{all}}>0.86$ by a numerical bound ($1.9\times0.45=0.855<0.86$). The upper bound $\rho_{\text{all}}/\rho_{\text{quark}}<2.0$ is equivalent to $\rho_{\text{all}}<2.0\cdot\rho_{\text{quark}}$, which follows from $\rho_{\text{all}}<0.87$ and $\rho_{\text{quark}}>0.44$ ($2.0\times0.44=0.88>0.87$).
\end{proof}
\lean{CrossCheck.density\_ratio\_near\_two}
\end{theorem}

\begin{theorem}
Within a sector, the relation has no content.  One equation and one unknown: given measurements of the beta drift and the threshold density this defines the weight $\kappa$, and substituting back returns the input.  Nothing has been tested.
\begin{align*}
\text{If } dN/dt \neq 0, \text{ then } (db/dt / dN/dt) \cdot dN/dt = db/dt
\end{align*}
\begin{proof}
Since $dN/dt\neq0$, division by it and multiplication by it are mutually cancelling field operations: $(db/dt \,/\, dN/dt)\cdot dN/dt = db/dt \cdot (dN/dt)^{-1}\cdot dN/dt = db/dt\cdot 1 = db/dt$.
\end{proof}
\lean{CrossCheck.sector\_relation\_is\_identity}
\end{theorem}

\begin{theorem}
Carried out for QCD between the charm and top thresholds: two flavours appear, the beta coefficient drifts by $2/3$ per flavour, and the extracted weight is $2/3$ — the input written backwards.
\begin{align*}
\text{If } dN/dt \neq 0, \text{ then } \left(\frac{2}{3} \cdot dN/dt\right) / dN/dt = \frac{2}{3}
\end{align*}
\begin{proof}
Since $dN/dt\neq0$, it is invertible in the field of reals, so $(\tfrac23\cdot dN/dt)/dN/dt = \tfrac23\cdot dN/dt\cdot(dN/dt)^{-1} = \tfrac23\cdot1 = \tfrac23$.
\end{proof}
\lean{CrossCheck.qcd\_extraction\_returns\_input}
\end{theorem}

\begin{theorem}
What a genuine test would require: two independent determinations of the same weight, so that the second can check the first.  The framework supplies only one, because it has only one coupling-like quantity per sector and no coupling that responds to every structure.
\begin{align*}
\text{If } dN/dt \neq 0 \text{ and } \kappa = db/dt / dN/dt, \text{ then } db/dt = \kappa \cdot dN/dt
\end{align*}
\begin{proof}
Substitute the definition $\kappa = db/dt / dN/dt$ into the goal, reducing it to $db/dt = (db/dt / dN/dt)\cdot dN/dt$, which is exactly the sector-relation-is-identity theorem above (using $dN/dt\neq0$).
\end{proof}
\lean{CrossCheck.no\_second\_determination}
\end{theorem}

\begin{definition}
Sum of the five quark log-gaps.
\begin{align*}
\text{quarkGapSum} = 11.270
\end{align*}
\lean{CrossCheck.quarkGapSum}
\end{definition}

\begin{definition}
Sum of their squares.
\begin{align*}
\text{quarkGapSqSum} = 31.583624
\end{align*}
\lean{CrossCheck.quarkGapSqSum}
\end{definition}

\begin{definition}
Mean of the five quark log-gaps.
\begin{align*}
\text{quarkMeanGap} = \text{quarkGapSum} / 5
\end{align*}
\lean{CrossCheck.quarkMeanGap}
\end{definition}

\begin{definition}
Variance of the five quark log-gaps.
\begin{align*}
\text{quarkVariance} = \text{quarkGapSqSum} / 5 - \text{quarkMeanGap}^2
\end{align*}
\lean{CrossCheck.quarkVariance}
\end{definition}

\begin{definition}
Squared coefficient of variation for the quark sector.
\begin{align*}
\text{quarkCvSq} = \text{quarkVariance} / \text{quarkMeanGap}^2
\end{align*}
\lean{CrossCheck.quarkCvSq}
\end{definition}

\begin{theorem}
$CV^2 = 0.243$ (so $CV = 0.493$) — the twelfth percentile of the constant-rate prediction for five gaps.  Low, and more regular than the prediction typically gives, but not excluded.  What survives is the per-sector prediction: each sector's thresholds should be a constant-rate process.
\begin{align*}
0.24 < \text{quarkCvSq} < 0.25
\end{align*}
\begin{proof}
Unfold the chain of definitions: $\text{quarkMeanGap}=11.270/5=2.254$, $\text{quarkVariance}=31.583624/5 - 2.254^2 = 6.3167248 - 5.080516 = 1.2362088$, and $\text{quarkCvSq}=1.2362088/2.254^2 = 1.2362088/5.080516\approx0.24333$. This lies strictly between $0.24$ and $0.25$ by direct numerical computation.
\end{proof}
\lean{CrossCheck.quarkCvSq\_value}
\end{theorem}

\begin{theorem}
The quark sector is not a geometric tower: its gap variance is positive.
\begin{align*}
\text{quarkCvSq} \neq 0
\end{align*}
\begin{proof}
By the previous theorem $0.24<\text{quarkCvSq}$, so in particular $\text{quarkCvSq}$ is strictly positive and hence cannot equal $0$.
\end{proof}
\lean{CrossCheck.quark\_not\_geometric}
\end{theorem}

\subsection{MassAudit}

The mass law $m_k = m_0 e^{|k|\Delta}$ was presented as a consequence of axiom A3 ("mass is inverse scale") and the scale stepping by $\Delta$ per unit winding.  Testing it against lattice glueball data falsified the constant mass ratio.  But the law was never derived — it was assumed.  The topology establishes only that winding is an integer, additive, conserved, and invariant under orientation reversal.  What the mass law additionally requires — how strained the scale field is by a defect of winding $k$ — is a dynamical question about energetics, which no theorem addresses.  This file exhibits two different mass laws, exponential and quadratic, both fully compatible with every proved property of winding, so the topology does not fix the mass law.

\begin{definition}
The law asserted in earlier work: a geometric tower.
\begin{align*}
\text{massExp}(m_0, \Delta, k) = m_0 \cdot e^{|k|\Delta}
\end{align*}
\lean{MassAudit.massExp}
\end{definition}

\begin{definition}
An equally available alternative: strain energy quadratic in the winding, which is what a naive energetic estimate for a topological defect gives.
\begin{align*}
\text{massQuad}(m_0, c, k) = m_0 + c \cdot k^2
\end{align*}
\lean{MassAudit.massQuad}
\end{definition}

\begin{proposition}
The exponential law satisfies particle–antiparticle mass degeneracy.
\begin{align*}
\text{massExp}(m_0, \Delta, -k) = \text{massExp}(m_0, \Delta, k)
\end{align*}
\begin{proof}
Unfold $\text{massExp}(m_0,\Delta,k) = m_0\cdot e^{|k|\Delta}$. Since $|-k|=|k|$ for any integer $k$, the exponents agree, so $\text{massExp}(m_0,\Delta,-k) = m_0\cdot e^{|-k|\Delta} = m_0\cdot e^{|k|\Delta} = \text{massExp}(m_0,\Delta,k)$.
\end{proof}
\lean{MassAudit.massExp\_neg}
\end{proposition}

\begin{proposition}
The quadratic law also satisfies particle–antiparticle mass degeneracy.
\begin{align*}
\text{massQuad}(m_0, c, -k) = \text{massQuad}(m_0, c, k)
\end{align*}
\begin{proof}
Unfold $\text{massQuad}(m_0,c,k) = m_0+c\cdot k^2$. Since $(-k)^2=k^2$ for any integer $k$ cast to $\mathbb{R}$, $\text{massQuad}(m_0,c,-k) = m_0 + c\cdot(-k)^2 = m_0+c\cdot k^2 = \text{massQuad}(m_0,c,k)$ by ring arithmetic.
\end{proof}
\lean{MassAudit.massQuad\_neg}
\end{proposition}

\begin{theorem}
The exponential law is positive for positive parameters.
\begin{align*}
\text{If } 0 < m_0, \text{ then } 0 < \text{massExp}(m_0, \Delta, k)
\end{align*}
\begin{proof}
Unfold $\text{massExp}(m_0,\Delta,k)=m_0\cdot e^{|k|\Delta}$: the exponential factor $e^{|k|\Delta}$ is always strictly positive, and $m_0>0$ by hypothesis, so the product of two positive reals is positive.
\end{proof}
\lean{MassAudit.massExp\_pos}
\end{theorem}

\begin{theorem}
The quadratic law is also positive for positive parameters.
\begin{align*}
\text{If } 0 < m_0 \text{ and } 0 \le c, \text{ then } 0 < \text{massQuad}(m_0, c, k)
\end{align*}
\begin{proof}
Unfold $\text{massQuad}(m_0,c,k)=m_0+c\cdot k^2$. Since $c\geq0$ and $k^2\geq0$, the product $c\cdot k^2\geq0$; adding this nonnegative quantity to $m_0>0$ preserves positivity: $0<m_0\leq m_0+c\cdot k^2$.
\end{proof}
\lean{MassAudit.massQuad\_pos}
\end{theorem}

\begin{proposition}
Both laws agree at zero winding.
\begin{align*}
\text{massExp}(m_0, \Delta, 0) = m_0 = \text{massQuad}(m_0, c, 0)
\end{align*}
\begin{proof}
For the first equality, unfold $\text{massExp}(m_0,\Delta,0) = m_0\cdot e^{|0|\Delta} = m_0\cdot e^{0} = m_0\cdot1=m_0$, using $|0|=0$ and $e^0=1$. For the second, unfold $\text{massQuad}(m_0,c,0) = m_0+c\cdot0^2 = m_0+0=m_0$.
\end{proof}
\lean{MassAudit.massExp\_zero}
\end{proposition}

\begin{proposition}
Both laws agree at zero winding (quadratic version).
\begin{align*}
\text{massQuad}(m_0, c, 0) = m_0
\end{align*}
\begin{proof}
Immediate from the definition: $\text{massQuad}(m_0,c,0) = m_0 + c\cdot0^2 = m_0+0 = m_0$ by ring arithmetic.
\end{proof}
\lean{MassAudit.massQuad\_zero}
\end{proposition}

\begin{theorem}
The topology does not determine the mass law.  Two laws that are genuinely different — their ratio at winding $1$ differs — both satisfy every property the topology has fixed: positivity, agreement at $k=0$, and invariance under orientation reversal.  Therefore the geometric tower is an additional hypothesis, and the constant mass ratio it predicts is a consequence of that hypothesis rather than of the framework.
\begin{align*}
\exists f, g : \mathbb{Z} \to \mathbb{R}, \quad &(\forall k, f(-k) = f(k)) \,\wedge\, (\forall k, g(-k) = g(k)) \\
&\,\wedge\, (\forall k, 0 < f(k)) \,\wedge\, (\forall k, 0 < g(k)) \\
&\,\wedge\, (f(0) = g(0)) \,\wedge\, (f(1) \neq g(1))
\end{align*}
\begin{proof}
Take $f=\text{massExp}(1,1,\cdot)$ and $g=\text{massQuad}(1,1,\cdot)$. Orientation-reversal symmetry for both is $\text{massExp\_neg}$ and $\text{massQuad\_neg}$ above; positivity of both for all $k$ is $\text{massExp\_pos}$ and $\text{massQuad\_pos}$ (with $m_0=1>0$, $c=1\geq0$). Agreement at $k=0$ is $\text{massExp\_zero}$ and $\text{massQuad\_zero}$, both giving $1$. For the final inequality, compute $g(1)=\text{massQuad}(1,1,1) = 1+1\cdot1^2=2$, while $f(1)=\text{massExp}(1,1,1)=1\cdot e^{1}=e$. Since $e>2.7$ (a standard numerical bound on Euler's constant) while $g(1)=2<2.7$, $f(1)\neq g(1)$.
\end{proof}
\lean{MassAudit.mass\_law\_underdetermined}
\end{theorem}

\begin{theorem}
Both the exponential and quadratic laws reproduce the one mass statement that is derived — the particle–antiparticle degeneracy — so that prediction survives the retraction while the tower does not.
\begin{align*}
\text{massExp}(m_0, \Delta, -k) = \text{massExp}(m_0, \Delta, k) \,\wedge\, \text{massQuad}(m_0, c, -k) = \text{massQuad}(m_0, c, k)
\end{align*}
\begin{proof}
The two conjuncts are exactly $\text{massExp\_neg}$ and $\text{massQuad\_neg}$ proved above, applied to the given parameters.
\end{proof}
\lean{MassAudit.both\_laws\_antiparticle\_degenerate}
\end{theorem}

\begin{theorem}
The quadratic law has a non-constant consecutive ratio, so a spectrum that fails the geometric-tower test is not thereby in conflict with a defect picture at all.
\begin{align*}
\text{massQuad}(1, 1, 1) \cdot \text{massQuad}(1, 1, 1) \neq \text{massQuad}(1, 1, 2) \cdot \text{massQuad}(1, 1, 0)
\end{align*}
\begin{proof}
Unfold and compute each factor: $\text{massQuad}(1,1,1)=1+1^2=2$, $\text{massQuad}(1,1,2)=1+2^2=5$, $\text{massQuad}(1,1,0)=1+0^2=1$. So the left side is $2\cdot2=4$ and the right side is $5\cdot1=5$, and $4\neq5$ by numerical computation.
\end{proof}
\lean{MassAudit.massQuad\_ratio\_not\_constant}
\end{theorem}

\subsection{Audit}

This file is the register — the single authoritative list of what was retracted, corrected, cited, assumed, or proved with no auxiliary hypotheses. Nothing here is new mathematics; it is the standing answer to ``Does this rest on the axioms alone?'' and it is meant to be read before every claim.  The retractions are real: a geometric mass tower (assumed, not derived), a colour identification (wrong by a factor of $3/2$), a cross-sector test (measuring two different quantities), and a codimension analysis (true of a circle-valued scale, not this framework). The corrections are more numerous than the retractions — most claims survived with their scope narrowed. An earlier version of this register was incomplete and omitted A6$^\prime$ (the scale response law) and the assumption that defects are uniform. The register grew under external audit, which is the point of having one.

\begin{definition}
How many times the scalar-for-directional substitution has been made and caught.  Recorded so the count cannot quietly drift.
\begin{align*}
\text{scalarForDirectionalCount} = 5
\end{align*}
\lean{Audit.scalarForDirectionalCount}
\end{definition}

\begin{theorem}
The scalar-for-directional substitution appears exactly five times in the development.
\begin{align*}
\text{scalarForDirectionalCount} = 5
\end{align*}
\begin{proof}
Immediate from the definition $\text{scalarForDirectionalCount}=5$.
\end{proof}
\lean{Audit.scalar\_for\_directional\_count}
\end{theorem}

\begin{definition}
Claims retracted outright.
\begin{align*}
\text{retractedCount} = 4
\end{align*}
\lean{Audit.retractedCount}
\end{definition}

\begin{definition}
Claims corrected in scope, where the theorem survives and the gloss did not.
\begin{align*}
\text{correctedCount} = 6
\end{align*}
\lean{Audit.correctedCount}
\end{definition}

\begin{definition}
External results cited and not proved.
\begin{align*}
\text{citedCount} = 2
\end{align*}
\lean{Audit.citedCount}
\end{definition}

\begin{definition}
Assumptions registered rather than smuggled: A6$^\prime$ (the scale response law), the assumption that defects are uniform in the symmetric space's volume, and A7 itself.
\begin{align*}
\text{assumedCount} = 3
\end{align*}
\lean{Audit.assumedCount}
\end{definition}

\begin{theorem}
The register is not empty.  A development with nothing retracted has not been audited.  The counts are carried so that a reader can weigh the results against what had to be taken back to get them.
\begin{align*}
\text{retractedCount} \neq 0 \,\wedge\, \text{correctedCount} \neq 0 \,\wedge\, \text{citedCount} \neq 0 \,\wedge\, \text{assumedCount} \neq 0
\end{align*}
\begin{proof}
Each conjunct unfolds to a concrete numeral inequality: $\text{retractedCount}=4\neq0$, $\text{correctedCount}=6\neq0$, $\text{citedCount}=2\neq0$, $\text{assumedCount}=3\neq0$, each true by direct computation on natural numbers.
\end{proof}
\lean{Audit.register\_nonempty}
\end{theorem}

\begin{theorem}
Corrections outnumber retractions: more claims survived with their scope narrowed than had to be withdrawn.
\begin{align*}
\text{retractedCount} < \text{correctedCount}
\end{align*}
\begin{proof}
Unfolding both definitions reduces the claim to $4<6$, true by direct computation on natural numbers.
\end{proof}
\lean{Audit.more\_corrected\_than\_retracted}
\end{theorem}

\begin{definition}
Structures carrying hypotheses the framework does not supply — true theorems whose physical reading is conditional.  These are theorems about anything meeting their hypothesis, not statements of what the axioms deliver.
\begin{align*}
\text{unwitnessedStructures} = 6
\end{align*}
\lean{Audit.unwitnessedStructures}
\end{definition}

\begin{theorem}
The register grew under external audit, which is the point of having one.  An audit that finds nothing has not been run adversarially.
\begin{align*}
0 < \text{unwitnessedStructures}
\end{align*}
\begin{proof}
Unfolding the definition reduces the claim to $0<6$, true by direct computation on natural numbers.
\end{proof}
\lean{Audit.register\_grew\_under\_audit}
\end{theorem}

\begin{definition}
Number of theorems put through axiom verification — every theorem in the development, generated from the sources rather than curated.
\begin{align*}
\text{auditedTheorems} = 650
\end{align*}
\lean{Audit.auditedTheorems}
\end{definition}

\begin{definition}
Occurrences of unresolved axioms in that audit.
\begin{align*}
\text{sorryAxCount} = 0
\end{align*}
\lean{Audit.sorryAxCount}
\end{definition}

\begin{theorem}
Nothing in the development rests on a gap, as a number rather than a mood.  Every audited theorem reduces to `propext`, `Classical.choice` and `Quot.sound` and to nothing else.
\begin{align*}
\text{sorryAxCount} = 0 \,\wedge\, 0 < \text{auditedTheorems}
\end{align*}
\begin{proof}
The first conjunct is the definition $\text{sorryAxCount}=0$ itself. The second unfolds $\text{auditedTheorems}=650$ to the numerical fact $0<650$, true by direct computation.
\end{proof}
\lean{Audit.clean\_count}
\end{theorem}


\end{document}
